% ============================================================
%  R. Nielsen, Natur og Aand.
%  Bidrag til en med Physiken stemmende Naturphilosophie.
%  Kjøbenhavn: Forlagt af J. H. Schubothes Boghandel, 1873.
%  (xiv + 549 pp.)
%  TRANSLATION (English) --- Hans Halvorson
%
%  SOURCE OF TRUTH is transcription.tex in this directory (Danish,
%  COMPLETE: 549/549 pages, 0 markers, builds 0/0 at 318 pp.).
%  Translate FROM it, never from the scan.
%
%  METHOD: ../../../TRANSLATION-PLAYBOOK.md, the standing method for every
%  book in this repo. Book-specific state and conventions live in
%  TRANSLATION-RESUME-NOTES.md in this directory. Read both before a batch.
%
%  STRUCTURE mirrors transcription.tex 1:1 --- three Deele, nine Kapitler,
%  thirty Afdelinger, in the same order, with \streg at the same breaks.
%  Every head carries its Danish in a % comment above it, so the mirror
%  stays checkable line by line.
%
%  MATHEMATICS. Roughly 83 pages carry real mathematics (Kepler's laws
%  derived in the footnote to p. 160; algebra through the mechanics
%  chapters). Formulas are COPIED from the transcription unchanged ---
%  Nielsen's notation, his \operatorname{tg}, his \div for minus at p. 46,
%  his mixed primes and raised numerals. Translate the prose around them;
%  do not modernise, retype or tidy the mathematics.
%
%  PRINTED DEFECTS. The transcription keeps the book's misprints as printed
%  and logs Nielsen's own errata in a % comment at the site. The translation
%  follows the PRINTED reading likewise and carries the comment across. Two
%  heads are affected: the division head on p. 233 prints "B." where the
%  sequence wants "C.", and the run-in head on p. 180 prints "c)" where it
%  wants "b)". Both stand as printed.
% ============================================================
\documentclass[12pt,a4paper]{book}
\usepackage[utf8]{inputenc}
\usepackage[T1]{fontenc}
\usepackage{libertinus}
\usepackage{libertinust1math}
\usepackage{amsmath}     % this book has real mathematics
\usepackage{textalpha}   % Greek copied verbatim from the Danish
\usepackage{graphicx}    % \reflectbox, for the old Danish ">o:" id-est mark
\DeclareUnicodeCharacter{0254}{\reflectbox{c}}
\usepackage[english]{babel}
\usepackage[protrusion=true,expansion=true]{microtype}
\usepackage{setspace}
\usepackage[margin=1.2in]{geometry}
\usepackage{enumitem}
\usepackage{fancyhdr}
\usepackage{hyperref}
\setlength{\headheight}{28pt}
\addtolength{\topmargin}{-13.5pt}

\hypersetup{pdftitle={Nature and Spirit},pdfauthor={R. Nielsen},hidelinks}

% Footnotes are marked *) in this book, as in the transcription.
\renewcommand{\thefootnote}{*)}

\pagestyle{fancy}
\fancyhf{}
\fancyhead[C]{\thepage}
\renewcommand{\headrulewidth}{0pt}

% ---------------------------------------------------------------------------
% Heading macros --- identical to the transcription's, so the two files set
% the same shapes. None forces a page break: two of the nine chapter heads do
% not open a page in the original (pp. 125 and 509).
% ---------------------------------------------------------------------------
\newcommand{\deel}[2]{%
  \par\vspace{2.2\baselineskip}%
  {\centering\large #1\par\vspace{0.3\baselineskip}\Large #2\par}%
  \addcontentsline{toc}{part}{#1 #2}%
  \vspace{1.4\baselineskip}\par\nointerlineskip}
\newcommand{\kapitel}[2]{%
  \par\vspace{1.8\baselineskip}%
  {\centering #1\par\vspace{0.25\baselineskip}\large #2\par}%
  \addcontentsline{toc}{chapter}{#1 #2}%
  \vspace{1.1\baselineskip}\par\nointerlineskip}
\newcommand{\afdeling}[2]{%
  \par\vspace{1.3\baselineskip}%
  {\centering #1\par\vspace{0.2\baselineskip}\normalsize #2\par}%
  \addcontentsline{toc}{section}{#1 #2}%
  \vspace{0.8\baselineskip}\par\nointerlineskip}
\newcommand{\streg}{\begin{center}\rule{2cm}{0.4pt}\end{center}}
% Run-in heads: a) b) c) print on the same line as the paragraph they open.
\newcommand{\runin}[1]{\emph{#1}\hspace{0.4em}}
\newcommand{\subrunin}[1]{\emph{#1}\hspace{0.4em}}

\setstretch{1.15}

\title{Nature and Spirit\\[0.4em]
  \large A Contribution to a Philosophy of Nature in Accord with Physics}
\author{R.\ Nielsen\\[4pt]
  \small Translated from the Danish by Hans Halvorson}
\date{Copenhagen: J.\ H.\ Schubothe's Bookshop, 1873}

\begin{document}
\maketitle
\tableofcontents

% ============================================================
%  FRONT MATTER
%
%  The Danish Indhold (printed V--X) is not reproduced: \tableofcontents
%  above generates the English one from the heads. The Indledning (printed
%  XI--XIV) is translated below.
% ============================================================

% --- p. XI ---
{\centering\large Introduction.\par}

\streg  % short, centred, SINGLE rule under the head, as printed

The influence of the science of nature upon the cultural life of the present
day is perceptible in every direction. To give a description of the great
practical results to which the discoveries of the latest age have led is
superfluous: steamships, railways, telegraphs are signs so eloquent that men
from the youth of Napoleon the First, were they to wake in their graves and
look up, would believe themselves set down in a world of wonders. A science
that reaches so mightily into all actual relations cannot, of course, confine
itself to serving merely finite interests and promoting what is useful in the
civil sense; its origin and its goal are of a higher kind, it wills a knowledge
of truth: its significance is essentially theoretical. It is natural science
that has opened men's eyes to what the wisest of antiquity still saw only as
through mists --- the laws of nature, the connection of things, the secret
kinship of the forces. The visible universe has become another, for the
illumination has become another. The arrangement of the solar system, the
development of the terrestrial globe, the natural descent of the living genera
and species, all this presents itself in broad strokes, with distinct outlines,
in a new light, reality's own daylight. There are, to be sure, still those who
harbour a secret mistrust of the reliability of research, but there is no one
able to withdraw himself from its influence. Those who hold to the old strive
against it in vain; the new picture of the world is actually there; it is
glimpsed in everything, they meet it everywhere, they see it even when they
close their eyes.

But precisely the circumstance that the natural sciences together make up a
spiritual great power, to whose authority, where
% --- p. XII ---
real insight into real things is at issue, the civilized world so willingly
bows, invites us to consider whether the common consciousness does not at the
same time attach to it an unwarranted afterthought, a misunderstanding that may
have grave consequences both for science itself and for life.

It is then, in the first place, a question whether the science of nature, so
long as it keeps within the bounds of experience, can own as its task: to solve
the riddle of existence by fathoming what nature is in itself, whether it has
any origin and temporal beginning or exists from eternity, whether it is matter
that brings forth Spirit or the reverse. If the science of experience, by its
character and method, neither can nor will concern itself with the solution of
such tasks, then it is surely put in a false position when one ascribes to it
an endeavour it does not have, and sets it up as judge over what it cannot
judge. Next, it is doubtful which of the many mutually combating world-views
--- materialism, pantheism, atheism, nihilism, etc.,
% The Danish prints the contracted "osv.", not "o. s. v."; read on the KB copy.
all of which appeal to ``Science'' --- can properly be said to have the full
support of the scientific study of nature. It is well known that the
investigators of nature themselves, where ``the theories'' are concerned,
express themselves in different spirits and proceed from highly different
hypotheses. Even though an investigator with a famous name holds it probable
that man is a modified ape, such a way of conceiving may not on that account be
confused, without more ado, with an infallible fact and exalted into a dogma.
In assumptions of this kind a distinction must surely be made between the
construing of the facts and the facts themselves, between the investigator's
individual opinion and his scientific proofs. The science of nature leans upon
hypotheses; but it treats its hypotheses critically, for it knows that their
nature is ambiguous. Every hypothesis has, so to speak, two faces: one severe
and earnest, turned toward the facts and aiming unswervingly at the discovery
of their hidden connection; another smiling, more ingratiating, which addresses
the general public and would win itself a place among its favourite doctrines.

% --- p. XIII ---
From whichever side one regards the matter, it is as impossible by the way of
hypotheses as by continued observation of facts --- with whose infinite
multiplicity research never comes to an end --- to get clear about those
fundamental relations to which all knowledge of nature must essentially be led
back; the investigation of
% MISPRINT, not corrected: the Danish sets ``tilhage'' where the sense requires
% ``tilbage'' (back). Both copies carry it, so the fault is the compositor's;
% the errata leaf is silent about the place. Translated by the sense.
these fundamental relations, so far as it can be carried through on the terms
of human knowing, falls not under the science of experience but under the
philosophy of nature.

The philosophy of nature is indeed dependent on the science of experience,
inasmuch as it must use the yield of the observations and keep as exactly as
possible in the track of research; but it has for all that its own tasks, its
own method, its own goal, and, in the knowing consciousness, its own
standpoint.

It is difficult in few words to give a generally intelligible account of the
course and the essence of the philosophy of nature; a closer understanding
depends on its being seen in the doing how its work is carried out. So much can
be said in advance, that it is to build neither upon hypotheses nor upon
abstract speculations.

As concerns the present contribution, the following propositions may serve as
an introduction. The collected impression of all contemplation of nature is an
impression of power: nature is a realm of power. The eternal power that reveals
itself in the order of nature is visible and invisible, visible in corporeal
things, invisible in the concepts of things. From the sensuous side power shows
itself as object-power, a visible power in the series of the phenomena, in a
concatenation of members of which the one incessantly presupposes the other and
points to the other, thus as otherness-power, outward power, \emph{power in
another}.
% Letterspaced in the Danish: „M a g t / i  A n d e t“, the spacing running
% across the line break; the preceding „udvortes Magt,“ is ordinary setting.
From the non-sensuous side, by contrast, power shows itself as power in the
inward, as the invisible power of thoughts, of concepts, of wisdom, as the
power of ideas, as \emph{self-power}.
% Letterspaced in the Danish: „S e l v m a g t“.
Out of this doubling in the essence of power the philosophy of nature develops
two fundamental concepts, rich in content and decisive for all existence: out
of power in another the concept \emph{nature}, out of the essence of self-power
the concept \emph{Spirit}.
% Letterspaced in the Danish: „N a t u r“ and „A a n d“; „Begrebet“ is ordinary
% setting in both places.

% --- p. XIV ---
It goes without saying that the philosophy of nature does not set out to
exhaust the doctrine of Spirit; it does not consider nature for the sake of
Spirit, but conversely, Spirit for the sake of nature. The science of
experience presents nature in nature's own light; the philosophy of nature
presents it in the light of Spirit. Will the altered illumination be able to
open to us new insights into the mechanical, elemental and organic spheres of
nature? That must depend upon an attempt.

% The Indledning closes with a short centred DOUBLE rule about two fifths
% down p. XIV, under blank setting; \streg serves, per the house rule.
\streg


% ============================================================
%  Part One. Mechanical Nature.
% ============================================================
% Danish: \deel{Første Deel.}{Den mechaniske Natur.}
\deel{Part One.}{Mechanical Nature.}
% A short centred rule stands between the Deel head and the Kapitel head.
\streg

% Danish: \kapitel{Første Kapitel.}{Grundlagsbegreber.}   (printed pp. 1--61)
\kapitel{Chapter One.}{Foundational Concepts.}

% Danish: \afdeling{A.}{Rum og Tid.}   (printed pp. 1--18)
\afdeling{A.}{Space and Time.}

Since all outward things exist in space and could not possibly exist if space
were not; since things arise and pass away in time and could not possibly either
arise or pass away were there no time --- it surely follows from this that space
and time must belong to the essence of natural things, that they must, in other
words, be regarded as the first universal, all-embracing forms of nature. But
although space and time are universally valid forms of nature, it will
nevertheless appear in what follows that the natural subsistence of these forms
has a spiritual ground; that they are object-forms only by being selfhood-forms,
objective only by being objectified by the eternal Spirit.

\runin{a) Space.} The most important contributions to a deeper understanding of
the nature of space one will, in more recent philosophy,
% --- p. 2 ---
chiefly seek in Kant and Hegel; but the two ways of seeing must be treated
critically, for, viewed superficially, the one contradicts the other.

Kant places himself, as is well known, upon an unconditionally subjective
standpoint and regards space as an a priori form of sense, a merely human form
of consciousness. Space is no concept borrowed from experience; in order to be
able to regard things in space, we must ourselves bring the intuition of space
with us. One cannot possibly represent to oneself that there is no space; but
one can very well represent to oneself that there are no objects in space. When
there is talk of many spaces, this is not to be understood as though there were
talk of parts out of which a whole was composed: the many parts flow together
into a single space given with our necessary intuition of sense. From this,
then, the conclusion is drawn that it is not the thing in itself but only the
phenomenon of the thing, that is, the thing as a thing of sight, that is
apprehended in space; which may also be expressed thus, that \emph{space is only
a selfhood-form, but no otherness-form valid in itself.}

In sharp and decisive opposition to Kant, Hegel stands upon an absolutely
objective standpoint. For him space is, in its ``abstract universality,'' the
first form of nature, the immediately necessary form of \emph{outward being}. In
order to make clear what lies in the concept of space's ``ideal
being-beside-one-another,''
% Misprint, carried over from the Danish: the second quotation there opens with
% the guillemet » and closes with “; both copies (Google and KB) set it so.
``being-outside-one-another,'' and so forth, he chooses to show how space looks
in our representation. In order to prove that space corresponds to our thought,
we must, he says, compare our representations of space with the determinations
that lie in our concept. The many \emph{Heres} and \emph{Theres} in space are
beside one another; the one \emph{Here} is completely the same as the other, and
this abstract plurality --- without true interruption and limit --- is precisely
outwardness. Space with its innumerable points is sheer discreteness; but this
discreteness is again a flowing
% --- p. 3 ---
unity, unity in ``outside-one-another,'' complete continuity. The point has
significance only insofar as it occupies a space, and so is an outward over
against itself and over against another; in the point itself there is then again
an above, a below, a right, a left, for a last \emph{Here} is not given. Beyond
its limit space is constantly ``at home with itself'': the continuous is
discrete, the discrete continuous. The unity of both moments, discreteness and
continuity, is ``the objectively determined concept of space.'' On this way of
seeing it is then so far from being the case that objective space should rest
upon any subjective intuition whatever, that on the contrary, as ``a
non-sensible sensibility'' resting upon its own concept, it furnishes the
foundation of natural things: \emph{space is otherness-form without being
selfhood-form.}

By holding the two conceptions sharply against one another one will easily be
convinced that both are one-sided.

That Kant rightly makes the form of space a human form of consciousness can be
seen expressly from Hegel's demonstration; for this demonstration supports
itself point by point upon human intuiting, human thinking and knowing. Our
insight into the validity of objective space depends absolutely upon its
validity for us. Space is determined as the ``\emph{abstract}'' form of outward
being; but it is we --- not nature --- who abstract the form. The many different
``\emph{Heres}'' and ``\emph{Theres}'' must be posited and cancelled; but it is
we --- not nature --- who posit them and cancel them. The unity of
``discreteness'' and ``continuity'' is to think in difference, and the
difference again in unity; but it is we --- not nature --- who must think this
unity of difference: for space to show itself as objective concept, it is we who
must conceive.

But that Hegel is nonetheless right in making space an objective concept
independent of us becomes recognizable, conversely, from the Kantian
demonstration. Kant holds fast
% --- p. 4 ---
that the phenomena, the sensible things, cannot possibly be apprehended without
space; herewith he has then at the same time conceded that space is a necessary,
objective form for the things of sight. Things in nature do not, after all,
become sensible by our sensing them; rather we sense them because they are in
themselves sensible. As little as it is our sight that makes things visible, or
our thinking that makes them thinkable, just as little is it our intuition of
space that gives them existence in space. If the partition wall Kant has raised
between the essence of things in themselves and their phenomena were more than a
critical fancy, then it would be altogether unthinkable that this unknown
essence, incompatible with space, should be able to furnish any content of
experience and to fill the empty intuition from without with actual objects.

Thus: every time the form of space is made exclusively an otherness-form, it
transforms itself and becomes a selfhood-form; and when it is held fast
exclusively as a selfhood-form, it transforms itself again and becomes a form
for another. This shows plainly enough that space could not possibly be a truly
real form of nature if this form of nature did not from eternity receive
objective being in divine intuiting. The Kantian error does not consist in his
making space a form of intuition, but in his arbitrarily restricting the form of
space to being a merely human selfhood-form, and precisely thereby coming to
deny its objective validity. Hegel's fault is not that he asserts the
objectivity of space; the fundamental fault is rather that in his philosophy of
nature, as in his logic, he constantly overlooks that nothing objective exists
except what is objectified, and that no objectification is possible without an
objectifying self. That human consciousness ascribes outward being to space in
nature rests at bottom upon the tacit presupposition that the world-space which
we in our contemplation of nature take to be objective must from the beginning
be ---
% --- p. 5 ---
objectified by the creating Spirit. The starry heaven makes a symbolic
impression; it is human intuiting that, in its self-knowing, is met in the
infinity of space by the divine, as eye meets eye.

How matters stand more closely, as regards the objectification of space, with
the reciprocal determination of selfhood and otherness can be seen by a
consideration of the dimensions of space. Length, breadth and height are in
space different and yet not different. A distinction must necessarily be made
between the principal directions just named; for without such a distinction
there would be no determinations in the intuition, and then no objectively
determined space either. But upon what does the difference rest? Length is
height, and height is again breadth: the one dimension is what the other is, and
can be exchanged for the other. They are posited simultaneously and cancelled
simultaneously; they presuppose one another to infinity. From this it is seen
that the three dimensions, in their outward reciprocity, are double-sided
determinations, determinations of selfhood and determinations for another. As
determinations of selfhood they are arbitrary, as determinations of otherness
accidental, as objectified --- and so objective --- determinations of reciprocity
mutually necessary. If we choose, e.g., a rectangular coordinate system with the
axes $X$, $Y$ and $Z$, the objectification is easy to illustrate. The three axes
must have a point of intersection. Now where is this point of intersection in
space itself? Everywhere and nowhere. The point is objective only by being
objectified. To be objectified, it must be posited; to be posited, it must be
posited arbitrarily, and so by a self-determination. The arbitrariness is
necessary; for whatever place we choose, we could just as well have chosen
another place. By being posited the point has become objective; it has received
a represented existence and now furnishes a determination for another, namely
for the axes. In this objective determinateness the point is to be taken as
given; now it is precisely here, not there.
% --- p. 6 ---
Its existence is accidental, for the axes could just as well have cut one
another at any other point whatever. But the accidentality is necessary, for
whatever point were given, another could with equal right have been given. What
holds of the point of intersection holds in part of the axes. It is arbitrary to
make $Z$ the axis of height, $X$ that of length and $Y$ that of breadth; but once
arbitrariness has given them this signification, then height, length and breadth,
as determinations for one another mutually, and so as determinations of
otherness, have thereby also obtained an objective, accidental existence. The
accidentality in this existence is just as unavoidable, that is, just as
necessary, as the arbitrariness. Neither the arbitrary nor the accidental does
any damage to necessity; on the contrary, these foundational determinations may
rightly be said to furnish presuppositions, indispensable conditions for an exact
objectification of conditioned necessity. If a point $M$ be assumed whose
coordinates in the given system are $x$, $y$ and $z$, and whose distance from the
point of intersection may be designated by $r$, then an equation such as:
\[
  r^2 = x^2 + y^2 + z^2
\]
will involuntarily recall that manifold of spatial determinations which, in their
exact necessity, are all conditioned --- subjectively by arbitrary, objectively by
accidental presuppositions.

The three coexisting axes are, in their outward, immediate accidentality, without
relations to one another; only with knowing do the relations come, and with the
relations necessity. Here, if we take all knowing away, there are not even three
axes, for threeness demands that a knowing of the two unite with a knowing of the
third. The representation of an abstract objective threeness derives from a
co-knowing which the objectifying consciousness, without thinking of it, itself
slips underneath. It is incontestably certain that the equation cited, answering
to the axes,
% --- p. 7 ---
has objective validity --- provided, be it noted, that it is objectified in
knowing; but when there is no knowing, the equation is meaningless. \emph{All
objective relation is subjectively grounded in reflexion}; that is a law of
objectification, as strict as any law of nature.

In the Schellingian philosophy of nature there is mystically profound talk of an
intuiting dwelling within nature, an objective phantasy of nature, an unconscious
knowing of nature, and the like. How groundless such turns of phrase are shows
itself every time the gaze is fixed upon something determinate. Were the three
dimensions of space posited by an unconscious knowing, then in this knowing there
would have to be a co-knowing of all three, a knowing in which each single one of
them was determined with respect to the two others. The unconscious knowing would
then have to be \emph{a knowing of respects}; here is the contradiction. A knowing
of respects is a knowing that hovers over the single determinations, i.e., a
self-knowing. In order that $X$ can be determined with respect to $Y$, and $Y$
with respect to $X$, the determining knowing must necessarily hover over both $X$
and $Y$. Such a hovering is possible only when the same self that knows $X$ also
knows $Y$. But since knowing of another is, by the nature of the case,
inseparable from self-knowing, profundities such as ``an unconscious knowing of
nature,'' ``a self-intuiting space'' --- on a par with ``a self-thinking air,'' ``a
self-conceiving matter'' --- are to be referred, not to ``strict science,'' but to
the utopian speculations.

To determine space is to posit limits in space. Here, then, a qualitative
difference shows itself between the limit, the negative, and the limited space,
the positive. The negative is for thought, the positive for intuition. The first
absolutely negative is the point. In the determination of the point a peculiar
conflict arises between intuiting and thinking. Thought demands that the point be
without extension; intuition demands that it have an extension, however
% --- p. 8 ---
small. If thought is not satisfied, the point becomes a surface, no point; if
intuition is not satisfied, the point can occupy no place: a pure point of thought
is no point in space. The resolution of the contradiction is a reflexion in which
intuition and thought alternately cancel and confirm one another. Insofar as
thought cancels intuition, the element of space, the positive, is consumed by the
negative and posited as point. Insofar as intuition cancels thought, the negative
is converted into the positive, and the point posited as element. For thought the
point is the limit of the line, but for intuition it is the element of the line;
for thought the line is the limit of the surface, but for intuition the element of
the surface; and so the surface too is again both the limit of the body and its
element. The element of the line can be integrated into a line, that of the
surface into a surface, that of the body into a body; but the limit itself, the
absolutely negative, cannot be integrated.

As the negative, the non-being, the limit can therefore set no barrier to the
being of space; there is space on both sides of the limit; the limit is itself the
negative unity in which the limited are united and separated. Everywhere that
thought posits a limit in space, intuition goes out beyond the limit and discovers
a new space. Space is full of limits, for the limit is everywhere, and yet
limitless, for the limit is everywhere overstepped. To exchange a limit that is
overstepped for a new limit that is again overstepped --- that is the secret:
infinite space is objective only for an infinite objectification.

\runin{b) Time.} The infinite extension of space is grounded in the infinite
becoming of the limit. This becoming, which posits a difference between the
vanished, the present and the coming limit, points to another form of nature,
distinct from space: the succession hidden in extension
% --- p. 9 ---
belongs to time. Indifferent toward its limit, space is indifferent toward its
becoming; the space that comes when the limit is overstepped was already there
beforehand, and that which has been is there still; the difference is only a
shadow of time. In time itself, by contrast, the difference between what has been,
what is, and what will come is an essential, though as yet outward, difference.

That the form of time is originally a selfhood-form can be seen from the reflexion
into which intuition is resolved. Time has only one dimension, \emph{length}; but
this is so broken in the tripartition of past, present and future that only a
thinking that runs through them is able to hold the unity fast. The past exists in
recollection, but recollection is an act of selfhood; the future exists in
expectation, but expectation is an act of selfhood. The part of time that might be
said to have an immediate being, independent of a recollecting and expecting
subjectivity, would then be the present. But the present is, when we look more
closely, just as reflected as past and future. If the present is apprehended as the
absolute instant, then it is apprehended not as time but as limit of time, the
limit in which past and future touch one another. This limit, the pure now, is a
non-being, a mathematical zero point. If the present is apprehended as relative
instant, as element of time, then it is a duration which intuition slips
underneath, in ceaseless flight from pastness to futurity, i.e., from non-being to
non-being. A time-element can indeed be apprehended as constant, $dt = dt$, when it
is applied as a measure for an element of space; but the proposition ``now is now''
becomes a self-contradicting judgment, for the now of the subject is past when we
come to that of the predicate. If the present be held fast as finite, existing
time, then it is solely the recollection and foresight of subjectivity which, in
reflected unity, are able to let intuition stretch itself along an arbitrary
present length of time.

Hegel's dialectic of what he calls ``the dimensions of time''
% --- p. 10 ---
is, from the objective side, fine and sharp. It is becoming, whose abstract
moments of unity are, each for itself, posited as the whole, but under opposite
determinations. The two determinations, past and future, are thus each of them a
unity of being and nothing; but they are also different. This difference can only
be that between arising and passing away. ``The one time, in pastness (Hades),
being is the foundation from which the start is made; the past has actually been,
as world-history, as events of nature, but posited under the determination of
non-being, which supervenes. The second time it is the reverse; in time to come
non-being is the first determination, being the later, though not in point of
time. The middle is the indifferent unity of both,'' and so forth. That is
excellent. But in glaring contradiction with these objectified conceptual
determinations --- determinations manifestly owing to a self-conscious,
objectifying subjectivity --- stands a declaration such as this, that the
difference between objectivity and a consciousness subjective over against it
does not concern the concept of time at all. ``Time,'' it is said,
% The Danish re-opens the quotation with „ after the interpolation „hedder det“
% and closes only once, at the end of the citation. Printed so.
``is the same principle as the I $=$ I of pure self-consciousness: but the same
principle, or the simple concept, still in its complete outwardness and
abstraction, as intuited mere \emph{becoming} --- pure being-in-itself as a
sheer coming-outside-itself.'' --- Under such turns of phrase the clear
objectification of time, an objectification possible only for the knowing Spirit,
is pushed aside, and the fixed idea of a self-intuiting, self-thinking form of
nature, a self-conceiving time, befogs consciousness.

As the form of becoming, time has something becoming as its presupposition. The
reciprocal determination of time's becoming and of what becomes in time throws
light upon the relation between the objective and the subjective. From the
objective side time goes together into one with what goes on in time: the
becoming content lays out its essence in the form of becoming and imparts to the
form
% --- p. 11 ---
its predicates. Time becomes, time hastens, runs; the times are changed (\textit{tempora
mutantur}), there are good times, bad times, etc. From the subjective side, by
contrast, the form is held fast in abstract universality for itself, in
distinction from the content. Everything finite begins in time and ends in time,
but time itself has neither beginning nor end. If a first, absolute point of
beginning be posited, thought easily discovers the contradiction that there must
then have been a time when there was as yet no time. If a last, absolute end
point be posited, the consequence must be that at last there will come a time
when there is no time. If Hegel does wrong in holding fast to objective time
without regard to an objectifying self-knowing, then Kant also does wrong in
letting the subjective and the objective fall apart into two pieces. How such a
parcelling out must become a snare in which the finite understanding catches
itself is shown most clearly in the ``antinomy'' of the origin of the world.

% The Danish re-opens the quotation with „ after the interpolation „hedder det i
% Thesis“ and closes only once, at the end of the citation. Printed so in both
% copies (Google and KB).
``\emph{The world},'' says the thesis, ``\emph{has a beginning in time}''; for if
one assumed the opposite, then at every given point of time a whole eternity must
have elapsed, and an infinite series of successive states have been completed.
\emph{The world}, says the antithesis, \emph{has no beginning in time}; for in the
contrary case there would have to be an empty time prior to the beginning of the
world, but in an empty time, destitute of all conditions, nothing can begin.

According to the thesis, the infinity of a series consists in its never letting
itself be ``completed by successive synthesis.'' What does this Never mean? That
the completion is postponed into an infinite time. The proposition can therefore
be expressed positively thus: to conclude an infinite series by successive
synthesis, an infinite time is required; or: by continuing the successive
synthesis through an infinite time, representation forms a series of infinitely
many members, an infinite series.
% --- p. 12 ---
% This page carries TWO footnotes in the Danish, marked *) and **).
The formation of the series is the same whether we begin in finitude and approach
infinity, or from infinity approach finitude.\footnote{% printed mark: *)
% The rows of dots in the two notes on p. 12 are NOT alike: here three SPACED
% points, in note **) four CLOSE ones. Both reproduced as printed in the Danish.
In the series $\dfrac{1}{1-x} = 1 + x + x^2 + \ .\ .\ . x^{n-1} + x^n +
\dfrac{x^n + 1}{1 - x}$, $x^n$ is the general member; setting $n = o$,
% „n = o“ is printed with the letter o, not the numeral nought; so reproduced.
we get $x^0 = 1$; setting $n = \infty$, we get $x^\infty$; but the way from 1 to
$x^\infty$ and from $x^\infty$ to 1 is equally long: it is infinite.} Between a
point of time given in temporality and a beginning in eternity there lies, then,
an infinite time. There is accordingly no contradiction whatever in an infinite
series of states having elapsed ``at every given point of time.'' If the proof of
the thesis were to signify anything, then it would have to be proved either that
an infinite series cannot possibly be formed in an infinite time, or that between
a beginning in eternity and a temporal point of time there cannot possibly lie an
infinite time. But to prove such impossibilities is impossible.

As concerns the antithesis, the emphasis must fall upon the relativity of the
beginning. Every finite beginning presupposes a preceding one and furnishes the
presupposition for one that follows. This, on Kant's train of thought, is supposed
to prove the impossibility of a \emph{first} beginning. But no. Let the finite
beginnings correspond to the members in an infinite series; there is nothing
whatever to prevent the series from beginning with a first finite member, for this
member has, precisely, the same general presupposition as all the following ones:
the infinite series is a successive unfolding of the function that encloses the
ground of infinity.%
\renewcommand{\thefootnote}{**)}%
\footnote{% printed mark: **)
The infinite series $1 + x^2 + x^3 + .\,.\,.\,.$ has its ground and
% Misprint in the series: the member $x$ is missing between 1 and $x^2$; so printed.
presupposition in the function $F(x) = \dfrac{1}{1-x}$; the infinite series
$1 + \dfrac{x}{1} + \dfrac{x^2}{1.2} + \dfrac{x^3}{1.2.3} + .\,.\,.\,.$
has its ground and presupposition in the function $f(x) = e^x$; there are examples
enough in mathematics.}%
\renewcommand{\thefootnote}{*)}
Were the
% --- p. 13 ---
proof of the antithesis to have any true validity, then it would have to be proved
that a beginning in time cannot possibly consist with an infinite succession of
finite beginnings, be it past or future. But this, for good reasons, has not been
proved.

That the conflict between the finite and the infinite on which the antinomy is
built rests essentially upon a critical rupture between consciousness itself and
objective being is easy to see. To what, then, does finitude answer? To that
objectivity which sensation can grasp and representation hold fast. In this
respect it is quite right that an infinite series of time can neither be sensed
nor represented. How do we then come to infinity? By the way of thought. We go, as
it is called in the language of mathematics, to the limit; that is to say: thought
dissolves finitude into its uttermost limit, posits the negation of the finite. It
happens by a leap. If to the infinity thus posited there is now also to answer an
objectivity of its own, then it is surely clear that objectivity in the infinite
must come into conflict with objectivity in the finite. But instead of laying the
blame on the intuition of time, we should on the contrary recognize that it is
precisely time which, in the natural sense, strives to reconcile the finite with
the infinite. Time is succession, and succession an outward transition which
satisfies representation and thought in equal degree. By succession the finite
follows upon the finite; each finite member comes to be successively for
representation. Representation can then demand no more than that the succession
shall further continue. But the continuance of the succession is precisely the
coming-to-be of infinity; the infinity that comes to be must then likewise satisfy
thought. Objectively regarded, succession is a becoming; but when becoming in time
is the outward manner in which the finite passes incessantly over into the
infinite,
% --- p. 14 ---
and conversely, then both, the finite and the infinite in time as time, have equal
objectivity.

To sharpen the sting of the antinomy Kant has thrown time and eternity together
into one and yet set them up in the sharpest opposition. Since the world, for the
sake of the infinite series of time, cannot be from eternity, it must have a
beginning in time; that is the thesis. Since otherwise an infinite, empty time (an
empty eternity) would have to lie prior to a beginning in time, the world must be
from eternity; that is the antithesis. Here there is an alternation between time
and eternity, and so time is after all eternity, and eternity time --- to infinity.
The fact is that at the ground of all becoming there must lie a being that does not
become. This being reflected into itself, the essence, is namely reflected in its
other, in the phenomenon. In virtue of this double-sided reflexion the essence has
its inner infinity in the form of selfhood, which is opposed to the outer infinity
of the otherness-forms. The outer infinity answers to time, the inner to eternity.
There is a qualitative difference between time and eternity.

Upon this difference rests in turn the difference between \emph{beginning} and
\emph{origin}. The beginning happens in time, but the origin issues from eternity.
The absolute essence is eternal, unoriginated, exalted above time; the relative
essence has, by contrast, its origin from the absolute and is by its relativity
subject to time. Origin points to the self-essence's objectification of its
essential other, to Spirit's objectification of nature, to God's objectification of
the world; origin is transition from eternity to time. But beginning is a category
of time; the beginning presupposes the origin, for it belongs to time, the infinite
as well as the finite: infinite time is a succession of nothing but finite
beginnings. Spirit is eternal in time, above time; but the temporal world is not
eternal, even though its becoming is from eternal time.

% --- p. 15 ---
% The Danish head prints a loose space inside the word: „Mellemvære nde.“ Both
% copies show it; set here as one word.
\runin{c) The Business between Space and Time.} If the concepts of nature, space
and time, were able in an unconscious knowing to conceive themselves, each for
itself, and one another mutually; if by a common conceiving effort they really
availed to bring forth something so tangible as what we call matter, then modern
nature-pantheism could undeniably boast of a sure footing in ``strict science'';
but the fixed idea becomes more recognizable with every step the development of
the concept takes.

The Hegelian deduction of place, motion and matter is instructive in this respect.
The words are obscure, but the thought is clear. Space, with its indifferent
discreteness and its equally indifferent continuity, is a contradiction that
resolves itself in time. But time for itself is also a contradiction; for the
present is an unarrestable conversion of the past into the future: the separated
moments dissolve into immediate indifference. The concrete point in which the
contradictions, converging from both sides, fall together is ``\emph{place}.''
Place is unity of ``\emph{here}'' and ``\emph{now},'' and thereby expressly posited
as unity of space and time. The place determined by here and now points, however,
out beyond itself to another place determined by here and now; this becoming of
place is ``\emph{motion}.'' But here and now are everywhere; the determinate place
is, in all its indeterminacy, only the universal place: the contradictory moments
of motion fall together in immediate unity; this immediately existing unity is
``\emph{matter}.'' ``Matter is \emph{impenetrable} and offers \emph{resistance},
is something tangible, visible, etc. These predicates are nothing else than that
matter is partly for determinate observation, in general \emph{for another}, but
partly also \emph{for itself}. Both are determinations which it has just as the
identity of space and time, of the immediate \emph{outside-one-another} and
\emph{negativity}, or of singularity that is for itself.''

% --- p. 16 ---
This deduction and its dialectic are characteristic of
% The Danish re-opens the quotation with „ after the interpolation „hedder det“
% and closes only once, at the end of the citation (after „Existens.“). Printed
% so in both copies (Google and KB).
the philosophy of immanence. ``Space,'' it is said, ``is not adequate to its
concept; it is therefore the concept of space itself that in matter procures
existence for itself.'' The doctrine of the objective concept's miracle-working
power cannot be uttered more plainly. But the sorcery vanishes, and the miracle is
called back again, by an explanation such as the following: ``One has,'' it is
added, ``often begun with matter, and then regarded space and time as its forms.
What is right in this is that matter is the real element in space and time. But
these must, on account of their abstraction, appear here as the first, and then it
must show itself that matter is their truth.'' The last of these statements we
subscribe to unconditionally; but it stands in screaming contradiction with the
first. If it is really the concept of space which, in union with that of time, has
availed to give matter existence, then matter, the real, cannot possibly be the
first; it must, in the subjective as in the objective sense, be the second, the
derived. If, on the other hand, matter is in itself to be the first, because it is
the concrete which takes space and time up into itself as abstract forms, then it
is surely an insane assertion that the abstract should have brought forth that from
which it is itself abstracted, that space, spurred on by its own restless concept
abstracted from matter, should have created matter.

Abstraction, contradiction, negativity, infinite negativity, etc., are categories
of selfhood which assuredly have validity in infinite objectification; but they are
misunderstood, misused and taken in vain when they are laid without more ado into
the objective, so that outward being is ascribed to them. Space and time are forms
whose abstract being must be integrated by a more complete existence: they must
have a \emph{business between them}, which is taken up into them and in which they
are taken up. But abstraction and integration are acts of thought, functions of
objectifying selfhood. What reality the
% --- p. 17 ---
required business between them is to have will then depend upon the
objectification. For abstract objectification, insofar as it is still only a
question of a general determining of place and mechanical motion, a represented
business between them is sufficient.

No motion is possible without a something that can move. Motion goes on in space,
but the moving thing is distinct from space, for space does not move. Change of
place is a change in space, but without space being changed. The change happens
with the moved thing, which passes over from place to place; and yet it is not the
moved thing itself, but only the place, that is changed. The place does not belong
to the moved thing, for it belongs to space; and yet it does not belong to space
for itself, for it first becomes place, in the motion, by being occupied by the
moved thing. The place is neither straightforwardly existent nor non-existent; it
is an objective determination of relation, objectively reflected. Motion goes on in
time, and what is moved in time is distinct from time; but time enters into motion
in another manner than space. The places of space are simultaneously existent,
namely according to possibility; every place traversed can for that reason be
occupied again. But the instants of time vanish, and every vanished instant is
irrevocably lost. With the movable as the business between them, space and time
each receive their reflected share in the motion. Determinate motion is velocity,
and velocity a quotient in which the element of space as numerator and the element
of time as denominator reflect one another objectively, $\dfrac{ds}{dt} = v$; but
on the terms of objectification: the objectifying reflexion is subjective.

What confusion of concepts can arise when, in the phenomenon of motion, one will
absolutely think the objective without considering its relation to the objectifying
subjectivity
% --- p. 18 ---
is well known. The arrow of the Eleatics, flying in infinite space, is now here,
now here, always now here! and does not get off the spot. It is an objective extreme
that cancels itself. For motion to come to be, the observer must recollect, gather
together and intuit --- namely the places the moved thing has traversed and goes on
traversing. No wonder that the objective extreme in the Eleatics struck over into a
subjective one, so that those inflexible thinkers declared all motion outright to be
semblance and deception. But that the opposite too can happen, that the subjective,
namely, can by a logical self-deception strike over into the objective, is evident
from the Hegelian dialectic. While Zeno, in order to get the arrow to stand still,
points to the infinite space in which all relativity has vanished, Hegel holds the
determination of place fast in thought. ``Here are three places: the one that is,
the next following, and the one abandoned. But it is at the same time only one
place, the universal place of those places, an unchanged amid all changes.'' One
learns from this how abstraction, contradiction and negation drive one another
forward. That the three places are one place rests upon an abstraction, in that
abstraction is made from that difference in outward being which is precisely
characteristic of space. It is at the same time a contradiction; but this
contradiction, conditioned by the abstraction, is resolved by the unity negated in
the threeness being negated again in the universal unity, the abiding, so that
``duration (\textit{die Dauer}) is according to its concept motion.'' That nature does
not institute such reflexions on its own account, but leaves it to subjectivity,
must surely be conceded. None the less the subjectively reflected concept of motion
is straightforwardly declared to be a not subjective but merely objective concept of
nature, and is moreover made one with motion itself.

As abstract forms, hovering between being and non-being, space and time are without
particular independence; their existence depends upon the business between them in
which they are united,
% --- p. 19 ---
and so upon the content with which they are filled: in this dependence their
ideality is seen. If space and time are filled with dream-images, they have only a
dreamt existence; if they are filled with a world of poetic figures, they attain an
aesthetic existence; only with a content of actual things do they become actual. If
it holds of the forms, as we have seen, that their objective being rests upon
objectification, then the same must necessarily hold of the content with which they
are filled.
% The A division runs on into p. 19; the B head stands part-way down that page.

% A short centred rule is printed here, between division A and division B.
\streg

% Danish: \afdeling{B.}{Materie og Kraft.}   (printed pp. 19--38)
\afdeling{B.}{Matter and Force.}

The actuality of space coincides with its actual filling. What fills space,
matter, and what is at work in space, force, belong to nature but are posited by
Spirit; the task is to show how the two foundational determinations --- origin
from Spirit and natural reality --- presuppose one another.

\runin{a) Matter.} Nature's first immediate real is not a concept, but a sensible
existent in which a concept exists. The concept of existence is a concept of
reality; it is then first and chiefly a matter of seeing into the relation between
concept and reality.

To presuppose reality and exclude the concept is characteristic of materialism.
The one-sidedness of this way of seeing does not consist in its ascribing reality
to matter --- for if matter were only a semblance, the ground of nature would be
hollowed out; the one-sidedness is rather that matter is made into something
% --- p. 20 ---
absolute in itself, into an eternal, unconditioned existent out of whose
fundamental parts everything is composed and formed. These fundamental parts, the
atoms, are apprehended at once both as parts of mass and as absolute real points.
That is a contradiction; for if the atoms are parts of mass, then they are
relative; if, on the other hand, they are absolute, then all actual relation
between them is impossible. To ascribe properties to the atoms is to make them
relative, for the property of the one atom acquires significance, in attraction and
repulsion, only through the property of the other. To presuppose the absolute unity
of the atoms with their properties is to concede their absolute relativity and to
reduce them to moments for one another mutually. With this consequence materialism
has dissolved itself.

To posit the concept as reality and reality as concept is characteristic of the
one-sidedness of idealism. Idealism is powerless; it mistakes the opposition of the
ideal and the real, of Spirit and matter, and does not see that to a mighty inner
there must answer a mighty outer. Idealism would transform the concept of matter
into actual matter; that is a barren speculation. The opposition of concept and
matter is an opposition in power, an opposition so hard and strong that the concept
of their unity coincides with the concept of omnipotence, of absolutely overreaching
self-power. Self-power's inner original opposition is object-power, power in
another; the first outer for power in another is matter. Hence the swarm of atoms:
nothing but points of otherness, points of reality. The qualitative opposition of
Spirit and matter is marked by the issuing of power. The issuing is not blind; it is
a fundamental act in virtue of an omnipotent conceiving: the concept is the ground
of reality, reality the existence of the concept.

An immediate object-power without a corresponding concept is just as impossible as
an existing concept of power without an object. One cannot divide the existent into
two absolutely different kinds: an existent \emph{with} a ground and an existent
% --- p. 21 ---
\emph{without} a ground; for the essence of existence is to have a ground. The real
ground is an existence dissolved into the concept of power; the existent that has no
ground does not exist, its being is a hollow phenomenon, a semblance.

Materialism's fundamental error now shows itself in this, that it makes the matter
``existing from eternity'' into an existent without a ground. The atoms, the absolute
existences upon whose separation and combination all finite, material formations
rest, then stand to actual things in their finitude as the groundless to the
grounded! A more unnatural separation of ground and existence cannot be given.
Materialism has a sharp eye for reality as object-power. It begins with a language of
power, points to the fact, takes the existent as something inviolably given; that is
sound and right. But to overlook the concept of power and to deny the essence of
existence, in order to build upon an existent without a ground, is unsound and false.

Objective idealism holds to the grounding, in that it holds to the concept; but its
fundamental fault is that it overlooks object-power and effaces the difference between
reflexion in knowing and reflexion in power. What this misstep signifies can be seen,
in matter, in atoms and mass. Existing mass has the atoms as its foundation; but the
existence of the atoms has in turn the presupposed mass as its foundation. To form a
foundation, an existent must go to the ground. The existent is represented; but the
ground and the foundation are thought. That the atoms exist means that they are
posited in representation, while the mass understood along with them is posited in
thought; if, on the other hand, the mass is in representation, the atoms being
understood along with it, then it is the mass that exists. The perishing of the atoms
in the mass and the dissolution of the mass into atoms are imaginary movements that go
on within thought and representation; it is a reflexion which, however many conceptual
determinations it takes up, never gets beyond a
% --- p. 22 ---
represented existence with a thought ground, never reaches actuality.

In the reflexion of power the start is expressly made from actuality, insofar as the
start is made in the existent with object-power; the ground is a concept of power, and
the foundational determinations are real moments of dissolved existence. The
transition from existence to ground, and from ground to existence, is an essential
change, an objective act of power, an event of nature. The mass exists, and the atoms
exist; and yet here there are not at one and the same time two sorts of existent or
two layers of existence. If existence falls in the mass, then the atoms have ceased to
exist; they are taken up, as moments of the foundation, into the existence of the
mass. But when the atoms perished, something actually took place: they approached one
another, gave up their independence in cohesion, and became moments in the mass. If
existence falls in the atoms, then the mass is dissolved: it has perished. But when
the mass perished, something actually took place: it dissolved itself and in the
dissolution lost its independence. That the mass lost its independence is the ground
of the atoms' acquiring independence and coming to exist. The existence of the atoms
has its real ground in the perishing of the mass just as surely as the existence of
the mass has its real ground in the perishing of the atoms. In this reflexion of power
the one is not an unconditioned first and the other a derivative; rather both
existences, that of the atoms and that of the mass, are equally original and equally
derived. Nor, then, is perishing annihilation; the existent that perished in another
existent comes to exist again when the existence resting upon its perishing perishes.

The reflexion of the concept of power and of object-power which furnishes the general
foundation of mechanical nature is everywhere recognizable by an abrupt conversion of
existing members into moments, and of moments into existing members. The concept of
power's
% --- p. 23 ---
reflected unity with object-power is expressed in the harshest opposition. The
existent is no concept; it is object, actual fact. The concept does not exist; it is
reflected in dissolved existence as ground and foundation. But the unity is in turn
just as mighty as the opposition: the existent has the ground in itself: the object is
the material expression of its realized concept. From the side of reflexion the
concept is what overreaches; from the side of the object it is fact, and fact first and
last. The infinite reciprocal determination of reflexion and fact, ground and
existence, is the form of actuality for \emph{power in another}. But power in another
presupposes self-power; the concepts dwelling within the facts must be conceived by an
overreaching sole power: \emph{the mechanical power of nature is borne by the
omnipotence of Spirit.}

In what sense existing matter must be said to be the first mechanical presentation of
power in another can be seen from the properties which physics has particularly
ascribed to it as matter. Matter is \emph{sensible}, not merely for us, but in itself.
To think matter is to transform it into a being of thought and to deny its existence.
If matter is for Spirit only a being of thought, then the whole of nature is a
semblance. But the truth is that matter's essence is an essence of power; its concept
is a concept of knowing converted into power, its existence an object-power, a
resistance for thought, an immediate object, for Spirit too. To be Spirit's immediate
object, the real other of intuition, a limit for thought, is to be sensible. Matter
\emph{fills space}. Limited matter occupies its place and offers resistance to other
matter that would occupy the same place; this power of resistance is otherness-power,
power to exclude another from another. Matter is \emph{divisible}. Divisibility rests
upon otherness. The mass is one thing, the atoms another. The unity of the mass is no
self-unity, but a
% --- p. 24 ---
unity of another with another; the mass subsists by its other, for it consists of
atoms. Were the atoms absolutely to subsist by themselves, it could still not possibly
be as mathematical points; it would have to be as small masses. But what holds of the
larger masses holds of the smaller, the smallest, the very smallest: the mass subsists
in another by another. Each single atom has the same subsistence as all the others. The
reality of one atom is the reality of all; but the reality of all is the mass: the atoms
subsist in the mass, and so in another. Matter is \emph{indifferent to rest and motion};
this indifference is quite a characteristic sign of otherness. That matter is subject to
the law of inertia is also expressed by saying that it is without will, which in other
words is to say that it is without self-determination. When it is at rest it cannot
determine itself to motion, but must be determined by another; when it is in motion it
cannot determine itself to rest, but must be determined by another. This doctrine of
matter's selfless independence (inertia) has decisive significance for the knowledge
both of matter's essence and of the laws of motion.

None the less the doctrine of immanence, as speculative philosophy of nature, will
still, in the teeth of empirical physics and rational mechanics, maintain the mystical
representation of a selfhood dwelling within matter, a silent longing for liberation, a
hidden capacity to set itself in motion. Matter, says Hegel, is removal in space. It
offers resistance, repels itself from itself: that is \emph{repulsion}, whereby it
posits its reality and fills space. But the single members repelled from one another are
all one. That it is the same one which, by repelling itself from itself, cancels the
removal of the single members --- in that consists \emph{attraction}. Both together
constitute, as gravity, the concept of matter. The unity of gravity is only ``an ought,
a longing, the
% The quotation opens here and closes only on the next page („naaes ikke.“); a
% quotation running across a page break, not a misprint.
unhappiest striving to which matter is eternally
% --- p. 25 ---
condemned; for the unity does not come to itself, it is not attained.'' Matter can,
however, console itself that ``it is not so stupid'' --- as the physicists --- ``as to
hold one and many apart from one another.'' It is only in the finite sphere, in the
circle of ``the selfless bodies,'' that matter suffers from inertia; in its ``free
concept,'' in the ``celestial sphere,'' matter has selfhood and moves itself.

In thus stepping forth polemically against physics, the doctrine of immanence at the
same time comes into a dubious position with respect to the doctrine of Spirit. To make
good the lack of Spirit's conceiving self-power, and to let nature conceive
% Misprint: the Danish has „sin egne Begreber“ in both copies. The book's own
% errata leaf corrects: „Side 25 Linie 11 f. o. sin --- læs sine“. Reproduced as
% printed; translated by the sense.
its own concepts, personifications are brought in: as a substitute for the excluded
Spirit, felicities of wit are introduced. ``Just as time is the simple formal soul of
nature, and space, according to Newton, God's sensorium, so motion is the concept of the
true world-soul. We are accustomed to regard it as predicate, as state, but it is in
actuality the self, the subject as subject, the abiding in all vanishing.'' Just as
motion is the essential self, so matter is, in the mechanics of the infinite, naturally
the actual self. The fact that matter moves matter to infinity is explained, by a
logical paraphrase, to the effect that matter accordingly \emph{moves itself}. By such
explanations the concept of nature's development is mystified from the ground up.

\runin{b) Force.} Matter moves and is moved; the one mass moves the other: matter has
the ground and cause of motion in itself, but under the form of otherness. Only when it
is clearly seen that being able to move and to let itself be moved necessarily belongs
to its space-filling reality is the concept of matter exhausted. It is not the abstract
determination of the form of motion, but actual motion, that is treated of here. We
must look about us for a
% --- p. 26 ---
fundamental phenomenon from which it can show itself most clearly what the moving thing
properly is.

From a finite mechanical point of view one has pointed to \emph{impact}. In the contact
of two colliding bodies the mutual resistance occasions a struggle over the place, which
seems to furnish sufficient ground for change of place, and so for motion. The ground is
so much the more natural in that it manifestly lies in the reality with which every mass
occupies its place. Immediate contact
% Greek, set in italic Greek type in the Danish: ἅπτεσθαι.
(ἅπτεσθαι) is accordingly posited by Aristotle too as the condition on which the one
matter can act movingly upon the other. But that impact cannot be the moving thing
itself is evident. For impact presupposes that the two colliding bodies must either both
have moved in opposite directions, or, if the direction is the same, that either the one
must have moved with greater velocity than the other, or the one have moved and the
other been at rest. It is seen then from all sides that impact cannot be what properly
moves; it is itself, after all, a result of motion.

To avoid the presupposed motion, one might try to place the moving thing in a pressure
and make \emph{pressure} a fundamental phenomenon. Pressure and counter-pressure now
express straightforwardly the power of resistance identical with matter's reality;
motion enters as a necessary consequence, whose presupposition is that the equilibrium
between pressure and counter-pressure is cancelled. Motion begins with a striving toward
motion; such a striving really lies in pressure, and so far the explanation has a
ground; but it is an insufficient ground: a something is here presupposed which is not
yet grounded. The phenomenon of pressure can be construed partly so that both bodies,
each from its own side, strive to press into one another while yet mutually resisting
one another; partly so that only the one is pressing forward, the other merely
resisting. In both cases the striving coincides with the pressing
% --- p. 27 ---
forward; and the question becomes whence this pressing forward then derives. Pressing
forward presupposes approach: \emph{the mutual approach of two separated masses} must
accordingly be the fundamental phenomenon sought.

Two masses separated in space, \textit{A} and \textit{B}, will, when extraneous
hindrances are thought away, approach one another, and so come into motion; but what is
the ground? Hegel appeals to gravity and places gravity in a common seeking after a
common centre lying outside the masses. ``The centre,'' he says, ``is not to be taken as
% The interpolation „siger han“ breaks the citation, but the Danish does NOT
% re-open with „ here; the quotation marks therefore balance on this page.
material; for the material is precisely this, to have its centre \emph{outside itself}.
Not the centre, but a striving after it, is immanent in matter. Gravity is, so to speak,
the confession matter makes of the nothingness of its outward being for itself, of its
dependence, its contradiction.'' But to ascribe to the masses such an inner lack, a
yearning toward an ideal existing outside them, is poetic enough. but
% Misprint: a full stop where the sense requires a comma --- „poetisk. men ikke
% videnskabeligt“. So printed in both copies (Google and KB); not corrected on
% the book's own errata leaf.
not scientific; just as it also expressly conflicts with the hard otherness, the selfless
independence, to explain the mutual approach of the masses by an obscure sympathy and to
invent a --- naturally unconscious --- longing in \textit{A} for \textit{B}, answered by
a --- naturally unconscious --- longing in \textit{B} for \textit{A}. Strict mechanics
has another construal. \textit{A}, it is said, attracts \textit{B} because \textit{B}
attracts \textit{A}: what attracts in both masses, indeed in every atom of both masses,
is \emph{force}.

Force is, mechanically understood, the cause of motion; but --- physics adds --- an
unknown cause. We shall now examine more closely how matters really stand with ``the
unknown cause we call force.'' When it is said that we do not know what force is, it is
tacitly understood that we do after all know well enough what matter is. But if matter
really is a known quantity, then it is impossible that force, in its peculiar connexion
% --- p. 28 ---
with matter, should be wholly unknown to us. Force, the non-sensible, answers to matter,
the sensible, as the reflexion of power answers to immediate object-power: matter is the
outer, force the inner. But an inner can never show itself straightforwardly; that it
shows itself means that it utters itself; it reflects itself in its outer and must be
known by its utterances. A proposition such as this --- we do not know what force is in
itself, we know only its utterances, its effects --- accordingly states nothing else than
what lies in the nature of reflexion, namely that the inner must necessarily be known by
its outer. The abstract logical essence reflects itself in a semblance of semblance, but
the reflecting of power is an efficacy, its phenomena are facts: as the categorical
expression of the reflexion of power, force is not a ``hypostatized thought'' in human
imagination, but a fact reflected in its effects. Force as inner is nothing absolute in
itself; the inner is a possibility which first becomes actuality by uttering itself and
working. The proposition cited above cannot then, in the objective sense, be regarded as
any sure designation of the absolute limit of human knowledge, but rather as a subjective
indication of unclarity about the essence of thinking, of reflexion, and especially of
the reflexion of power.

How unclear and hovering the usual representation of force is comes to light from the way
in which its relation to matter is more closely determined. On the one side one speaks as
though matter were one thing and force something quite other. It is then not the mass
\textit{A} that by attraction sets \textit{B} in motion; no, it is the force, the force
of attraction present in \textit{A}, whereby the motion of \textit{B} is caused, and
conversely: it is not \textit{B}, but the force in \textit{B}, that moves \textit{A}. The
masses are moved, but the forces move: matter is \emph{passive}, force \emph{active}. On
the other side one speaks again as though force were only a property of matter.
% --- p. 29 ---
Then it really is the mass \textit{A} that in its capacity as mass attracts
\textit{B}, and conversely. The latter way of seeing stands in manifest conflict
with the former. By making force a property of matter and, as regards universal
attraction, a fundamental property at that, one lets the concept of force merge
without more ado into that of matter. On the first way of seeing there would here
be two essences, but on the second there is only one essence. As a property of
matter, force comes into a
% „høist“: both OCR witnesses read „heist“, but the ø stands plainly on the KB
% copy. The reading „heist“ rejected.
most questionable position with respect to matter's other properties, notably
inertia. By inertia matter can only be moved, not move; by force it can only move,
not be moved. It must then seem as though matter were a self-contradicting reality.
As what fills space, matter has besides a number of properties in common with
space, not to speak of a quantity of immediate properties of sensibility which do
not well admit of being reconciled with force as force. For one will surely be as
little able to speak of a broad force, a round force, a four-cornered force as of a
thin or thick or grey force.

The difficulties cited fall away when we hold to the reciprocal determination of
the reflexion of power and immediate object-power set out above. Matter is in its
own way reflected, namely in the mutual grounding of atoms and masses; force is in
its own way
% „iøinefaldende“: both OCR witnesses read „ieinefaldende“; the ø stands plainly
% on the KB copy.
immediate, namely in its utterances. And yet the difference is conspicuous; for
with matter immediacy, but with force reflexion, is what is decisive. The reflexion
of power is something more than a logical mirroring of the one in the other; it is
an actual determining, an efficacy in which it shows itself what the one can do
over the other. It is precisely such a capability that mechanically utters itself in
motion. The motion of \textit{B} reflects the force in \textit{A}, in that the
motion of \textit{A} reflects the force in \textit{B}. The mutual motion is a
phenomenon of the essence of the reflexion of power: the reflexion that overreaches
in essence
% --- p. 30 ---
proves that the universal in the relation of the members governs the separate
members. Physics is so far within its rights when it apprehends force, the active,
as an essentiality for itself, in distinction from inert matter, the passive. But
in this physics does not after all overlook that the two essences are at bottom
side-determinations of the same essence. The unity of the reflexion of power with
object-power is expressly brought out by force being apprehended as matter's own
reflected determinateness, and so as its property. As matter's pervasive property,
force is something other than an immediate quality or constitution; it is a
categorical expression, not for matter's sensible reality, but for its reflected
essence, its indwelling energy, its inner. Herewith the contradiction between force
as a reflected property and matter's other, immediate properties disappears. The
immediate properties belong to \emph{the extensive}, but the determinations
attaching to force belong to \emph{the intensive}, for they belong to reflexion ---
be it noted: reflexion in power. ``One can,'' it is said of the centre of gravity in
mechanical natural science, ``represent things to oneself as though the whole
weight of the body were united in the centre of gravity, and all the remaining parts
had no weight.'' This representation is more than a representation; it is an actual
truth. The weight distributed over the parts, weight in the extensive, is
transformed in the reflexion of power into a unity of weight, a weight in the
intensive: the centre of gravity \emph{is} the intensive point of
unity.\footnote{% printed mark: *)
Cf.\ \textit{Mathematik og Dialektik}, pp.\ 129--131.}

The validity of the above remains to be tested more closely, so far as a final
construal of the otherwise so enigmatic phenomenon of attraction is concerned. The
masses \textit{A} and \textit{B}, which at a distance each fill their own space,
seem to have nothing whatever to do with one another; and yet they have to do with
one another, for they approach one another. It is a
% --- p. 31 ---
strange opposition of essence and existence, reflexion of power and object-power,
that lets itself be seen in the phenomenon; existence says one thing, essence
another. The solution of the riddle lies in the insight that material actuality has,
through essence and existence, uttered two opposed but equally irrefusable demands,
and that the phenomenon of attraction is precisely the fulfilment of both demands.
The first demand is that the material members must have immediate independence;
this independence is made intuitable in the sensible existence of the masses.
Existence answers to place. Each mass occupies its place, and it may easily appear
to the beholder as though space-filling reality were the complete expression of
matter's actuality. That, however, is a deception of sense; what is thus seen in
immediate existence is only half an actuality. Existence conceals something in the
existing members that can come into view only in the form of reflexion, as essence.
The second demand is that the essence must reflect itself in the relation of the
members, in the mutual relation of the masses. Terms of a relation without a
relation is just as abortive a thought as a relation without terms. If the members
are actual, the relations must be so too; their actuality is in their activity, and
this activity is the reality of the relation. When power is reflected in a relation
of actual members, the essence shows itself as force. The distance-relation of the
masses becomes, through force, an active and so an actual relation. In the masses
matter's existence is limited, restricted to the place; the masses are as actual
members moments in actuality. As essence, by contrast, force goes from mass to mass
out beyond the limit, and so out beyond the place, and the moment of actuality is
seen in the mutual change of place. That force goes through space, that masses act
at a distance, is a mechanically characteristic expression for that reciprocal
determination of the reflexion of power and object-power which makes actuality
complete.
% --- p. 32 ---
The apparent contradiction, that a mass has force to move another mass but no force
to move itself, has with this at once found its solution. Matter's actuality is an
independence in the form of otherness, a selfless independence. Selflessness is
dependence; material actuality is a dependent independence. The dependence is
matter's inertia, the independence its force. Inertia is that matter is moved by
another; force is that it moves another. In masses that attract one another,
inertia and force are united on both sides: force is the complement of inertia,
inertia is receptivity for force. Thus independence and dependence too are united on
both sides. The mass \textit{A} has independence, it sets \textit{B}, from its own
place, in dependence; but \textit{B} also has independence and sets \textit{A}, from
its own place, in dependence. The relativity is that both masses are, in their
independence, dependent. Here no outward observer is needed to hold the two places
together in representation and determine their distance; for so pervasive is the
relativity that it is in actuality the masses that mutually determine one another's
places: matter's relativity is, in the motion, actual by means of force.

\runin{c) The reflected unity of matter and force: Spirit and matter.} The
development carried thus far has the ideality of power, reflected power, as its
foundation; its significance stands and falls with the answer to this question:
\emph{how is the reflexion of power possible?} That gravitating masses come forward
as holders of power, exercising in proportion to their magnitude a certain dominion
over one another, is only a figurative turn of speech so long as the possibility of
the fact is not seen into. It is, as already shown, the reciprocal determination of
members and moments that matters. Outwardly, in the form of parcelling out and
dividedness, the
% --- p. 33 ---
moments are actual masses; their reflexion is relation, and the relation reciprocal
action. But in the concept, inwardly, the actual masses are moments of the same
unity, reflexes of one and the same conceiving. To conceive is to master. It is then
clear that the omnipotent conceiving which unites all separate concepts of power in
the unity of the universal concept must go on in an all-conceiving self-power. Only
an infinite self-power, a knowing self that objectifies in power, avails to give the
system of objective concepts actuality in the existence of the masses and to raise
itself above the patchwork. The reflexion of power in another has its ground in
reflecting self-power; herewith matter's absolute dependence upon Spirit is known in
the concept.

The history of science teaches that all attempts in the philosophy of nature to
clear up the relation between matter and motion, matter and force, have hitherto
foundered upon a total misapprehension of the relation between Spirit and matter.
Matter's unity with force is, namely, reflected in such a way that, besides being a
logical unity in essence, it is at the same time an actual unity in fact. The
reflexion in which the double-sided unity is grounded is itself double-sided; it is
a reflexion of the essence of matter and of force; but it is at the same time a
reflexion of existing matter and of force's utterance. If the unity be posited in
essence, as when one says: matter and force are at bottom the same, then it is
broken in existence; for existing matter is not force, and acting force is not
matter. If the unity be posited in existence, as when one says: force is material
(the atomistic conception) or: matter is a product of forces (the dynamical
conception), then the difference in essence comes into view in the reflexion; for
matter is an outer, force an inner. If the unity is to be known in a double
reflexion of essence in essence, of fact in fact, then it is possible only by
self-power and power in another,
% --- p. 34 ---
essence-power in ideality and object-power in reality, furnishing extremes so
strong, so hard, that Spirit's omnipotence can show itself as absolutely
overreaching sole power.

Without selfhood no reflecting, and conversely; without a reflecting self no
reflexion; without reflexion no relation; without reflexion in power no active
relation, no cause and effect, no reciprocal action. In the form of conceiving the
reflexion of power is subjective; in the form of the concept, of the conceived, it
is objective; in a reflexion of subjective and objective reflexion the
objectification is completed.

\subrunin{α) Subjective reflexion of power.} Self-power is not power's empty,
% The book's first Greek sub-head. The letter is read on the page itself (KB
% copy): an unambiguous α, not the „a“ the Indhold's OCR gives. The head is
% run-in, set on the same line as the paragraph it opens, and letterspaced;
% a full stop after „Magtreflexion“, not a comma.
merely formal likeness to itself without inner difference and opposition, as when
infinite power is assumed without more ado to be existent and is determined by a
tautological proposition: power $=$ power, $A = A$. No, self-power must expressly
posit its own likeness and thereby make it self-likeness. In actuality this can
happen only by the unity of power showing itself different from itself, doubling
itself through another and confirming itself by opposition to another. Only that
confirmation in which the same essence is in the result its own presupposition and
in the presupposition its own result can in truth be called \emph{self}-confirmation,
a confirmation of the \emph{self}.
% Only the first element „Selv“ is letterspaced in „Selvbekræftelse“;
% „bekræftelse“ is set in ordinary type. The word is one word, not two.

The greater the independence the other has, the mightier is the self thereby
confirming itself, for the self reflects itself subjectively by mastering its other:
the stronger the resistance, the greater the power by which it is cancelled. From
this it is seen that self-power cannot be apprehended as an immediate existent, an
absolute substance indifferent to everything else. Self-power's being is efficacy,
and this efficacy an infinite reflecting of itself in another, of another in itself.
The inner motion in which it
% MISPRINT: „bvori“ printed for „hvori“ — a wrong sort, b for h. The bowl is
% closed and answers exactly to the b in „bestaaer“ in the same line; both
% copies have the same shape, so the fault is the compositor's, and the errata
% leaf does not correct it. Reproduced as printed; translated by the sense.
eternally subsists is a circular motion. Absolute self-power does not leave itself
because it goes out from itself; it does not reach
% --- p. 35 ---
its other as an outwardly given, for it has its other in itself; it is not changed,
for it masters its other; the motion is no unrest, for the self is in its motion at
home with itself, its motion is rest; it rests in itself upon itself: \emph{absolute
self-power is Spirit.}

For Spirit as Spirit the whole material world is a nothing and yet an actual
something. Both determinations are reconciled in the fundamental determination that
the material world is Spirit's other. As other for Spirit, the material is a moment
for the omnipotent self-reflecting, but nothing, nothing at all in itself: only
Spirit, the one, is; all else is nothing. As other for Spirit, the material is
actual; it is, after all, object-power, the inner power of resistance belonging to
self-confirmation, the measure of Spirit's omnipotence. Were matter a semblance and
the world a dream, then Spirit itself were a dream, then there were no reality,
either inward or outward, no actuality, either ideal or real.

\subrunin{β) Objective reflexion of power.} The immediate and yet
% The letter read on the page (KB copy): β, not the „B“, „ß“ or „P“ the
% Indhold's OCR gives. Run-in head, as α) above.
reflected unity of matter and force is a unity of power. But the objective unity of
power is at the same time a broken power, power in another, power of multiplicity;
upon this rests its peculiar reflexion. The reflexion is recognizable in gravity.
Physics apprehends gravity from three points of view: as a \emph{property} of
matter, as a particular \emph{force}, and as a \emph{power} spanning the universe.
As property, gravity is reflected in the independence of the masses; as force, in
their mutual reciprocal action; as world-power, in the great laws. It is from all
three sides the same essence, the same power, the same substantiality; each side of
the power reflects the two others in such a way that it must itself be said to be
the whole power. If we have before us the actual masses, then we have in their
common property force, in force attraction, and in attraction gravity as universal
power. If the attracting forces are given, then we have with them at once
% --- p. 36 ---
the masses, for the masses are, in this light, nothing other than the existence of
the forces, and force and mass again moments of gravity as power. If we have,
finally, gravity as power, then the existence of the masses with the corresponding
distribution of the forces is precisely its own reality, an issuing of the
particular in which it subsists as the universal. From all sides, then, objective
power is reflected.

But objective power cannot reflect itself by itself; only self-power can do that.
Let us look again at the two masses. The moving $A$ is cause, the motion of $B$
effect. The cause, which is reality in $A$, is ideality in $B$, for the ideality of
the cause is to be --- reflected --- in its effect. Now since $A$ is removed from
$B$, it has as cause its effect outside itself, is as reality removed from its
ideality. And yet the ideality is in the cause, and the cause is in $A$. $A$ cannot
then possibly be removed from what it itself is; as surely as it has the cause in
itself, it must necessarily also have the ideality in itself. This contradiction ---
that each mass, namely, has its ideality in itself by having it outside itself ---
is a contradiction in the objective reflexion, which plainly proves that it cannot
bear itself, but must be borne by Spirit. From Spirit, in and by the inspiriting,
the reality of the mass receives an addition in which the ideal unites itself
objectively and essentially with the real; this ideal-real addition is force. By
means of force, matter in $A$ is
% MISPRINT: „Formeldelst“ printed for „formedelst“ — one letter l too many. The
% ascender between „Forme“ and „delst“ stands plainly in both copies. The errata
% leaf does not correct it. Reproduced as printed; translated by the sense.
according to possibility more than matter, it is cause. That the activity sets in,
that the cause becomes actual and gets its effect, is objective reflexion in power.
If objective power were not reflected in Spirit, then gravity would not be the
universal power, then there would be no universal powers in nature, then the laws
would be impossible. Objective, universal power is not an abstraction from masses
and forces, but their own substantial essentiality, their indwelling foundation.
Objective power reflects itself --- by being
% --- p. 37 ---
reflected; it is in actuality objective, but --- by being objectified.

\subrunin{γ) Objectification in power.} The concepts of nature have a double-sided
% The letter read on the page (KB copy): γ, not the „y“ or „7“ the Indhold's OCR
% gives. Run-in head, as α) and β) above. The Danish prints „Naturbegrebene“ here
% and „Naturbegreberne“ further down in the same section.
subsistence, subsistence in actual things and subsistence in the self's conceiving.
The latter conditions the former; the concepts could not give themselves an imprint
in existing objects if they were not themselves stamped by an infinite conceiving.
It is not, after all, sensible natural things that in the objective sense determine
the concepts; on the contrary, it is the concepts of nature that non-sensibly
determine sensible things. To be the determining essence of the thing, the concept
must necessarily have determinateness in itself, must be determined in advance to
furnish determination. The determinate concept of mass does not first come to be by
means of existing masses, nor the concept of force by forces already existing. If
masses and forces had an existence independent of concepts, then things could at
once dispense with all concepts; there would be no reason at all in nature. If the
coming-to-be of the concepts coincided immediately with that of things, it would be
altogether impossible to say whence the determining came. It could not then be the
concept that determined the thing, for the concept was precisely to be determined by
the thing; but that neither could it be the thing determinate in itself that
determined the concept goes without saying, since what determines a thing is just
its concept. It stands fast, then, that the concepts of nature must have in them a
determinateness independent of the existence of things and decisive for things
themselves. Whence do the concepts now get this inner determinateness, indispensable
for all outward existence? From a conceiving. To determine concepts is to conceive;
to present the concepts, determinate in themselves and for that reason determining,
in existing and existence is not to \emph{conceive}, it is to \emph{work}. Nature
presents concepts, forms concepts out, but it does not conceive, it works. Spirit's
ideal working
% --- p. 38 ---
is conceiving, nature's rendering of what is conceived is not conceiving, it
% A comma, not a full stop, after the first „Begriben“; verified on the KB copy.
% The sequence of clauses is comma-joined, as printed.
is real working, material activity, the unfolding of the determinations of the
concept in sensibly actual facts.

The opposition of Spirit's conceiving and nature's working is grounded in the
reflexion of subjective and objective reflexion of power, an all-sided expression for
objectification in power. In the subjective reflexion of power, the ideal
objectification, the concepts of nature have being as spiritual determinations of
power, decisive for all material existing. But under the form of ideality, in
self-power's conceiving, the concepts do not themselves have existence; their being,
a being in knowing, is the infinite reciprocal shining of the determinations of power
in one another; the single and the particular reflect themselves here in the clear,
free universal: power in another is a moment for self-power. From the side of
existence, and so from the side of nature, the emphasis falls by contrast upon power
in another, upon object-power, upon reality, and the real objectification deposits
concepts in corresponding objects. The objects are things, a swarm of single members,
atoms and masses. The universal that determines as concept here merges immediately
into the single, is bound in power by single members, by outward objects. So it is
said in the language of physics: only the single members, the fundamental parts with
the things composed of them, have existence in nature; the universal does not exist.
And yet it is precisely the universal that comes to existence; it is, after all,
objectified in power precisely through the issuing of single members. The single
members work, take on through their reciprocal action the stamp of the particular,
and fall due to the universal; then it shows itself that the single and the
particular are subject to the universal. Objectively reflected power, the universal
power of nature, of objectivity, comes to existence in the forces of things and the
laws under which the forces work.

% A short centred rule stands after the last line of p. 38, closing division B.
% Present in both copies.
\streg

% Danish: \afdeling{C.}{Kræfter og Love.}   (printed pp. 39--61)
\afdeling{C.}{Forces and Laws.}

In objective reflexion universal power is both existence itself and the essence of
existence; but the difference between essence and existence is just as mighty as
the unity. Power in another, which is one with itself in essence, is different from
itself in existence; it subsists with this difference as object-power in the masses
of the members and as acting forces in the reciprocal action of the members; but it
also has subsistence in the universal form of essentiality, as reposing,
all-determining power, namely in the laws. The utterances of forces are facts, but
the laws look like thoughts, mighty thoughts, thoughts of nature. To know the acting
forces and the conditions of their activity is to know the laws; and yet the forces
are not laws, and the laws not forces. To come to an insight into the universal
unity-in-difference of forces and laws, we shall now investigate, from a mechanical
point of view, the laws'
% Letterspaced setting, verified on the KB copy: the whole phrase from
% „Oprindelse“ to „Sammenhæng“ is letterspaced. A full stop, not a comma, after
% „Sammenhæng“; tesseract reads a comma here.
\emph{origin, fulfilment and mutual connexion}.

\runin{a) Represented forces: the origin of the laws.}%
``The understanding,'' says Kant, ``does not borrow its (a priori) laws from
nature, but prescribes them to it.'' This proposition points to the subjective
origin of the laws and is justified by countless examples from analytical mechanics.
If two equally great forces act upon a material point in opposite directions, they
must cancel one another, and the point remain unmoved; that is a law. That this law
has objective validity is just as evident as that its validity is independent of all
experience. The law's universal validity lies in the necessity with which the
posited presupposition yields its consequent: the law is, in logical form, a
% --- p. 40 ---
hypothetical judgment, in mathematical form an equation. No weight attaches to
whether the assumed case occurs in outward nature or not, for the law is the nature
of the matter; it has its nature in itself. It is the thinking self that sets up the
presupposition; it is the thinking self that draws the consequence; it is the
thinking self that gives itself the law and assures itself of its objectivity.
There is neither in heaven nor on earth any outward power that avails to overturn
the law given in the train of thought or to make it doubtful; the only holder of
power able to cancel the law is the legislating subjectivity. Subjectivity cannot,
indeed, contradict the law without contradicting itself; but it can cancel it by
cancelling the presupposition, it can alter it by altering the presupposition.

That a law must alter when its presupposition is altered is itself a law. If the two
forces acting in opposite directions are unequal, the point must come into motion,
and the difference of the forces will then be the moving force; but if the two
forces form any angle whatever with one another, it will at first glance look as
though we had got a new law. The necessity with which the law holds is, namely,
always bound to a presupposition of one or more conditions, and consequently is not
an absolute but a hypothetical necessity. By introducing now fewer, now more
conditions into the presupposition, and letting the conditions vary, the legislating
subjectivity can let the one law proceed from the other, can in this way lay out and
develop systems of laws. We see it in the parallelogram of forces. One should think
that here a law of its own was set up for a special case of its own; but the thinker
soon discovers that the apparently special case is precisely the general case, for
two forces acting at one
% --- p. 41 ---
point must always form an angle with one another.%
% The footnote mark is *), as elsewhere in the book. One footnote on this page.
% The formulas are set from the image, not from the OCR. The radical's bar
% covers the whole expression, „cos (P. Q)“ included; in the KB copy the bar is
% set in two pieces, which is not reproduced.
\footnote{Let the two forces be $P$ and $Q$ and their resultant $R$. For the
parallelogram of forces the equation then holds:
$R = \sqrt{P^2 + Q^2 + 2\,P\,Q \cos (P.\,Q)}$. If now the angle
$(P.\,Q) = 0^{\circ}$, then $\cos = 1$, and $R = P + Q$; if the angle
$(P.\,Q) = 180^{\circ}$, then $\cos = -1$, and $R = P-Q$. The law of the
parallelogram is a fundamental law for the composition and resolution of mechanical
forces, from which an infinite system of particular laws lets itself be developed.}

Motion along the diagonal of the parallelogram is, essentially regarded, a composite
motion, for the force by which it is determined is a co-working of two component
forces; but just as the difference between simple and composite force is only
relative, so too the difference between simple and composite motion is relative. The
resultant marked by the diagonal can in its continuation itself become a component
force, whose angle with another corresponding one yields a new resultant in a new
direction, so that the track of the motion becomes, with the finite alteration, a
broken line. If the alteration happens continuously, so that the component forces
(the composants) are incessantly converted into a resultant, and the resultant again
into a component force, then a curved line comes forth.

The true concept of motion in the curved line, whose law mechanics exactly
determines, rests upon a resolution of two layers of misunderstanding. The ancients
supposed that it belonged to the peculiar essence of certain bodies to move in
curves, that notably the celestial bodies, in accordance with their higher nature,
moved in circular orbits. That was a misunderstanding. By resolving this
misunderstanding physics became an exact science. But to make composite motion clear,
one had to resort to parcelling out. Under the parcelling the curved line becomes a
broken line: let the fragments be as small as they will, they remain
% --- p. 42 ---
fragments nevertheless. Hereby a new misunderstanding came forth. The method of
parcelling has to such a degree given offence to the more recent philosophy of
nature that speculation has openly broken with mechanical physics and sought the
truth in a modified conception of the ancient representation. Hegel is so angry at
the parcelling that he not only forbids physics --- what he concedes to logic in
ample measure --- any use of abstraction of its own, but goes so far as to want the
representation of forces wholly rooted out. ``Newton,'' he says, ``assures us that a
leaden ball \textit{in coelos abiret et motu abeundi pergeret in infinitum}, if
(quite so: if) one could only impart to it the requisite velocity. Such a separating
of outward from essential motion belongs neither to experience nor to the concept,
but to abstracting reflexion alone.'' By introducing ``the representation of
forces,'' Newton is supposed, for all his reservations, to have done considerable
damage to the true knowledge of nature. These are speculative misunderstandings.
These misunderstandings resolve themselves and disappear when the relation between
continuity, parcelling and infinite discreteness is sufficiently cleared up.

If a body be thrown upward in a slanting direction, then the force answering to the
impulse can be regarded as a composition of two forces, of which the one would carry
the body vertically upward, along the Y-axis, the other drive it horizontally
forward, along the X-axis. If now every other action be thought away, the body will
go along the diagonal of the rectangle which the coordinates x and y form with one
another, and the path s traversed in the time t will be equal to $\sqrt{x^2 + y^2}$.
The angle $\alpha$ formed by s and x will be constant, and its tangent
$\frac{y}{x}$ of course likewise. In such a case it is then indifferent whether one
brings out the continuity or the parcelling or the infinity of the
% --- p. 43 ---
discreteness. Every finite piece of s is a continuum, and the arbitrarily posited
limits do no damage; but s can also be resolved into elements answering to elements
of x and y:
\[
ds = \sqrt{dx^2 + dy^2};\quad \frac{y}{x} = \frac{dy}{dx}.
\]

It is otherwise when the body is, during the motion, represented as subjected to a
vertically downward action; here an essential difference between parcelling and
continuity shows itself.

If the action be represented as setting in at certain intervals, then s will separate
into finite pieces, $s_1$, $s_2$, $s_3$ and so forth; the angle $\alpha$ will alter
each time, since each piece of s gets an inclination of its own to the axis of
abscissas, and $\frac{y}{x}$ will obtain different values. Every piece is a diagonal
in a parallelogram of forces, one side of which is a continuation of the previous
resultant, and the other side of which is formed by the supervening action: the
body's path becomes a broken line whose parts can be measured and exactly determined.

If the partial action be exchanged for a continuous one, such as gravity g, then s
alters its direction continuously, and the path becomes a curved line. Straightforwardly
in the form of continuity such a path cannot be determined, for the infinity of the
alteration cannot be determined for itself. Straightforwardly in the form of parcelling
the path cannot be determined either, for the parcelling is a finite discreteness which
breaks the continuity. Only that reciprocal determination of finite and infinite which
is conditioned by infinite discreteness --- whereby an operation with elements can
become just as exact as a calculation with finite magnitudes --- can resolve the
parcelling into a finite determining of the infinite.

% The book writes „tg“, not „tang“ or „tan“; reproduced as printed. The period
% after α is the book's own multiplication sign, likewise as printed.
The equation: $y = f\,(x) = \operatorname{tg} \alpha .\, x -
\frac{g}{2\,a^2 \cos^2 \alpha}\, x^2$
% --- p. 44 ---
shows how the smallest alteration in x continuously brings with it a corresponding
alteration in y; from which again follows a continuous alteration in s. The manner
in which y depends upon x expresses the law. We have in the law the universal essence
determining for the infinity of the phenomena --- the infinite quantity of single
cases. But in this infinity we have at the same time finitude. In an essential manner
finitude is determined by the determinate function; in the outward form of the
appearances, laid out for division and parcelling, the infinity of the finite comes
into view in the swarm of single values. The transition from finite to infinite and
from infinite to finite is determined by the mutual relation of the elements. Every
element of the curve coincides with a corresponding element of the curve's tangent;
how, i.e., by what law this tangent alters its direction in the infinite, is seen from
the equation: $\frac{dy}{dx} = \operatorname{tg} \alpha -
\frac{g}{a^2 \cos^2 \alpha}\, x = \varphi\,(x)$.

What holds of the law holds also of the forces. The resolution of forces is no
straightforward parcelling; it is an analysis under which each component force is
brought out for itself by an abstraction necessary for science. From the side of
essence every one of the component forces is a moment of the resultant. Where a
parcelling really takes place, it is nature's own parcelling, existence itself, that
in the form of otherness makes the one force independent of the other. Nor is the
composition of forces a simple addition; the force acting in the resultant is the unity
of the component forces, a unity in which their independence is wholly cancelled. Here
there is not, as Hegel nevertheless asserts, the least ground for entering, on the
concept's behalf, any objection whatever against the analyses of physics.

But is it now, strictly speaking, the laws that determine
% --- p. 45 ---
the forces, or might one not perhaps reverse the relation and say that it is really
the forces that determine the laws? A superficial conception of the manner in which
rational mechanics presents its problems seems at first to speak for the latter. An
example. A certain velocity is imparted to a body in any direction whatever, only not
vertical, whereupon it is left to the action of gravity: of what character will the
path be? The answer is that the path, if the motion be assumed to go on in a vacuum,
must be a parabola. The law of the parabola is accordingly the law of the motion. The
represented forces and the posited conditions seem here to be the first, and the law
the second, derived from them; and so: the forces are determined by representation,
the law by thinking, and what is thought depends upon what is represented. How
one-sided this way of seeing is comes to light from the relation between
representation and thinking. The merely representing consciousness does not know the
law at all. It cannot even, without the aid of thought, set up what is represented as
a presupposition; for the presupposition is only in the reflexion of its consequent,
but the consequent is only for thought. A determining of the forces by mere
representation is an act of imagination, a fleeting, preliminary determining that
loses all significance when thought is wanting. The determining representation is a
function of thought; it is sent on ahead by thought itself and stands from the
beginning in thinking's service.

From the standpoint of mechanics the matter easily takes on an appearance as though
thought's task were exclusively to discover the laws, as though the laws were in and
by themselves objective without needing objectification. The objectification would
accordingly be merely for subjectivity's sake, not for the laws'. It is an illusion
which strict science will instantly reject once it becomes aware of it. If the
law-discovering subjectivity were not itself legislating, it
% --- p. 46 ---
could not discover the laws at all; it could not judge of the conditions, not
distinguish the possible from the impossible, not determine the hypothetically
necessary, not assure itself of the a priori validity of the alleged laws. And what
becomes of the objective laws when they are separated from the legislating,
objectifying subjectivity? Dead, meaningless formulas! A law of equilibrium for two
points connected by a line and acted upon by the forces $P_i$ and $P_m$, whose angles
with the X-axis are $\alpha_i$ and $\alpha_m$, the unknown force in the connecting line
being $R$, its angle with the X-axis $\delta$, and $\varphi_i$ and $\varphi_m$
arbitrary angles, is preformed, for a system of several mutually connected points, by
the equations:
% The second equation has ÷ where the first has +. Verified on the KB copy at high
% magnification: the sign is unmistakably ÷, and both copies agree. ÷ was the usual
% minus sign in Danish setting at this period, so this is presumably Nielsen's own
% notation and not a compositor's slip; the errata leaf does not mention the place.
% Reproduced as printed. The rows of dots stand as four spaced points in both
% equations.
\begin{align*}
P_i \cos (\alpha_i - \varphi_i) + R \cos (\delta - \varphi_i) + \ .\ .\ .\ . &= 0\\
P_m \cos (\alpha_m - \varphi_m) \div R \cos (\delta - \varphi_m) + \ .\ .\ .\ . &= 0
\end{align*}
But as little as these equations would have significance in themselves if no one
understood them, just as little would the law preformed by the equations have in itself
the least validity if there were no subjectivity, human or divine, that could think the
law, discover the law with self-legislating assurance, and objectifying maintain the
law's objectivity.

\runin{b) Actual forces: the fulfilment of the laws.} From the subjective standpoint,
so far as it is merely a question of represented forces, the agreement between the
divine and the human reason is pervasive; for the laws of nature are laws of thought,
and the laws of thought at once both human and divine. But although all laws of nature
are laws of thought, the laws of thought are in their free ideality not yet laws of
nature. Apprehended for themselves, the thought laws have only an ideal objectivity,
articulated by the representations; in nature itself, by contrast, the laws have a
material actuality answering to the forces; this actuality is their fulfilment.

From the objective standpoint, so far as one looks expressly
% --- p. 47 ---
to natural actuality, the dualism, the difference in principle between divine and
human objectification, comes into view again. From the divine side all facts must
surely be derivable from the determining laws; but from the human side the actual laws
must be drawn out of given facts.

The laws uttered in factual propositions about nature as in a divine language of power
cannot, however peculiar they may be, be thought to stand in conflict with the general
laws of reason; but it makes an essential difference whether the laws are treated as
abstractly universal determinations of form, and so \emph{mathematically}, or are shown
in intuitable unity with the content of actuality in which they have their fulfilment,
% Misprint: „alsaa“ for „altsaa“. The letter t is missing. Verified on the KB copy,
% where the word likewise stands „alsaa“, while the same word three lines above
% stands „altsaa“; the fault is therefore the compositor's and not this copy's. The
% book's own errata leaf does not correct it.
and so \emph{physically}.

``Just as mathematics,'' says H.\ C.\ Ørsted, ``must treat the objects of nature in its
own way, so also must physics. The ideal of the former is to accomplish everything
without borrowing anything from experience, and when, owing to the impossibility of
wholly attaining the ideal, it does make a loan from experience, this must be something
as general, as little individualized, as possible. It belongs to its artistic perfection
to take the most circuitous ways rather than to take up the alien stuff of experience
into itself. Quite otherwise does it stand with the science of nature. Nature is its
scientific property; it does not borrow from it, but draws from it as from its proper
source. Its ideal is to present the objects to intuition'' --- and so the laws in their
fulfilment --- ``and with a view to this it works up forms and magnitudes.''

The scientific significance accorded to physics as a science of experience is naturally
dependent upon the manner in which the essence of experience is conceived as a whole.
True experience rests upon an even co-working of sensation and thinking; but this
co-working can be perverted, according as
% --- p. 48 ---
it is either sensation that overwhelms thinking, or the reverse. If sensation is
one-sidedly predominant, then the gaze fastens exclusively upon the facts in their
singularity. Wading in heaps of accumulated facts, the observing consciousness gets so
busy sorting, combining, adding and subtracting the many singularities that it can
hardly recognize the thought in the universal representations it has itself abstracted.
If a priori thinking is one-sidedly predominant, then the facts are involuntarily
reduced to fleeting, transient appearances of the universal. Without entering upon any
exhaustive investigation of the material in its singularities, the observer, tempted by
one or another favourite hypothesis, will all too easily spin out a theory and give it
out for a truth grounded on experience. Against both aberrations physics, the true
science of experience, must wage a constant struggle.

The fundamental scientific concept upon which the science of experience builds is the
concept \emph{induction}.

That the induction must become incomplete when the observations are defective goes
without saying. But that an outward heaping together of merely sensible observations,
even if they be multiplied on an ever so enormous scale, can likewise give no true
knowledge of nature, is then also evident.

One of the most thorough proofs of a superficial conception of induction's logical
validity we have in Stuart Mill's inductive
% The citation of Stuart Mill runs across the page break and closes only on p. 49
% („Aarsag og Virkning.“). No compositor's fault --- merely a citation over two
% pages; hence one „ too many on p. 48 and one “ too many on p. 49.
logic. ``We learn,'' in this philosopher, ``from experience that there is in nature an
invariable order of succession, and that every fact is always preceded by another fact.
What we call cause is the invariable antecedent, what we call effect is the invariable
consequent. The concept of necessity, of which so much is made, here finds its
explanation. Necessary is that which will occur, whatever we may assume with respect to
all
% --- p. 49 ---
other things. Necessity is present when the antecedent is sufficient and complete, when
it contains all the conditions sought, and no new condition is required.
Unconditionally to be about to follow --- that is the whole concept of cause and
effect.'' So far as such a manner of explanation can be supposed to have something
dazzling about it, this lies in the unembarrassment with which thought as thought is in
one and the same breath invoked and dismissed. Seen in the light of thought, it is
incontestably true that the relation between cause and effect encloses necessity, that
the consequent must unconditionally occur when all the conditions are present. But
thought is not allowed to count as thought; it is not thought's a priori assurance, but
the testimony of the sensations, of the repeated sensations, that here, in the name of
experience, gets to signify something. That a sum of sensible observations, however
ingeniously they be sifted and weighed, cannot possibly secure us an ``unconditionally
to be about to follow,'' is evident to such a degree that it has not even been able to
escape the author's own logical attention. Stuart Mill ``is convinced that a man
accustomed to abstraction and analysis'' can very well ``conceive that in certain
places --- for instance in one of those firmaments into which astronomy now divides the
universe --- events might follow one another by pure chance without any fixed law, and
there is nothing, either in our experience or in our mental constitution, that gives us
sufficient ground to believe that this nowhere takes place.'' According to such a
logic the laws do not belong to the essence of things, but are in certain regions of
nature accidentally pasted on the outside of things, much as one pastes a device on a
bottle. If Hume rose from his grave, he would soon convince Mill that he has no need
whatever to trouble himself with the long journey to ``one of those firmaments'' in
order to get the laws away and necessity dissolved;
% --- p. 50 ---
that it happens far more easily when he stays where he is and bethinks himself of what
he has himself said.

There is reason in nature; this central thought, which Ørsted has so beautifully uttered
and maintained in our literature, gives induction another and deeper scientific
significance than the Millian explanation; here one starts from the premiss that the
laws lie in the essence of things. This central thought must have stood clear before
Newton when the idea of universal gravitation dawned upon him. The eye that saw in the
falling apple an utterance of the force that also draws the moon down toward the earth
belongs not to sensation but to thinking. It was thought's intimation of a universal
essentiality indwelling in all bodies that gave the impulse to the great discovery. But
with the thinking, sensation too came into its right. The universal essentiality
indicated in one fact had again to be more closely determined by facts. Only so could
the thought law confirm itself as a law of nature, a universally valid law of
experience; in other words: the law of gravitation had to be sensed and seen in its
fulfilment.

The reciprocal action between empirical physics and rational mechanics shows itself
here in the fairest light. For the observations to begin and experience to be
introduced, the problem had to be put into an equation; but the equation that was to
give the observations significance had itself to be built upon observations.

That the moon's fall must stand to the apple's as the force that draws the moon to the
force that draws the apple goes without saying; but in what relation do these forces
stand to
% The fraction stands in the setting as a stacked fraction in mid-line; the figure
% 15 5/8 is the distance of fall in feet.
one another? The force that draws the apple is known from the law of fall ---
$15\tfrac{5}{8}$ feet in the first second; the force that draws the moon --- the moon's
fall in 1 second --- can, since the moon's mean distance and mean period of revolution
are known magnitudes, likewise be determined. If the moon's distance be denoted
% --- p. 51 ---
% Θ is capital theta, ϱ the curled rho; both read on the KB copy. The printed
% variables stand in upright antiqua; per the house rule this is not reproduced,
% and the ordinary mathematical italic is kept.
from the earth's centre by R, and the angle answering to the arc (in one second of
time) by $\Theta$, then the piece denoting the moon's fall toward the earth is
$\varrho = 2 R \sin^{2} \tfrac{1}{2} \Theta$. The relation of the forces is
accordingly determined by the quotient
$\dfrac{2 R \sin^{2} \tfrac{1}{2} \Theta}{15\tfrac{5}{8}}$;
for R to be determined, the earth's radius r must also be known. Here is a circle of
presupposed observations, and yet it cannot without special preparation be enriched by
new observations. It lay near, indeed, to seek the ground of the difference in
attraction in the difference of distance, but with this general surmise no experience
could yet be made; only with the assumption that the attracting forces stand in inverse
ratio to the squares of the distances was the equation determined:
\[
\frac{2 R \sin^{2} \tfrac{1}{2} \Theta}{15\tfrac{5}{8}}
  = \frac{\dfrac{\mu}{R^{2}}}{\dfrac{\mu}{r^{2}}}
  = \frac{r^{2}}{R^{2}}.
\]
The equation thus determined was a question put to actuality in so exact a form that
the fact itself necessarily had to answer either yes or no.

The first answer that Newton received to his question was, as is well known, an
unfavourable one.
Since actual nature could not err, he drew, disheartened, the conclusion that the error
lay in his thought; that, curiously enough, was his real error. In the matter itself his
original thought was as faultless as nature. The error lay in the value then current for
the earth's radius. When this error was afterwards corrected by a new survey, Newton
found his thought confirmed; and only then had the turning point arrived at which the
law found for a particular case could, by an induction answering to continued
observations, be extended to a universal world-law.

The Newtonian law has in its time furnished a magnificent example of how the fulfilment
of a law can, by means of the
% --- p. 52 ---
supervening conditions, alter a law's appearance so that it becomes almost
unrecognizable. In occupying himself with the perturbations, particularly with the
peculiar alterations to which the moon's orbit is subject (the motion of the apogee),
Clairaut came to raise doubts about the universal validity of the Newtonian law. He
remarks, on the occasion of a dispute with Buffon, who would defend the law with a
priori grounds, that the ways in which the Newtonian law had hitherto been supported
were so encumbered with exceptions and piecemeal taking into account of various
influences that a law different from $\dfrac{a}{x^{2}}$ could give just as good an
agreement between the results of calculation and of observation. In solving ``the
problem of three bodies'' Clairaut came provisionally to the result that the expression
for the law would have to be, not $\dfrac{a}{x^{2}}$, but
$\dfrac{a}{x^{2}} + \dfrac{b}{x^{4}}$. Further investigations led, however, to the
result that the final solution of the problem precisely confirmed Newton's formula. The
reciprocal action between the three bodies --- the central body, the perturbed and the
perturbing body --- goes on, namely, in each and all according to the same law, so that
the varying conditions bring with them only modifications: the particular is in this
example a
% The footnote mark is printed *) --- the book's usual mark.
modification of the universal unchanged in itself.\footnote{That an alteration of the
fundamental formula may nevertheless perhaps be possible must be conceded, since
distinguished astronomers, such as Bessel, still harbour some doubt.}

But the conditions of actuality under which a law is fulfilled may also be of such a
character that they introduce one or more new determinations of law, heterogeneous with
the given universal form, and thereby effect an essential alteration in the law. One may
think of the alteration that goes on with the law of fall when the fall happens on an
inclined plane and the modifying circumstances are taken into consideration;
% --- p. 53 ---
% The footnote mark is printed *), set after „Pendul“ and before the semicolon.
% In the note the literal calculation is set in mathematical italic; the printed
% variables stand in upright antiqua, which per the house rule is not reproduced.
% Where a letter carries a prime or an index, the whole enumeration is set in
% math, so that m, m', m'' does not change face part-way along the list. The rows
% of dots stand as spaced points.
of the alteration that happens with the law of pendulum oscillations, according as one
considers either ``the simple'' or ``the compound'' pendulum\footnote{At equal angle of
displacement and force of gravity, the times of oscillation of two simple pendulums
stand to one another as the square roots of the pendulums' lengths. With the compound
pendulum the law cited holds, by contrast, only for the pendulum's reduced length, i.e.,
for the distance from the point of suspension to the centre of oscillation. In
determining the centre of oscillation the moment of inertia is introduced. If $m$, $m'$,
$m''$ $\ .\ .\ .\ .$ are the parts of the mass, and their distance from the point of
suspension $x$, $x'$, $x''$ $\ .\ .\ .\ .$, then the moment of inertia is
$mx^{2} + m'x'^{2} + m''x''^{2} + \ .\ .\ .\ .$ If the whole mass be called M, and the
distance of the centre of gravity from the point of suspension a, then the reduced
length
$l = \dfrac{mx^{2} + m'x'^{2} + m''x''^{2} + \ .\ .\ .\ .}{Ma}$
is the designation of a law of its own. That the laws fused together in the fulfilment
do not contradict one another goes without saying.}; of the difference between the
parabola and ``the ballistic curve,'' etc. The deviation is in the last case
particularly so pervasive that in actuality a wholly new law shows itself.

The fulfilment of a law happens, then, in essentially different ways, according as the
supervening conditions either merely alter the phenomenal expression, or bring with them
a transformation of the law, or even change it in such a way that it is taken up into
another law and disappears in it. It stands with the laws as with the forces. How the
particular laws, in mutual reciprocal action, partly confirm, partly alter, partly
mutually cancel one another can first be fully elucidated when we consider the
corresponding forces from the standpoint of totality.

% Run-in head, read on the page itself (p. 53) and not in the contents; colon and
% wording agree with the Indhold.
\runin{c) Totality of forces: the realm of laws.} The laws of nature are unalterable;
this fundamental proposition, on which physics stands fast, and without which it would
not be the strict, exalted science it is, seems, when it is
% --- p. 54 ---
uttered in all unconditionality, to strike down at one blow everything set out above
about alterations in forces and laws. And yet the philosophy of nature knows itself, on
this point, in full agreement with physics. Physics has, namely, a right to set up its
proposition --- that the laws of nature are unalterable --- only because logic has a
right to set up its own: that the laws of thought are unalterable. When, then, logic
proves that the unalterability of the laws of thought consists precisely in their
alterability, it has \textit{eo ipso} shown that the same must hold of the laws of
nature. The unalterable cannot be held fast for itself in finite separation from the
alterable, any more than essence can be set up finitely over against semblance and
phenomenon. Essence for itself is semblance and nothing, when semblance and phenomenon
are thought away. Essence is itself involved in the alterations; its unalterable
subsistence shows itself precisely in this, that through all its alterations it
preserves itself unaltered. The unity of unalterable and alterable does not efface the
difference; no, it confirms the difference: here is essential difference, because here
is essential unity.

The forces of nature are unalterable; and yet they are altered, in that they work and
enter into reciprocal action: in their essence they are unalterable, but in their
utterances they are highly alterable. So too with the laws: they are in their essence
unalterable, but in their fulfilment alterable. When the proposition about the
unalterability of the laws of nature is not understood in this way, but is proclaimed
outright as a dogma, an article of faith in which essence and phenomenon, possibility
and actuality run out into one, then protest must be made. In spite of the catholic
semblance of infallibility with which the dogma is preached, sound sense will easily
smell a rat and ask: what does accidentality mean? what becomes of freedom?

That necessity and accidentality are everywhere united in the great whole of nature,
that the necessary in the world of facts is accidental and the accidental necessary, has
its
% --- p. 55 ---
ground in this, that the unalterable subsists in alterations. Every actual law must have
actual conditions and show itself determining for active forces. Let a law of nature be
as universal as it will, it can still never be fulfilled in the merely universal; a
universal fulfilment is an abstraction: the actuality of the fulfilment is in the actual
cases. The cases all fall under the law; but they do not for that reason lose their
accidentality, do not cease to be cases. Were there no bodies that could fall, the law
of fall would not be fulfilled. If there were no other bodies in celestial space, the
force with which the earth's body acts upon other bodies could not utter itself, could
not become actual. In essence and possibility the law of fall is the same unalterable
law, whether in given cases it is fulfilled or not fulfilled; but in actuality it is not
the same, for when it is not fulfilled it is unactual. In essence and possibility the
earth's force of attraction is unalterably the same, whether it works or does not work.
But actuality introduces an alteration, for only the working force is actual; the force
that does not work is only possible.

It will now easily be seen how law and case enter into a many-sided reciprocal relation.
That a stone falls to the earth is, with respect to the single stone, only a single
case, but it expresses at the same time the necessity with which all stones in all like
cases must fall. Bodies are, as regards gravity, under the same law because they are
under the same necessity, a necessity that expresses their indwelling, universal
essence. But, curiously enough, there are also laws which permit the accidentality in
each single case to withdraw itself from the
% The footnote mark is printed *). The note begins on p. 55 and runs over onto
% p. 56, where it continues without a new mark.
universal necessity.\footnote{For because one man marries at a certain age, it is not
necessary that others marry at the same age; because one man hangs himself, it is not at
all necessary that others hang themselves too; because one man in distraction posts a
letter without an address, it is not for that reason necessary that others must do
likewise. And yet the calculus of probability shows that large groups of such cases,
from whatever sphere of existence they be drawn, let themselves be brought under
corresponding laws.}
% --- p. 56 ---
There are, then, in the world of cases two essentially different circles of laws,
\emph{laws of essence} and \emph{laws of accident}. In the first circle accidentality is
necessary, in the second necessity is accidental. The two circles form, in their sharply
stamped opposition, two extremes which must not be confused. But the extremes do not
stand over against one another separated by a yawning chasm; innumerable determinations
of transition form a web in which the threads of necessity are, through manifold
crossings, as it were woven into the motley carpet of the accidentalities.

There is a pervasive connexion in the whole of nature, no leap, no interruption, no newly
created matter, no interpolation of new forces, no essentially new laws: everywhere
events engage in one another as links in a closed chain; the new that is formed is formed
out of the old, out of what already was in primeval time. The Hellenes, says Anaxagoras,
wrongly assume becoming and passing away, for no thing becomes or passes away, but things
are composed out of subsisting things, and thus they would more correctly call becoming
composition and passing away separation. What Anaxagoras divined, mechanical natural
science has from its standpoint sufficiently endorsed. Just as there is in the universe a
given system of ``eternal forces,'' so there is also a system of ``eternal laws'' closed
in itself; this system is \emph{the realm of laws}.

To have discovered and conceived nature's infinite unity of connexion, and thereby opened
to itself a magnificent view into
% --- p. 57 ---
the realm of laws, is ``modern science's'' pride. The philosophy of nature acknowledges
what is true and warranted in this. But in the hymns that ring out in honour of deified
necessity, unresolved dissonances are not seldom heard, jarring tones which recall that
modern science, out of excessive joy over necessity, is on the way to denying freedom.

Can freedom, one asks, really consist with necessity? This question presents, as regards
the knowledge of nature, two principal sides, according as one reflects rather upon
accidentality or upon the pervasive unity of connexion.

Just as all material actuality has a moment of accidentality, so freedom's ideal
actuality has a moment of arbitrariness. What accidentality is in the outward, the
factually given, arbitrariness is in the inward, what originally determines. If necessity
can be reconciled with the accidental outwardly, then it must surely also be reconcilable
with the arbitrary inwardly. But arbitrariness in nature? The mere representation of this
is an offence to physics. The offence will pass off, however, once the physicist comes to
himself. It is, after all, precisely physics itself which, without in any way having to
do with the arbitrary, supplies philosophy with the most important presuppositions and
most reliable points of support for the recognition of the moment of arbitrariness.

Physics has broken with speculation in order to hold unconditionally to the fact. That is
a decisive turning point in the history of science. To hold to the fact is to take it as
it is, without reasoning about why it is now just so, whether it might not perhaps just
as well have been otherwise. The fact is because it is; but what only is because it is
is, in its given necessity, immediately to be apprehended as
% --- p. 58 ---
something accidental. Observation sees in the many peculiar facts nothing but peculiar
cases; whence, then, do the given cases have their original determinateness? From other
cases. And the other cases? From other cases. And these again? From other cases. The
endless reasoning that incessantly exchanges case for case drowns in accidentality
without getting so far that it can even so much as glimpse the grounding necessity.

Only by holding to reason and starting from the a priori presupposition that nature's
order, grounded in itself, rests upon an all-embracing ground-unity that determines
everything from within, can the observing investigator discover necessity. Every
particular case must, after all, have its particular ground: as many cases outwardly, so
many particular grounds inwardly. But the separation of the ground-unity into particular
grounds cannot possibly be accidental; to make what grounds accidental is to annihilate
all grounding. The ground-unity must therefore itself determine itself from within to a
manifold of different grounds, connected in their
% Misprint: „Forkjellighed“ for „Forskjellighed“. The letter s is missing.
% Verified on the KB copy, where the word likewise stands „Forkjellighed“, while
% „forskjellige“ one line above stands in full; the fault is therefore the
% compositor's and not this copy's. The book's own errata leaf does not correct it.
difference; the ground-unity must, for the sake of the connexion, necessarily be a
determining self-unity.

If the observer will now, from this standpoint, take aim at accidentality, he shall
assuredly discover arbitrariness. So as not to be dazzled by extraneous considerations,
one may choose such general, uncompounded facts as: the force of cohesion, the force of
gravity, magnetic forces, electrical attractions, etc. None of these forces lets itself be
derived from the others; nor has physics any prospect of being able to deduce the
corresponding properties of things either from their other properties or from the general
essence of matter. The difference of the fundamental forces must have a determining
ground. But what is the determining ground of accidentality? It is arbitrariness. As what
is
% --- p. 59 ---
determined is, so must what determines be. The greater the emphasis that falls upon the
determinate force in its singularity and immediate accidentality, the greater the emphasis
that must fall, from the side of the grounding, upon the exclusive utterance of will, the
act of arbitrariness by which it has been determined. Accidentality says that the given
force is because it is; arbitrariness says that it is because it shall be. If the gaze
fastens upon the single facts in their singularity, then everything inwardly is arbitrary,
just as everything outwardly is accidental. What accidentality is for the system of facts,
arbitrariness is for the system of grounds.

Rational mechanics is from first to last systematically conditioned by arbitrary and
accidental presuppositions, and yet an articulation of necessity so sure and firm in
thought that it truly does honour to human reason. Without a reciprocal determination of
accidentality and arbitrariness it is impossible to set oneself any mechanical problem
whatever. Let the mathematician choose whatever problem he will, he sees at once that he
must either posit something himself, and so begin arbitrarily, or presuppose
% The footnote mark is printed *). This note is the book's longest so far: it
% begins on p. 59, fills the whole of p. 60 beneath the two first lines of text,
% and ends on p. 61 --- throughout without a new mark. It is gathered here at the
% mark. The rows of dots stand as spaced points, now three, now four; the number is
% reproduced as printed.
something as given, and so begin accidentally.\footnote{How arbitrariness and
accidentality stand in principle to the necessity of connexion can, as regards exact
analysis, be seen for instance from \emph{interpolation}. A number of scattered points,
lying on a certain determinate curve, is given. Whether the given curve answers to an
actual fact and so far falls under the accidental, or is expressly posited by a free
self-determination and belongs under the arbitrary, is indifferent. The only thing that
can be known from the beginning, because it lies straightforwardly in the necessity of
connexion, is that the ordinate y and the abscissa x must for every interval answer
exactly to one another; $F(x,y) = 0$ or $y = f(x)$. The problem to be solved by the
interpolation calculus is then the following: from the \emph{given} values of a function,
answering to \emph{given} values of the independent variable, the argument, to calculate
the value of the function answering to any value whatever of the argument lying between
the given ones.

If it can now be presupposed that the curve answering to the given function has a
continuous and even course, then the problem will be soluble by an approximation which is
the greater, the closer the given values of the argument lie to one another; the
interpolation formula becoming:
\[
y = a_{0} + a_{1} x + a_{2} x^{2} + \ .\ .\ .\ a_{n} x^{n}. \qquad (1).
\]
To determine the (n + 1) constants entering into this, (n + 1) corresponding values of the
argument and the function must be given. If these are respectively:
\[
x_{0},\ x_{1},\ x_{2},\ \ .\ .\ .\ .\ x_{n}\ \text{and}\ y_{0},\ y_{1},\ y_{2}\ \ .\ .\ .\ y_{n},
\]
one can write down (n + 1) equations of the form:
\[
y_{r} = a_{0} + a_{1} x_{r} + a_{2} x_{r}^{\,2} + \ .\ .\ .\ a_{n} x_{r}^{\,n}.
\]
By the solution of these (n + 1) equations of the first degree the constants in equation
(1) can be determined. Geometrically regarded, the problem can accordingly be solved by
drawing through the given points a parabola of the n\textsuperscript{th} degree.

The solution of the (n + 1) equations of the first degree is avoided by bringing equation
(1) into the form:
$y = X_{0} (x-x_{1}) (x-x_{2}) \ .\ .\ .\ (x-x_{n}) + X_{1} (x-x_{0}) (x-x_{2}) \ .\ .\ .\ (x-x_{n}) + \ .\ .\ .\ X_{n} (x-x_{0}) (x-x_{1}) \ .\ .\ .\ .\ (x-x_{n-1})$,
in which every term contains all the factors $(x-x_{0})$ etc. but one, the first term
namely lacking the first factor, the second the second \ .\ .\ .\ . the
n\textsuperscript{th} the n\textsuperscript{th}. The constants $X_{0}$, $X_{1}$,
$\ .\ .\ .\ .\ X_{n}$ are determined from the corresponding given values of x and y; thus
$X_{r}$ from
$y_{r} = X_{r} (x_{r}-x_{0}) (x_{r}-x_{1}) \ .\ .\ .\ .\ (x_{r}-x_{r-1}) (x_{r}-x_{r+1}) \ .\ .\ .\ (x_{r}-x_{n})$.
The interpolation formula sought (Lagrange's) accordingly becomes:
% The errata leaf corrects p. 60, line 1 from the foot: „sidste Led i Tælleren
% (x-x_n) --- læs (x-x_{n-1})“. Per the house rule the text stands here as
% printed, with (x-x_n) in the last fraction's numerator.
\begin{align*}
y = {}&+ \frac{(x-x_{1}) (x-x_{2}) \ .\ .\ .\ .\ (x-x_{n})}%
          {(x_{0}-x_{1}) (x_{0}-x_{2}) \ .\ .\ .\ (x_{0}-x_{n})}\ y_{0}\\
      &+ \frac{(x-x_{0}) (x-x_{2}) \ .\ .\ .\ (x-x_{n})}%
          {(x_{1}-x_{0}) (x_{1}-x_{2}) \ .\ .\ .\ .\ (x_{1}-x_{n})}\ y_{1}
          + \ .\ .\ .\ .\ .\\
\ .\ .\ .\ .\ .\ {}&+ \frac{(x-x_{0}) (x-x_{1}) \ .\ .\ .\ (x-x_{n})}%
          {(x_{n}-x_{0}) (x_{n}-x_{1}) \ .\ .\ .\ .\ (x_{n}-x_{n-1})}\ y_{n}.
\end{align*}

This so richly articulated necessity of connexion can, as said, have arbitrariness just as
well as accidentality for its foundation.

If we now let the given curve with the even, continuous course be an image of given nature,
then it will be possible to see how speculation and empiricism, each from its own side, are
tempted by the necessity of connexion to apprehend this actual world, this present order of
things, as the only possible one. Speculation holds to the universal, stops at $y = f(x)$,
and for want of mathematical insight cannot get further. Empiricism starts from the
particular, the given single members. The more scattered the observed single members are,
the looser the whole connexion of nature seems to be. But as the sum of the observations
increases, the single members move closer to one another: the countless possible curves
approach coinciding with a single actual one. In his restless labour at finding out and
calculating the new intermediate values, the empiricist gets an ever stronger impression
that this order of nature, this curve with the even course, must be the only possible one,
not merely under the given presuppositions but in unconditional universality. To lift the
illusion, however, the empiricist need only go back to pure mathematics and bethink himself
of the general foundation of interpolations. The \emph{constants} in equation (1) must for
every \emph{different} curve be \emph{different}; but an infinite quantity of different
curves, answering to an infinite quantity of different series of (n + 1) corresponding
constants, are after all possible, and just so are infinitely many different worlds
possible.} There is, in the circle of conditions,
% --- p. 60 ---
% Only two lines of body text stand on p. 60; the rest of the page is taken up by
% the note from p. 59, which is set here at its mark.
the same relation between the necessary and the arbitrary as between the necessary and the
accidental: the
% --- p. 61 ---
consequence is equally necessary whether the presupposition is arbitrary or accidental.

Just as the moment of accidentality merges into the unity of connexion, so the moment of
arbitrariness merges into self-unity. The reciprocal determination of self-unity and unity
of connexion shows itself in the solution of any mechanical problem whatever.

% The chapter closes with a short centred DOUBLE rule. Per the house rule the form
% of the ornament is not reproduced: \streg serves throughout.
\streg

% Danish: \kapitel{Andet Kapitel.}{Mechaniske Aarsager.}   (printed pp. 61--125)
\kapitel{Chapter Two.}{Mechanical Causes.}

% Danish: \afdeling{A.}{Mechanisk Idealitet.}   (printed pp. 62--86)
\afdeling{A.}{Mechanical Ideality.}

That natural necessity unites in itself the arbitrary and the accidental as cancelled
determinations has its ground in this, that the mechanical cause is at once to be
apprehended both as concept and as object. As object the cause is real, as concept it is
ideal. To make the ideality of the cause recognizable, we shall now show in particular how
matters stand in mechanics with
% Letterspaced setting: „Kraft, Betingelse“ and „Aarsag“ are letterspaced, but the
% intervening „og“ is not; hence two emphases and not one.
\emph{force, condition} and \emph{cause}, and upon what it essentially rests that the unity
of the cause masters the outward manifold and takes up into itself a plurality of
\emph{heterogeneous moments}.

\runin{a) Force and cause.} When mechanics makes ``what we call force'' an ``unknown''
cause, it thereby indicates both a difference and a unity of force and cause. The
difference is that the unknown lies in the force, not in the cause. That the concept of
cause has validity, that every motion of matter must necessarily have its cause, cannot be
doubted; in the contrary case rational mechanics would have to cease being a science. But
one is in doubt about what the active element in the cause really is, how
% --- p. 63 ---
it is best designated, and so makes a kind of apology by saying: ``we call it force.'' It
has been shown in our foregoing that this apology is, from the standpoint of essence,
quite superfluous, that the concept force is in its kind just as trustworthy a concept of
nature as the concept cause. And yet the apology has an acceptable ground when we look to
actuality. Precisely by the cause being determined as force, the concept goes out beyond
itself; it is not exchanged for another concept, but is set over into activity and
expressed in fact. The active element is a concept and yet more than a concept; it is
force. It is precisely through force that the mechanical concept of cause passes over into
actuality and becomes a working cause.

How conceiving and working stand to one another in this respect can be found out by seeking
the origin of the peculiar designations in which mechanics --- though it is otherwise
indifferent to the question of the forces' ``physical origin'' --- holds fast a qualitative
difference among the forces themselves, between accelerating force, moving force and living
force. It is in actuality the same force --- the active element is force, and force the
active element --- and yet, on account of the determinations of content with which the
concept of cause is enriched in a progressive development, qualitatively different forces;
it is in essence the same cause, and yet, on account of the identity of the active forces,
different causes.

% The Greek letter is read on the KB copy: α, not a. The working copy's pencil
% underlinings under „instantan“ and „accelererende“ are an earlier reader's marks;
% on the KB copy both words stand in ordinary antiqua.
\subrunin{α) Accelerating force.} As the cause of motion, force can either work suddenly,
e.g. in impact, and then gets the name of instantaneous force, or continuously alter the
velocity and is then called --- positively or negatively --- accelerating. In the concept
itself the difference is vanishing. The instantaneous force requires a finite time, however
small, to effect a finite velocity, and is so far
% --- p. 64 ---
essentially accelerating; conversely, the accelerating force can be resolved into
infinitely many impacts and is in each impact to be regarded as instantaneous. Essential
for the concept of cause is, by contrast, the peculiar circumstance that force in and for
itself as force is not the cause of the acceleration, that it is so only by means of
inertia. If we think inertia away, then force would, in order to accelerate or retard a
motion continuously, itself have either to be increased or diminished continuously; but in
that way the cause of the acceleration would come to lie not in the force, but in what
caused the alteration in the force. If the cause of the alteration of force be again
regarded as a force, then this force must again have a cause of alteration, and so on to
infinity.

% MATHEMATICAL SETTING, pp. 64--70. Per the house rule, single standing letters in
% the running text are set as ordinary text (k, g, t, v, M, m, P, A, B), while
% equations and compound expressions are set in math; where a single letter is put
% directly against a compound expression in the same clause (g and mg, m and mV),
% both are set in math, so that the face does not change in mid-contrast.
Only by working in unity with inertia is force accelerating. In acceleration we have, then,
the sought unity of reflecting and working, a categorical expression for the mechanical
ideality of the cause. Analysis shows exactly how matters stand with this unity. A constant
force k works upon a material point; in the relative instant dt, $kdt = dv$. The material
point receives in each instant an infinitely small but constant increment of velocity; but
what brings it about that the velocity, which in the first instant is only $dv$, becomes in
the next instant $2dv$, in the next again $3dv$ and so forth, so that in the
$n^{\text{th}}$ instant it has grown to $ndv = v$, for $n = \infty$, is precisely inertia.
It is inertia that takes up, and it is inertia that preserves every increment of velocity,
and thereby makes it possible for the velocity to grow with the time, to be a function of
the time: $v = f(t) = kt$. From this it is seen that the passive is just as much
determining for the cause as the active. The active is passive, for the force's action upon
its point of application is from first to last dependent upon the inertia that receives and
preserves the action; the passive is active, for inertia is on its side just as
co-operative in the acceleration as force.
% --- p. 65 ---
So pervasive is the reflexion: the accelerating cause is the reflected, active unity of the
active and the passive.

In the form of reflexion the cause is concept; but the reflexion here spoken of is the
facts' own determining reflexion, a reflexion in power. Every reflex in logical conceiving
has on the side of actuality its antilogically existential moment in the conceiving power;
infinite reflecting has activity for its actual expression: point for point this reflecting
goes together into one with a corresponding working. In nature itself the cause is concept
and yet more than concept. The concept does not work; it is the force reflected in inertia,
the physical cause, that works. But the physical cause is determined to activity by its
concept; everything mechanics can elucidate about force and inertia must in every respect
borrow its validity from the concept of cause. If we take away the active element in the
cause, then there is only a logical concept, no actual cause; but if we take away the
concept of cause, then force, inertia, the active element, and whatever else one will call
it, are only idle representations; there is then no cause at all, either actual or unactual.

% The Greek letter is read on the page itself: β, not B or ß.
\subrunin{β) Moving force.} Since it is precisely the moving cause that is called force, it
might at first glance seem superfluous to introduce a category of its own and set moving
force alongside accelerating force. The accelerating force is moving, and the moving force
accelerating; upon what, then, does the difference rest? Upon a particular development of
the concept of cause. In determining acceleration and the accelerating cause, the material
element in the matter was tacitly understood, but not expressly brought out. True enough,
force issues from matter, inertia is in matter, the reflexion of inertia and force has
matter for its foundation; but the material moment first comes into its right when the
concept of mass enters into the concept
% --- p. 66 ---
of cause, and the quantity of matter becomes decisive for the cause. If we choose for
example gravity at the earth's surface as accelerating force g, then, as regards the
acceleration, it is altogether indifferent whether g works upon a small or upon a large
mass; every smallest part of the body will --- in empty space --- fall with the same
velocity, whether the parts be connected or separated: $gdm$ and $\int_{0}^{m} gdm = mg$
give in this respect one and the same result.

It is otherwise with the moving force, which is cause on account of the mass. If the mass M
is greater than the mass m, then the moving force in $Mg$ has undeniably a greater strength
than in $mg$; but while between $Mg$ and $mg$ there is only a quantitative difference,
since they both belong under the same concept of cause, there is by contrast a qualitative
difference between $g$ and $mg$, for here the concept of cause is different. For the moving
force as cause it is not, strictly speaking, necessary that it cause an actual motion. The
cause is indeed not without effect, but the effect may be cancelled by a counter-effect, a
pressure, and the cause has its reality none the less. So it is with weight. Weight is a
product of mass and acceleration without preceding motion: $Mg = P$. The weight P is a
reflected unity of the extensive in the mass and the intensive in gravity, an immediate and
yet reflected unity of reality and concept; that is weight's mechanical ideality as cause.
Mass in its immediate being is no concept; it may easily appear to a superficial observer
as though mass might well be cause without entering into the concept of cause, but that is
a contradiction. From weight it is seen that the concept of mass is an essential concept of
nature, for from weight it is seen that mass is in itself only what it is, in its reflexion
with gravity, for another. If gravity alters its intensity, then the mass gets another
weight and, as the concept uttered in the equation $M = \frac{P}{g}$
% --- p. 67 ---
shows, another reality. The mechanical equations are in this respect more than formulas
with an equals sign, they are judgments. The equation $M = \frac{P}{g}$ can indeed, in the
abstract quantitative sense, be resolved into a tautology: $M = \frac{Mg}{g} = M$, but,
apprehended as a judgment, it means something quite other. That $Mg$ is expressed by $P$
means, namely, that the factors M and g must not be held fast in their product as
heterogeneous units, but apprehended as moments in the real concept P, the concrete concept
that gives existing mass its peculiar value as working cause.

As cause, the moving force is a mass-cause. If now the resistance to the effect be
cancelled, then motion sets in, time gets its significance, and the moving force displays
itself in the product of mass and velocity, i.e., in the \emph{quantity of motion}: again
an alteration in the concept of cause. If we call the moving force $mk$, the time $t$ and
the velocity $v$, then the equation $mkt = mv$, and so $mk = \frac{mv}{t}$, shows that any
constant force whatever can be measured by the quantity of motion it brings forth in a unit
of time. That the moving force $mk$, working in the time $t$, brings forth $mv$, means then
that $mkt$ is cause, $mv$ effect. Insofar as the quantity of motion gained is expressly
posited as effect, the working cause has taken the determination of time up into its
concept and has become a function of time.

Here now arises the question: is the concept of working cause completely expressed in that
concept of time and force which has the quantity of motion for its measure? Cartesius
answers yes, but Leibnitz no. To appraise a force, one must know exactly what it can
accomplish. It is, says Leibnitz, not enough that the effect be proportional to the cause;
the complete effect must be
% --- p. 68 ---
perfectly equal to its cause. To lift a body A of 1 pound to a height of 4 ells, the same
force is required as to lift a body B of 4 pounds to a height of 1 ell; it is precisely the
force the two bodies attain by falling, each from its own height. But the velocity A
attains by the fall stands to that B attains as 2 to 1.
% The multiplication sign is printed as a cross; set with \times.
The quantity of motion in A will then be $1 \times 2$, while that in B is $4 \times 1$. The
quantity of motion in A is only half as great as the quantity of motion in B, although the
force is the same; consequently force cannot be measured by the quantity of motion. Abbé de
Conti objected that the law set up by Cartesius required regard to simultaneous forces, or
to motions which two bodies attain in the same time; but --- asks Leibnitz --- ``when one
sees a body of a determinate magnitude move with a determinate velocity, should one not be
able to appraise its force without knowing how long it has taken to attain this velocity?
And when two equal bodies have the same velocity, but the one has obtained it by a sudden
impact, the other by falling for a longer time, will one then say that their forces are
different? That would be the same as to say that a man was richer because he had taken
longer to win his money.''

% The Greek letter is read on the page itself: γ, not y or 7.
\subrunin{γ) Living force.} The ground why the quantity of motion cannot be the completely
uttered effect of the moving force, of the moving cause, lies in this, that the opposition
of cause and effect is as yet unactual. The equation $mkt = mV = Mv$ states that the moving
force $mk$ can in the time $t$ awaken a quantity of motion in $m$ which is $mV$, and in $M$
which is $Mv$. If we now ask how great the cause is, we are referred to the effect. If we
ask how great, then, the effect is, we are referred to the cause. The cause is equal to the
effect; that is a merely formal judgment, a tautological proposition, so long as it is not
yet elucidated
% --- p. 69 ---
wherein this greatness actually consists. The question is whether the moving cause in the
form of $mV$ can really accomplish the same as in the form of $Mv$; for if it cannot, then
by assuming two different forms it has come under two different concepts and has thus
become two different causes.

To test the working cause it is necessary to see it in activity. The activity cannot then
be stopped by an immediate resistance; the resistance must be overcome, and the activity
each time determined by the manner in which it overcomes its resistance. To this
corresponds the distinction Leibnitz has introduced between
% The letterspacing verified on the KB copy: „døde“ is letterspaced, but the
% preceding „den“ is not, while the following „den levende“ is letterspaced in its
% entirety. Reproduced as printed.
\emph{dead} and \emph{the living} force. Among dead, elementary forces he reckons those
whose effect is immediately cancelled by a counter-effect; the living force, by contrast,
is ``always connected with actual motion.'' During the motion the resistance gives way,
while the force is gradually expended.

If $mkt = mV$ be posited, the activity in the first element of time will be $mkdt = mdV$.
What is now accomplished in the running time? The force moves the burden forward in space,
but in dragging the burden the velocity decreases, and the decreasing velocity is again an
expression of the decreasing force. How long the force can hold out must then be seen from
the length of way through which it avails to carry its burden. The equation $mdV = mkdt$
must accordingly be transformed into $mdV = mk\frac{ds}{V}$,
% MISPRINT: the Danish has „kvad der udrettes i Tiden“ for „hvad“. The reading is
% verified on the KB copy, which has the same letter; the fault is the compositor's
% and not this copy's. The errata leaf does not correct it. Translated by the sense.
so that it is known from the space what is accomplished in the time. The old saying that
time is equally long but not equally useful here finds a peculiar confirmation; dt is
constant, but V varies, from which it follows that $ds$ also
% The rows of dots are printed with three points after ds_2 and four after V_2; the
% difference is reproduced. The points stand SOLID, without spacing between them ---
% closer than anywhere else in the book. Set in the house rule's close form.
varies: $ds > ds_{1} > ds_{2} .\,.\,.$, since $V > V_{1} > V_{2}
.\,.\,.\,.$. From $mdV = mk\frac{ds}{V}$ follows $mVdV = mkds$; by integrating this equation
we obtain $mV^{2} = 2mks$ and with it the real opposition of cause and effect
% THE PAGE ENDS HERE, without a full stop after „og Virkning“ and with about a third
% of the type page standing empty; p. 70 begins with a new paragraph. Both copies
% agree at this point, so the lack is the printing's and not this copy's; in all
% probability a displayed equation has dropped out of the setting. The errata leaf
% does not correct it. Reproduced as printed.
% --- p. 70 ---

But when to the quantity of motion $mV$ there answers a living force $mV^{2}$, then to the
quantity of motion $Mv$ there answers a living force $Mv^{2}$. Since now $mV^{2}$, by the
presupposition $mV = Mv$, has another value than $Mv^{2}$, it is seen from this that the
force in $mV$, when it is set to work, accomplishes something other than the force in $Mv$.
The living force, which is measured by the work, is accordingly qualitatively different from
the force moving in time, which is measured by the quantity of motion. Here there is no leap,
either in nature or in thought; it is force, the active element in the cause, that in the
course of the development of the concept takes up and gathers step by step the one-sidedly
emerging determinations and thereby reaches its complete expression in the form of living
force.

\runin{b) Condition and cause.} When the determinations of existence belonging to the concept
of cause are held fast in representation as existing, each for itself, they become
conditioning, and the working cause resolves itself into a sum of conditions. But it lies in
the nature of the condition to point out beyond itself to a conditioned and to other
conditions; as the conditions reflect themselves in one another mutually, they lose their
apparent independence, and the unity of cause, undivided in itself, comes into view again as
the unconditioned. This unconditioned, which accordingly does not fall outside its conditions
but goes together into one with them, is the conditioned cause.

% The Greek letter is read on the page itself: α.
\subrunin{α) The condition.} The cause is a concept, but a concept of actuality, and a
concept of actuality always has reflexes of actual things in its content. The reflexes of
things necessary for actuality are, on causality's behalf, \emph{the condition} and \emph{the
conditioned}. Not even force can in this connexion be held fast against the conditions, for
force for itself can merely be represented; what is represented, ``we call force,'' sinks to
being a condition. The condition is a presupposition which the representing consciousness
% --- p. 71 ---
in the train of thought partly puts in, partly finds in actuality, a demand for an
independent which for the dependent's sake must be satisfied. The condition, the independent,
encloses the conditioned, the dependent, and yet they stand over against one another as
member against member. The reflexion of condition is a reflexion of represented members, an
outward reflexion hovering between essence and existence, concept and actuality,
characteristic of mechanical activity. Thus force and inertia are conditions for the
acceleration; they are represented each for itself, and yet point incessantly to one another.
Mass and acceleration are again conditions for weight; the same mass can be subjected to
another acceleration, and the same acceleration applied to another mass, and so far the
conditions must be separated; and yet they cannot be separated, for the altered product of
mass and acceleration alters the weight. If time be drawn into the circle of conditions, then
weight comes into motion, shows itself as moving force, and comes to be determined as
quantity of motion. If, finally, the condition be added that the given quantity of motion
shall be applied to overcoming a resistance successively, and be measured by the space the
resistance traverses, then we have the living force, though not as cause; for now the cause is
dissolved, it is altered into a sum of conditions.

From the categorical peculiarity of the condition it follows that several particular
conditions can be enclosed in one fundamental condition, and this again resolved into a
quantity of single conditions. Rational mechanics is laid out so as to show how resolution and
composition happen systematically. For forces acting upon a point to be in equilibrium, their
resultant must be zero; this is the universal condition of equilibrium, its fundamental
condition. But what a swarm of particular conditions and varying single conditions can enter
into this!
% The working scan's leaf has two black blots on this page; the KB copy is clean at
% both places, so it is ink or dirt in that copy and no compositor's fault. Nothing
% reproduced.

% --- p. 72 ---
% NOTATION: the letters P, Q, R, A, B, C, O, a, b, c are on this page members of the
% equations set up, and are therefore set throughout in math, so that the face does
% not change in mid-series (the house rule from batch 5).
If two forces act in parallel directions toward the same side upon two points connected by an
inflexible straight line, then one can determine the magnitude and position of their
resultant, and thereby of a force opposed to it which holds them in equilibrium. One sees at
once how the condition set up comprehends several conditions: that there are two forces, that
they are parallel, connected, etc. If the forces be called $P$ and $Q$, the resultant $R$,
then we have the conditioned: $R = P + Q$. If it be further made a condition that $A$, the
point of application of $P$, and $B$, the point of application of $Q$, shall have from an
arbitrary point $O$ in the prolongation of the connecting line the distances $a$ and $b$, then
the distance between the arbitrary point $O$ and the resultant's point of application $C$, a
distance we may denote by $c$, is expressly conditioned thereby. The composition of the
conditioned, the relation between $Rc$, $Pa$ and $Qb$, answers exactly to the composite
condition, and this conditioned can by a turn in the reflexion again be set up as a condition.
For since $Rc = Pa + Qb$, the equilibrium between the three forces $P$, $Q$ and $R$ is
precisely conditioned by $R$ being applied at the point $C$ and in the direction opposite to
$P$ and $Q$.

Progress in analysis happens on the whole in such a way that what in a preceding connexion was
the conditioned passes over in a following one into being a condition, and that new ground is
continually won by partly generalizing, partly specifying a given fundamental condition. To
% The minus sign is here a true minus, not Nielsen's ÷ (verified on the KB copy under
% magnification); ÷ stands in this book only where he expressly prints it, cf. p. 46.
the above circle of conditions belong further the equations $Q = R - P$,
$b = \dfrac{Rc - Pa}{Q}$; if the conditions now be specified so that $R = P$, but not
$c = a$, then $Q = 0$, and $b = \infty$. From this it is seen that two equal parallel forces
acting in opposite directions ($R$ and $P$) cannot possibly be held in equilibrium by a single
force ($Q$), for an infinitely small force ($Q = 0$) at an infinitely great distance
($b = \infty$) are conditions that cannot be brought about. Thus we are led
% --- p. 73 ---
by the way of the conditions from the doctrine of parallel forces to the doctrine of couples
% „Svingkræfter“ in the Danish --- the forces of a couple, the case just reached, where
% two equal opposed parallel forces admit of no single balancing force.
--- an example in little which shows what it means in the large that the way of the conditions
is the way of mechanics, that exact science with its equations of condition and its final
equations has ground to resolve the cause into conditions. But philosophy of nature is not
mechanics; it goes the reverse way and resolves the conditions into

% The Greek letter is read on the page itself (KB copy): β, not „B“, „ß“ or „P“, which
% is what the OCR gives. The head is letterspaced and run-in: the sentence above ends
% without a full stop and is continued in the words of the head („resolves the
% conditions into β) the cause“). The letters A, B and C below are loose letters in the
% body text on a page without equations, and therefore stand as ordinary text.
\subrunin{β) The cause.} The conditions are not to be effaced, and yet they are to disappear.
What is characteristic of the condition is that in its represented independence it is so
dependent. The hypothetical judgments --- if A is, then B is; if an aberration of the fixed
stars takes place, the earth must move about the sun; if B is, then C is; if the earth moves,
the fixed stars must have parallax --- show us the dependence of the self-subsistent, the
dependence of the independent, when the judgments cited furnish premisses for
corresponding inferences. One infers: A is, therefore B is; the fixed stars display an
aberration, therefore the earth must move, although we cannot immediately sense it; B is,
therefore C is; the earth moves, therefore the fixed stars must have parallax, even if on
account of its smallness we are not able to measure it. In both cases the inference is by
\textit{modus ponens}, and the inference is a proof of the conditioned's dependence upon its
condition. But in actuality the relation can also be reversed; it is precisely the condition
that is dependent upon its conditioned --- which is expressly seen when we infer by
\textit{modus tollens}: B is not, therefore A is not; C is not, therefore B is not. Which of
the two modes of inference has actual validity in a given case must of course depend upon
facts and be determined by observations; but that both have logical validity proves that the
condition, dependent as it is, stands and falls with its conditioned.

Where a fundamental condition is resolved into particular conditions,
% --- p. 74 ---
the same dependence shows itself. Motion in a straight line has another fundamental condition
than motion in a curved line. If we analyse the latter, then the condition for the straight
line will repeat itself here and deposit single conditions, for motion along the tangent gives
a straight line, and so does motion in a direction forming an angle with the tangent. If a
fixed centre be introduced as a particular condition, then the position of the two directions
--- along the tangent and toward the centre --- will likewise come under particular conditions
and be again conditioning. If the one direction is constantly perpendicular to the other, the
curve will be a circle and the velocity constant; if, by contrast, the position varies so that
the two directions form different angles, now a right, now an acute, now an obtuse one, the
velocity will vary, and the curve cannot be a circle.

It will now be possible to see from this how the conditions cancel one another and disappear in
an unconditioned. None of the single conditions cited is particularly fulfilled; the one
incessantly prevents the fulfilment of the other. Motion along the tangent is at each instant
prevented by the motion toward the centre, and this again by the former. Motion along the
tangent disappears into two conditioned results; the one is a flight from the centre --- for
representation the so-called centrifugal force --- which is cancelled by a corresponding
striving, the normal force; the other is the advance in the orbit, the only result that is not
cancelled.

The unity thus conditioned, the motion in the orbit, now depends upon a corresponding unity of
the conditions which does not let itself be resolved, since it contains the conditions not as a
sum but as cancelled conditions. This actual unity of the cancelled conditions is the
unconditioned, the matter in itself, the original matter, \emph{the cause}.
% The Danish turns on a word that English cannot keep: „Sagen i sig, den oprindelige
% Sag, Aarsagen“ --- Sag (matter, case) and Aar-sag (cause, literally ur-matter). The
% etymology is Nielsen's argument here, not an ornament.
But when the unconditioned is determined as cause, then it is thereby
% --- p. 75 ---
more than a conditioned, it is something caused, an effect.

That there is another fundamental relation between cause and effect than between condition
and conditioned is easy to see. If the conditioned is not, then neither is the condition,
for the conditioned is a content of the condition, represented as a member for itself; the
condition with its conditioned belongs solely to outward reflexion. But the cause embraces
actuality, and thereby both inward and outward reflexion. The cause has a being in
possibility, in inward reflexion, even when it does not work, just as force has a being as
an inner, even when it does not utter itself. A condition must supervene for a force to be
able to utter itself, for force and condition are coordinate categories. But the cause is
above them both, for it is in them both. The cause does not scrape by, borrowing here a
force, there a condition; the cause itself owns the forces, it masters the conditions; its
essence is power, the unconditioned of the conditions.
% „hutler“ (from „at hutle sig igjennem“, to scrape by) stands so in both copies. In
% the working scan the word is underlined in pencil by an earlier reader; the pencil is
% not text.

% The Greek letter is read on the page itself (KB copy): γ, not „y“ or „7“, which is
% what the OCR gives. The head is letterspaced and run-in.
\subrunin{γ) The conditioned cause.} As the pervasive unity of the conditions the cause is an
unconditioned; but since this unity is itself determined by the conditions taken up into it,
the cause determined in itself is conditioned. If one reflects upon the difference between
cause and condition, then the same cause can, by taking up different conditions, have
different effects; if one reflects upon the unity of condition and cause, then the
differently conditioned cause becomes different causes, and each effect answers exactly to
its cause. That this double-sided reflexion is by no means idle, but on the contrary decisive
for mechanical ideality, we learn from rational mechanics.

The cause of rest is the equilibrium of forces; the cause of motion is the cancelling of the
equilibrium: it is at bottom the same cause and yet in actuality a different cause. If we
% --- p. 76 ---
reflect upon the difference, then statics and dynamics have each its own circle of problems
and fall so far apart from one another; if we reflect upon the unity, then the static problem
becomes dynamic, and the dynamic static. The reciprocal determination is grounded on two
corresponding principles, the principle of virtual velocities and D'Alembert's principle of
the equilibrium of the lost forces.

% „Tale“ is set with a loose space before the last letter („Tal e“), but the letters
% are not letterspaced among themselves. It is therefore loose setting and not emphasis;
% set as one word without \emph. Verified on the KB copy.
Everywhere that there is talk of a principle, one must necessarily, however the matter is
otherwise to be understood, think of an unconditioned. The general equation of condition for
equilibrium makes in this respect no exception; for it does indeed receive such an extension
that it enters into all particular conditions and finds application to any system whatever of
material points; but precisely the manner in which the particular conditions go together into
unity gives the equation the stamp of an unconditioned. The fundamental thought uttered in
the characteristic equation of equilibrium, the formula of virtual
% The note's mark is printed *) and stands after the comma. The note begins on p. 76
% and is continued without a new mark at the top of the note field on p. 77
% („Totalvirkningen af alle Kræfterne ...“); per the house rule the whole note is set
% here at its mark.
velocities,\footnote{% Misprint in the note: „af Punktet. som kan tænkes“ --- a full
% stop where the connexion requires a comma. Both copies show the full stop, so the
% fault is the compositor's and not this copy's; the book's errata leaf does not correct
% it. Translated by the sense.
``By a point's virtual velocity is understood any infinitely small displacement of the point
that can be thought to take place, without regard to the acting forces, solely in accordance
with the conditions to which the point is subjected.'' If the force be called P, the
projection of the virtual velocity upon the direction of the force $\delta p$, and the
reciprocal determination of the universal and the single --- the logical and the antilogical
--- be indicated by an appended index k, then the equation of condition for the equilibrium
of a single point is $\Sigma\, P_{k}\, \delta p_{k} = 0$. The forces are reckoned positive,
but the projections positive or negative, according as they fall either along the direction
of the force or along its opposite prolongation. The formula holding for a system of points,
in which all conditions are taken up, can be expressed by
\[
P_{1}\, \delta p_{1} + P_{2}\, \delta p_{2} \ .\ .\ .\ .\ + P_{n}\,
\delta p_{n} + \lambda_{1}\, \delta u_{1} + \lambda_{2}\, \delta u_{2}
\ .\ .\ .\ \lambda_{r}\, \delta u_{r} = 0.
\]
The terms $P_{k}\, \delta p_{k}$ are called the virtual moments of the forces, the terms
$\lambda_{s}\, \delta u_{s}$ the virtual moments of the resistance, and the formula presents
the principle. The sum of the virtual moments
$P_{1}\, \delta p_{1} + P_{2}\, \delta p_{2} + P_{3}\, \delta p_{3} +
\ .\ .\ .$ indicates at every instant
% Here the note passes over onto p. 77.
the total effect of all the forces. That this sum must be zero at the instant the position of
equilibrium is passed is a straightforward consequence of the concept of equilibrium, since
the forces then neither promote nor counteract the motion.}
% --- p. 77 ---
--- a fundamental thought which comprehends all conditions in the conditioned unity of the
determinate cause, the unconditioned --- can with full right be called a principle.

According to this principle equilibrium is, to be sure, tested upon a motion which is
nevertheless no motion, for a virtual velocity is a possibility that does not become
actuality. But this possibility, the merely represented displacement of the points, serves
meanwhile to let motion shine through in the equilibrium, to let the equilibrium show itself
as cancelled motion, and to make the static problem dynamic.

In the formula for the principle of virtual velocities one has a sum of virtual moments in
which each term is a product of the force $P_{k}$ and the virtual velocity $\delta p_{k}$
projected upon the direction of the force $p_{k}$, and so $P_{1}\, \delta p_{1} + P_{2}\,
\delta p_{2} + P_{3}\, \delta p_{3} \ .\ .\ .$ What is peculiar about this sum is that it
indicates at every instant the total effect of all the forces, however different its value
may be. What for statics is therefore the universal --- namely that the sum of the virtual
moments shall expressly be zero --- can then, with respect to dynamics, be regarded as a
special case, a conditioned cause which goes together with another condition of cause and
forms another, more comprehensive unconditioned, a principle which at once determines the
cause both to equilibrium and to motion --- D'Alembert's
% This note too is marked *) --- the page therefore carries the end of the note from
% p. 76 (without a mark) and then this, p. 77's own note. The note is continued without
% a new mark through the whole note field on p. 78 and is per the house rule set here at
% its mark. The second differential quotient is printed with the exponent 2 standing
% OVER the letter (d with 2 above x); that manner of setting is not reproduced, but the
% term is set as d^{2}x_{k}/dt^{2}.
principle.\footnote{Let a system of points $m_{k}$ be determined by rectangular coordinates
and acted upon by the forces $P_{k}$, whose components along the axes are $X_{k}$, $Y_{k}$,
$Z_{k}$. The actual motion of $m_{k}$ in the time element $dt$ can be assumed to be the same
as if this point were free and acted upon by the final velocity $v_{k}$ and by an accelerating
force $Q_{k}$ which in magnitude
% Here the note passes over onto p. 78.
and direction is different from $P_{k}$. If now every point $m_{k}$ be assumed subjected,
besides the given force $P_{k}$, also to the force $Q_{k}$ and to another equal and opposite
to it, namely $Q'_{k}$, then the given system of forces is not thereby altered. The whole
system moves as though the points were free and subjected to the force $Q_{k}$ alone. That
part of the forces which goes to overcoming pressures and tensions in the system itself gets
no influence upon the motion and must in this respect be regarded as lost forces. The two
forces $P_{k}$ and $Q'_{k}$ acting on $m_{k}$ have a resultant $p_{k}$; this resultant is
\emph{the lost force}, which by composition with $Q_{k}$ gives $P_{k}$ as resultant. If now
the components of
\[
Q'_{k}\ \text{be denoted by}\ -\frac{d^{2}x_{k}}{dt^{2}},\
-\frac{d^{2}y_{k}}{dt^{2}},\ -\frac{d^{2}z_{k}}{dt^{2}},
\]
$u_{1} = c_{1}$, $u_{2} = c_{2} \ .\ .\ .\ .$ indicating equations of condition between the
coordinates, then one has the characteristic equation:
\[
\begin{aligned}
\Sigma\, m_{k} \biggl[ \biggl( X_{k} &- \frac{d^{2}x_{k}}{dt^{2}}
   \biggr)\, \delta x_{k}
   + \biggl( Y_{k} - \frac{d^{2}y_{k}}{dt^{2}} \biggr)\, \delta y_{k}\\
   &+ \biggl( Z_{k} - \frac{d^{2}z_{k}}{dt^{2}} \biggr)\, \delta z_{k}
   \biggr] + \lambda_{1}\, \delta u_{1} + \lambda_{2}\, \delta u_{2}
   \ .\ .\ .\ .\ \lambda_{r}\, \delta u_{r} = 0,
\end{aligned}
\]
which shows how the problems of dynamics are reduced to problems of statics. Cf.\ C.\ Ramus,
\textit{Anal.\ Mech.}, p.\ 210.}
% --- p. 78 ---
% NOTATION: the letters of cause and force on pp. 78--79 (m_k, P_k, Q_k, Q'_k, p_k and
% the causes A, A', A'' with the effects B and C) belong to one and the same series; per
% the house rule from batch 5 the whole series is set in math, so that the face does not
% change in the middle of it.
Let a system of material points $m_{k}$, acted upon by the forces $P_{k}$, be given. Since now
the forces with their conditions fall together, according to our foregoing, in the cause, we
have then in this system one single conditioned cause, the cause $A$. But analysis shows that
a part of the forces $P_{k}$ is cancelled by the reactions lying in the system, while another
part is applied exclusively to the motion. The cause $A$ accordingly yields two mutually
different effects, of which the one, $B$, is seen in the system's maintenance, the other, $C$,
in its motion. Here is the same cause and yet two different causes. If we reflect upon
$P_{k}$, then it is the same cause; for it is $P_{k}$ that is decomposed, so that a
% --- p. 79 ---
% The errata leaf corrects p. 79, line 1 from the top: „Q_k læs Q'_k“ --- and the leaf
% prints, in this entry alone, no dash before „læs“. The correction concerns the SECOND
% $Q_{k}$ in the line, the one that is „equal and opposite to $Q_{k}$“ and forms with
% $P_{k}$ the resultant $p_{k}$; cf. the note to pp. 77--78, where $Q'_{k}$ stands in
% full. Verified on the KB copy under magnification: $Q_{k}$ does indeed stand there
% without the prime. Per the house rule the text stands here as printed and the
% correction is not introduced.
part of $P_{k}$, namely $Q_{k}$, yields the motion, while $Q_{k}$, which is equal and opposite
to $Q_{k}$, forms with $P_{k}$ a resultant $p_{k}$ that is lost to the motion; $P_{k}$ is
accordingly decomposed into $p_{k}$ and $Q_{k}$. If we reflect upon the opposition between
$p_{k}$ and $Q_{k}$, then we have two different causes, namely in $p_{k}$ the cause of the
system's maintenance ($A'$) and in $Q_{k}$ the cause of the motion ($A''$). How intimately the
causes thus separated engage in one another is seen from this, that an alteration in the one
brings with it point for point an alteration in the other. During the system's motion $Q_{k}$
varies, and consequently $p_{k}$ too. The sum of the virtual moments of the lost forces and of
the resistances must of course be zero at every given instant; but as $p_{k}$ and $Q_{k}$ are
determined as functions of time, and with them the cause $A'$ answering to the effect $B$, so
$Q_{k}$, with the cause $A''$ answering to the effect $C$, is determined at every instant, and
conversely; the problem is at once both static and dynamic, because the cause conditioned by
the system is itself both a cause of equilibrium and a cause of motion.

% The run-in head c) is printed in SEMIBOLD antiqua, not letterspaced like the Greek
% sub-heads α) β) γ). Verified on the KB copy. The house macro \runin serves all the
% same; the difference of face is not reproduced.
\runin{c) The unity of the moments of cause.} When in a mechanical phenomenon one is expressly
to look to what is immediately active, one names force and, insofar as this is in matter, also
mass. If, by contrast, one will particularly bring out everything which, without itself being
active in and for itself, gets an influence upon the activity and thereby becomes a
co-operating element, then one speaks of conditions. If now, as we have seen, the cause has
all forces and conditions in its power, then it is intelligible from this that the mechanical
cause takes up and unites highly heterogeneous moments. In its essence the cause is concept;
in this lies that all moments of cause drawn from actuality, whether they be otherwise called
forces or conditions, must be \emph{ideal}. But in its working the cause is existent; in this
lies that the
% --- p. 80 ---
active moments of cause must all be \emph{real}. In the naturally determined fundamental
relation between the ideal and real moments of the mechanical cause the unity-in-difference of
ideal and real cause must accordingly be clearly uttered.

% The Greek sub-head is run-in and letterspaced; the letter is read on the KB copy
% (α, not a).
\subrunin{α) The ideality of the moments.} Just as forces and conditions merge into the causes
thereby determined, so the causes merge into the existing bodies; bodies are from this
standpoint to be regarded as the only causes working in space and time: to decompose bodies
into corporeal elements is to decompose causes into moments of cause.

It is instructive to see how rational mechanics, though its gaze is fixed most immediately
upon forces, conditions and laws, precisely for the laws' sake transforms corporeal causes
into concepts of cause and, without rightly knowing it, sets a whole world of mechanical
causes forth in the light of ideality as a system of embodied concepts of cause.

To penetrate properly into the corporeal relations of cause, the mathematician begins by
resolving the body into its elements: $M = \int dm$. But he does not stop at the infinitely
empty explanation that every body consists of infinitely many small parts; he takes note of
the manner of limitation; he takes aim at the concept of corporeality and is thereby put in a
position to determine punctually both a body's inner structure and its outer form, its figure.
% The series of symbols ϱ, V, dm, M is kept in one face (the house rule from batch 5),
% although V and M stand as upright letters in the printing.
If $\varrho$ be the density, $V$ the volume, then $dm = \varrho\,dv$ and $M = \int
\varrho\,dv$; if the body is of unequal density, and so $\varrho$ variable, then it is a
matter for him of determining the $\varrho$ answering to every element of volume as a function
of that element's coordinates. But how does he get the function determined? By presupposing
determinate corporeal relations in the body in question, and so by starting from the concept
of the represented body. He does not have in advance a given
% --- p. 81 ---
body to which he adjusts the concept --- the representation of this is from this standpoint
only a disturbing fancy; no, he has by means of the intellectual intuition the body completely
in his power; he determines it in virtue of his concept and derives the law from the concept.
Whether he says ``if the body is of unequal density'' or says ``a body which is to be of
unequal density'' is, with respect to the concept of body, one and the same. It is precisely
what is rational in rational mechanics that its functions and laws do not lean upon empirical
inductions, but are developed out of concepts. The corporeal elements are not objects of
experience; they are reflexes of representation which the conceiving understanding receives
from intuition, ideal moments in images of corporeal reality.
% „ideelle“: both OCR witnesses read here „ideella“/„ideell&“; the KB copy shows plainly
% „ideelle“.

In indeterminate universality the one $dm$ is like the other, an element that can again be
resolved into elements; but in virtue of the concept the one $dm$ can become altogether
heterogeneous with the other. An element of a sphere and an element of a parallelepiped can, as
regards the infinite vanishing in infinite smallness, certainly pass for homogeneous; whether
the infinitely small lines of limitation be apprehended as elements of straight lines,
$dxdydz$, answering to a rectangular coordinate system, or whether by the use of polar
coordinates curvilinear elements be taken up into the limitation,
% The element is here printed with θ, ω, θ; on p. 84 Nielsen writes the same element with
% α, α, θ. The inconsistency is his own and stands in both copies; reproduced as printed.
$r^{2}\sin\theta\,d\omega\,d\theta\,dr$, makes in this respect no difference. But when the
finite is to be brought forth out of the infinite, when an integration, $\int\int\int dxdydz =
\int\int\int r^{2}\sin\theta\,d\omega\,d\theta\,dr$, is to be carried out, partly so that the
body becomes a parallelepiped, partly so that it becomes a sphere, then it is recognized that
the two bodies have heterogeneous elements: the moments of the concept are different because
the concepts themselves are different.

Just as every seed has in it the disposition to a plant of a determinate species, answering to
a determinate genus, so
% --- p. 82 ---
too every element of body determined by its concept has in it the disposition, under the
integration, to grow a peculiar growth and become a determinate body.
% The fractions stand in the printing as stacked fractions in mid-line; \tfrac follows the
% house rule from batch 4. The limit „r = o“ is set with the letter o, not with the
% numeral nought --- as in the note to p. 21; so reproduced.
That the mass of a spherical shell $dm = d\tfrac{4}{3}\pi r^{3} = 4\pi r^{2} dr$, when
integrated between the limits $r = o$ and $r = a$, must give a sphere $= \tfrac{4}{3}\pi
a^{3}$, is just as necessary as that the seed of a rose must, when it comes to growth and
development, bring forth a rose. The categories \emph{genus}, \emph{species} and
\emph{individual} have just as much validity in mechanics as in botany.

That a homogeneous sphere must, on the presupposition that attraction stands in inverse ratio
to the square of the distance, attract an external point as though the whole mass of the sphere
were gathered at the centre --- this actual relation of cause is not determined by the
observation of any outward fact, but solely by insight into the essence of the sphere and of
attraction.
% The working scan shows a small raised mark after „Afstanden,“ which tesseract read as a
% quotation mark. The KB copy has here a clean comma and nothing else: the mark is a blot
% in the one copy, not the type.
Even when empirical determinations enter into the calculation, it is still always a system of a
priori corporeal concepts of cause that furnishes the rational foundation. Physics has weighed
the planets, by empirical observations no doubt, but on the scales of ideality.

\subrunin{β) The reality of the moments.} One should now suppose that the analysing
mathematician would have to be fully assured of his concepts of body having physical validity,
since he is fully assured of the physical validity of the laws he finds; but no! About the
laws' validity for nature he harbours no doubt; they confirm themselves for him through all the
modifications which their fulfilment under particular conditions brings with it in actuality.
But with the relation between the concepts of body and actual bodies it is another matter; here
the mathematical physicist is inclined to regard the fundamental determinations necessary for
the analysis as arbitrary designations, subjective ways of representing, which have nothing to
do with bodies in nature,
% --- p. 83 ---
and consequently nothing to do with the corporeal causes themselves either. To rely on the laws
and yet doubt the concepts from which the laws are derived is assuredly a most singular
inconsistency; but whence does it come? From an unclear conception of the reality of the
moments.
% „høist“: the Google layer reads „heist“ (the familiar ø→e misreading); the KB copy shows
% the slash of the ø plainly.

An element such as $dm$ does not exist in nature: \textit{un infiniment petit est une grandeur
moindre que toute grandeur donnée de la même nature (Poisson)}; it is the problem of division
that here presses forward.\footnote{% printed mark: *)
% The title is letterspaced here; in the corresponding note to p. 30 it stands
% unletterspaced. Reproduced as printed on each page.
Cf.\ \emph{Mathematik og Dialektik}, Copenhagen 1859; the problem of division,
pp.\ 116--131.}
It goes with the element of body $dm$ as with the element of space $ds$ and with the element of
time $dt$; the mathematical elements do not exist and yet do exist. They do not exist as
particular subsisting parts, but they are moments in actual existence. The material body is just
as surely an integral of material $dm$ as the mathematical body is of mathematical $dm$; the
actual event is just as surely an integral of moments of event as the time in which it goes on
is an integral of moments of time. The question does not turn at all upon the relation between
the absolutely simple and what is composed of it, for the element is just as little in the real
as in the ideal sense anything absolutely simple. The fundamental question is: how can that
which is without magnitude bring forth magnitude? how can that which is itself without extension
bring forth something in actuality extended?

To derive the real from the ideal, Leibnitz introduced the monads and called them metaphysical
points:
% The compositor had no room for the last t in „soient“ and set it dropped at the end of
% the line. Both copies show the same, so it is the compositor's shift and not damage in
% the one copy; the word is reproduced whole.
\textit{il n'y a que les points metaphysiques ou de substance, qui soient exacts et réels, et
sans eux il n'y auroit rien de réel, puisque sans les véritables unités il n'y auroit point de
multitude.} But out of the monadic points, the soul-like unities, he was
% --- p. 84 ---
still not able to deduce corporeal actuality; as regards matter, it remained with ``obscure,
confused representations.'' In his exposition of ``the amphiboly of the concepts of reflexion''
Kant showed very rightly that there is a difference between matter as a concept of the
understanding built upon ``the power of representation'' and the natural phenomenon matter,
which is an object of sensation and whose apprehension is conditioned by the a priori form of
our intuition of sense. From the standpoint of intellectual philosophy matter is only a complex
of representations, a phenomenon of the representing being, a semblance of immediacy belonging
to the intellectual; for critical philosophy the chasm between the subjectively ideal and the
objectively real is so inconceivably great that matter can neither be brought forth by
representation nor have anything to do with ``the thing in itself.'' Both ways of seeing are
one-sided. The actual unity of the ideal and the real is a unity in opposition, a unity of
power. The ideal in $dm$ does not bring forth the material, any more than the reverse; but both
mutually opposed fundamental determinations are brought forth in and by the objective, universal
power which in its doubling is both ideal and real.

% The hyphen in „Ideel-reel“ really stands in the printed head (checked on the KB copy);
% cf. „den ideel-reelle Aarsag“ in the body text on p. 87.
\subrunin{γ) Ideal-real cause.} To determine still further how matters stand in the mechanical
sense with the concept of cause and the working cause, we choose, with a view to the foregoing,
a determinate example: the formation of a sphere.

By starting from the prismatic element
\[ r^{2}\sin\alpha\, d\alpha\, d\theta\, dr \]
% The lower limits of the integrals are set with the letter o, not with the numeral
% nought, like „r = o“ on p. 82; so reproduced.
and carrying out the integrations $\int_{o}^{\pi}\int_{o}^{2\pi}\int_{o}^{r}
r^{2}\sin\alpha\, d\alpha\, d\theta\, dr$, one sees how the one element arises out of the other,
and the whole sphere out of the system of the elements.\footnote{% printed mark: *)
Cf.\ \textit{Forelæsninger over Phil.\ Propæd.} 1861--62, pp.\ 31--33.}
The sphere's
% --- p. 85 ---
concept in unity with the productive intuition is the sphere's ideal cause. But what is thus
produced is likewise only something ideal, only the sphere's form, not the actual sphere. The
sum of conditions of form to which the integrations answer is arbitrarily chosen with particular
regard to known functions, and stands in no exclusively necessary connexion with that sum of
material constituents and outward conditions in which the formation of an actual sphere must
have its real cause. The moments of the concept and the elements of actuality fall into two
circles separated from one another: here is a yawning chasm between the ideal and the real.

This chasm is, in the matter itself, filled out by the divine concept of power, whose energy
embraces the ideal and the real evenly. But a human insight into the divine concept is attained
only on human terms; that is why the example is chosen.

From the ideal side the mechanical essence of the concept of cause can be made clear by means of
the required integrations. Out of pure nothing nothing can be formed. Every formation must begin
with something; the something with which the beginning is made is itself a beginning formation,
an element of formation. The element's possibility lies in conceiving intuiting; its definitive
determinateness is posited arbitrarily as the first presupposition of the formation: the
prismatic element with its three infinitely small dimensions is only a general condition of
formation. What can now be formed on such a condition? According to the formula set up, the
successive formations are particularly determined by the particular concept. If we integrate
first with respect to r, then an elementary pyramid comes forth, an element of one finite and
two infinitely small dimensions; if we integrate next with respect to $\theta$, then a conical
shell comes forth, an element with two finite and one infinitely small dimension; if we
integrate finally with respect to $\alpha$,
% --- p. 86 ---
we obtain the form sought, the spherical form with three finite dimensions. Thus, according to
the governing concept which is the ideal cause of the sphere's formation, the prismatic element
will condition the pyramidal element, the pyramidal element the conical shell, and the conical
shell the sphere.\footnote{% printed mark: *)
By carrying out the integrations in six different orders, new elements will indeed come in and
the course of the conditions be modified, but since the governing concept is the same, these
modifications are inessential. Only when the concept of the sphere is exchanged for another
concept must the elements be employed in another manner, and the formation becomes another.}

This series of conditions shows what the mechanical concept, apart from finite operations of the
understanding, is in itself as a concept of cause. It is not abstract logical universal forms,
but punctual conditions, that make up the concept's content. What is singly seen as arbitrary
becomes necessary in the connexion. No single condition is anything for itself, but all of them
are moments for one another mutually. Outside the concept the inward, ideal conditions are
meaningless representations; only in being penetrated and inspirited by the governing concept do
they have, in the unity of conceiving, each its peculiar significance.

To the conditions in the ideal there answer the conditions in the real point for point, when the
concept is a concept of power. Power's dialectic is, as we have seen, this double-sidedness,
that the moments are existing members, and the existing members moments. Outside the concept of
power the outward, real conditions are meaningless representations. Every condition must
necessarily be a condition for something, but the something whose condition it is cannot be
immediately determined by another condition. For the conditions to be able to be determined by
one another, they must, in the matter itself,
% --- p. 87 ---
% THE SEAM: the sentence from p. 86 („... maae de i Sagen selv“) continues
% immediately here. There must be NO blank line before „reflecteres“, which
% would create a false paragraph break.
be reflected in one another; the matter in which they are mutually reflected is the cause, the
concept in power, the ideal-real cause.

If we now return with these presuppositions to our original question --- \emph{how can that
which is itself without extension bring forth something in actuality extended}? --- a punctual
answer can be given. The greater the mass of the sphere, the greater the force with which it
attracts. If the sphere be resolved into its material elements, then the force is resolved into
nothing but corresponding moments: each moment of force is just as dependent upon its
peculiarly determined element of mass as the collected force is upon the whole mass. That the
moments of force can be integrated into a finite force is conditioned by the elements of mass
being integrable into a finite mass. The integration presupposes that in each element of mass
there is a disposition to actual extension, immediate reality in the extensive, just as
original as the disposition in each moment of force to actual strength, reflected reality in
the intensive. The ideal-real cause both of matter and of force is power, Spirit's power, the
divine object-power.

% A short centred rule in the middle of p. 87: here --- and only here --- division A,
% Mechanical Ideality, closes, having run two paragraphs onto this page.
\streg

% Danish: \afdeling{B.}{Mechanisk Realitet.}   (printed pp. 87--105)
\afdeling{B.}{Mechanical Reality.}

In mechanical reality the cause is to be apprehended as immediate, natural object-power. It
then comes to look as though the concept cause were only a name, for the object steps into the
cause's place: the investigation answering to this moves about \emph{power and cause}. Further,
object-power is power in another. It then comes to look as though the existing causes fell
outside one another
% --- p. 88 ---
in all accidentality: the investigation moves about \emph{the parcelling out of causes}.
Finally, the outward manifold of causes presupposes a foundation for many-sided reciprocal
action, so that the members in their mutual connexion come forward as caused causes, and the
question of reality becomes how matters originally stand with the infinite \emph{series of
causes}.

\runin{a) Power and cause.} An exactly determined separation of what belongs to the object, the
outwardly actual, and what the contemplating consciousness itself brings with it to the
apprehension of the object, is not so easy as it might seem at first glance. Declarations such
as this --- ``there are neither bodies nor causes, but only systems of facts, groups of present
or possible motions and of present or possible thoughts'' --- show an endeavour to fix the
separation, but at the same time furnish proofs of the superficiality with which an immature
reflexion can fancy itself raised above more penetrating investigations. If one says: the object
\emph{stone} is not a body; \emph{the object}, the subject in the judgment, is a \emph{fact}
lying outside consciousness, but \emph{body}, the judgment's predicate, is a word, a universal
concept falling within consciousness --- then we have the contradiction. In the one judgment it
is the word body that is predicate; but in the other it is the word fact. If the predicate body
is not to concern the object as object, but solely the predicating consciousness, then the same
must necessarily hold of the predicate fact. If one will insist that it is exclusively the fact
that is object, then it becomes still worse, for in the judgment ``the fact is object'' the
object is itself predicate and enters into consciousness. If one says: the fact is not cause,
for the cause is a concept, but the fact is only fact --- then it is manifest that
% --- p. 89 ---
one is playing with words. In the one case one lets the word point out beyond itself to what is
not a word but an existing thing; in the other case one lets the thing coincide with the word
and makes the thing a word. To end this play on words it is necessary to bethink oneself of the
relation between the logical and the antilogical.\footnote{% printed mark: *)
Cf.\ \textit{Grundideernes Logik} I, pp.\ 146--166; \textit{Grundideernes Logik i kort Begreb},
p.\ 28 ff.}

In the antilogical sense there are no thoughts or concepts whatever in nature: everything is
actuality, everything is object, thing, fact, in the smallest as in the greatest. The words
object, thing, fact will not be understood as words; they point out beyond themselves, they
speak to the one who senses, they say to the beholder: open your eyes, come and see! But in the
logical sense the fact is image, representation, concept. The investigating consciousness finds
in the whole of nature nothing but images, nothing but representations, thoughts and concepts;
and that for the good reason that consciousness, however busy it may be discovering objects,
things and facts, meets in everything its own apprehensions and cannot possibly at any point
leap out beyond itself. The earth attracts the moon. In the antilogical sense the earth is
earth, the moon moon, and the motion a fact that is observed; but the cause is not observed as
cause, the effect not as effect, for concepts are not objects of immediate observation. In the
logical sense, by contrast, the earth is cause and the moon's motion effect, for the actual
relation between earth and moon is impossible without a concept of relation belonging to
omnipotent conceiving.

The relation of experience between object and consciousness is well suited to clear up the
relation in the matter between power and cause. The object's reality is its power; power is
known
% --- p. 90 ---
by reality, and conversely. That reality is power is brought to consciousness through the
immediate impression of experience. If the impression could not with certainty be assumed to
issue from the object as effect from cause, then an experience of the object would be
impossible; one could not, after all, know that outside consciousness there really was any
object. In the relation of experience the object is cause and its apprehension effect. The
object is cause because it is power, and it is for consciousness power because it is the cause
of the impression. Just so the object's apprehension, its image and corresponding concept, is an
effect determined by the impression: the effect is determined by being referred to the cause. To
doubt the cause and rely on experience is accordingly meaningless, since the relation of
experience is a relation of cause and effect. In the knowledge of the real cause thought goes
together with sensation, the logical with the antilogical. From the logical side the object is
cause, from the antilogical side the cause is object: the language of experience is, in the
style of fact, a language of power that speaks for the cause's reality.

But whence does the observer now know that the relation between bodies themselves is also an
actual relation of cause? Solely from this, that the fundamental relation of experience is
itself a relation of cause. That the sense organs are acted upon by the sensible object is an
unmistakable proof from experience that body acts upon body, that, in other words, the relation
of cause has corporeal validity. And yet --- as follows straightforwardly from the extremes in
the essence of power --- the opposition of the logical and the antilogical is so abrupt that an
observer can be fully assured of the concept of cause being applicable as a concept of
experience and still harbour doubt about its a priori worth as a trustworthy expression for the
essence of things.

The mechanical phenomenon in which the relation of cause utters itself
% --- p. 91 ---
as a relation of power in a quite immediate and intuitable manner is impact. For centuries
impact was acknowledged as an actual cause before human knowing got so far as to be able to fix
the laws of its effect or, what is the same, to get the still indeterminate concept of cause
exactly determined. From the six propositions Cartesius set up it is seen that he misses the
problem, in that he directs himself more by his own self-made concept than by the facts. His
fundamental thought, that the quantity of motion in the whole corporeal world must always remain
equal to itself, misleads him, namely, into the assumption that two motions opposite in
direction do not cancel one another as positive and negative magnitudes, but that bodies
colliding in this way must go back, each with the same velocity it had before the impact ---
% The Latin citation is set in italic against the page's antiqua. „æqvalia“ and „æqve“
% are the book's own qv spelling, checked on the KB copy.
\textit{si duo illa corpora essent plane æqvalia et æqve velociter moverentur, B a dextra versus
sinistram, C a sinistra versus dextram, cum sibi mutuo occurrerent, reflecterentur et postea
pergerent moveri: B versus dextram, C versus sinistram, nulla parte suæ celeritatis amissa.} The
formula holding for inelastic bodies,
% The sign in the numerator is ÷, Nielsen's minus (cf. p. 46), checked on both copies.
$x = \dfrac{MH \div mh}{M + m}$, which for the mass $M = m$ and the velocity $H = h$ gives the
velocity after impact $x = 0$, now shows that Cartesius has not become aware at all of the
difference between elastic and inelastic bodies.

Mistakes of this nature indicate that the determining power lies not in the human concept but in
the fact. The investigator's knowledge is conditioned by trials; he must experiment. And yet the
fact, immediately regarded, can teach him nothing; it can be resolved into a manifold of facts,
but each of the parts is like the whole --- a fact, no concept. To learn anything from the fact
he must himself think and judge, and so bring concepts with him. But when it is
% --- p. 92 ---
the investigator who must add the concepts out of his own consciousness, the special ones as
well as the universal, then all his facts are, as regards the conceiving, only occasions, not
working causes: then he learns only occasionally from the fact what he originally learns from
himself. So after all the determining power comes to lie in the concept.

The elastic body, for example, is in antilogical determinateness a fact, but elasticity is a
concept. How does the observer now get from fact to concept and from concept to fact? If B and
C, he says, are elastic, they will in striking against one another be compressed to a certain
maximum. In this maximum of compression they are, if only for an instant, inelastic, and the
common velocity comes, as for inelastic bodies, to be expressed by
% The sign in the numerator is here ±, not the ÷ that stands in the same formula on
% p. 91 and in the following parentheses on this page: a plus over a stroke, without
% dots. Checked on both copies, which agree; reproduced as printed.
$x = \dfrac{MH \pm mh}{M + m}$. B has accordingly lost the quantity of motion $M(H \div x)$, and
C the quantity of motion $m (h \div x)$. But now it lies in the essence of elastic bodies that
after the compression they strive to return to their original form; effects similar to those
that took place during the compression then also take place during the expansion, whereby B and
C again lose the same quantity of motion as before. The whole loss in quantity of motion is
accordingly, for B's part, $2 M (H \div x)$ and for C's $2 m (h \div x)$. If the observer could
give a punctual account of how much in this explanation he has from the fact and how much from
his own concept, a problem of great significance for the science of experience would be solved.

The explanation given is built upon the observation of facts; and yet it goes far out beyond the
observed facts, for it is built upon concepts and judgments. Every fact is antilogically a
\emph{this}; only
% --- p. 93 ---
this, the single, can each time be observed; the universal cannot be immediately observed.
\emph{This} body can be observed, but a body in general, or the universal in the body, cannot be
observed; \emph{this} compression and \emph{this} elasticity can be observed, but elasticity in
general, or the universal in the tension, cannot be observed. None the less the observer may
rely unconditionally upon his concept as determined by the fact. That there are found in nature
neither absolutely elastic nor absolutely inelastic bodies never holds him back from laying the
concepts elasticity and inelasticity at the foundation of a mechanical analysis and, in spite of
all particular conditions, letting the consequences of the concept hold good for the universal
essence of the facts. Led by his concept he observes even in the fact itself not only the
visible but also the invisible.

How it goes when two bodies, be they more or less elastic, momentarily alter one another's form
during the collision cannot be seen with sensible eyes, and yet the observer sees mathematically
what happens. Before the impact the body B with the mass M and the velocity H had a living force
$MH^2$; after the impact, that is, at the maximum of compression, it has a living force $Mx^2$.
The body has accordingly lost in living force $M (H^2 \div x^2)$, and why? Because in compressing
the body C it has performed a work. If the pressure answering to the penetration s be called P,
then $P = f(s)$, and $M (H^2 \div x^2) = 2 \int\limits_{s_2}^{s_1} Pds$. It lies in the nature
% Note the signs on this page: ÷ (Nielsen's minus) in the parentheses
% $M (H^2 \div x^2)$, but an ordinary minus stroke in „$s_1 - s_2$“. Both copies agree;
% reproduced as printed.
of the case that the interval $s_1 - s_2$ must be greater when the impression is made, under
otherwise equal conditions, in a softer than in a harder body. If in both cases the left side of
the equation is to be the same, and so $M (H^2 \div x^2)$ to have the same value, then the right
side, regarded as a sum of $Pds$, must contain, in the midst of infinity, a greater
% --- p. 94 ---
% p. 94 opens with six lines set in SMALLER type than the text; the change of size falls
% exactly at the turn of the leaf, in mid-sentence, and the smaller type runs to the end
% of the paragraph („... ved det blødere.“). Checked on both copies, which agree.
% Reproduced with \small.
{\small number of addends for the softer than for the harder. The sum is namely equally great in
both cases, for, as the equation shows, it depends only on the impact's initial and final
velocity; but the single addends must, in equal sums, necessarily be greater where there is a
smaller number; in other words: the pressure with the harder body is greater than with the
softer.}

When the investigator in virtue of observations goes in this way out beyond what admits of being
observed, then he argues not from facts but from concepts. By laying concepts at the foundation
of the analysis he is put in a position to derive the laws for elastic bodies from the laws for
inelastic ones, and to see the necessity that the latter must always lose in living force upon
collision, while the former must not.

But there still remains the answer to this question: what entitles the observer thus to exchange
the single that he sees in the fact for the universal that he does not see? He abstracts, it is
said, the concept from the fact; we must then take note of what such an abstraction essentially
presupposes. The single is not the universal; but if the single in the fact could not present
the universal, then it would be impossible to abstract the universal from the single. The object
is not a concept; but if no concept were uttered in the given object, the investigator would
deceive himself in supposing that he could abstract a concept from the object: to abstract a
concept from something in which there is no concept is a contradiction; to discover and pursue
consequences of the thoughts uttered in natural things, when there are no thoughts in the
things, is a contradiction; to make the thing concept and the concept thing is a contradiction.
As an object of scientific experience the thing is at once the concept's effect and its cause.
It is effect; for the divine concept of power is its ideal cause. It is
% --- p. 95 ---
cause, namely cause of experience in reciprocal action with the observing consciousness, which
thereby comes to a concept of the thing, and by the thing's concept to a concept of the
reciprocal action of things, and so to a concept of the cause's reality. As the object's concept
power is cause; as the concept's object the cause is power.

\runin{b) The parcelling out of causes.} The real cause lets itself be parcelled out; the whole
of outward nature is a magnificent example of the parcelling out of causes. Now if every real
cause is an embodied concept, then it should surely seem that the embodied concepts must have
the same fate as the bodies; and so: when corporeal causes are parcelled out, the concepts of
cause must be parcelled out; and when corporeal causes are pieced together, the concepts must be
pieced together. But concepts can neither be parcelled out nor pieced together; the concept is as
concept an indissoluble totality of determinations belonging together: how, then, does it stand
with the mechanical parcelling?

The gravity indwelling in all bodies works as cause; the concept of gravity is a concept of power
pervading the universe; and yet here is a division. The planets in space have in actuality each
its own share of universal gravity, the earth its own and the moon its own. But divisibility in
the real is precisely grounded upon indivisibility in the ideal. If the concept of gravity, and
with it gravity's mighty essence itself, were divided and by the division had gone to pieces,
then the piece of gravity dwelling in the earth would find itself at a mean distance of 51,032
Danish miles from the piece dwelling in the moon; but with this the bond of gravity would be
broken: gravity could not be gravity. The secret of the cause is precisely this, that the real
can be divided while the ideal is and remains an indivisible power. In the real sense the earth
is cause, and the moon is cause, and the two
% --- p. 96 ---
% p. 96 begins WITHOUT indentation, in mid-sentence from p. 95. No new paragraph here.
% The head b) Udstykning af Aarsager does NOT stand on p. 96.
real causes really find themselves at a distance of 51,032 miles from one another; but in the
ideal sense gravity is cause, and gravity's ideal bond cannot burst. Between the two causes,
ideal for earth and moon, there is no distance: they are, each for itself, the whole cause, and
yet moments of one cause.

That the indivisibility of the ideal can consist with a parcelling out of the real has its ground
in this, that ideal indivisibility itself consists in a fundamental division, but a fundamental
division is no parcelling out.\footnote{% printed mark: *)
\textit{Grundideernes Logik} II, pp.\ 347--458.} The same concepts of power can be individualized
in nature in different ways, just as the same laws can be fulfilled in countless ways. It must
constantly be held fast that the laws of nature in all facts, just as much as the facts
themselves, rest upon the concepts of nature. Physics does indeed discover manifold laws whose
validity it secures by ordinary induction, without yet being able to ground them upon a clear
insight into the nature of the real causes; but induction itself would lack a true scientific
ground if the investigator did not dare presuppose that it is everywhere concepts --- if hidden,
as yet undiscovered concepts --- that furnish the foundation of the laws. The laws for the
% COMPOSITOR'S FAULT: the Danish has „Virkniuger“ with an inverted n (u) --- checked on
% the KB copy, where the third letter of the syllable is plainly a u with an open top,
% while the n in the same word is closed. Both copies show it, so it is the compositor's
% fault. Reproduced as printed; read „Virkninger“ (effects).
effects of electricity cannot indeed yet be derived from the concept of electricity in the same
penetrating manner as the laws for the collision of elastic bodies can be derived from the
concept of elasticity; but there is surely no physicist who doubts on that account that the
fundamental relation between concepts and laws must constantly be the same.

How the laws stand to the concepts of cause, and these to the real causes, can already be seen
from the so-called mechanical powers. Of such simple machines as the lever, the inclined plane,
etc., one may certainly
% --- p. 97 ---
say that each of them is a piece of cause, an arbitrarily formed implement; but that the whole
concept of cause is just as much realized in the lever as in the inclined plane, that every one
of these implements presents a whole concept for itself, is evident. That these restricted
concepts, in all their one-sidedness, are not fragments but moments of the great mechanical
concept of cause embracing the universe follows from this, that the laws of equilibrium and
motion individualized in these implements are laws of nature, universal world-laws.

The law as law is one side of the concept and itself the whole concept, namely in the form of
universality, while the corporeal causes in their mutual difference present the concept in the
particular, and each body for itself the antilogically existing single concept. From the side of
reality the single concept is, however, to be seen not as concept but as matter, as body. The
observer does not notice that bodies are in their essence concepts just as much as the law is,
and that the law is for that reason sensible in its fulfilment just as much as bodies are. He
senses bodies and abstracts the law; the law thus abstracted therefore becomes for him not a
concept but a categorical form for an unexplained necessity. The one-sidedness of the standpoint
may, however, be said to be justified by the peculiar character of the science of experience; it
does not lie in the investigator's interest to develop the concepts of cause as real concepts of
nature: he holds to facts and laws.

But what has the philosophy of immanence accomplished? Hegel inveighs against the physicists and
complains of their parcellings; but the physicists are surely excused when nature's ``being, as
it is, does not answer to its concept.'' How should the investigators of nature be able to hold
to ``the Idea'' when nature cannot even do so, but must in the Hegelian sense be regarded as
``the Idea's falling away from itself''? In its
% --- p. 98 ---
endeavour to set nature, fallen into ``outwardness,'' on its feet again by means of ``the
concept'' --- or rather, to make it conceivable how ``the Idea'' works its way up step by step
out of its degradation --- speculative philosophy of nature has in a sorry manner had to
overstrain logical necessity.

As surely as it is ``the absolute Idea'' that begins its absolutely necessary journey upward from
below, the philosophy of nature must quite rightly furnish the proof that the stepwise ascent
happens with immanent necessity, and that this present world-order is the only conceivable one in
which the Idea could come to outward, actual self-revelation. If such a proof had really been
furnished or could be furnished, then we should certainly be bound to acknowledge the proof and
bow to the necessity; we should, as H.\ C.\ Ørsted says, have to ``be ashamed in our inmost self
if we caught ourselves wanting a truth other than the actual one.'' But when such a proof has
neither been furnished nor can ever be furnished, then we offend against the actual truth in
bowing to a fixed idea.

It is true that everything that happens must have its cause and so far its necessity; but it is
untrue that any relative actuality whatever should be absolutely necessary. In abstract
universality necessity is an empty thought; to give the thought a content we must give it a form
and take it up into a determinate concept. Only by the concept can anything necessary come out of
necessity. Necessary is that whose opposite is impossible. Conceptual necessity everywhere marks
the limit of the possible and the impossible, but never reaches the limit of the actual.

The limit of necessity for an actual inclined plane, for example, is a limit of its possibility.
All determinations of logical necessity concern actuality only so far as
% --- p. 99 ---
they lie in possibility. It is necessary that every inclined plane must have a base, a height and
a length; why? Because an inclined plane without height, length and base is an impossibility. But
within
% COMPOSITOR'S FAULT: „Mulighedeu“ with an inverted n (u) as the word's last letter.
% Checked on the KB copy against „kan der“ in the same line, where the n is closed at the
% top; here stands an open u. Both copies show it. Reproduced as printed; read
% „Muligheden“.
possibility there can be an infinity of actual inclined planes, different in their actuality; to
the actual differences necessity is indifferent: in actuality it is just as content with the one
specimen as with the other. The same holds of the laws. ``The tendency to fall on the inclined
plane must stand to the pressure against it as the inclined plane's height to its base.'' That is
a law, for it is necessary; its opposite is impossible. But this necessity decides not the least
thing concerning the actual height, the actual base, the actual pressure. Conceptual necessity is
hard and strict against the impossible, but supple, compliant, extremely liberal toward the
actual. Speculation may widen its horizon as much as it will; its logical necessity never reaches
out beyond possibility.

It might perhaps appear at first glance to one or another that the doctrine of this world's
absolute necessity had a reliable foundation in rational mechanics; but not to speak of the fact
that mechanics, in its presentation of the laws for the motion of the heavenly bodies, lays the
Newtonian law at the foundation without an a priori proof of that law's exclusive, absolute
necessity, it goes on the whole with necessity for the world's planets as it goes in particular
with necessity for the inclined plane. If we represent a planet with the mass m moving about a
central body with the mass M, and decompose the accelerating forces along three mutually
perpendicular axes, so that
% The setting has a short hyphen without spacing: „og for m-x, y og z:“. The sense
% requires a mark of division (colon or dash); both copies show the same short stroke, so
% it is reproduced as printed.
the coordinates for M become X, Y and Z and for m-x, y and z: then the distance r between M and m
will be
% ÷ is Nielsen's minus sign (cf. p. 46 and the house rule in the header); it stands so in
% both copies under the radical, whose angle covers the whole expression.
$r = \sqrt{(x \div X)^{2} + (y \div Y)^{2} + (z \div Z)^{2}}$; that is necessary, for its opposite
is impossible. If we reflect solely upon
% --- p. 100 ---
the X-axis and call the angle r forms with this axis α, so that $\cos\alpha = \frac{x \div X}{r}$:
then the attraction the central body exerts upon the planet, and so the latter's acceleration,
will be
% After the equals sign stands an independent ÷ before the fraction: it is Nielsen's minus
% sign, not a division sign. The second equation has no sign. Both copies agree.
$\frac{d^{2}x}{dt^{2}} = \div\frac{M(x \div X)}{r^{3}}$, and the attraction the planet exerts upon
the central body, and so the latter's acceleration, will be
$\frac{d^{2}X}{dt^{2}} = \frac{m(x \div X)}{r^{3}}$; that is necessary, for its opposite is
impossible. If we let the central body be immovable, then its acceleration must be added
negatively to the planet's, and the equation for the latter's motion becomes
$\frac{d^{2}(x \div X)}{dt^{2}} + \frac{(M + m)(x \div X)}{r^{3}} = 0$; that is necessary, for its
opposite is impossible. That under such analyses images of actual planets, an actual central body,
actual planets and actual distances may float before the mathematician signifies nothing with
respect to actuality itself and the actual parcelling. Where actual motions take place it is
certainly necessary that the general laws must be fulfilled; but from this logical necessity to
the absolute necessity of an antilogically determinate planetary system there is an infinite leap.

The absolute Idea, one says, has all relative moments, all essential conditions for its realization
in itself. The whole outwardly cast finitude is accordingly a parcelling out in which the absolute
must be absolutely omnipresent. The more one widens the concept of the whole, the more conditions
one gradually takes up into the general presupposition of the analyses, the more one will be able,
by the crossing ways of inner necessity, to approach the outwardly actual. We have in this way of
seeing one of the many immanent ambiguities with which ``modern'' philosophy has become so
overloaded that only a reference to exact determinations of form can be supposed likely to exert a
% --- p. 101 ---
% NOTATION: the series m, m', m'' carries primes, and per the house rule (batch 5) the whole
% enumeration is set in math, so that the face does not change in mid-series. The isolated
% letters R, x, X, r in what follows stand, by contrast, as printed, in ordinary setting.
purifying effect. Let us then add to the presuppositions given this further one, that besides the
planet $m$ other planets $m'$, $m''$, etc., move about the central body. What does this now lead to?
On the last presupposition the sum
$\frac{d^{2}(x \div X)}{dt^{2}} + \frac{(M + m)(x \div X)}{r^{3}}$
or, when $x \div X$, m's abscissa with respect to M's centre, is denoted by $\xi$, the sum
$\frac{d^{2}\xi}{dt^{2}} + \frac{(M + m)\xi}{r^{3}}$ cannot indeed be zero; here a perturbing
function R comes forth, and the equation becomes
$\frac{d^{2}\xi}{dt^{2}} + \frac{(M + m)\xi}{r^{3}} = \left(\frac{dR}{d\xi}\right)$. Such an
equation certainly gives the expression of necessity a more articulated form; but the limit is only,
in possibility, a limit for the impossible. Even if we were, guided by the method of ``the variation
of the arbitrary constants,'' to wander through the whole primeval forest of equations, series and
formulas that lies behind the indications above, we should still not reach the antilogically
determinate facts, should not get the material actuality of the parcelling pressed into absolute
necessity.

% The head is printed SEMIBOLD, not letterspaced like the other a)/b)/c) heads (the same as
% with c) Aarsagsmomenternes Eenhed, p. 79); „c)“ itself stands in ordinary antiqua. The
% edition does not reproduce the difference, so \runin is used here too.
\runin{c) Series of causes.} The parcelling out of the mechanical causes must necessarily have a
connexion of things of its own for its foundation. From the side of universality the connexion is
realized in gravity, which under the form of particularity varies with the corporeal states. Physics
sees, in the resistance bodies offer to the separation of their parts, a force of cohesion different
from the force of gravity, and surely
% THE NOTE RUNS OVER TWO LEAVES: it begins beneath the rule on p. 101 and fills the whole
% note field on p. 102; the mark is not repeated on the continuation. The whole note is set
% here at its mark on p. 101.
% NOTATION: the series of letters T, r, μ, ϱ, ε mixes Latin and Greek signs; per the house
% rule about enumerations the whole series is set in math, so that the face does not change
% in mid-enumeration.
rightly.\footnote{% printed mark: *)
The attraction of cohesion holds, namely, only for the small parts among themselves; if this
attraction were an utterance of the force of gravity, it would, on account of the smallness of the
parts and their great mutual nearness, have to become extraordinarily slight. According to the law
of gravity $T = \frac{\mu}{r^{2}}$ --- $T$ being the attraction, $r$ the distance and $\mu$ the
intensity --- the intensity for a spherical mass with radius $r$ and density $\varrho$ will be
$\mu = \tfrac{4}{3}\pi\varrho r^{3}$. If now the sphere shrink to an atom, the intensity for the atom
will be $\varepsilon = \tfrac{4}{3}\pi\varrho(dr)^{3}$ and consequently
$\frac{\varepsilon}{(dr)^{2}} = \tfrac{4}{3}\pi\varrho\, dr$, a vanishing magnitude. The objection
that the small parts are not infinitely small signifies nothing in this connexion; for the formula
shows that the force of cohesion must become the weaker, the more intimate the joining is ---
something which is both at variance with experience and in itself a contradiction.}

% --- p. 102 ---
% The body text on p. 102 begins here; the whole note field on the leaf is taken up by the
% continuation of the note from p. 101, set above.
But however little the force of cohesion may be regarded as a mere offshoot of the force of gravity,
it will yet be intelligible that the two heterogeneous forces must be referred, not indeed to one
fundamental force, but still to one fundamental concept, the comprehending concept of power whose
inner articulation is the ideal cause of all forces. Just as the force of gravity, as has been shown
in its place, rests upon reflexion in power, so too does the force of cohesion. The particular
concepts in this reflexion are \emph{whole} and \emph{parts}. That the body is a whole real in
parcelling out is synonymous with its consisting of parts. These parts are in the mechanical sense
dependently independent members. Their dependence is an expression of the real concept's power in
the given whole; but their independence shows the whole's being pieced together, for the real concept
is power in the whole only by being particularly power in the parts: the expression of cause
reflected in existence for the dependence of the independent parts is the force of cohesion. From the
forms of state resting upon particular relations between ``force of expansion'' and ``force of
cohesion'' it is seen how the properties of the parts reflect themselves in the whole, and from this
again how matters stand, on the foundation of the parcelling, with the concepts of body and existing
bodies.

It is a fundamental proposition in mechanics that action and reaction must answer exactly to
% --- p. 103 ---
% NOTATION: the series k, M, v, M', v' carries primes, so the whole enumeration is per the
% house rule set in math. ÷ before kt is again Nielsen's minus sign; the other signs in the
% section are true plus and equals signs.
one another. The proposition holds for gravity; for the attraction a mass exerts upon a point is
always equal to that which the point exerts upon the mass. A falling grain of sand attracts the whole
earth with a force equal to that with which the whole earth attracts the grain of sand --- that the
acceleration is highly different is quite another matter. --- The proposition holds for quantities of
motion. If a constant force $k$ act for a certain time upon the mass $M$, then the acceleration is
$\frac{k}{M}$, and the quantity of motion in a certain determinate direction $Mv = kt$. But the force
acting upon the mass $M$ must have issued from another mass, e.g. $M'$, and the mass $M$ acted upon by
$k$ then reacts upon $M'$ with a force which is also $k$; only that the activity goes in the opposite
direction, and the quantity of motion becomes $M'v' = \div kt$. From this follows $Mv + M'v' = 0$, an
equation showing that the sum of the quantities of motion for both bodies continues unaltered. --- The
proposition holds for living force. It cannot indeed be said that the work given off by one body is
always received by another, for the activity can assume different forms; heat, for example, can give
off work, and work again be stored up. But if the activity be regarded under all forms as the same
function of force, and be measured by the same measure, then this
% COMPOSITOR'S FAULT: „Functiou“ with an inverted n (u). Checked on the KB copy against
% „Kraftfunction“ two lines above, where the n is closed at the top. Both copies show it.
% Reproduced as printed; read „Function“.
function too will remain unaltered in passing from body to body: action and reaction answer to one
another.

The fundamental proposition teaches us that forces always come forward two by two, parallel, equal and
opposite: pressure answers to counter-pressure, impact to counter-impact, A's attraction of B to B's
attraction of A. The mechanical cause is, in other words, never any immediate single cause, but always
a reciprocal action of two causes. That every outward cause is, essentially regarded, a reciprocal
action of two outward causes spreads a peculiar light over all mechanical relations of cause.

% --- p. 104 ---
To begin with gravity, the most comprehensive of all outward causes: it is not merely outward in the
relation between the masses upon which it works, but it is in itself outward. The mass $m$ attracts
$m'$ insofar as $m'$ attracts $m$; that the attraction is equally great from both sides proves that
the cause in both masses is equally original and equally derived; both causes are in their activity
\emph{caused}: they cause one another mutually. What holds of two members holds of all; here an
infinite series of caused causes forms itself; no member marks the first cause. The contradiction that
gravity is everywhere equally original and equally derived is resolved in nature in such a way that
gravity, which presupposes itself to infinity, comes outside itself and concentrates its effects in
other causes, different from it: \emph{pressure}, \emph{impact}, \emph{quantity of motion}.

Where the distance between the heavy masses is cancelled, gravity works as \emph{pressure}. According
to the fundamental proposition of the equality of pressure and counter-pressure, a body's weight is to
be determined by gravity; and yet pressure as cause is different from gravity. By gravity the parts of
mass strive to unite in one point, the centre of gravity, but pressure counteracts this striving and
is, though itself an effect, nevertheless cause --- cause, namely, of the parts keeping outside one
another, so that masses fill space. Only what is pressed presses; pressure as cause is not merely
outward with respect to its object; it is in itself an outwardly, infinitely caused cause. Series of
causes formed of pressure and counter-pressure cannot possibly have any first member in which pressure
as pressure was original cause. The contradiction that every particular pressure in its immediacy is a
derived, caused cause resolves itself into a transformation of pressure. By this transformation it
becomes a moment in
% --- p. 105 ---
another cause, \emph{impact}. Impact presupposes motion, and motion must have its cause.

Where a distance between bodies is given, gravity may indeed cancel the distance and become the cause
of the impact; but gravity, which is not originally the cause of the distance, cannot be the original
cause of the impact. Just as pressure is caused by a striking together, the distances of bodies would
accordingly have to be caused by a striking apart; impact would then from this standpoint be
apprehended as the concrete and so far original cause, having gravity for its universal and pressure
for its particular moment. Older mechanical philosophy of nature (Cartesius) has not let references be
wanting to a first, original impact, a fundamental impact, from which all motions and world-phenomena
were to be explained. But that the representation of a first original impact must resolve itself into a
conceptless concept is clear from the fundamental proposition: only from impact comes impact. But when
only what has previously been struck can strike, then every impact is a cause caused by an impact. The
infinite series of causes in which impact unbrokenly presupposes impact cannot then possibly begin with
a first impact. If a beginning be made with a first impact, i.e., with an impact whose cause is not an
impact, then the beginning of the series falls outside the series: the beginning is then made with a
cause which does indeed contain the possibility of impact, but is nevertheless different from impact;
the beginning is made with a \emph{quantity of motion}.

So the quantity of motion would then be the cause rich in content which, by taking gravity, pressure and
impact up into the circle of its real moments, could in spite of the parcelling be the principle of
motions and come forward as fundamental cause; but this too is impossible. Quantities of motion can,
says physics, even without regard to the unalterability of force, be propagated from body to body
without altering magnitude or direction, so that what the one

% --- p. 106 ---
% p. 106 does NOT open with division C. Division B (Mechanical Reality) runs on to
% here with TWO paragraphs: the first continues the sentence from p. 105, the second
% ends with „mechanisk Realitet.“ Only then comes the short centred rule, and only
% then the C head. The same shape as at the A -> B turn on p. 87.
body receives, the other loses. For a world of bodies the quantities of motion in every
direction will accordingly be unalterable magnitudes, whatever forces work between the bodies or
between the corporeal parts themselves; but in the form of quantity of motion too the cause is
outward to itself.
% The sign before M'v' is Nielsen's MINUS (÷), read on the KB copy; cf. the house rule
% from batch 4 (p. 46).
The cause of an $Mv$ is in a $\div M'v'$; the quantity of motion that communicates itself must
everywhere have been communicated: the infinite series of causes is a series of nothing but
communicated, \emph{caused} causes. Nor does the relation of causality alter its nature by the
quantity of motion pointing out beyond itself, being taken up into a still more concrete cause
and converted into \emph{living force}.

% „levende Kraft“ is letterspaced at the end of the paragraph above, but NOT in this
% paragraph, where the words stand three times in ordinary antiqua (checked on the KB
% copy).
A body's increase in living force is, it is said, equal to twice the work received, and a body's
loss in living force equal to twice the work given off; the equation
$M(v^{2} \div c^{2}) = \pm\, 2\,MGs$ shows both: the conversion of work into living force and of
living force into work can be continued to infinity. With the co-operation of the different
solid, liquid and gaseous states in which bodies find themselves, the outward causes can be
transformed in countless ways and enter into manifoldly entangled reciprocal actions; but that
these outward causes everywhere, in the great whole as in the smallest parts, necessarily
presuppose a ground-cause and yet keep the ground-cause hidden, is and remains the mark of
character of all mechanical reality.

% Short centred rule between division B and division C, in the middle of p. 106.
\streg

% Danish: \afdeling{C.}{Grundaarsagen: det første Bevægende.}   (printed pp. 106--125)
\afdeling{C.}{The Ground-Cause: the First Mover.}

According to mechanical ideality the working causes are reflected in corresponding concepts of
cause, and these
% --- p. 107 ---
holders of power in causality then show themselves to be the ideal causes determining from
within. According to mechanical reality, by contrast, the concepts of cause are, by means of
object-power, hidden under the outward material existence in which they have their actuality; the
ground-cause has disappeared, and the beholder of the infinite series looks about him in vain for
a first mover. If it is to be possible to discover the ground-cause in the series of motions,
then the difference between the two heterogeneous motions --- the \emph{ideal} in the concept and
the \emph{real} in the fact --- must be expressly shown. It is under the ideal motion that
Spirit, the self-power that causes everything, determines the particular concepts of power to be
the material causes of object-power. But the ideal motion presupposes precisely what it itself
determines, namely the real. Not immediately, but through actual causes in nature and real
motions, is self-power \emph{the original cause of the whole mechanical parcelling out}. The
parcelling rests as little upon an absolute necessity as upon a blind accidentality: the
ground-cause itself, omnipotent freedom, is \emph{the first mover}.

% Run-in head, read on the page itself (KB copy): set in bold antiqua, not letterspaced;
% the „er“ before it is ordinary type.
\runin{a) Ideal and real motion.} The doctrine of immanence has done its utmost to make ideal and
real motion look like one and the same thing. In Hegel's logic the motion of the concept is to
pass for the motion of the matter; in his philosophy of nature the motion of the matter is in
return to pass for the motion of the concept. One need only bethink oneself of the fundamental
contradiction lying in this to come to certainty about how hollow the absolute is upon which
``immanence'' must go on building its ``pantheistic,'' ``naturalistic,'' ``atheistic'' castles in
the air.

If it is to hold in earnest that the logical Idea has an eternal being resting in itself and
independent of nature, then there must necessarily be presupposed a thinking, reflecting and
conceiving exalted above space, time and
% --- p. 108 ---
matter; for without thinking, reflecting and conceiving there is not a single logical
determination, far less an absolute logical Idea. In declaring for the logical Idea's originality
and holding fast its supernatural being, the philosophy of the absolute goes out beyond its own
vaunted immanence and makes a questionable loan from transcendence. From this standpoint the
ideal motion in the concept is then also essentially distinct from the real motion of nature's
development, and yet the two motions are supposed to be one.

To evade the contradiction in the consequence just drawn, one has recourse to another
explanation. That the absolute is in its essence logical, and the concept its original form, must
not then be understood as though the Idea of logic had really thought through its categorical
determinations of content before it altered itself, struck over into the real and began afresh as
the Idea of nature. The absolute's eternal being is on the contrary an original natural being;
its immanent determinations of essence then transform themselves as at a stroke of magic; they
are not in the absolute in an absolute manner; their existence comes to be apprehended as the
sensible, material forms of nature, of outwardness, inadequate to the Idea: the motion is real.
For the absolute nevertheless to be able to remain absolute in the midst of the dispersion, the
thoughtless must itself be thinking. It is the universal, the all-thinking, the great thinking
neuter, upon which immanence rests. This universal is indeed a form lost in things, bound to
nature, which can first be freed into conscious thinking in the human spirit; but the unfree
state of nature makes no breach in the universal's infinite thinking; consciousness is quite
superfluous; conscious thinking is only a lower kind of thinking. Absolute thinking is a thinking
without consciousness, just as absolute knowing is ``an unconscious knowing,'' a clarity without
light, a conceiving without insight,
% --- p. 109 ---
an understanding without understanding. On these terms it is not at all unreasonable that time
can conceive space, and space time, and that the matter brought forth by the mutual conceiving of
space and time can again conceive and move itself: the real motion is ideal, for the ideal is
real.

What is true in all this is that the concepts of nature are indwelling in natural things. It is
Hegel's immortal merit to have developed Spinoza's barren doctrine of immanence into so fruitful
an all-sidedness that ``the concept'' is seen to hold a world-content. Hegel's philosophy of
nature deserves in this respect just as full appreciation as his logic. But the fundamental error
is that immanence, however much is said of the ``overreaching'' --- and so transcendent ---
concept, everywhere dissolves and denies the original presupposition upon which it itself rests:
immanence is impossible without transcendence. The concepts of nature have a double medium; they
are indwelling in things: actual nature is their real medium; they proceed from infinite
conceiving: the all-knowing, self-powerful Spirit is their ideal medium. The unity of the two
media is a unity in opposition which grounds the essential difference between ideal and real
motion.

% The Greek letter is read on the KB copy: α, not the „a“ or „&“ the two OCR witnesses
% give. The head is run-in and letterspaced.
\subrunin{α) Ideal motion.} For the mechanical cause to be concept, it must be conceived. How the
conceiving goes on in human knowing has been shown above. Thinking and representation are
functions of one another mutually. The cause cannot be thought without a something representing
the active element, namely force, and the active element cannot be thought without a something
upon which it works, namely matter, whose inertia is determined precisely by its receiving and
preserving the action. Force and inertia reflect one another mutually: the phenomenon answering
to this in time and space is acceleration, and the cause represented in the concept is the
% --- p. 110 ---
accelerating force. In acceleration the mass, limited matter, is understood along with it but not
yet taken up into the concept itself. The concept of cause demands that acceleration and mass
expressly enter into the reflected unity of the moving force, the more determinate expression of
the cause conceived in time, whose measure is quantity of motion. But quantity of motion in exact
form is a one-sided measure. The relation of causality demands that the effect be opposed to the
cause, namely in such a way that the cause reflected in its effect answers exactly to what is
effected. The completely determinate effect is work; the cause is set to work and assumes the
form of living force. However many determinations representation may have brought to thinking in
the course of this development of the concept, it is nevertheless evident that none of these
determinations of representation is drawn in under the cause from without, either arbitrarily or
accidentally, but that they all develop from within out of the mechanical cause's own concept. If
the beginning be made in the abstract with acceleration, then the motion cannot stop before the
concrete is reached in living force; if the beginning be made in the concrete with living force,
then analysis will discover in it space, time, quantity of motion, and in the quantity of motion
acceleration: that is ideal motion.

And yet the development of the concept is not yet completed. The mechanical cause is by character
a cause of otherness; its activity is in the outward; the concept of outwardness belongs
originally to its essence. The conceiving reflexion is determined by reflexes of outward things,
and what causes spreads itself into a circle of conditions. But the independence of the
conditions is only apparent; the conditions, cancelling one another, disappear as moments in the
cause's unity, the unconditioned. By taking the conditions up in the form of unconditionedness
the cause is conditioned in itself, and in such a way that it can develop its own difference from
within out of its conceptual content: one cause can,
% --- p. 111 ---
on the foundation of the conditions, become several, and several again one. It is no real but an
ideal motion that is seen in these transitions. The concept of cause is as yet only a concept;
but the conceiving shows that the ideality of the moments reflects itself point for point in a
corresponding reality.

% spacing.py reported „Aand“ and „vanvittigt“ here as possible letterspacing; both
% rejected after inspection on the KB copy --- ordinary antiqua, loose justification in
% the working scan.
As a concept of nature the cause has its ideal medium in the all-conceiving Spirit; to the
infinite conceiving there must then also answer an ideal motion. As unreasonable, indeed insane,
as it would be to represent the absolute Spirit after the likeness of the human one ---
observing, deliberating, abstracting, calculating, differentiating, integrating --- just as
unreasonable, indeed insane, is it to want to think a cause without a concept, a concept without
a development of the concept, a development of the concept without inner, free, ideal motion. The
metaphysics of immanence, with its selfless
% „Grund“: both OCR witnesses read „Gruud“ (the familiar n/u confusion). The KB copy
% does not settle it positively at 300 ppi, and per the playbook the sense's reading
% stands here as „Grund“.
thinking neuter, ought at least to be able to learn this much from the analysis of infinity that
lies at the foundation of rational mechanics: that it should at last come to consciousness of how
matters absolutely must stand, within conceiving itself, be it human or divine, with the
articulation of the infinite by the finite.

% The Greek letter is read on the KB copy: β, not the „B“, „ß“, „6“ or „ØB“ the OCR
% witnesses give.
\subrunin{β) Real motion.} Real motion stands in sharp opposition to ideal motion, because the
existing cause stands in sharp opposition to the concept of cause. A body's motion from place to
place has, outwardly regarded, no likeness to a motion of the concept. It is a language of power
that dictates the opposition; but this language of power is the inborn natural language of the
real concepts. The concepts' transition from the ideal into the real medium is no finite
natural-historical occurrence; it is a fundamental motion, an infinite event belonging equally to
time and eternity. The transition does not happen in such a way that the real concepts, untouched
by all realities, are preserved in self-power's eternity so that object-power may afterwards
supervene in time and fill the empty forms with
% --- p. 112 ---
stuff-like content. Self-power and object-power are extremes of the same omnipotence in an
eternal-temporal unity of doubling. It has been made clear by analyses that the same differential,
$dm$, which is ideal in the body's concept of nature, is real in the existing body --- that, in
other words, object-power, which posits existence, is infinitely penetrated and inwardly moved by
the concept of existence. But precisely by expressing its concept the object comes to conceal the
concept; the fact veils the ideality, and the intellectual semblance is so faint that the observer
can hardly discern whether it is the thing's concept or his own that he has to do with. This
opposition of the ideal and the real, stemming from the moment $dm$, is objectivity's opposition
in power; but the observer objectifies in powerlessness.

The real concepts' transition from the ideal into the real medium is grounded in an omnipotent
objectification. The objectification is double, inward and outward. Insofar as a determinate system
of world-concepts is thought through in self-power, the motion is logical and the act of
objectification ideal; insofar as the same system of concepts is translated into object-power, the
motion is --- now \emph{here}, now \emph{there} --- antilogical, and the act of objectification
real: the real motion has the ideal for its presupposition.

% The Greek letter is read on the KB copy: γ, not „y“ or „7“.
\subrunin{γ) Reciprocal determination of ideal and real motion.} For real concepts to be formed, an
outward must unbrokenly reflect itself in its inward; reflexes of space and time, matter and force
with all the determinations of outwardness belonging to them, belong essentially to the formation
of the concept. There is no thinking in space, for space is no form for abstract thinking; and yet
when material actuality is to be conceived, a gleam of spatial relations must necessarily fall into
the ideal medium of conceiving. The moments of the real concepts are dissolved
% MISPRINT, not corrected on the book's own errata leaf: the Danish has „Forstillingerne“
% for „Forestillingerne“ --- a dropped e. The word „Forestillinger“ stands correctly two
% words earlier in the SAME line, so the difference is immediately measurable; both copies
% have the fault. Translated by the sense.
representations, and the representations dissolved images; there is in every image an intuiting,
however
% --- p. 113 ---
% MISPRINT: the Danish has „Rummct“ for „Rummet“ --- a c set for an e. Read on the KB copy
% under high magnification and calibrated against the certain e's in „Materien“, „er“,
% „ingen“ in the same line, which all show the horizontal cross-stroke this letter lacks.
% Both copies have the fault. The errata leaf does not correct it.
fleeting, of space. Matter is no thought, but thought's object, its resistance; the material in
matter is not thought, it is conceived as thought's other in a reflexion of intuiting and thinking:
the image of sense, whose significance is that thought's other has existence outside thought, brings
the reflex of the material with it. There are accordingly in the ideal motion of the real concepts
indications everywhere of real motion.

Real motion now has what the ideal presupposes: an actual matter in actual masses, an actual space
with determinate places, an actual time within given limits. The two heterogeneous motions cannot,
to be sure, be parallel; the intellectual way from concept to concept is incomparable with the
outward way from place to place; but that real motion is impossible without ideal motion is the
pulsing central thought of the philosophy of nature.

Ideal motion belongs to eternity; real motion goes on in time. The great world-system of mechanical
concepts of cause does not come to be in time; it is eternally finished. Infinite conceiving is
simultaneous, and the succession inseparable from ideal motion is only the indication of time
reflected in eternity itself. The real causes in which the concepts are realized belong by contrast
to time as time, and the motions go on with different velocity in space. In the eternal sense all
motions are determined in advance, for the ideal must be completed in itself before the real can
come to be
% The book's own errata leaf corrects p. 113, line 8 from the foot: „blive til --- læs
% blive til.“ --- Nielsen adds the missing full stop. Per the house rule the text stands
% here AS PRINTED, and so without the stop; the correction is not carried out.
Motion along the diagonal of the parallelogram of velocities must be completed in the concept, for
real motion can only in an imitative manner express what is originally in the concept. The laws,
self-power's thoughts, which are eternally thought and have been thought, are determined by
\emph{ideal}, but the fulfilment of the laws by \emph{real} motion.

% Run-in head at the FOOT of p. 113, set in bold antiqua. The signature „8“ stands beneath
% it and is not text.
\runin{b) The cause of the parcelling out.} In the reciprocal determination of
% --- p. 114 ---
ideal and real motion lies this, that the mechanical distribution of masses and forces in space
must be referred to a ground-cause working through outward causes. From the ideal side the
ground-cause is to be sought in the all-determining self-power, from the real side by contrast in
object-power, the power of existence answering to outwardness. If now the emphasis falls
particularly upon self-power, then the fundamental determinations point to nothing but acts of
will, and the necessity uttered in the connexion rests upon \emph{arbitrariness}. If conversely the
emphasis falls exclusively upon object-power, then the division proceeds from particularly existing
causes, and the necessity of connexion rests upon \emph{accidentality}. Only by absolute self-power
being known to have all power in another as the reality of its unconditioned supremacy can the
moment of arbitrariness from inner and the moment of accidentality from outer causality be
reconciled with necessity; and the plan of wisdom in the given world-order then presents itself as a
work of divine \emph{freedom}.

% The Greek letter is read on the KB copy: α, not „a“ or „&“.
\subrunin{α) Arbitrariness.} Every original determination, be it otherwise of whatever kind, is a
determination of the self, a self-determination. If the self-determinations be apprehended in part,
singly, they take on the stamp of arbitrariness. Yet it is easily seen that a plurality of
determinations of arbitrariness must either fall apart from one another, contradict one another and
lose themselves in the meaningless, or by mutual agreement presuppose one another, lie in a plan and
form a unity of connexion.

By means of the connexion the determinations of arbitrariness come under necessity: the arbitrary
becomes necessary, and the necessary arbitrary --- in the infinite reflexion. But in the finite
reflexion it looks as though arbitrariness and necessity fell apart from one
% --- p. 115 ---
another, as though the arbitrary belonged to one part and the necessary to another part of the plan.

Let it lie, for example, in the plan of omnipotence that two planets in world-space shall stand to
one another in such a way that the one furnishes a fixed centre about which the other moves at a
certain distance: then certainly, when it is omnipotence that absolutely wills it, each of the
planets can be arbitrarily posited for itself, the one be given a motion along the tangent, inertia
be added, and a law of gravity be fixed so as to preserve the relation; here in this parcelling out
there are no fewer than five omnipotent acts of arbitrariness. But now comes necessity.

The acts of arbitrariness presuppose one another mutually in the plan itself. Rational mechanics
maintains with full right that to the solution of such a problem there belong equations of
condition,\footnote{% printed mark: *)
% This note begins on p. 115 and continues beneath the text on p. 116; the mark is NOT
% repeated on the continuation. The whole note stands here at its mark, as the house rule
% prescribes.
If the forces be resolved along three perpendicular axes into X, Y and Z, and the running coordinates
be denoted by x, y and z, then the equation $X : Y : Z = x : y : z$ states a condition that must
necessarily be fulfilled if the resultant of the forces is to pass through the origin of the
coordinates and have in it a fixed centre. Since $\dfrac{d^2x}{dt^2} = X$,
$\dfrac{d^2y}{dt^2} = Y$ and $\dfrac{d^2z}{dt^2} = Z$, the equations
$\dfrac{xdy - ydx}{dt} = a$, $\dfrac{zdx - xdz}{dt} = b$ and $\dfrac{ydz - zdy}{dt} = c$ show that
all the constants a, b and c cannot possibly be zero if the moved body is not to plunge straight
toward the centre, but its radius vector to describe an area. From the equations it further follows
with necessity that the areas projected upon the planes XY, XZ and YZ must stand to one another as
the times employed, and that the orbit must be a plane curve, since the equation derived from the
presupposition, $az + by + cx = 0$, is precisely the equation of a plane passing through the origin
% „0“ is here the numeral nought, not the letter o; checked on the KB copy.
of the coordinates. Finally, since from the equations set up it can likewise be derived that
$\dfrac{dx^2 + dy^2 + dz^2}{dt^2} = \dfrac{ds^2}{dt^2} = 2F(x, y, z) + C$
(s denoting a piece of the orbit, and $F(x, y, z)$ the integral of $Xdx + Ydy + Zdz$), and that the
velocity $\dfrac{ds}{dt}$ --- as surely as $Xdx + Ydy + Zdz$ is a complete differential containing
neither the time nor the velocity --- depends solely upon the coordinates, it follows from this that
the living forces are maintained, that the moved body has the same velocity every time it comes to
the same place in the orbit, and that the motion is accordingly continued to infinity in the same
orbit.} which must inviolably be fulfilled, that
% --- p. 116 ---
% At the foot of p. 115 stands the signature „8*“; it is not text.
the principles of the conservation of areas and of living forces have, under the presupposition
given, a necessity from which not even omnipotence itself can withdraw.

He who wills the end must will the means; that means, in the language of infinity, that the end
is willed in the means and the means in the end. The conceiving self-power, which in unclouded
clarity sees the whole determined by the parts and conversely, chooses the conditioned for the
conditions' sake just as much as the conditions for the conditioned's sake. At no point in the
chosen connexion does the arbitrary fall outside the necessary. In infinite conceiving there can
accordingly be no thought of arbitrariness either weakening or breaking the laws of necessity;
but there may well be a danger here that the real causes should from this standpoint become
empty phenomena of power belonging to the ideal ones, that the abstract ideality of the
determinations of arbitrariness should undermine natural reality, that all-determining self-power
should in such a one-sidedness devour object-power.

% The sub-head is run-in and LETTERSPACED (the Greek α) β) γ) are letterspaced, while
% a) b) c) are semibold); the Greek letter is read on the KB copy: β, not „B“, „ß“ or „P“.
\subrunin{β) Accidentality.} The real cause coincides with object-power in corporeal existence.
If one now looks away from the question of ideal causality and takes things as they show
themselves, then the laws are laws of nature, and
% --- p. 117 ---
the conditioned necessity suffers no damage. Indeed it might almost seem as though the reciprocal
action between the natural causes furnished a surer guarantee of lawful necessity than a formal
reflexion of divine acts of arbitrariness. To keep to our example: the two bodies will now take
their place and attract one another not by divine order, but in virtue of their own material
nature. Every side-thought that it might perhaps one day please absolute self-power to transform
matter, impart new properties to it, introduce new laws, a new connexion with a new necessity ---
all such representations of arbitrariness fall away, and the factually given necessity is as
though built of solid masonry.

By this conception something is certainly won; the cause's mechanical reality is maintained, and
nature so far come into its right. But those who suppose that necessity itself, rational
necessity, would be better secured upon the real than upon the ideal foundation are remarkably
mistaken. Of the concept of necessity it does indeed hold that, in the character of a
selfhood-form, it has the conditions as moments in its essence; but of mechanically real necessity
precisely the opposite holds; as an otherness-form it has the conditions outside itself: the
existing members by which it is conditioned everywhere present a side of existence that lies under
accidentality. We learn it from our example. The two planets, of which the one moves about the
other, must indeed be acknowledged to stand in a real relation of necessity and to presuppose one
another in outward reflexion; but the one is nevertheless not the cause of the other: their being
together is, in the form of the parcelling out, a pure accidentality. If one now tries to lift the
accidentality by carrying both masses back to a common original mass, then the accidentality is
there again; it is in the mass's existence. If one lifts the accidentality of this existence by
letting the originally rotating
% --- p. 118 ---
mass come forth as a necessary result of atoms scattered in space, then the accidentality is there
again; it is in the atoms' existence.\footnote{% printed mark: *)
% p. 118 carries TWO footnotes, marked *) and **); the second is set after the pattern
% from p. 12 with a local \renewcommand. The book titles stand in ordinary antiqua, not in
% italic; „Über“ in the first title and „Ueber“ in the second are printed just so.
In his work \textit{Über die Natur der Cometen, Beiträge zur Geschichte und Theorie der
Erkenntniss}, 2\textsuperscript{te} Auflage, Leipzig 1872, J.\ C.\ F.\ Zöllner --- under the heading
\textit{Ueber die Endlichkeit der Materie im unendlichen Raume} --- has by reasoning and calculation
sought to show that, since the assumption of an act of creation is an \emph{arbitrary} limitation of
the series of causality, it is --- on the presupposition that the elements of matter are subject to
the Newtonian and the Mariotte law --- necessary to assume a finite quantity of matter, existing
from eternity, in infinite space. --- However interesting these investigations are in their details,
and however willingly we concede that physics too may profit from ``starting out from universal
concepts,'' that the progress of knowledge ``shall not be hampered by inherited prejudices,'' it
still stands fast that a physical series of causes cannot possibly begin with an outwardly given, a
first member that is not caused. To avoid ``arbitrariness'' one throws oneself into the arms of
accidentality; finite matter in infinite space comes, from the Zöllnerian standpoint, to be taken as
the first, uncaused cause, the great accident in eternally groundless accidentality.}

% The Greek letter is read on the KB copy: γ, not „y“ or „7“. The head is run-in and
% letterspaced.
\subrunin{γ) Freedom.} Necessity in and for itself is a powerless being, a brooding night in the
womb of possibilities, too heavy and formless to help itself forward. To win shape and come to
standing, necessity must be supported from within by arbitrariness, from without by accidentality.
It is Spirit, self-power in eternal reciprocal action with object-power, that breaks the seal of the
possibilities and brings dark necessity to the light.%
\renewcommand{\thefootnote}{**)}%
\footnote{% printed mark: **)
On the concept of necessity, cf.\ \textit{Grundideernes Logik} II, pp.\ 253--275;
\textit{Grundid.\ Logik i kort Begreb}, pp.\ 78--83.}%
\renewcommand{\thefootnote}{*)}
Every complete actuality, every mechanical, relatively closed world-whole must originally be
determined from within as a totality of concepts
% --- p. 119 ---
satisfying the demand of selfhood. This is the plan, the arbitrarily formed but content-rich,
comprehending and pervasive universal concept in which the particular concepts reflect and determine
one another with rational necessity; this is the ideal cause of the parcelling out.

And yet here there is as yet no parcelling out. Actual parcelling points everywhere, in the whole as
in the single case, to real causes. With the real causes the facts come forth, with the facts
accidentality; accidentality is woven into necessity, and outwardly conditioned necessity, restless
and yet at rest, reposes upon accidental presuppositions.

If now there could be discovered in the whole universe a real cause so free of presuppositions, so
factual, that in its given existence it fell outside every determining concept of cause, then we
should have here a genuine representative of facticity's unconditioned power, eternally blind
accidentality. But so long as such a discovery has not yet been made, we can freely let ``modern''
science's ``free thought'' play at necessity with blind
% The two quotations checked on the KB copy: ordinary antiqua, NOT letterspaced.
accidentality, while we inviolably hold fast the fundamental proposition that every fact, as an
outward cause, must be determined from within by its concept of cause, its ideal cause.

Upon the original unity of ideal and real cause all knowledge of nature must be grounded.
Philosophical criticism has plainly shown that matter, as little as its properties and forces, admits
of having natural existence imparted to it by any power outside nature. This is a decisive thought
for the philosophy of nature; but the thought is misconstrued and distorted when it is used to defend
nature's unconditioned autonomy, its absolutely necessary self-grounding. The truth is that all
development of nature rests upon natural presuppositions: nature begins with nature, however far back
we go. But if this truth is not to
% --- p. 120 ---
be perverted, it must necessarily be understood along with it that the causes in nature are real, and
that every real cause is in essence determined by an ideal one. The cause of the parcelling out is as
little to be referred to an absolute necessity as to facticity's accidentality; the natural cause is
in object-power --- an all-embracing system of outward causes.

The system of natural causes is a world-plan uttering itself in a categorical language of power.
Arbitrariness, accidentality, necessity are one-sided concepts. Accidentality is cancelled when it is
seen that the
% MISPRINT: „insees“ for „indsees“ --- the d is missing. The same page sets „indsees“
% correctly three lines further down. Both copies agree, so the fault is the compositor's;
% the errata leaf does not touch p. 120. Translated by the sense.
accidental is because it \emph{shall} be; and so: accidentality is cancelled when it is reflected in
arbitrariness. Arbitrariness is cancelled when it is seen that the arbitrary is because it \emph{must}
be, that the closed system is necessary for the plan; and so: arbitrariness is cancelled when it is
reflected in necessity. Necessity is cancelled when, in the unity of connexion of the members, it is
reflected as the form for absolute causality's content. Absolute causality, ideal in self-power, real
in object-power, is overreaching self-determination, is freedom; and so: necessity is cancelled when
it is reflected in freedom.

% The run-in head c) is printed SEMIBOLD (not letterspaced), after the same pattern as a)
% and b) earlier in the book. This is not reproduced.
\runin{c) The first mover.} The relation between the motion of thought and motion in space has been
developed in a peculiar manner by Aristotle. Motion in nature is transition from possibility to
actuality, and conversely; motion presupposes two members, a moving and a moved, an actual being in
the form and a potential one in the matter: motion is eternal, as are form and matter. If motion had
a beginning, then prior to this beginning there must either have been a moving and a moved, or not.
If they have not been before, then they must have come to be, and a motion has then preceded the
first motion. If they have been before, then it is unthinkable
% --- p. 121 ---
that they should not have brought forth motion; in the contrary case an activity must have set in
whereby they came to move, and we should then in this case too have a motion prior to all motion.
Just as motion in a finite manner presupposes itself to infinity, so in the absolute sense it
presupposes a first mover that is not moved (τὸ πρῶτον κινοῦν ἀκίνητον ὄν); this first mover is
% The Greek is read on the KB copy. The printing uses the ϰ form of kappa and the ϑ form of
% theta; these are variants of the face and are reproduced with κ and θ.
the Godhead. Everywhere in the order of things the moving is itself a moved; but what is moved is
altered and cannot work any constant, uniform motion. The first mover must be unalterable; but
unalterable it can be only by excluding, according to its essence, the possibility of being otherwise;
and to exclude the possibility of being otherwise is to exclude matter. Matterless, incorporeal,
indivisible, outside space, without motion, passivity or alteration (ἀπαθὲς καὶ ἀναλλοίωτον), the
first mover is the perfect form, eternal actuality, pure energy --- in a word: God.

With this metaphysical concluding thought Aristotle remains standing on the border between two
mutually irreconcilable world-views, pantheism and theism.

In a theistic direction Aristotle maintains Spirit's absolute originality and lays a decisive weight
upon the opposition between God and the world. The concept of the Godhead coincides, in perfect
thinking, with the perfect good. As the eternal, perfect good the Godhead moves everything in being
desired by everything; for everything in nature is laid out toward the good as toward a highest
common end (πρὸς μὲν γὰρ ἕν ἅπαντα συντέτακται). As the good perfect in itself the Godhead can move
% „ἕν“ is printed with an acute, not with a grave; the mark is the same as the one over the
% α in „ἅπαντα“. The rejected reading is „ἓν“, which the editions have.
everything without itself being moved; for that which is desired and thought moves without itself
being moved. What \emph{shows itself} beautiful is desired; what \emph{is} beautiful is willed. We
desire the beautiful because it shows itself, but it does not show itself because we desire it. The
% „viser sig“ and the two-syllable „er“ are both letterspaced on the KB copy; spacing.py
% found the first as a RUN and, as expected, overlooked the second.
% --- p. 122 ---
end moves because it is loved, and what is moved by the end moves in turn the rest. Since the divine
consists in eternally completed energy, it consists in eternal living, for life is energy (ὥστε ζωὴ
καὶ αἰὼν συνεχὴς καὶ ἀΐδιος ὑπάρχει τῷ θεῷ). This life is blessed; such a joy is granted to us only
in brief instants; but all joy (ἡδονή) is a draught of the divine principle.

That Aristotelian theism nevertheless stands on weak feet, criticism has no difficulty in showing. A
Godhead that does nothing but eternally
% Three points in the Greek citation, read on the KB copy. The spacing is 15--16 px against
% a word space of 22--24 px in the same line: a BORDERLINE case between the book's two
% justifications, here carried to the close form.
think the thinking of thinking (αὑτὸν ἄρα νοεῖ .\,.\,.\ καὶ ἔστιν ἡ νόησις νοήσεως νόησις), and which,
so as not to stain its perfection with anything imperfect, cannot possibly have a thought to spare for
anything in the world, is an empty, natureless being whose absolutely necessary content of thought
resolves itself into pure tautologies.

Where the working of the self-subsistent God upon the world is treated, the turns of phrase are
ambiguous. There was, it is said, neither eternal chaos nor eternal night when motion began; but the
same has always been,
% These lines are merely loosely justified, not letterspaced; checked on the KB copy.
either in constant alternation or in another way, as surely as actuality is prior to possibility. If
now the same is always to work in constant alternation, then there must be something that always works
in the same manner; if there is to be arising and passing away, then another must incessantly work
otherwise and again otherwise. The first mover, which is always like itself, thus comes into a
questionable intercourse with what works otherwise and is always unlike itself; and yet the eternal
thinking of thinking was supposed to be untouched by everything else. What is said of the need and
desire with which things move toward the perfect one is only a teleological, figurative expression for
the negativity indwelling in everything finite, which shows that the relative is a moment for the
absolute.
% --- p. 123 ---
In a far more immediate manner thinkers such as Cartesius and Gassendi have made God the first mover.
However much these two philosophers otherwise disagree about the causes in nature, they both start
from a dogmatic separation between the absolute principle, the first cause of everything, the Creator,
and the secondary causes, among which is reckoned ``everything which besides God has a force to effect
and accomplish something in the world.''

The cause of motion, says Cartesius, ``is twofold; partly a universal and original one, which causes
% The quotation checked on the KB copy: ordinary antiqua, not letterspaced.
all the motion there is in the world at all; partly a particular one, which brings it about that the
single parts of matter obtain motions they did not have before. The universal cause is God, who
created matter both with motion and with rest, and who moreover by his natural co-operation maintains
all the motion and rest in matter that he laid into it from the first.''

What is characteristic of this way of seeing is that, starting from an abrupt opposition of the
supernatural and the natural, it is laid out so as to explain nature's coming-to-be by a finite act of
creation, and so by an act that is not natural. First matter is created; the Creator implants in it
certain properties, a certain quantity of motion, etc. With a created matter, with created forces and
properties, created nature is then to begin working in a natural manner. Such a beginning is, on the
presupposition set up, impossible. As the supernatural effect of a supernatural cause, created matter
with its properties and forces gets at the beginning a supernatural existence of its own; its first
motion is accordingly supernatural: whence are the natural causes with the natural motions now to
come? They cannot lie in matter's essence, for matter's created being is supernatural; just as little
can they be ascribed to the forces, since the forces are themselves likewise supernatural. As the
cause is, so must the effect be; cause and effect
% --- p. 124 ---
must, in spite of the separation in the outward, absolutely have the same essence. That created matter
with its properties and forces is wisely distributed by the Creator and becomes, by the supernatural
act of distribution, ``secondary causes,'' will then mean that the secondary causes are not natural
either.

The dogmatic assurance that a supernatural beginning must necessarily be presupposed here, though we
human beings are not able to derive the natural from the supernatural, is only a convenient evasion for
theologians, but signifies nothing in science. One of two things! Either
% „Videnskaben.“ with a full stop; both OCR witnesses read a comma, the KB copy settles it.
nature as nature has a supernatural beginning, and then it must necessarily continue as it began: all
so-called natural causes, effects and reciprocal actions must in being and working be supernatural; or,
if there really are natural causes and naturally necessary motions, then the beginning must answer to
the continuation: nature, natural in being and working, must have a natural beginning.

It is an unclarity resting upon the concept of omnipotence that from all sides breeds the
misunderstandings. While Aristotle apprehends the absolute in the form of thinking as a metaphysical
self-being without actual omnipotence, the absolute is in Gassendi and Cartesius represented as an
omnipotent self, but without object-power, without power in another; while Aristotle gets a first mover
which in actuality can move nothing and must alongside this let nature be nature and move itself on the
foundation of a motion smuggled in under a negative form, Gassendi and Cartesius make absolutely ideal
self-power into an immediate power of reality and presuppose an omnipotence which must consistently
devour the nature it itself brings forth.

The fundamental contradictions find, as we have seen, their solution in power's original self-doubling.
The mechanical causes in nature belong to object-power and move with
% --- p. 125 ---
% p. 125 opens with the three last lines of the SECOND chapter, division C. They continue
% without a new paragraph from p. 124.
\emph{necessity}; but the ideal causality, the inner determining of the natural causes, belongs to
self-power, the first mover, which moves with \emph{freedom}.

% Short centred rule between the second and third chapters. It is printed as a DOUBLE line
% in both copies; neither the length nor the doubling is reproduced --- \streg serves
% throughout, per the house rule.
\streg

% The chapter head DOES NOT OPEN A PAGE: it stands in the middle of p. 125, beneath the
% close of the second chapter and above division A.

% Danish: \kapitel{Tredie Kapitel.}{Himlens Mechanik.}   (printed pp. 125--196)
\kapitel{Chapter Three.}{The Mechanics of the Heavens.}

% Danish: \afdeling{A.}{Verdensbygningen.}   (printed pp. 125--146)
\afdeling{A.}{The Structure of the World.}

There is in the world nothing objective except what is objectified: space and time, matter and its
forces, everything in outward nature, in the realm of power, has objectivity --- not immediately of
itself, but from Spirit's omnipotence. From the standpoint of finitude it does indeed hold that
visible things are in their outward actuality independent of the eye by which they are seen;
% „sees“ without an accent, read on the KB copy; both OCR witnesses give „sées“, which is
% the working scan's speckle and not the type.
but if the question is of knowledge's truth, then the law of objectification teaches that the
validity which can belong to the object as an existing object is dependent upon the validity which,
according to scientific insight, can belong to the object's objectification. The objectivity of the
world-structure is a magnificent and eloquent example of this. The starry heaven is indeed an object
of immediate contemplation, and it might so far seem as though everyone who looks up toward the high
vault must involuntarily see in what style this building is raised. And yet it is an illusion; the
unsurveyable building with its many rooms does not stand open to naively
% --- p. 126 ---
curious spectators. With earthly buildings it is the actual erection that costs labour; the
contemplation comes of itself. To get the heavenly building raised costs no labour; that work is
completed; but to see into the building and, if only approximately, to acquaint oneself with the plan
in its peculiar construction is conditioned by labour. It has cost the human spirit millennia of
labour to get the world-structure objectified in all the main features answering to actuality; this
ideal work of building is of infinitely greater worth than all the material toil expended on Egypt's
pyramids. We follow the progress of the objectification from \emph{the stage of illusion} through
\emph{reflexion to the final testing}.

% Run-in head, read on the page itself (KB copy) and not in the Indhold: set in semibold, on
% the same line as the paragraph it opens. The distinction between semibold and letterspaced
% is not reproduced.
\runin{a) The illusion: the Ptolemaic world-structure.} Antiquity's consciousness of nature is
essentially mythical. What is brought to the senses from without and must be referred to outward
nature melts involuntarily together with what comes from within, from the ground of the mind, under a
mood-laden co-operation of feeling, fantasy and thinking, and transforms the outward; how it comes
about that the world-picture and the actual world, in spite of their mutual difference, meet and fuse
in consciousness itself --- this secret of objectification is not there for a naive contemplation. The
sense of sight is objective, but by reason of an illusion; in the apprehension of the visible heaven
the illusion rules so much the more as the earthly standpoint is fixed of itself, and none of the
other senses reaches with its testimony out into infinite space. It is then a matter of course that
the first picture of the world-structure must be illusory. But the illusion is not a simple deception.
The revolution of the celestial sphere about the earth, the sun's rising and setting, the moon's
particular motion and changing appearance --- all such things have from an earthly station their truth
just as surely as
% --- p. 127 ---
the separation of light and darkness, morning and evening, day and night.

The Ptolemaic world-structure rests upon ancient foundation stones; the first material was gathered at
a time when a sustained beholding of the phenomena in the heaven was less a concern of science than of
religion, of popular faith. The seven planets --- the moon, Mercury, Venus, the sun, etc. --- were of
course not regarded as lifeless masses of rolled-together matter, but as high powers in radiant
splendour, as gods. In this illusion Spirit's unity with nature was held fast. The gods wandering
across the heaven did not merely, by their measured motions, give men a standard for calculating civil
time; they were themselves the rulers of time. In order to survey events and govern earthly affairs
these lords of time divided time among them; it was then no accident that each of them got his own day
in the week.\footnote{% printed mark: *) --- the only note on p. 127.
By arranging the planets in the following series --- Saturn, Jupiter, Mars, the sun, Venus, Mercury and
the moon --- and starting from the presupposition that each of them has indeed its hours in the day,
but in such a way that the one that reigns in the first hour is overlord and gives the day its name,
one will understand --- in the sense of the illusion --- how the seven days of the week can be
subjected to the seven planet-gods. If the first hour of the day begins, for example, with Saturn, then
we have a Saturday, \textit{Dies Saturni}. By sharing the dominion with the six co-regents Saturn will
further come, in his own person, to rule the eighth, the fifteenth and the twenty-second hour; then
Jupiter will get the twenty-third, Mars the twenty-fourth, and the regency in the first hour of the
following day falls to the sun, so that Sunday becomes the sun's day, \textit{Dies Solis}. As it went
with Saturn, so it now goes with the sun; when it has last ruled in the twenty-second hour, Venus rules
in the twenty-third, Mercury in the twenty-fourth, so that the first hour of the following day falls to
the moon and becomes a Monday, a moon's day, \textit{Dies Lunæ}; and so forth.}

% --- p. 128 ---
If the illusion is the naturally necessary form under which the human spirit, at the stage of naive
intuition, avails with the gods' assistance to see a reflected splendour of its objective essence, then
it cannot surprise us that antiquity was not only involuntarily dazzled by the illusion, but we
understand that it also, in an arbitrary and yet instinctive manner, half consciously, half
unconsciously, held the illusion fast as the truest form for its richest and deepest content. Thus the
constellations are by no means formed by merely arbitrary grouping with arbitrary names; and just as
little can it be supposed that the ancients really fancied they saw the heroes of the past alive among
serpents and monsters in the heaven. The arbitrary and the involuntary melt together in a need of the
mind for exalted symbolism, for an illusory form for Spirit's self-objectification in nature. The
constellations are, as J.\ L.\ Heiberg has so aptly expressed it, ``true pictures which antiquity has
hung in the great gallery where all generations shall behold them. Here it rightly shows itself that
the starry heaven, far from being foreign to the earth, is a reflex of earthly destinies, earthly
thoughts and dreams.''

The ground-plan for the Ptolemaic world-structure was supplied by Greek philosophy; it is a plan that
lies open to the illusion from first to last. The fantasy filled with the world-idea, the beautifully
rounded cosmic whole, knew how to form the stuff-like content to thought's liking so ingeniously that
the phenomena of the illusion received a most profound rational grounding.

In mythically ideal images Plato's \textit{Timaeus} describes how the god, the highest master-builder,
arranged everything from the beginning for the best, in forming the world-soul into that intermediate
being in which reason and sensibility, idea and corporeality alone admit of being united. Everything
corporeal is moved by another; only the soul, self-motion's own
% --- p. 129 ---
% The Greek is read on the KB copy under strong magnification. The circumflex in „κινεῖν“
% is in the printing set over ε, not over ι; that is the face's habit throughout the book's
% Greek and cannot be reproduced in Unicode, where the mark must sit on ι.
force (ἡ δυναμένη αὐτὴ αὑτὴν κινεῖν κίνησις), moves itself and everything. Even before the corporeal
elements were formed, the god mixed the indivisible substance, like to itself, and the corporeally
divisible in such a way that a third being came forth, namely the substance of soul, in which the
self-like and the other are joined together. This whole was now divided, according to the fundamental
numbers of the harmonic and astronomical system, in such a way that the circles of the fixed-star
heaven and of the planets came forth by a division of length of its own; for all relations of number
and measure are contained in the soul, its essence is nothing but number and harmony. Plato lets the
seven parts of the world-soul stand to one another as the numbers 1, 2, 3, $2^2$, $3^2$, $2^3$, $3^3$;
upon this division rests the planets' distance from the world's centre; the sun is 2 times, Venus 3
times, Mercury 4 times, Mars 8 times, Jupiter 9 times and Saturn 27 times as far from the earth as the
moon. As the soul, in the circular motion, returns to itself, it says through its whole being
everywhere to itself what it is that it touches in its revolution; whether it is something divisible or
indivisible, with what it is of one kind, from what it is different, and how it in general stands in
its becoming and being. This speech propagates itself soundlessly in the self-moving and begets
knowledge. From the circle of the other, the ecliptic, proceed right representations and opinions; from
the fixed-star circle, the equator, the circle of the self-like, stem rational knowledge and
understanding knowing. The illusorily profound central thought is that the harmoniously moved
world-whole is a living, ensouled being.

Aristotle takes a step forward. He disdains the mythical clothing; he subjects Plato's doctrine of the
world-soul to a strict criticism and holds to actuality. One should now suppose that the Aristotelian
building, freed from the dazzling light-mists of myth and illusion, would have to stand fast on the
ground of reason, and the actual
% --- p. 130 ---
outlines keep themselves in the clear day of science. And yet the whole objectification is a complete
illusion. How little human reason is able, by a few general observations, to penetrate so objectively
into nature that it can in virtue of its own world-idea construct the actual world --- of this
Aristotelian physics is a grand example. Everything from first to last has this philosopher reasoned
through, grounded, defined, demonstrated; the system, worked through with enduring acuteness, makes the
impression of a comprehensive, irrefutable necessity of thought; but the necessity rests upon an
illusion.

How Aristotle has thought the first mover has been shown in the preceding chapter; here it shall now be
touched upon briefly how he apprehended the motions in the corporeal world. All motion in space is
either rectilinear or circular or composed of both. Rectilinear motion is, according to the nature of
the elements, again either \emph{toward} or \emph{away from} the centre.
% The letterspacing of „mod“ and „fra“ is confirmed on the KB copy; the later „mod
% Midtpunktet“ in the same paragraph is by contrast not letterspaced.
The four elements move along the straight line --- earth, which is absolutely heavy, toward the centre;
fire, which is absolutely light, away from the centre; water and air form transitions; but circular
motion belongs to the aether. Since now every element has its natural motion and its natural place, the
fundamental shape of the world-structure is determined with necessity. Immovable, in the middle of the
whole, rests the spherical earth. It is impossible for the earth to move; it conflicts with its nature.
The earthly can move only insofar as it moves toward the middle; but the terrestrial globe is the
middle; every body comes to rest when it has reached the place to which its motion carries it:
everything earthly accordingly has its rest in the middle. If it stands fast that the earth with its
sphere of air and fire is spherical, it follows that the heaven surrounding the earth must also be
spherical.
% --- p. 131 ---
But it is seen besides on other grounds. The heaven must, namely, as the most perfect, have the most
perfect form; but the most perfect corporeal form is the sphere. As the measure of all motions the
heaven must move with the greatest velocity, and so follow the shortest way; but the shortest way from
the same to the same is the circle. All heavenly bodies move daily from east to west; but seven among
them move besides, in shorter or longer periods, from west to east. Since it is now settled that the
earth rests, either the whole heaven must move, or only the stars, or both the heaven and the stars. But
that both the heaven and the stars should move is unthinkable; for how should it come about that the
stars kept pace in velocity, as of themselves, with the velocity in their own circles? That the stars
should move and their circles rest is also unthinkable; it would have to be an inexplicable accident if
a star's motion should answer exactly to the size of the circle. There remains, then, only to assume
that the circles move while the stars fixed upon them rest (τοῦς μὲν κύκλους κενεῖσθαι τὰ δὲ ἄστρα
ῆρεμεῖν).
% The Greek citation (Aristotle, De caelo II, 8) is read on the KB copy under strong
% magnification and is reproduced AS PRINTED. It departs at three points from the correct
% form (τοὺς μὲν κύκλους κινεῖσθαι, τὰ δ' ἄστρα ἠρεμεῖν): τοῦς with a circumflex over υ;
% κενεῖσθαι with ε in the second syllable; ῆρεμεῖν with a circumflex over η. No comma is
% printed between κενεῖσθαι and τὰ.
The famous astronomer Eudoxus of Cnidus had already sketched a system with 27 spheres, of which 26 fell
to the planets. To these Callippus added a further 7 spheres, so that sun and moon each got two, and
Mercury, Venus and Mars one each. From the manner in which Aristotle takes up and grounds the intricate
theories of spheres it becomes quite evident that he, far from cancelling the illusion, on the contrary
works it so deeply into the thinking consciousness that nearly two millennia had to pass before science
got it purged out from the ground up.

On an illusory foundation an understanding objectification can meanwhile go on upon the prudent
condition that the observer will not undertake to give grounds for everything, but holds to facts and by
empirical induction
% --- p. 132 ---
sets about finding out the laws of motion. Herewith scientific astronomy begins.

The understanding objectification now runs through three principal stages. The first thing necessary is
to draw circles and mark off points, whereby it becomes possible to orient oneself on the celestial
sphere. To this answer the familiar great circles --- the horizon with zenith and nadir, the equator with
the poles of the world, the ecliptic --- in short, what belongs, from different points of view, to
indicating a star's coordinates, that is, determining its spherical or angular place. Although such
points, poles and circles have no material actuality, they have nevertheless an undeniable validity for
intuition and thinking; they have an essential objectivity.

As the sum of the first conditions of objectification is gradually widened, and continued observation of
the phenomena of motion leads to ever more exact and more special results, it will dawn upon consciousness
that the world-harmony, the wonderful connexion of the whole, is far more intricate than one had
represented it from the beginning. Here a contradiction sets in between the original presupposition of the
earth as the centre of all circles, of the heavenly bodies' even, uniform motion, etc., and the manifest,
at least apparent irregularities which the observations display; the endeavour to solve this great problem,
indicated by the contradiction, is characteristic of the next stage. If the sun's orbit is a circle and the
motion even, how can it then come about that the seasons are not of equal length, that the sun takes a very
different time to traverse each of the ecliptic's four quadrants? To explain the phenomenon Hipparchus set
up two hypotheses which, however --- as Ptolemy afterwards saw --- are different only in form, not in
actuality. According to one of these hypotheses the sun describes with even velocity a circle whose centre
lies outside the earth;
% --- p. 133 ---
according to the other the sun is ascribed an even motion on a circle --- the epicycle --- whose centre
itself moves on a circle --- the deferent circle --- which has the earth for its centre.
% Misprint: between „Centrum“ and „selv“ both copies show a vertical inked mark which has
% neither the comma's hook nor its position below the line --- probably a risen space that
% has taken ink. Reproduced as no mark; tesseract reads a comma.
Hipparchus is the real founder of the epicycle theory. Finally --- and this is, as regards method, the
last, completing step of the objectification --- the whole theory of the planets is brought to a relative
conclusion by the continued application and extension of the epicycle hypothesis. Precession, the retrograde
motion of the equinoctial point, was shown by Hipparchus; the unevenness in the moon's orbit which
afterwards got the name of evection was discovered by Ptolemy, whose \textit{Μεγάλη σύνταξις} set the
garland upon the ancient edifice of doctrine. In illusory objectivity, with a few Arabic repairs, the
building stood unshaken for fourteen hundred years.

% Run-in head, read on the page itself (KB copy): semibold.
\runin{b) Reflexion: the Copernican world-structure.} Fantasy can breathe in thinking, set thoughts in
motion and reflect; conversely, thinking can build upon an intuition and, when the intuition is hovering,
fantasize. Illusion and reflexion form in and for themselves no sharp opposition. There can, as we have
just seen, be much reflexion brought to bear in developing and grounding an illusory theory, and likewise
there can, as we shall see in what follows, be illusions of various kinds that creep into an otherwise
thoroughly reflected fundamental intuition answering to actuality. But when the talk is of what on the
whole characterizes a standpoint of objectification, the difference becomes recognizable.

There is a deep-going difference between a view of the world whose starting point belongs unconsciously
and straightforwardly to the illusion, and an investigation of the connexion of the phenomena which begins
by subjecting the foundation of the view, the first, general presupposition, to a strict scientific
testing. In this sense we can, as regards
% --- p. 134 ---
the doctrine of the world-structure, rightly say that Copernicus stands in the history of astronomy as a
heroic figure, insofar as he, with an opened eye for a higher world-harmony, is the first who had the
superiority of genius and the intellectual courage to break with the power of habit, place himself on the
standpoint of reflexion and combat the illusion. ``The authors,'' he says, ``are usually so agreed that the
earth rests in the middle of the world that they consider it meaningless, indeed even laughable, to assume
the opposite. But the question is not yet settled and cannot be dismissed with
% The book's habit: the quotation is re-opened with „ after the interpolation „siger han“.
% Here, however, it is also closed after „Forfatterne“, so the marks balance on the page
% (two „ and two “). Reproduced as printed.
contempt.''\footnote{% printed mark: *) --- the first of THREE notes on p. 134.
De Revolutionibus orbium cœlestium. Norimbergæ MDXLIII. Lib.\ I.\ Cap.\ V.}

To choose one's standpoint in the sun, when all the adherents of tradition are agreed on keeping to the
earth, must in many eyes look like a whim, a foolhardy notion of a high-flying fantasy. But Copernicus will
not fantasize; he will reflect in earnest. It is the actual relation between earth and heaven --- a
relation that only a carried-through reflexion, not an immediate judgment, can appraise --- that he will
have investigated from the ground up.%
\renewcommand{\thefootnote}{**)}%
\footnote{% printed mark: **) --- the second of three notes on p. 134.
Qvam ob causam ante omnia puto necessarium, ut diligenter animadvertamus, qvæ sit ad cœlum terræ habitudo,
ne dum excelsissima scrutari volumus, qvæ nobis proxima sunt, ignoremus, ac eodem errore qvæ telluris sunt
attribuamus cælestibus. De Revol.\ Lib.\ I.\ Cap.\ IV.}%
% The compositor mixes the ligatures: „cœlum“ with œ but „cælestibus“ with æ in one and the
% same note. Read on the KB copy. Reproduced as printed.
\renewcommand{\thefootnote}{*)}
It is the difference between apparent and actual motion, and so the objectivity of motion, upon whose
determination it turns.%
\renewcommand{\thefootnote}{***)}%
\footnote{% printed mark: ***) --- the third of three notes on p. 134.
% THIS NOTE RUNS OVER THE LEAF: it begins beneath the text on p. 134 and is continued at the
% top of the note field on p. 135 without the mark being repeated. Per the house rule the
% whole note is set here at its mark; the word-break „ali-/qvis“ is resolved.
Omnis enim, qvæ videtur secundum locum mutatio, aut est propter spectatæ rei motum, aut videntis, aut certe
disparem utriusqve mutationem. Nam inter mota æqvaliter ad eadem non percipitur motus, inter rem visam dico
et videntem. Terra autem est, unde cœlestis ille circuitus aspicitur, et visui reproducitur nostro. Si
igitur motus aliqvis terræ deputetur, ipse in universis, qvæ extrinsecus sunt, idem apparebit, sed ad partem
oppositam, tanqvam prætereuntibus, qvalis est revolutio cotidiana imprimis. Lib.\ I.\ cap.\ V.}%
\renewcommand{\thefootnote}{*)}

% --- p. 135 ---
That the heaven's daily motion from east to west can be explained as a deceptive rendering of the earth's
actual motion from west to east, Copernicus is so far from wanting to give out as any new thought that he
on the contrary recalls that among the ancients, for instance, the Pythagoreans Heraclides and Ecphantus
and Nicetas of Syracuse\footnote{% printed mark: *) --- the only note of p. 135's own.
Should no doubt be Hicetas.} were, according to Cicero, of this opinion. But, he adds, when this is
conceded, then another no lesser doubt follows about the earth's place, although it is now assumed by almost
everyone that the earth is the middle of the world. From this doubt reflexion then begins in such a way that
the rejected way of seeing is set against the authorized one, and the grounding happens by reasoning against
reasoning. Ptolemy finds it unreasonable that the earth should move. By the rapid rotation the connexion of
things would have to be dissolved, all earthly bonds burst --- indeed, what is quite a laughable
representation, the heaven itself would have to fall down. But, asks Copernicus, why does one not rather
apply this to the world, whose motion must be so much the more rapid as the world is greater than the earth?
Has the heaven perhaps become so infinitely great because it is driven away from the middle with
indescribable velocity (\textit{ineffabili motus vehementia})? In truth, if this reasoning holds good, the
heaven's greatness must be infinitely increasing; the velocity must bring forth the greatness and the
greatness again the velocity, without measure or limit. But according to the physical axiom that what is
infinite can neither be overstepped nor moved, the heaven must necessarily stand still (\textit{stabit
necessario cœlum}). But when the heaven stands
% --- p. 136 ---
still, it must be the earth that moves. If the earth were the world's motionless middle, then it
would have to be the common centre for all the orbits; but that it is not the fixed common centre
comes to light both from the manifest unevenness in the planets' motion and from their changeable
distances from the earth. Since there now demonstrably exist several centres, there is every ground
to look about for another world-centre, another centre of
gravity.\footnote{% printed mark: *) --- the first of TWO notes on p. 136.
Pluribus ergo existentibus centris, de centro qvoqve mundi non temere qvis dubitabit, an videlicet
fuerit istud gravitatis terrenæ, an aliud. Eqvidem existimo, gravitatem non aliud esse, qvam
appetentiam qvandam naturalem partibus inditam a divina providentia opificis universorum, ut in
unitatem integritatemque suam sese conferant in formam globi coëuntes. Qvam affectionem credibile
est etiam Soli, Lunæ, cæterisque errantium fulgoribus inesse, ut ejus efficacia in ea qua se
repræsentant rotunditate permaneant, qvæ nihilominus multis modis suos efficiunt circuitus. Lib.\ I.\
cap.\ IX.}
% The note mixes qv and qu: „qvoqve, qvis, qvam, qvandam, qvæ“ against „integritatemque,
% cæterisque, qua“. Read on the KB copy and reproduced as printed.
All the planets' testimonies (\textit{tot tantaqve errantium syderum testimonia}) agree
(\textit{consentiunt}) that the earth moves forward in space. Copernicus ascribes to it a threefold
motion: a \emph{daily} one about its axis, an \emph{annual} one about the sun through the ecliptic,
and --- superfluously%
% The letterspacing of „daglig“ and „aarlig“ is confirmed on the KB copy; the following
% „til“ is by contrast not letterspaced.
\renewcommand{\thefootnote}{**)}%
\footnote{% printed mark: **) --- the second of two notes on p. 136.
As Rothmann had already noticed.}%
\renewcommand{\thefootnote}{*)}
--- yet another motion, whereby the axis's inclination remains constantly parallel
% The errata leaf corrects p. 136, line 14 from the top: „Hældning --- læs Heldning“. Per
% the house rule the text stands AS PRINTED.
with itself (\textit{tertius declinationis motus}).

We see the ancient philosophers starting from the fundamental consideration that bodies moving with
equal velocity must seem to move the more slowly in proportion as they are at a greater distance ---
that the moon, with the smallest orbit, therefore has the shortest period of revolution, and Saturn,
with the greatest orbit, the longest. But about Venus and Mercury opinions are divided. Whether one
places them with Plato above the sun or with Ptolemy below, the matter is equally
% --- p. 137 ---
confused. As a ground why Venus and Mercury ought to be placed below the sun, it has been urged that
otherwise there would be an all too enormous space between the sun and the moon. That is a most
unreliable ground; for precisely the measures of distance upon which the assumption is to be
grounded show that the epicycle of Venus, if the earth is to be the deferent's centre, encloses a
space so great that this empty space could more than hold the earth, the air, the aether, the moon
and Mercury together. Another argument is still unluckier. Ptolemy considers it fitting that the sun
should get a middle place between the planets which move away from it on all sides
(\textit{omnifariam}) and the planets which do not; but how false the proof is can be seen from the
moon, which in this case would have had to have its place above the sun. Accordingly, Copernicus
concludes, either the earth is not the centre, or the order in which one arranges the planets is
without significance.\footnote{% printed mark: *) --- the first of TWO notes on p. 137.
Oportebit igitur, vel terram non esse centrum, ad qvod ordo syderum orbinmqve referatur: aut certe
rationem ordinis non esse, nec apparere cur magis Saturno qvam Jovi seu alii cuivis superior debeatur
locus. Lib.\ I.\ cap.\ X.}
% Misprint in the note: „orbinmqve“ for „orbiumqve“. The n is unambiguous on the KB copy
% under strong magnification, so the fault is the compositor's. Reproduced as printed.

To change the geocentric standpoint into a heliocentric one is possible only for one who has become
quite clearly conscious of the difference between appearance and actuality. From the beholder's actual
standpoint all motions are apparent; only from an unactual, reflected-to standpoint are the motions
actual; it is a reversal which only a superior reflexion has been able, at first hand, to grasp and
carry through. Whether the Pythagorean Philolaus%
\renewcommand{\thefootnote}{**)}%
\footnote{% printed mark: **) --- the second of two notes on p. 137.
Cf.\ De Revol.\ Lib.\ I.\ cap.\ V.}%
\renewcommand{\thefootnote}{*)}
or Aristarchus of Samos were of the opinion that the earth was a planet makes in this connexion no
difference; for what it turns upon is an astronomical-geometrical reflexion which, with due insight
into
% --- p. 138 ---
the manifold of entangled relations, avails to carry the appearances back to actuality and then again,
starting from actuality, to get the apparent exactly explained and grounded on all sides in all its
details. It is this problem that Copernicus, in relation to the astronomical knowledge of his age,
completely solved. An exalted central thought led him, an idea about which his objectifying fantasy was
from first to last at one with his reflexion: the idea of an ``admirable symmetry of the world''
(\textit{admiranda mundi symmetria}) and ``a harmonious connexion between the motions and the size of
the orbits.'' The fixed-star heaven, ``the highest of all,'' is immovable, for it encloses everything;
but the planets Saturn, Jupiter, Mars, the earth with its moon, Venus and Mercury all move about the
common centre, the sun. ``Who could in this beautiful temple set this lamp in another or better place
than there, whence it can at once irradiate everything? Some call it, quite aptly, the world's lantern,
others its soul, others its governor; Trismegistus a visible god, Sophocles's all-seeing Electra. Yes,
thus the sun sits as on a royal throne, governing the whole family of encircling heavenly
bodies.''\footnote{% printed mark: *) --- the only note on p. 138.
De Revol.\ Lib.\ I.\ cap.\ X.}

Copernicus has not enriched astronomy with any new art of observation or with the discovery of new
particulars; his discovery belongs not so much to the facts as to reflexion; it is a secret of
objectification that he has discovered; by this discovery he has transformed not the world, but our
consciousness of the world: he has, for our knowing, remade heaven and earth.

% Run-in head, read on the page itself (KB copy): semibold, and here genuinely run-in --- the
% text continues on the same line. The head stands at the foot of p. 138, not on p. 139.
\runin{c) The discussion: three hypotheses.} The Ptolemaic system was dogmatic; dogmatism, with its
naive reliance
% --- p. 139 ---
on the objective validity of human objectification, is on good terms with illusion. About particulars
doubts might arise, but the foundation as a whole stood fast, for it rested upon immediate intuition. The
Copernican system is thetic in the form of reflexion, and so far dogmatic too, though in a way of its own.
The senses' immediate testimony is indeed in each and all subordinated to the understanding's testing; but
the standpoint reflexion has chosen must for the time being hold to the dogma of harmony. Even though
Copernicus now and then expresses himself with caution, as though it were enough for him if the reader
merely found his view probable, or at any rate not unreasonable, it is nevertheless clear from the whole
impression that he himself considers it irrefutable; the world-structure he has raised is for him not one
of several possible ones, but the only possible one, because it is the only reasonable one.

Now it lies in the nature of the case that two mutually contradictory theories, both coming forward with
equal title, since both of them, each in its own way, satisfy the observations in about equal degree, must
make one another problematic. But when the foundations are problematic, then either one of the theories
must be laid down as an article of faith, or both must remain standing as hypotheses.

% „charachteristisk“: so printed, read on the KB copy (ch...cht). The book elsewhere writes
% „charakteristisk“ (pp. 142, 145); the errata leaf does not correct it.
The unnamed author of the preface to the work \textit{De Revolutionibus} has expressed himself in this
respect in a quite characteristic manner. Since it is impossible, he says, to discover the true causes,
the astronomer has only to devise and invent such presuppositions or hypotheses as can give him a
convenient and satisfactory foundation for the calculations. In several turns of phrase it is then put
before the reader most emphatically that the Copernican theory neither is nor will be anything but a plain
and simple hypothesis. It is surely more than considerations of
% --- p. 140 ---
prudence that stamps this presentation; it is honest
% The note's Latin is set in upright antiqua, NOT in italic --- checked on the KB copy.
% The errata leaf: p. 140, line 2 from the foot: „quiqvam --- læs quidqvam“. Per the house
% rule the text stands as printed (quiqvam). The ellipses stand in the printing as FOUR
% spaced points, so reproduced.
conviction.\footnote{% printed mark: *)
Est enim Astronomi proprium, historiam motuum cælestium diligenti et artificiosa observatione colligere.
Deinde causas earundem, seu hypotheses, cum veras assequi nulla ratione possit, qualesumque excogitare et
confingere, quibus suppositis, iidem motus ex Geometriæ principiis, tam in futurum quam in præteritum
recte possint calculari.\par
Neque enim necesse est, eas hypotheses esse veras, immo ne verisimiles quidem, sed sufficit hoc unum, si
calculum observationibus congruentem exhibeant .\,.\,.\,. Sinamus igitur et has novas hypotheses, inter
veteres, nihilo verisimiliores, innotescere .\,.\,.\,. Neque quisquam, quodad Hypotheses attinet, quiqvam
certi ab Astronomia expectet, cum ipsa nihil tale præstare queat.}

The two hypotheses stood meanwhile on so strained a footing with one another that the adherents of the one
showed themselves irreconcilable toward the adherents of the other; it might then, for the calming of
minds, be quite to the purpose to attempt a mediation by setting up a third hypothesis, a hypothesis which
removed what was offensive in the new one and in part made good the imperfections of the old: this third
hypothesis was set up by Tycho Brahe. The old standpoint of illusion is maintained, insofar as the earth
continues to be the immovable centre: about the earth move the moon and the sun; but the sun --- that is
the concession made to the new theory --- is the centre for all the planets. One can from astronomy's
present standpoint apprehend the Tychonic hypothesis as a modification of the Copernican, arising from the
intersection point of the coordinate axes being transferred to the earth's centre; but one can with just as
much ground regard it as a special case of the Ptolemaic, namely the one obtained by taking the sun's orbit
as the common deferent circle for all the planets. On this hypothesis one then knows the relation between
% --- p. 141 ---
all the planets' mean distances, but the value of the distance in radii of the sun's orbit is, like the
relation between the moon's and the sun's distance, still unknown.

With the three hypotheses the world-structure's objectivity has been lost; actuality itself has gone missing
for consciousness. But that the actual connexion between earth and heaven should not admit of being known is
an intolerable representation, equally unnatural for scientific research and for naive contemplation. If the
reflecting understanding has, in an ingenious, combined play of hypotheses, squandered its original dowry,
the certainty of knowledge's first immediate reality, then neither does it find peace until it has been
sated anew with actual content and has conquered back the lost objectivity. Doubt and groping led to
comprehensive \emph{discussion}; contradictory considerations crossed one another; the points of view from
which they proceeded are characteristic of the scientific condition of the time.

% The first of the three points of view. The letterspacing on this page covers ONLY „kirkelig
% dogmatiske“ --- „Det“ and „Synspunkt“ are set in ordinary antiqua. On p. 142 the
% corresponding letterspacing does cover „Synspunkt“; the inconsistency is the compositor's
% and is reproduced as printed. Not a head: no a)/b)/c) and no Greek letter stands before it.
The \emph{ecclesiastically dogmatic} point of view was at bottom subjective, but, as usual, mistook itself
and fancied that it was in all infallibility objective. It was some time before the theologians became
properly aware of the danger with which the Copernican theory threatened Protestant as well as Catholic
orthodoxy; but hardly had they got their eyes open before they, from both churches, defied the danger and
went into the fight with contempt for death; history gives them the testimony that in this matter
``important for salvation'' they conscientiously used every occasion on which they could in a fitting manner
come to prostitute themselves in a properly scientific way.
% The Greek stands in italic and with an acute on ό; read on the KB copy. The citation „tomme
% Ordgydere ... detortum)“ runs over the turn of the leaf and closes only on p. 142. It is not
% a misprint and gets no note.
What Copernicus foresaw in the dedication to Pope Paul III --- that ``empty word-pourers
(\textit{ματαιόλογοι}), destitute of all mathematical knowledge, would arrogate to themselves a judgment
upon the work by wilfully distorting one or another passage of Holy
% --- p. 142 ---
Scripture (\textit{propter aliqvem locum scripturæ male ad suum propositum detortum})'' --- was amply
fulfilled.

% The second point of view. Here the letterspacing covers „Synspunkt“ as well --- confirmed on
% the KB copy; cf. the note to p. 141. „spiritus“ is likewise letterspaced and set in upright
% antiqua, not italic.
The \emph{empirically philosophical point of view}, represented by Baco, was undeniably better suited to a
prudent discussion of the contending hypotheses, grounded on the method of induction; but, remarkably
enough, Baco had, in open contradiction with his own theory of induction, so thoroughly settled himself in a
certain conception of nature's fundamental qualities and of the \emph{spiritus} indwelling in bodies that he
not only held to the immovable earth --- without, however, either settling down with Ptolemy or accepting
Tycho Brahe's modification --- but was not even able to enter upon exact investigations. He sees in
Copernicus a man who takes up notions of every kind and introduces them into nature whenever they can in any
way be brought into agreement with his calculations. By such an appraisal Baco has, however, characterized
not Copernicus but, as regards the astronomical problems,
% „Thema Coeli“ stands in italic AND letterspaced; confirmed on the KB copy. This edition
% reproduces the italic and notes the letterspacing here, since \emph inside \textit would turn
% the type upright.
rather himself; for an assertion in \textit{Thema Coeli} such as this --- that the planets' so-called
retrograde motion from east to west does not on the whole exist at all, but is an empty semblance arising
only from the fixed vault of heaven advancing further on the one side, so that the planet on the other side
seems to be left behind --- must surely be reckoned among notions of the kind no astronomer finds it worth
the trouble to refute. By boldly breaking with scholastic formalism, laying bare the barrenness of
syllogistic, maintaining induction and pointing to experience, Baco gave the thinking of his age a mighty
push forward. But the fundamental relation between the empirical and the philosophical method he was unable
to see through. For a penetrating analysis of the peculiar reciprocal action of empirical and a priori
thinking he lacks on the one side insight into the mathematical-geometrical, on the other side an eye for the
% --- p. 143 ---
critical-dialectical. He has the general propositions in the form of wit; but, unable to let the idea itself
penetratingly dissolve and confirm the peculiar facts of concrete actuality, he remains standing at the
merely witty and does not reach the spirit: Baconian empiricism is only a somewhat improved edition of the
Aristotelian.\footnote{% printed mark: *)
In vehement polemic against the Copernican hypothesis Baco adduces: that Copernicus burdens the earth with a
threefold motion and tears the sun loose from the planets, with which it is nevertheless so nearly akin;
that his system introduces too much that is immovable into nature, in that it deprives the sun and the fixed
stars --- and so precisely the bodies that shine and radiate most clearly of all --- of motion, etc.; these
arguments are in scientific worth altogether Aristotelian. The truly significant central thought that the
form is the matter's inmost essence, \textit{ipsissima res, natura naturans, fons emanationis}, is
Aristotelian too. Baco does indeed make the form-giving element in bodies, the principle of motion and
alteration, the spiritus indwelling in every body, into a material substance, and will on no account have
this material form-being confused with an Aristotelian energy; but by this unclear conception of the nature
of bodies he is so far from raising himself in the concept above Aristotle that he on the contrary sinks down
to a kind of witty dalliance with magic and sorcery. That Baco, with all his zeal for a unity of experience
and thinking, has not been able to reconcile the empirically investigating science of nature with the
philosophy of nature proper is explicable from this, that he has not been able, for his time --- as little
as the English to this day --- to become determinately and clearly conscious of the opposition between the
two qualitatively different sciences. In Baco the interest in the empirical is preponderantly philosophical;
therefore he has no sense for mechanics and looks down loftily upon mathematics; in Newton the interest rests
decisively on the mathematical-mechanical, but for all that Newton considers his physically exact principles
to be principles of the philosophy of nature.} Those who in the name of the philosophy of experience
% --- p. 144 ---
still consider Baco the non-plus-ultra will on closer testing assuredly make the discovery that there is, for
the principle itself, too little experience in his philosophy and too little philosophy in his experience.

% The third point of view. „exact empiriske“ is letterspaced --- both words, across the line
% break --- confirmed on the KB copy. The same expression twice on p. 145 is by contrast NOT
% letterspaced.
So there remains only a third point of view which could, and then really did, acquire a true scientific
significance in this discussion, namely the \emph{exactly empirical}. Its great representatives are Gilbert
and Galilei. It is not the final result but the way to it, not the verdict upon a hypothesis --- whether it
ought to be accepted, rejected or doubted --- but the grounding of the verdict, upon which everything in the
discussion first and last turns. As regards doubts and misgivings, Gilbert stands close to Baco; but in
grounding reflexion they are widely separated. Gilbert does not speculate, he investigates; when he cannot
convince himself that the earth has an annual motion, but wavers constantly between Tycho and Copernicus, it
is not the ideas but the facts that give him trouble. In his work De Magnete
% The title „De Magnete“ stands here in ordinary antiqua, not in italic; the Latin
% interpolations in parentheses by contrast in italic. So printed.
he does indeed assume that the earth turns about its axis,\footnote{% printed mark: *)
Guilielmi Gilberti Colcestrensis, Medici Londinensis. De Magnete Magneticisque corporibus et de Magno magnete
tellure, Physiologia nova. Londini MDC. --- De terrestris globi diurna revolutione magnetica. Lib.\ VI.\
cap.\ III. Terram circulariter moveri. cap.\ IV. Terræ motum negantium rationes et earum confutatio. cap.\ V.}
but sets the doctrine of this in an unclear connexion with the doctrine of the magnet and of electrical
attraction (\textit{Globus telluris per se electrice congregatur et cohæret}), in assuming that the electrical
striving aims at a binding heaping-together of matter (\textit{motus electricus est motus coacervationis
materiæ}). Although his proofs in this respect are unsatisfactory, there is nevertheless something that shows
he stands on a scientific foundation in quite
% --- p. 145 ---
% This page carries THREE footnotes, marked *), **) and ***) at the foot.
% „Ingenium“ is letterspaced --- confirmed on the KB copy; upright antiqua, not italic.
another sense than Baco. He does not decide by any speculative explanation what nature can and what it cannot
do, what suits its \emph{Ingenium} and what does not
% MISPRINT in the note: „hnmanum“ for „humanum“ --- an inverted u. The letter is unambiguously
% an n on the KB copy, and both copies show it, so it is the compositor's fault. Likewise
% „docenque“ stands for „docentque“. Reproduced as printed.
suit it:\footnote{% printed mark: *)
Theologi etiam curiosi, mysteria divina ultra hnmanum sensum posita, per magnetem et succinum illustrant, ut
vani metaphysici, cum inutilia phantasmata fundunt docenque, magnetem habent tanqvam Delphicum gladium,
exemplum semper ad omnia accommodandum. De Magnete. Lib.\ II cap.\ II.} he takes the facts as they show
themselves to him; he can be mistaken in determining the premisses, but in this he is not mistaken, that the
grounding of physical phenomena must be derived exclusively from
% The Latin in this note alternates between qv and qu in the same word: „æqvinoctiorum ...
% obliqvitatis“ against „inæquali ... obliquitatis“, and „Qvo magis“, „qvam“ against „Quare“,
% „neque“, „aliquas“. All read on the KB copy and reproduced as printed.
physically determinate data;%
\renewcommand{\thefootnote}{**)}%
\footnote{% printed mark: **)
The last chapter --- Lib.\ VI cap.\ IX. --- De præcessionis æqvinoctiorum et obliqvitatis Zodiaci anomalia
closes in a quite characteristic manner with the following words: Qvo magis hæc omnia de inæquali motu tam
præcessionis qvam obliquitatis incerta et incognita sunt. Quare neque nos illius causas aliquas naturales
proferre et certo statuere possumus. Quare etiam et nos magneticis rationibus et experimentis hic finem et
periodum imponimus.}%
\renewcommand{\thefootnote}{*)}
he holds fast to the exactly empirical method even in his missteps.

% MISPRINT: the reference mark in the text at „Dialog“ is printed **), while the note at the
% foot is marked ***). Both copies show it, so it is the compositor's fault. Reproduced as
% printed, with \footnotemark/\footnotetext.
The weightiest contributions to the great trial were made, however, by Galilei. Not content with letting one
of the speakers in his Dialogue%
\renewcommand{\thefootnote}{**)}%
\footnotemark%
\renewcommand{\thefootnote}{***)}%
\footnotetext{% printed mark at the foot: ***)
% „gjornate“ for „giornate“ and „Firense“ for „Firenze“ stand so in the printing; read on the
% KB copy. Reproduced as printed.
Dialogo di Galileo Galilei, dove ne' Congressi di quattro gjornate si discorre de' due massimi sistemi
Tolemaico e Copernicano. Firense per Landini 1632.}%
\renewcommand{\thefootnote}{*)}
conduct a defence of the Copernican hypothesis --- which afterwards drew upon him the persecution of the
Inquisition --- he aims higher, he arms the eye and widens the horizon; with the telescope he himself
constructed he makes astonishing discoveries in celestial space. With Galilei the exactly empirical method
makes its triumphal entry into the history of
% --- p. 146 ---
% This page carries TWO footnotes, marked *) and **); both begin and end here.
science.\footnote{% printed mark: *)
It is of course of Galilei as the representative of his age, not of Galilei as an isolated investigator, that
this can be said. That the telescope was first invented in Holland, and that it was the news of the invention
made in Holland that gave Galilei occasion to construct his telescope, is well known.} Of all the senses sight
is the most objective; with the widened sight, the widening thinking gets a widened foundation; so actuality
is overtaken, and the objectivity that had fled is called back again. The mountains of the moon, the phases of
Venus and above all the moons of Jupiter%
\renewcommand{\thefootnote}{**)}%
\footnote{% printed mark: **)
``Jupiter's moons, the first secondary planets of all found by the telescope, were, as it seems, discovered
almost simultaneously by Simon Marius in Ansbach on 29 December 1609 and by Galilei in Padua on 7 January
1610, without these men standing in the least connexion with one another. With respect to the publication of
this discovery, on the other hand, Galilei with his Nuntius Sidereus (1610) forestalled Simon Marius's Mundus
Jovialis (1614).'' A.\ v.\ Humboldt, Kosmos, translated by C.\ A.\ Schumacher, second volume, Copenhagen 1849,
pp.\ 289--90.}%
\renewcommand{\thefootnote}{*)}
belong to a cycle of facts which, by the new and surprising in what they offered to the sharpened sight,
necessarily had to make a strong and magnificent impression upon the age, shake confidence in the old way of
representing, open the eye to the Copernican world-structure and awaken thought-rich intimations of a
hitherto unknown mechanics of the heavens.

% DIVISION A ENDS HERE, on p. 146, and not first on p. 147. Beneath the last line of text
% stands a short, CENTRED rule (checked on both copies).
\streg

% Danish: \afdeling{B.}{Solsystemet og Universet.}   (printed pp. 147--171)
\afdeling{B.}{The Solar System and the Universe.}

Both at the stage of illusion and at that of reflexion there shows itself a peculiar relation between
the planetary spheres and the fixed-star heaven; but it is impossible to form a reasonable concept of
the universe so long as the representation of the solar system is still unclear. The immediate task
accordingly becomes: by sharpened observation and corresponding calculation to bring about the most
exact possible agreement between the orbits and the laws, and thereby to secure \emph{the objectivity
of the formal foundation}. Next, when the phenomena of motion are so far cleared up that they show
themselves as law-determined effects of determinate causes, the causes can be investigated, and the
causal connexion of the facts lead to an assurance of \emph{the objectivity of the real foundation}.
Finally, when the foundation, the real as well as the formal, is secured on all sides, the method's two
elements --- the art of observation and mathematical analysis --- will mutually drive one another
forward, and the objectification of the mechanical world-structure, in its details as in its great
outlines, complete the picture of \emph{an outwardly objective universe, infinite in finitude}.

% Run-in head, printed SEMIBOLD on the same line as the paragraph; the letterspacing in the
% three emphases above is by contrast true Sperrsatz --- both checked on the KB copy.
\runin{a) The objectivity of the formal foundation: the Keplerian laws.} One essential hindrance to a
real and true insight into the plan and structure of the solar system still remained to be removed: it
was the epicycle theory. However resolutely independent Copernicus was where it was a question of
making the sun the centre, he was nevertheless still dependent upon the old presupposition of the
necessity of circular orbits. ``We must,'' he says,\footnote{% printed mark: *)
De Revol.\ Lib.\ I Cap.\ IV.} ``assume that the celestial
% --- p. 148 ---
motions are circular, for their inequalities follow a determinate law and return in unalterable periods,
which they can do only when they go on in circles.'' We shall now see, from the standpoint of the
philosophy of nature, that the epicycle method must, by the nature of the case, make a final solution of
the great problem of objectification impossible, noting particularly wherein what is confusing for the
objectification really consists.

What dazzles and confuses cannot be placed in this, that the observer is tempted to take the many
circles, great and small, for so many actual orbits and to ascribe to each of them a material existence;
for the unreasonableness of such a way of representing is too conspicuous to be able to dazzle the
knowledgeable. ``Even if dogmatic philosophers such as Aristotle have regarded the celestial spheres as
really subsisting, solid bodies, Ptolemy nevertheless speaks of them only as of merely imaginary things,
and already from the proof he gives for the identity of an epicycle with an eccentric circle it clearly
follows that he has not held the epicycle to be anything but a geometrical conception by which he
attempted to present the apparent motions of the heaven in agreement with the observations.'' (Whewell.)
The mystical reveries of the Middle Ages and the Reformation period about solid heavens and celestial
spheres were most sharply distinguished by the astronomers from the mathematically geometrical epicycle
theory.\footnote{% printed mark: *)
% This note begins on p. 148 and ends on p. 149; the mark is not repeated at the
% continuation. The whole note stands here at its mark.
\dots\ ``in the middle of the sixteenth century, when the learned polymath Girolamo Fracastoro's theory
of the 77 homocentric spheres met with so much approval, and when later the opponents of Copernicus
sought out every means for the maintenance of the Ptolemaic system, the representation --- favoured
especially by the Church Fathers --- of the existence of solid spheres, circles and epicycles was still
very widespread. Tycho Brahe expressly prides himself on the merit of being, by his considerations of
the orbits of comets, the first to have proved the impossibility of solid spheres and shattered their
artificial scaffolding.'' (A.\ Humboldt.)}

% --- p. 149 ---
Nor can what misleads the objectification be sought in this, that the epicycle theory, by piling circle
upon circle, had at last to become too intricate; for in the first place a part of the superfluous
hypotheses did gradually admit of being reduced. One had begun by laying the deferent circles and first
epicycles in the plane of the ecliptic, and had therefore to explain the latitudes by introducing a
motion on small circles perpendicular to the ecliptic. Copernicus, however, helped himself by assuming
certain librations or oscillations of the epicycles; and Tycho Brahe avoided the perpendicular circles by
giving the planetary orbits a slight inclination to the ecliptic. Next, and this is what is decisive, the
growing complications hung most closely together with the multiplication of the observations and the
imperfection of the mathematical methods at the time. If the intricate\footnote{% printed mark: *)
% p. 149 carries TWO notes in the setting: the end of the note from p. 148 (without a mark)
% and then this one, marked *).
When the intricate arises from clumsiness in setting and treating a problem, an improved method will be
able to remove all difficulties to an astonishing degree. Thus the higher mathematics now treats with the
greatest ease a quantity of problems whose solution is difficult, cumbersome, indeed impossible, if it is
to be done solely by the means the elementary mathematics offers. It is another matter when the intricate
lies in the problem's own nature; for in this case the increasing entanglement of the problems will stand
in direct relation to the progress of the science. Entanglement presupposes development. In this sense it
must be said that the problems in the newer astronomy are far more intricate than in the older.} were in
and for itself misleading, then the infinite series of the newer analysis would be misleading in quite
another degree, for in comparison with these the complications of the epicycle theory may surely be said
to be very simple.

% --- p. 150 ---
No: what misleads the objectifying consciousness lies in this, that the alternating hypotheses make it
impossible to distinguish objectively between actual and feigned orbits. That the sun's orbit, apprehended
abstractly and geometrically, must be the same whether one assumes that the sun with even velocity
describes a circle whose centre lies outside the earth, or ascribes to the sun an even motion on a circle
whose centre moves evenly on a circle having the earth for its centre, is
incontestable.\footnote{% printed mark: *)
% The indeterminate magnitude is written in the printing as y with a raised 1 (a numeral, not
% a prime: the figure has both flag and foot serif and is identical with the 1 in „1+e cos v“;
% checked on the KB copy). The otherwise isolated letters in the running text are here
% likewise set in math, so that the face does not change in mid-enumeration y, y^1.
% MINUS is here a true minus, not Nielsen's ÷; read off the page.
% MISPRINT: the second time the Danish has „detererende“ for „defererende“ --- the same word
% stands correctly two lines above. Both copies have it; the errata leaf does not mention it.
If, on the latter hypothesis, the intersection point of the coordinate axes be laid in the earth's centre,
while the axis of ordinates coincides with the line of apsides, and the ordinate for the deferent circle
be denoted by $y$, and that for the sun on the epicycle by $y^{1}$; then $x^{2} + y^{2} = r^{2}$ is the
equation for the deferent circle, and $y = y^{1} - \varrho$; whereby one obtains
$x^{2} + (y^{1} - \varrho)^{2} = r^{2}$, which is precisely the equation for the eccentric circle answering
to the first hypothesis, the eccentricity being $= \varrho$.} But although the mathematical result is the
same, the mechanical connexion of the matter is nevertheless highly different. On the first hypothesis
there is here in the objective sense only one orbit; on the second there are, by contrast, two. In the
first case one will for the time being be able to explain the orbit by pointing to the nature indwelling
in all heavenly bodies of moving in circles; the only metaphysical difficulty that remains must then turn
upon the ground why the earth has not become the actual centre of the sun's orbit. In the second case it
is really not the sun but the sun's epicycle that moves about the earth. Here there must then in the
objective sense be a double ground for the double motion: a particular ground for the motion of the
epicycle's centre about the earth's centre, and alongside this a particular ground for the epicycle's
motion with the sun or --- what for the grounding gives a new problem --- for the sun's motion about the
epicycle's centre. If two alternating hypotheses, however
% --- p. 151 ---
simple each of them may be in itself, can already give the connexion of the matter so subjective and
doubtful an appearance, then it follows of itself that along this way one comes with every step ever
further from the objective. With the orbit of Mars, for example, the mechanical connexion and its
grounding is still more entangled, since account must be taken of the sun's orbit lying between the earth
and Mars. If one lets Mars move on an epicycle whose centre moves evenly, with a period of revolution of
687 days, on a deferent circle about the earth in such a way that Mars's place on the epicycle is
determined by the sun's place, the epicycle radius to the planet being constantly parallel with the line
connecting earth and sun: then the objective element in the orbit of Mars is already doubtful in itself,
and becomes still more doubtful in that the objective character of the sun's orbit is doubtful. For motion
in circular orbits speculative reason could still think a ground; but for eccentric orbits and motions of
epicycle centres not even the most ingenious speculation has been able to find out anything that could
resemble a ground in nature. Grounds without causes can be thought, but causes without grounds cannot be
thought; if the phenomena are so entangled that one cannot glimpse a determining ground, how should it
then be possible to discover a working cause? If the laws of motion are to be known, the epicycle theory
must go: here is the turning point at which Kepler's genius breaks through.

% The Latin interpolation is printed in ITALIC, not letterspaced; the exclamation mark with it.
To follow the discoverer of the great world-laws on his thorny way with its bewildering windings ---
\textit{per errores ad veritatem, per ardua ad astra!} --- is not the task of the philosophy of nature; an
indication of the movement must suffice for a presentation of the course of the objectification.

In the year 1600 Kepler attached himself to Tycho Brahe, whom he found in Prague eagerly occupied with a
theory of the planet Mars answering to the observations he himself
% --- p. 152 ---
had made. It was these new and improved observations that led Kepler onto the new path and put him in a
position to publish, nine years later, his epoch-making work
% The title is printed in ITALIC AND LETTERSPACED at once (the same case as „Thema Coeli“
% p. 142): \textit, with the letterspacing only noted here, since \emph inside italic would
% turn the type upright again. The full stop is set outside the italic.
\textit{De motibus stellæ Martis}. Had it been possible by immediate observations to determine an
arbitrary number of places in the orbit of Mars about the sun, then the problem of finding out the orbit's
character would undeniably have been considerably simplified; but the problem was entangled to such a
degree that the relations in the objective had to pass through a quantity of testings in subjective
reflexion before the solution could be attained. The reciprocal determinations of the determinate and the
indeterminate with the known and the unknown are manifold. There must indeed in the objective sense be a
system of points lying in Mars's heliocentric orbit; but these points had first to be determined in
relation to the earth's orbit. First the period of revolution had to be found by means of the
oppositions\footnote{% printed mark: *)
% p. 152 carries TWO notes, marked *) and **); both begin and end here.
At the instant of opposition the planet's heliocentric longitude is the same as the geocentric; the
geocentric longitudes of planet and sun differ by 180 degrees from one another. The Tychonic series of
observations of Mars covers a period of 20 years and gives 12 oppositions.}; next, one and the same place
at which Mars had twice been in its orbit%
\renewcommand{\thefootnote}{**)}%
\footnote{% printed mark: **)
If one chooses, among the geocentrically angular places a series of observations can give, two which are
separated by exactly an interval of 687 or more exactly of 686.98 days, then one knows that Mars has both
times been at the same point of its orbit.}%
\renewcommand{\thefootnote}{*)}%
, had to be seen in relation to the two corresponding places for the earth in its orbit: only by two lines
of direction from two known stations toward the same point did it become possible to find out the planet's
distance from the sun and determine the completely heliocentric place in the orbit of Mars. Only when one
point had been found could the determination be repeated for an arbitrary number of points.
% --- p. 153 ---
When Kepler now, with these points, set about establishing by calculation the nature of the line, he found
that the orbit of Mars could not, as had hitherto been assumed, be any circle, but that it must have an
oval form. By assuming the line to be an ellipse he found, not indeed complete agreement with the
observations, but nevertheless a significant approximation. Why only an approximation? Dazzled by the
theory of the circle\footnote{% printed mark: *)
% p. 153 carries THREE notes, marked *), **) and ***). The Latin Kepler text in this note is
% printed in ANTIQUA (upright), not italic; so too the reference to the title --- unlike the
% italicized title in the text on p. 152.
% MISPRINT: „metaphycicæ“ (with c) for metaphysicæ. The letter is positively a c on the KB
% copy; both copies have it, and the errata leaf does not mention it.
Primus meus error fuit, viam planetæ perfectum esse circulum, tanto nocentior temporis fur, qvanto erat ab
auctoritate omnium philosophorum instructior et metaphycicæ in specie convenientior. De motibus stellæ
Martis. Pars III. cap.\ XI.}, Kepler had started from the presupposition that the earth's orbit, whose
eccentricity is comparatively insignificant, had to be a circle; but when this error was corrected, the
so-called \emph{first} law was also discovered: \emph{the planetary orbits are ellipses, in one focus of
which the sun is situated.}

The polar equation for the ellipse shows that the \textit{radius vector} is a function of its angle with
the major axis%
\renewcommand{\thefootnote}{**)}%
\footnote{% printed mark: **)
% The note consists solely of the displayed formula. Minus is a true minus, not ÷; the
% exponent is a 2, not a 3 (checked on the KB copy).
$r = f(v) = \dfrac{a\,(1 - e^{2})}{1 + e \cos v}.$}%
\renewcommand{\thefootnote}{*)}%
; this leads in turn to a closer determination of the motion in the orbit. Hitherto it had counted as an
astronomical dogma that the planet must absolutely move with even angular velocity --- in other words,
traverse equal angles in equal times. In honour of the dogma one invented the \textit{punctum æqvans} with
corresponding epicyclic artifices. Kepler at length saw the unreasonableness of this. There was indeed
something even in the motion; but it was not the angles traversed, it was rather the elliptical sectors
described by the \textit{radius vector}%
\renewcommand{\thefootnote}{***)}%
\footnote{% printed mark: ***)
% This note begins on p. 153 and ends on p. 154; the mark is not repeated at the continuation.
% The whole note stands here at its mark.
From the law of the areas it follows that the planet's angular velocity at every point of its orbit stands
in inverse ratio to the square of the distance. Hereby a new means is won for determining the planet's
distances from the sun, when one of these distances is assumed as given.}%
\renewcommand{\thefootnote}{*)}%
. Hereby the so-called
% --- p. 154 ---
% „Radii vectores“ is printed in ITALIC AND LETTERSPACED inside the letterspaced statement of
% the law. \emph is therefore broken around it, so that the italic is preserved.
\emph{second} law is determined: \emph{the areas which the} \textit{radii vectores} \emph{describe stand to
one another as the times employed in the motion.}

Like Copernicus, Kepler too has a magnificent conception and a living assurance of an objectively subsisting
world-harmony. There was in this respect no thought so bold that this wonderfully cosmic genius did not dare
take it up. In his
% this title too is italic AND letterspaced; see the note at p. 152.
\textit{Mysterium Cosmographicum} (1596) he expresses himself about the relations of the planetary orbits in
the following manner: ``The earth's orbit is a circle. If one describes about the sphere to which the earth's
orbit answers as great circle a dodecahedron, then the circumscribed sphere gives the orbit of Mars. If one
describes about this orbit a tetrahedron, then the corresponding circle presents the orbit of Jupiter. If one
describes about Jupiter's orbit a cube, then the circle answering to it will give the orbit of Saturn,'' and
so forth. These Pythagorean-Platonic speculations found at length a fruitful scientific clarification,
insofar as they underwent a transformation and after 27 years' exertions led to the discovery of the
\emph{third} law, which says that \emph{the squares of the different planets' periods of revolution stand to
one another as the cubes of their mean distances from the sun.}

These laws, whose objective validity no one can draw into doubt without at the same time drawing the whole
validity of the science of experience into doubt, furnish the formal foundation for a scientific
objectification not only of the solar system but of the universe. That the formal element in these laws'
objectivity points out beyond itself with a promise of the real is manifest. Instead of empty geometrical
centres, as in the epicycle theory, the elliptical orbits have a full corporeal
% --- p. 155 ---
centre, an actual central body in the focus; instead of the imaginary circles of subjective reflexion we find
here objective relations between the central body and the peripheral bodies moving in the orbits. As
attention fastens upon these fundamental relations between the bodies themselves, the question of moving
causes gets sure points of support and a new life that promises a new and rich yield: formal knowledge has
prepared the real; abstractly empirical mechanics is on the point of becoming rational physics.

% This head too is printed SEMIBOLD and runs into the paragraph; it stands on p. 155, not on
% p. 156.
\runin{b) The objectivity of the real foundation: gravitation.} The first objections of any significance
against the Copernican system stemmed partly from unclarity about the distance between the planetary system
and the fixed-star heaven, in that one demanded that the Copernicans should point out the motions which the
fixed stars, on the new doctrine, must necessarily have\footnote{% printed mark: *)
Remarkably enough, the question of the fixed stars' distance was afterwards, when the law of light's
propagation had been discovered, to lead to one of the strongest proofs for the Copernican system, through
Bradley's discovery of aberration.}, and partly, and chiefly, from ignorance of mechanics. If the earth moved
from west toward east, then a stone falling from a high tower must necessarily --- so it was objected --- fall
in a westerly direction, whereas the fact is that it falls in an easterly one: inertia was not yet known. By
Galilei's discovery of the law of fall the reciprocal determination of inertia and attraction was shown in
exact form, and mechanical physics founded. Meanwhile, as regards the mechanics of the world, several
introductory steps and groping attempts still had to be made before Newton came to utter the fundamental
propositions, work the fragmentary into a rational whole, and mark the new
% --- p. 156 ---
epoch by a work such as \textit{Philosophiæ naturalis Principia mathematica.}

It is surely more the work's title than its content that has in more recent times given Hegel and the Hegelian
school%
\footnote{% printed mark: *) --- the first of two notes on p. 156.
Jul.\ Schaller. Geschichte der Naturphilosphie. 1ster Theil.
% „Naturphilosphie“: the o is missing in „philos[o]phie“. Compositor's fault; both copies show
% the same. Reproduced as printed; read „Naturphilosophie“.
Leipzig 1841. S.\ 353--400.}
occasion for objections against the Newtonian principles as unclear as they are unwarranted. If one holds
strictly to the designation philosophy of nature (\textit{philosophia naturalis}), and takes the word as a
category in the Hegelian sense, then objections can undeniably be made against the manner in which Newton has
defined and set up his fundamental propositions; but such a conception rests upon a colossal
misunderstanding of the Newtonian standpoint. To avoid unnecessary wrangling about obscure words, Newton
declares straightforwardly that the expressions he has chosen --- inertia, attraction and several others ---
are not to be taken in a physical but only in a mathematical sense.%
\renewcommand{\thefootnote}{**)}%
\footnote{% printed mark: **) --- the second of two notes on p. 156.
Voces autem attractionis, impulsus, vel propensionis cujuscunque in centrum, indifferenter et pro se mutuo
promiscue usurpo; has vires non physice, sed mathematice tantum considerando. Princ.Mathem.\ Amstælodami
MDCCXXIII. pag.\ 5.}%
% „Amstælodami“ with æ (for Amstelodami) stands so in both copies. Reproduced as printed.
\renewcommand{\thefootnote}{*)}
From this speculative philosophy has taken occasion to show how ungrounded his separation between the
physical and the mathematical is in itself, and how little he has been able to hold the separation
consistently. But how foreign this whole conception of physics and mathematics is to Newton's train of
thought one will easily be able to see on a little reflexion. Newton guards against one's taking his
expressions physically, that is, metaphysically, and the meaning is that he will not at all enter upon
speculative investigations
% --- p. 157 ---
concerning the inner How of the physical activities%
\footnote{% printed mark: *) --- the only note on p. 157.
Vocem attractionis hic generaliter usurpo pro corporum conatu qvocunque accedendi ad invicem; sive conatus
iste fiat ab actione corporum vel se mutuo petentium vel per spiritus emissos se invicem agitantium, sive is
ab actione ætheris aut aëris mediive cujuscunque seu corporei seu incorporei oriatur corpora innatantia in se
invicem
% Three CLOSE points.
utcunque impellentis .\,.\,. l.\ l.\ p.\ 172, confer.\ p.\ 483.
% The two letters after the points are alike and have no flag at the upper left (unlike the
% figure one in „172“): two small l's, therefore, and not „1. 1.“ nor „l. I.“ (liber I), which
% was considered and rejected. The note mixes qv and cu in one and the same sentence:
% „qvocunque“ beside „cujuscunque“ and „utcunque“ --- as printed.
The fundamental thought of the whole standpoint can hardly be uttered more plainly.}%
; but to tear the physical loose from the mathematical could surely not possibly occur to the man who from
first to last is laid out so as to make clear how the quantitative in the effect points point for point to the
qualitative in the cause.

Mechanics is originally a unity of mathematics and physics; space and time, matter and force are in equal
degree physical and mathematical factors. To develop the mathematical Newton must lay weight upon the
physical, and conversely. Mechanics expressly demands that the relation between the two be apprehended in a
double-sided way, partly from the standpoint of cause and effect, partly from the standpoint of law and
phenomenon. As regards law and phenomenon, it is intelligible that Newton can exchange categories such as
attraction and impact for one another, for what attraction is in the form of continuity, impact is in the
form of discreteness. As regards cause and effect, there is by contrast an essential difference between
attraction and impact. But that, mechanically understood, there is no conflict between the first and the
last conception is seen straightforwardly from this, that a body is acted upon \emph{from without} just as
much when it is attracted as when it is struck.
% The letterspacing of „udenfra“ is confirmed on the KB copy.
The contradiction between ``\textit{vis attractiva}'' physically apprehended
% The Latin is set in italic inside the quotation marks; read on the KB copy.
and the law of inertia, to which Schaller appeals, does not in actuality take place. What Newton brings
% --- p. 158 ---
out --- that the force of attraction cannot have a physical significance for the very reason that the centres
from which it issues are only mathematical points --- must be understood from the connexion in this way: that
the mechanical theory which makes a beginning by concentrating the attracting causes in mathematical points
thereby gives notice that it will not investigate the inner unity of essence of matter and force; but from
this it can in no way be derived either that the physical is excluded, or that the mathematical in attraction
should cease because the centre is laid into a corporeal mass.

(The relation between Kepler and Newton is, at the stage of mechanics, a
% The parenthesis opens here on p. 158 and closes only on p. 161. It is not a misprint and
% gets no note.
relation between objectification in empirical reflexion and objectification in pervasive conceiving. The
Keplerian reflexion abstracts the laws from the cases before it and secures them by empirical induction; the
Newtonian method, by contrast, aims at knowing the causes from their effects, determining the particular on
the foundation of the universal, and conceiving the laws in their fulfilment. That is, on the way of
objectification, a great advance.

Kepler speaks only of the areas described in the planetary orbits; but Newton proves that the law of the
areas must hold for any motion whatever of a material point about a fixed centre, since it is in general
determining for the objective concept of central motion.

Kepler holds to the factual --- that the planetary orbits are ellipses --- but, bound to the fact, he is not
able to state the cause; from Newton's principles, by contrast, the cause can be seen in the effect, insofar
as it can be seen that not merely the ellipse but the conic section generally is a naturally necessary
consequence, an express consequence of cause, of attraction's peculiar law:
% --- p. 159 ---
$T = \dfrac{\mu}{r^{2}}$, and that conversely this law
% The formula stands within the line of text, not displayed on a line of its own.
is a necessary consequence of motion in a conic section.%
\footnote{% printed mark: *) --- the only note on p. 159; it fills the whole leaf beneath the
% two lines of body text.
J.\ Schaller remarks that Newton has not been able to ``prove generally'' that the curve must be a conic
section when the attraction stands in inverse ratio to the square of the distance. It must in this connexion
be expressly brought out that what essentially matters is not the provisional manner of proof, but what
admits of being derived from the fundamental thought itself. Everyone who makes himself acquainted with the
fundamental propositions of rational mechanics will have ample occasion to convince himself of the universal
validity of the Newtonian principles. For central motion in general --- r being the radius vector, c an
arbitrary constant, θ the angle r forms with a fixed axis, and φ the accelerating force --- one finds (cf.\
Duhamel, Cours de Mécanique, Paris 1846, II, p.\ 3) the formula:
\[
\varphi = \frac{c^{2}}{r^{2}} \left( \frac{1}{r} +
    \frac{d^{2}\dfrac{1}{r}}{d \theta^{2}} \right)
\]
If now the orbit is a conic section, the angle between r and the polar axis θ, that between the major axis
and the polar axis ω, and the equation of the orbit
\[
\frac{1}{r} = \frac{1 + e \cos (\theta - \omega)}{\dfrac{p}{2}},
\]
then it follows from this that:
\[
\frac{d^{2}\dfrac{1}{r}}{d \theta^{2}} = \div \frac{e \cos (\theta - \omega)}{\dfrac{p}{2}},
\]
% The sign before the fraction is Nielsen's MINUS (÷), read on the KB copy. The sign between θ
% and ω is by contrast an ordinary minus in both formulas on this page, while p. 160 sets ÷ in
% the same place. Not harmonized.
which, in connexion with the equation of the orbit, gives:
\[
\frac{1}{r} + \frac{d^{2}\dfrac{1}{r}}{d \theta^{2}} = \frac{1}{\dfrac{p}{2}}
\]
and consequently:
\[
\varphi = \frac{c^{2}}{\dfrac{p}{2}} \cdot \frac{1}{r^{2}} = \frac{k}{r^{2}}
\]
From this it is seen that the attraction stands in inverse ratio to the square of the distance when the orbit
is a conic section and the central force issues from the focus.}
% --- p. 160 ---
Kepler confines himself to setting up the three laws alongside one another without any deduction
whatever; but on Newton's principles the one law can, if one starts from the universal and introduces
conditions, be deduced from the other.%
\footnote{% printed mark: *) --- the only note on p. 160. The note EXTENDS OVER TWO LEAVES: it
% fills the foot of p. 160 and is continued at the top of the note field on p. 161 (from
% „c = ...“), without the mark being repeated. Per the house rule the whole note stands here
% at its mark. The formulas stand in the printing WITHIN the lines of text, not displayed on
% lines of their own; so reproduced here.
If the force of attraction be put $\varphi = \dfrac{\mu}{r^{2}}$, then the equations in the plane (xy)
become $\dfrac{d^{2}x}{dt^{2}} = \div \dfrac{\mu x}{r^{3}}$,
$\dfrac{d^{2}y}{dt^{2}} = \div \dfrac{\mu y}{r^{3}}$ (I).
% The sign is Nielsen's MINUS (÷), read on the KB copy, not division. The numbering is a
% capital I with serifs (not the figure one) in both places.
If one introduces polar coordinates, so that $x = r \cos \theta $, $y = r \sin \theta $, and the area
described by the radius vector $\tfrac{1}{2} c\, dt = \tfrac{1}{2} r^{2}\, d\theta $, then by
multiplying the equations (I) by $cdt = r^{2}\, d\theta $ one gets the equations
$c\, d \dfrac{dx}{dt} = \div \mu \cos \theta d\theta $ (α) and
$c\, d \dfrac{dy}{dt} = \div \mu \sin \theta d\theta $ (β).
If these equations be integrated so that the constant for (α) becomes $\div \gamma \sin \omega $, and
for (β) $\gamma \cos \omega $; if the first be then multiplied by $y = r \sin \theta $ and the second by
$x = r \cos \theta $ (remembering that $xdy \div ydx = cdt$ and
$\sin^{2} \theta + \cos^{2} \theta = 1$), one obtains:
$r = \dfrac{\dfrac{c}{\mu}}{1 + \dfrac{\gamma}{\mu} \cos (\theta \div \omega)}$
% THE ERRATA LEAF corrects p. 160, line 13 from the foot: „c/μ in the numerator --- read
% c²/μ“. Per the house rule the text stands here AS PRINTED, with c (without exponent) in the
% numerator. The KB copy is unambiguous: c alone stands there. The working scan shows a raised
% 2 beside the c --- that is an earlier reader's PENCIL, not the printing, just like the forged
% full stop on p. 113.
that is, \emph{the orbit is a conic section with the centre of attraction in the focus.}

If the equation $r^{2}\, d\theta = cdt$ be integrated so that $\theta = o$ when $t = o$, then
$t = \dfrac{\displaystyle\int_{o}^{\theta} r^{2}\, d\theta}{c}$,
% The limits and values are set with the LETTER o, not with the numeral nought; read on the KB
% copy.
which put in words means that \emph{the areas described by the radius vector are proportional to the
times in which they are described.}

If, finally, the special case be chosen in which the conic section is an ellipse, and the integral
$ct = \int r^{2}\, d\theta $ be taken from o to $2\pi $, then this becomes twice the ellipse's area,
$2 ab\pi = 2a^{2} \sqrt{1 \div e^{2}}\, \pi $, and
$T = \dfrac{2 a^{2} \sqrt{1 \div e^{2}}\, \pi}{c}$ the time in which the whole orbit is traversed. If
a, e and T are known, c can be determined, since $c = \sqrt{a (1 \div e^{2}) \mu}$, whereby one obtains
$t = \dfrac{\int r^{2}\, d\theta}{\sqrt{a (1 \div e^{2}) \mu}}$, and consequently
$T = \dfrac{2 a^{2} \sqrt{1 \div e^{2}}\, \pi}{\sqrt{a (1 \div e^{2}) \mu}}$;
so that $\dfrac{T^{2}}{a^{3}} = \dfrac{4\pi^{2}}{\mu} = K$.
% In the two last formulas the printed radical bar does not reach past the μ; the sense requires
% μ under the root, and it is so set here.
That the ratio between the square of the period of revolution and the third power of the major axis is
a constant magnitude shows that this ratio holds for different material points moving in ellipses about
a common focus; in other words: \emph{the squares of the periods of revolution stand to one another as
the cubes of the semi-major axes, or mean distances.}

On these conditions the Keplerian laws can be deduced from Newton's principles.}
That the deduction must constantly be tied to
% --- p. 161 ---
given conditions recalls that the proofs, as regards the physical, are neither merely a priori nor
merely empirical, but that they must, by the nature of the case, rest upon an exactly articulated
reciprocal determination of experience and thinking.)
% The parenthesis opened on p. 158 closes here.

It has been asserted from the standpoint of ``the concept'' that ``the more universally attraction is
apprehended, the more the mathematical proof loses in significance and shows itself insufficient to
strike the essence of the phenomenon.''%
\footnote{% printed mark: *) --- p. 161's own note. In the printing it stands beneath the
% continuation of the note from p. 160.
Schaller, Anfr.\ Skr.\ p.\ 367.}
% „Anfr.“ (not „Anf.“) stands there; read on the KB copy.
The opposite of this is precisely what is true. The history of science teaches that insight into the
essence of the phenomena has advanced steadily with the knowledge of the laws, and the knowledge of the
laws with the mathematically mechanical determination of the objective relations, from which alone the
causes and their conditions admit of being found out.

The doctrine of gravitation, the doctrine of the planets' mutual attraction, is grounded evenly upon
concept and fact, and is in equal degree a priori and empirical. It can be shown from experience and
fixed by induction how the masses act upon one another; but the mutual action that admits of being
observed is in turn grounded by an activity
% --- p. 162 ---
that does not admit of being immediately observed, but must be determined by thinking --- namely the
reciprocal action whereby the smallest parts in the one mass attract the parts in the other. The
attraction of the parts is explained from that of the masses, and that of the masses again from that of
the parts. The mutual attraction between two masses is proportional to their product, while the force
with which the one mass is accelerated in its motion about the other is proportional to the sum of the
masses.%
\footnote{% printed mark: *) --- the first of two notes on p. 162.
If M be the sun's and m the planet's mass, and the distance r, then the mutual attraction is to be
expressed by $\dfrac{f.\,Mm}{r^{2}}$; the accelerating force for the planet by $\dfrac{f.M}{r^{2}}$, for
the sun by $\dfrac{f.m}{r^{2}}$; for the planet's \emph{relative} motion about the sun the driving force
must by contrast become $\dfrac{f\,(M + m)}{r^{2}}$. By f is denoted the mutual
% The period after f is Nielsen's multiplication sign (cf. p. 46); the last fraction prints no
% period. Reproduced as printed.
attraction of two units of mass, reduced to simple points and placed at a distance that gives a unit of
length.}
When the masses, the actual planets, are introduced in place of material points, the Keplerian laws do
indeed undergo manifold modifications; but these modifications are so far from weakening their validity
that they are on the contrary proofs of their actual fulfilment.

If the planets were perfectly spherical, formed of homogeneous concentric layers, they would act upon
one another as though each of them had its whole mass compressed into its centre; but since they deviate
somewhat in form from the sphere and, what is more essential, stand in relations of attraction not merely
to the sun but to one another, the reciprocal action becomes infinitely entangled, and the perturbations,
the periodic as well as the secular, make in this respect great demands%
\renewcommand{\thefootnote}{**)}%
\footnote{% printed mark: **) --- the second of two notes on p. 162. The note EXTENDS OVER TWO
% LEAVES: it begins at the foot of p. 162 and is closed at the top of the note field on p. 163,
% without the mark being repeated. Per the house rule the whole note stands here at its mark.
The formula for the moon's longitude presented by Delaunay in connaissance des temps for 1865 takes up no
fewer than 45 printed pages, and yet this formula takes account only of the perturbations due to the sun.}%
\renewcommand{\thefootnote}{*)}
upon the observing
% --- p. 163 ---
and calculating science. The perturbations are not, however, what the word one-sidedly suggests,
disturbances of something actual; on the contrary! they belong essentially to the character of mechanical
totality; they are phenomena of many-sided reciprocal action; it is precisely the perturbations that give
the actual the stamp of outward material objectivity. For a scientific objectification of the
world-structure it is not sufficient to consider the single planets in the single orbits, each for itself,
or to stop at a general representation of the various supervening actions to which a planet in its orbit
may be subject. The concept of actuality must be developed so exactly and completely that the reciprocal
determination of the universal and the single, the logical and the antilogical
% THE ERRATA LEAF corrects p. 163, line 13 from the top: „det Antilogiske --- læs det
% Antilogiske,“. Per the house rule the text stands here AS PRINTED, and so WITHOUT the comma.
% The KB copy confirms that no comma stands in the setting; the comma in the working scan is an
% earlier reader's PENCIL.
the laws and their modifications in the particular cases, stands clear before thought.%
\footnote{% printed mark: *) --- the only note on p. 163. The note EXTENDS OVER TWO LEAVES: it
% begins at the foot of p. 163 and fills the whole note field on p. 164, without the mark being
% repeated. Per the house rule the whole note stands here at its mark.
If one denotes the sun's mass by M, its coordinates by X, Y and Z, the planets' masses by $m$, $m'$,
$m''$ .\,.\,.\,., their coordinates
% The enumeration is kept in one face: since the terms carry primes, the whole series is set in
% math (the house rule from batch 5).
by $x$, $y$, $z$, $x'$, $y'$, $z'$, $x''$, $y''$, $z''$ .\,.\,.\,., and if one further lets the sun's
distances from the planets be $r$, $r'$, $r''$ .\,.\,., and the perturbed planet's distances from the
perturbing ones $s'$, $s''$ .\,.\,.\,., then the equations for the motion of the single planet, m, about
the sun as a fixed
% „Ceutrum“ for „Centrum“: an inverted or wrong sort. The letter is unambiguously a u on the KB
% copy. Compositor's fault; reproduced as printed.
% THE ERRATA LEAF corrects p. 163, line 10 from the foot: „være med --- læs være, med“. Per the
% house rule the text stands AS PRINTED, WITHOUT the comma.
centre will be, with respect to the X-axis:
\[
\frac{d^{2}\,(x - X)}{dt^{2}} + \frac{(M + m)\,(x - X)}{r^{3}} = 0;
\]
% THE ERRATA LEAF corrects p. 163, line 8 from the foot: „$=0$; --- læs $=0$, o. s. v.“ Per the
% house rule the text stands AS PRINTED: „= 0;“ with a semicolon and without „o. s. v.“
but if the planet m is perturbed by $m'$, $m''$ .\,.\,., the equation for the X-axis will be:
\begin{align*}
\frac{d^{2}\,(x - X)}{dt^{2}} + \frac{(M + m)\,(x - X)}{r^{3}}
  = {}& \frac{m'(x' - x)}{{s'}^{3}} + \frac{m''\,(x'' - x)}{{s''}^{3}} + \\
.\,.\,.\,. {}& - \frac{m'\,(x' - X)}{{r'}^{3}}
        - \frac{m''\,(x'' - X)}{{r''}^{3}} - .\,.\,.\,. \quad (I)
\end{align*}
% The signs here are ordinary minuses, not Nielsen's ÷; read on the KB copy.
and the corresponding ones for the Y- and Z-axis have the same form. That we have in such equations a
system of objective relations is certain; but by what is the objectivity secured? By the objectification!
And by what is the objectification secured? By an articulating unity of observation and thinking. It is
thinking that sees that when the law for the motion of the single planet, m, is
$\dfrac{d^{2}x}{dt^{2}} = - \dfrac{M\,(x - X)}{r^{3}}$ etc., and for the sun's motion
$\dfrac{d^{2}X}{dt^{2}} = \dfrac{m\,(x - X)}{r^{3}}$ etc., then the law for the perturbed planet, m, must
become:
% THE ERRATA LEAF corrects p. 164, line 15 from the top: „pertuberede --- læs perturberede“.
% Per the house rule the Danish text stands as printed.
\[
\frac{d^{2}x}{dt^{2}} = - \frac{M\,(x - X)}{r^{3}}
  + \frac{m'\,(x' - x)}{{s'}^{3}}
  + \frac{m''\,(x'' - x)}{{s''}^{3}} + .\,.\,.\,.
\]
and for the sun:
\[
\frac{d^{2}X}{dt^{2}} = \frac{m\,(x - X)}{r^{3}}
  + \frac{m'(x' - X)}{{r'}^{3}}
  + \frac{m''\,(x'' - X)}{{r''}^{3}} + .\,.\,.\,. \quad (II)
\]
On the left side the concept of (perturbing) force is determined in exact logical universality; on the
right side it is shown how the concept expresses itself in existing cases. From the equations (II) for the
absolute motions the equations (I) for the relative ones are then derived. The objectifying reflexion
articulates the logical, the universal, by making the law clear in its different forms and formulas; but
the language of signs also shows that the general formulas are only empty symbols so long as the moments
of actuality, with number and measure in corresponding antilogical determinations, are still wanting. The
facts are objective, but the analyses subjective. What do we learn from this? That no science, not even
the most objective, can take the subjective away and keep the objective --- that, in other words, nothing
in the world is objective except what has been objectified.}
Only from the standpoint of objectification
% --- p. 164 ---
does it become fully evident what it means that realistic sciences such as mechanics and physical astronomy
stand in Spirit's service. It is something far higher than a petty craving to get the manifold phenomena of
motion punctually calculated that drives the investigations forward; it is the investigating Spirit that
will have actuality theoretically in its power; it is the thinking subjectivity that desires the objective
and works restlessly with the conditions of objectification. A magnificent example, showing what
% --- p. 165 ---
science can attain by building on Newton's foundation, the present age has admired in Laplace's
\textit{Mécanique céleste}.

% Run-in head, read on the page itself (KB copy) and not in the Indhold, whose OCR has lost the
% letter prefix: the prefix is „c)“. Printed SEMIBOLD and on the same line as the paragraph it
% opens --- the pattern from pp. 79, 101 and 107.
\runin{c) The universe's mechanical objectivity: a system of systems.} When Hephaestus was of old cast out
of Olympus's window he had to fall for 9 days and 9 nights before he reached Lemnos. A fall from such a
height may well awaken a shudder; but when one considers that it is the way from the highest heaven down to
the earth that had to be traversed on this journey, one will not be able to say that the length is
exaggerated. In the Greeks' view of the world limitation is so essential a moment that Aristotle can even
in a manner be said to make the Godhead itself
% THE GREEK is read on both copies under high magnification. The last word stands as οὐρανοῦ,
% quite regularly. In the third word, by contrast, the printing sets the perispomene over the
% OMICRON and not over the upsilon --- „τõυ“ --- in both copies, so it is the compositor's (or
% the face's) form and not damage in this copy. Unicode/LaTeX cannot conveniently carry that
% placement, so the word is reproduced here as τοῦ, and the deviation noted here alone.
the limit of the heaven (τὸ πέρας τοῦ οὐρανοῦ). Observing astronomy had indeed to come upon the question of
the distance-relations of the heavenly bodies from many sides, but the observations were imperfect, and the
Ptolemaic epicycle theory had to content itself with an imagined standard. Tycho Brahe himself still assumed
the sun's parallax to be so great that the earth became a considerable magnitude in the universe. Only after
the discovery of the telescope and the introduction of the new methods of observation did science get so far
as to obtain measure and standard. Not only did the radius of the earth's orbit become a known magnitude;
even the fixed stars' infinitely slight, generally vanishing parallax led to insight into immeasurable
measures of distance. When Ole Rømer, by means of the eclipses of Jupiter's moons, had discovered the
velocity of light, thought went with the ray through celestial space; the immeasurable was measured, the
incomprehensible understood. Fantasy gave up, but insight triumphed. To traverse the distance from the sun to
the earth, 20 million miles, light needs 8 minutes and 13 seconds; to traverse the distance from the 61st
star in the Swan, whose parallax Bessel has determined,
% The fraction stands in the printing as a stacked fraction within the line. \tfrac per the house
% rule (cf. „15\tfrac{5}{8} Fod“ on p. 50).
light needs $9\tfrac{1}{4}$ years; such a
% --- p. 166 ---
length of way is incomprehensible for intuition, and yet it is exactly apprehended by the understanding --- it
is, after all, measured. But when one must further conclude ``from the character of the telescope'' that the
light from the most distant stars in the Milky Way has been 2 or 3000 years on the way --- indeed, that there
are stars we see by pre-Adamite light --- then the immeasurable makes even thought giddy.

For the classical intuition of fantasy the world was limited; for the modern --- we may say: romantic ---
reflexion the universe is unlimited. Each of these ways of seeing is one-sided. If the actual universe really
extends as far as space, then it must be unlimited like space. But the unlimited is something unactual;
everything actual is limited: insofar as the universe is actual, it is limited. The apparent contradiction
that the universe is both limited and unlimited is, however, as already shown in the doctrine of space, no
insoluble antinomy. Astronomy itself has indicated
% The Humboldt citation: the quotation marks balance on this page --- three „ and three “. The
% interruption „, siger Humboldt (Kosmos I, S. 68),“ closes and re-opens, but also closes at the
% end; so NOT the book's familiar habit (pp. 10, 11, 16) of a single closing.
the solution of the riddle. ``If one compares,'' says Humboldt (Kosmos I, p.\ 68), ``world-space with one of
our planet's island-rich seas, then one may think of matter as distributed in groups: now as irresolvable
nebulae of various ages, condensed about one or several nuclei, now
% „sammenboldtet“ stands so in both copies; reproduced as printed (read „sammenholdt“).
as clotted together in star-clusters or isolated sporades. Our star-cluster, the \emph{world-island} to
which we belong, forms a lens-shaped, flattened layer separated on all sides, whose major axis is estimated
at seven or eight hundred, its minor by contrast at a hundred and fifty Sirius-distances.'' We are thereby
led to apprehend the universe not merely as an unsurveyable manifold of worlds and world-islands, but as
\emph{a system of systems}.

The \emph{physical} double stars, which by mutual attraction describe orbits about one another according to
the Keplerian laws, contain a confirmation that Newton's principles
% --- p. 167 ---
hold good outside our solar system too. Now if attraction is a property indwelling in all matter, then it
follows that not merely the planets in each single system must attract one another, but that the systems must
also act upon one another mutually and show themselves as members which, by the infinite roundabout way of
mediation, can be regarded as realizations of a single, all-embracing fundamental system. This way of seeing
is further confirmed by there being nothing absolutely at rest in world-space; everything is in motion, not
merely planets and comets but the sun itself and the suns: there are, properly speaking, no fixed stars. The
question presses forward: is there here a common centre for all these motions? Corporeal this centre surely
cannot be. As is well known, the moon does not move about the earth's centre, but earth and moon about a
common centre; and the same holds of the sun and the planets. When, then, our whole solar system has a
translatory motion, this does not exactly indicate that there is a single fixed body about which the motion
happens; it rather presupposes that the system is acted upon by another, or rather by other systems, so that
the common centre of gravity is determined by the systems' mutual attraction. In the universality of the
concept the matter is accordingly to be apprehended thus: the universe is limited, for every actual system
has its limit; the universe is unlimited, for outside the one actual system there is another actual one, and
outside this again another: every system is a whole for itself, and yet a part, a member of a more
comprehensive, still more composite system: the actual limit is everywhere and is everywhere overstepped. The
finite does not end here or there, any more than the infinite begins here or there; the infinite is not the
ceasing of the finite, does not mean that the finite ceases,
% THE LETTERSPACING RUNS OVER THE TURN OF THE LEAF: „d e t  E n d e l i g e  e r“ stands as the
% last words on p. 167 and is continued on p. 168. Confirmed on the KB copy on both leaves. The
% colon after „overalt“ is set outside the emphasis here.
but on the contrary that it repeats itself incessantly; \emph{the finite is
% --- p. 168 ---
everywhere, because the infinite is everywhere}: this is, in short, \emph{the universe's infinity}.

How must the objectification of a universe infinite in finitude now necessarily go on? As the object is, the
objectification must be, and conversely. That the infinite is in actuality finite, and the finite infinite,
comes to light most clearly of all from the newer science's art of observation and analytical methods. The
articulating reciprocal determination of finite and infinite, which is in equal degree characteristic of
observation and of calculation, points to the doctrine of \emph{infinite approximation}.

Although with the Keplerian laws corporeal centres were introduced, the centres of gravity in the doctrine of
gravitation do not always have material existence. They are determined by the co-working of given forces, and
are thereby distinguished from the empty designations of place in the epicycle theory. But besides such ideal
realities the newer astronomy has a quantity of feigned magnitudes, especially mean magnitudes --- mean sun,
mean anomaly, mean motion, mean time and so forth. To see what an infinite approximation signifies in this
respect it will be sufficient to fasten attention upon a single problem, for example the so-called Keplerian
one: the equation of the centre.

It is something peculiar about the mathematically empirical sciences --- and astronomy is a mathematically
empirical science --- that they demand the exact and yet must incessantly struggle with the approximate. The
stricter the method, the more plainly it becomes evident that there is a disagreement between the actual and
its apprehension, between the object and its objectification, which never quite admits of being lifted. The
old astronomy stopped in all naivety at the approximate. On the hypothesis that the sun with even velocity
describes a circle whose centre lies outside the earth, it was then a rather simple problem to
% --- p. 169 ---
% p. 169 carries TWO footnotes, marked *) and **). The body text on the leaf is only four lines;
% the rest of the page is the notes.
find the equation of the centre.\footnote{% printed mark: *)
Let the diameter passing through the earth be the line of apsides, its end points P and A, perigee and
apogee; let the eccentricity, the distance from the earth J to the circle's centre C, be denoted by
$\varepsilon$; the angle which the lines from J and C to the sun S form with one another, by E; the angle
PCS, the mean anomaly, by nt: one then has, for determining the equation of the centre,
\[
\frac{\sin E}{\sin (nt + E)} = \frac{JC}{CS} = \varepsilon
\]
or, in all approximation, $E = \varepsilon \sin nt$.}
If we start, by contrast, from the Keplerian laws and carry the analysis through in a peculiar direction, then
we get a sine series with multiplied arcs which is indeed strongly convergent, but nevertheless infinite:%
\renewcommand{\thefootnote}{**)}%
\footnote{% printed mark: **)
% THIS NOTE RUNS OVER TWO LEAVES. It begins beneath the rule on p. 169, fills the rest of that
% leaf, and is continued beneath the rule on p. 170 from „erholdes saa Ligningen“ to the end. The
% book does not repeat the mark on the continuation. Per the house rule the whole note stands
% here at its mark, and the page mark for p. 170 is set at the text's, not the note's, turn.
The ellipse's major axis is the line of apsides, and
% MISPRINT WHICH NIELSEN HAS NOT CORRECTED: the printing has „Fxcentriciteten“ with a capital F
% where E is meant. Read under high magnification on the KB copy and confirmed on the working
% scan --- both copies show an F. The compositor's wrong sort. Translated by the sense.
the eccentricity e. The earth, not the sun, is assumed to be in the focus. E is the angular distance between
the mean sun, a feigned sun, and the actual sun; m the mean anomaly. If we compare the expression from the
previous solution, $E = \varepsilon \sin nt$, with the first term of the infinite series opposite, it is seen
that the hypothesis of the eccentric circle, in that
% THE ERRATA LEAF corrects p. 169, line 19 from the top: „E = 2e, --- læs ε = 2e,“. Per the house
% rule the text stands here AS PRINTED, and the correction is not applied. The sort was checked on
% the KB COPY, precisely because the heavy 1-bit working scan blurs the difference between capital
% E and Greek ε: a capital, three-stroke E stands there.
$E = 2 e$, gives the equation of the centre a value agreeing with the assumption of the elliptical orbit, so
long, be it noted, as one is content with the approximate. But if the approximate is to be lifted, and the
demand for the exact satisfied, then the limit of finitude must be overstepped within finitude, to infinity.
That is, as regards the analyses, an enormous step. What transformations the mathematical method has had to
undergo before it became possible to get from the first term in the series opposite to all the following ones
--- of that one gets a collected impression by casting a glance at the elegant solution of Kepler's problem to
be found
% „Cours de Mécanique“ stands in italic; „Duhamel“ and „II. S. 63--69“ in ordinary antiqua. The
% page numbers read on the KB copy: 63--69.
in Duhamel (\textit{Cours de Mécanique} II, pp.\ 63--69). As a general foundation, reference is made to
Lagrange's formula $z = x + \alpha f(z)$, from which one gets:
% The numerator is printed „d.(f(x))²“ with the argument (x) in smaller, lowered type; hence
% $f_{(x)}$.
\[
z = x + \alpha f (x) + \frac{\alpha^{2}}{1.2}\ \frac{d.(f_{(x)})^{2}}{dx}
+ \ .\ .\ .\ .\ ;
\]
from this, with respect to the formulas for the problem:
\[
u = nt + e \sin u;
\]
\[
r = a (1 - e \cos u), \quad \operatorname{tg} \tfrac{1}{2} \theta =
\sqrt{\frac{1 + e}{1 - e}}\ \operatorname{tg} \tfrac{1}{2} u,
\]
% Here the note passes over onto p. 170.
one then obtains the equation:
\[
u = x + e \sin x + \frac{e^{2}\, d\, .\, \sin^{2} x}{1.2\ dx}
+ \ .\ .\ .\ \frac{e^{m}\, d^{m-1} \sin^{m} x}{1.2 .\,.\,. m\ dx^{m-1}}
+ \ .\ .\ .
\]
To carry out the differentiation, however, one must know in advance what is given in equations such as:
\[
(2 \sin z)^{m} \cos \frac{m \pi}{2} = \cos mz - \frac{m}{1} \cos (m-2) z
+ \frac{m (m-1)}{1.2} \cos (m-4) z \ .\ .\ .
\]
\[
(2 \sin z)^{m} \sin \frac{m \pi}{2} = \sin mz - \frac{m}{1} \sin (m-2) z
+ \frac{m (m-1)}{1.2} \sin (m-4) z \ .\ .\ .\ .
\]
% The source stands in the printing CENTRED on a line of its own beneath the two formulas. The
% placement is not reproduced.
(Ramus, Algeb.\ og Functionsl.\ p.\ 202.)
and so forth, before one (Duhamel, p.\ 69) reaches the definitive form of the series in question. For the sake
of the objective relations the analysing thinking must sink itself in operating, combining
reflexions. The thoughts in the objective must be thought \emph{subjectively} in order to be thoughts; an
unconscious objectification of thoughts is just as meaningless a turn of speech as a shining darkness, a
seeing blindness, an unconscious knowing.}%
\renewcommand{\thefootnote}{*)}
% --- p. 170 ---
\[
E = 2e \sin nt + \frac{5}{4}\ e^{2} \sin 2nt + \frac{1}{4}\ e^{3}
\left( \frac{13}{3} \sin 3nt - \sin nt \right) + \ .\ .\ .\ .
\]

Here we have, in the form of calculation, an example of infinite approximation.

If we fasten attention upon the task of observation, then the same shows itself from another side. No single
observation is exact; one can repeat an observation as often as one will, one commits an error every time,
though not the same one. It is not clumsiness on the observer's side, not an accidental teasing by fate's
caprice; it lies in the nature of the conditions, in the objective necessity of accidentality --- it is a law.
How characteristic that it is precisely
% MISPRINT WHICH NIELSEN HAS NOT CORRECTED: the printing has „dan exacte Videnskab“, where „den“
% is meant. Read on the KB copy under high magnification: a clean a, not an e; both copies show
% it. Translated by the sense.
the exact science that has become aware that there is a law of the errors, a \emph{law of error}. The calculus
% --- p. 171 ---
of probability, with the law of error and the method of least squares, is a telling testimony to the
persistent struggle which mathematically empirical knowing must wage in all directions against the approximate
by means of infinite approximation.

Astronomy's objectification of the world-structure culminates in approximation, and a culmination of
that kind is infinitely dialectical. From the one side it looks as though the actual, the objective in itself, withdrew
itself from the observer's measure to infinity; for if exact, definitive determinateness could really be
attained, then the approximation would cease of itself. From the other side actuality with its determinate
limits comes constantly nearer, for if there were a finite fixed limit to the approximation, then it would not
be infinite. The exact cannot be attained by a restricted number of single calculations and carefully verified
observations; for the one series of observations points back to the other, and of course the one calculation to
the other as well; but in this lies also that observations and calculations can be corrected and secured by one
another mutually. The more exact the various series of observations become, the more means the analysing
investigator obtains for correcting them further and approaching the exact; the residue of indeterminateness
that still remains does not disturb him by any foggy approximateness, for now he knows exactly between what
limits it falls, and under what conditions
% The mark stands after „forsvindende“ and BEFORE the full stop, as printed.
it can be regarded as vanishing\footnote{% printed mark: *)
Cf.\ \textit{Mathematik og Dialektik}, pp.\ 34--56.}.

Observing science must have instruments and set store by their perfecting; telescopes, refractors, universal
instruments and so forth play an important part in the exact investigations. It is
% The word-break runs over the turn of the leaf: „Iagt-“ on p. 171, „tageren“ on p. 172.
necessary for the observer

% --- p. 172 ---
% The leaf opens MID-WORD: „Iagt-“ stands on p. 171, „tageren“ here. No new
% sentence; division B runs on and closes only further down this page.
to know his instrument and understand how to use it. For the sake of the objectification the
philosophy of nature has yet to draw attention to one indispensable instrument --- in truth a
universal instrument --- whose share in determining the objectification is otherwise so easily
overlooked, and that is consciousness itself. Outside consciousness there is indeed a presupposed
actuality, but no knowing of its objectivity. Not merely in the determination of the objective, but in
the essence of objectivity --- be it the objectivity of the earth or of the sun, of the Milky Way or of
the comets --- the objectifying consciousness is an essential factor. However gigantic the steps exactly
empirical science has taken with the objectification of the actual world-structure, however much it has
in theory widened its ideal power by continued conquests from the immeasurable realm of the starry
heaven: the building itself it has not conquered --- the actuality of nature, the objective in itself, is
in small as in great a barrier to consciousness, a barrier which at every step reminds it that, infinitely
far from being the first objectifying agent, which objectifies in power, it must be content with its lot:
to objectify in powerlessness.

% Short centred rule between division B and division C, in the middle of p. 172.
\streg

% Danish: \afdeling{C.}{Solsystemets Oprindelse.}   (printed pp. 172--196)
\afdeling{C.}{The Origin of the Solar System.}

As an infinite system of worlds the universe is itself without a finite beginning. The becoming, arising
and passing away of the single worlds takes place in the fundamental system; but this
neither is nor becomes. To investigate what science at its present stage avails to elucidate about the
formation
% --- p. 173 ---
of our solar system, we bring out: from the standpoint of the facts the physical-astronomical hypothesis
of \emph{the primal nebula}, and from the standpoint of ``the concept'' the system's self-formation in
virtue of \emph{the immanent Idea}. By placing the two sharply opposed views under a common illumination
of objectifying omnipotence, we can then convince ourselves that both are insufficient, while they
nevertheless, each from its own side, furnish essential contributions to the clearing up of the matter,
when the question is of the system's origin and the fundamental relation between \emph{nature and Spirit}
is fully known.

% Run-in head, read on the KB copy: set in semibold antiqua (not letterspaced), on the same line
% as the paragraph it opens.
\runin{a) The primal nebula: the physical-astronomical hypothesis.} The ground why no particular weight
can be laid upon the supernatural presuppositions of creation and creating omnipotence, with which the
older physicists and philosophers, from Gassendi and Cartesius to Newton and Leibnitz, willingly equip the
foundation of their view of nature, has been treated in the preceding chapter. The only yield which the
doctrine of gravitation might be supposed to afford the supernaturalistic view would in any case consist
in this, that a more deeply penetrating insight into the mechanics of the world-structure gave the
beholder renewed occasion to trace and admire the divine wisdom. That the Creator must be as great a
mathematician as physicist followed from this standpoint plainly enough from the whole plan. What good
would it have been that the Almighty had ``implanted'' inertia and attraction in matter, if he had
forgotten to distribute the masses, to place sun, planets and secondary planets each in its determinate
position, and to calculate the first impulse so that the orbits could become ellipses? Nor is this all.
There is, namely, not only a mutual attraction of the sun and each single planet, but so manifoldly
crossing a reciprocal action among the planets themselves that the perturbations thereby caused manifestly
threaten the system
% --- p. 174 ---
with destruction. This the Creator has foreseen and from the beginning taken his precautions against. He
has indeed not wished to make the planets perfectly spherical, just as on the whole he has not shunned the
intricate, so that his mechanical wisdom might be properly revealed to the investigator of nature. But he
has seen to it that the orbits got only a slight inclination, that
% The formula is read on the KB copy: T over t = A^{3/2} over a^{3/2}, the exponents set as
% stacked fractions.
the periods of revolution became incommensurable:
$\frac{T}{t} = \frac{A^{\frac{3}{2}}}{a^{\frac{3}{2}}}$,
and, what is more, that the sun got so great a mass that it more than outweighs all the
rest\footnote{% printed mark: *) --- the only note on p. 174.
The sun may, even after Neptune's discovery, be assumed to be about 720 times greater than all the masses
together that move about it.}. The last especially is of extraordinary significance; for the one planet
cannot now act disturbingly upon the other without the sun in majestic superiority, so to speak, placing
itself between them and keeping order. The more one thus sinks oneself in the mathematical analyses, in
the manifold of infinite series through which a Laplace has developed the laws of perturbation, the more
the admiration of the world-creator's divine mechanics grows.

Amid such edifying considerations it must then be highly surprising to hear Laplace himself exclaim:
\textit{Philosophe, montre moi la main qui a jeté les Planètes sur le tangent de leur orbite!} What does
this challenge mean? Is the mathematician then, for all his wisdom, unable to discover the Creator's hand
and know his wisdom?

There is a fundamental proposition which says that \emph{everything in nature must go on naturally}. By
departing from this fundamental proposition natural science would make itself guilty of a contradiction as
dangerous to life as the peasant's in the tree when he sawed through the branch on which he himself was
sitting. But that
% --- p. 175 ---
everything in nature goes on naturally means, in other words, that all causes working in nature must
without exception be natural. There must then just as necessarily be a natural cause for the distribution
of the planets in the solar system as for their motions this very day; just as much a natural cause for
the preponderance of the sun's mass over the rest as for the difference between sidereal and solar time.
It is not an arbitrary, presumptuous craving out beyond the limit of what can be known, but a scientific
need for a coherent account of given facts, that lies at the ground of the world-hypothesis by which Kant
first and in detail\footnote{% printed mark: *) --- the first of two notes on p. 175.
Allgemeine Naturgeschichte und Theorie des Himmels. Kant's Werke. Achter Band. Leipzig 1838. S.\ 224 flgd.}
and afterwards Laplace in all brevity%
\renewcommand{\thefootnote}{**)}%
\footnote{% printed mark: **) --- the second of two notes on p. 175.
% The note is set in italic as far as „Tome second.“; „S. 293 flgd.“ is antiqua.
% „francaise“ is printed WITHOUT the cedilla --- both copies show it. Reproduced as printed.
\textit{Exposition du Système du Monde. A Paris. L'an IV de la republique francaise. Tome second.} S.\ 293
flgd.}%
\renewcommand{\thefootnote}{*)}
have sought to explain to themselves the natural coming-to-be of the planetary system.

Physical astronomy presupposes that matter with its forces is ``eternal'' like space, and then sets about
conceiving, out of the first existent, the becoming and stepwise formation which the present must in the
course of the ages have passed through. The first existent is to be apprehended as a state in which
everything existing can be assumed to have found itself from the beginning. By a magnificent reduction we
are then carried back to a matter evenly spread in space and as yet indifferent, an extraordinarily thin
but enormous mass of vapour whose diameter is estimated by some at 12,000 billion miles.

The formation now begins with the parts attracting one another, whereby condensation, rotation and heat
come forth. The mass of vapour becomes spherical, raised at the equator; the rotation increases, and the
centrifugal force grows: all of it natural! In the raised belt of the spherical mass there will
% --- p. 176 ---
then gradually form rings, which loosen, are separated off and roll themselves together again; the
rolled-together masses move about their axes and about the main mass: here we have a first sketch of the
system, in all naturalness. What happened with the great mass can repeat itself with the smaller ones;
some among these can again deposit one or more rings, and these rings be separated off and become moons.
Saturn's ring, which in actuality consists of three rings plainly separated from one another, is an element
arrested in its development, from which one may learn how it has gone with the formations on the whole.
Since the sun's mass is precisely the remaining main mass, we easily see the natural ground of its
preponderance; only concerning the closer character of the planetary orbits does a doubt remain which the
theory, if it is to be accepted, must be able to lift. It must indeed in general be conceded that we no
longer need a supernatural cause to explain the first impulse along the tangent, since the force of rotation
is already in the rotating mass; further, it follows from the law of attraction that all the orbits must
become conic sections. But of the three conic sections only one, the ellipse, is suited to give a closed
system of motions; the other two would by contrast make the system meaningless, since the planets, moving
either in a parabolic or in a hyperbolic branch, would gradually remove themselves from the sun without
ever returning. If one will not set it down to chance that there is reason in the solar system, then it
undeniably seems as though precisely the particular initial velocity on which the elliptical orbits are
conditioned must presuppose a particular arrangement of wisdom, and so after all point to an invisible
hand, an arbitrary
% The errata leaf corrects p. 176, line 4 from the foot: „Ingriben --- læs Indgriben“. The
% printing has „Ingriben“ without the d --- read on the KB copy, because the working copy carries
% a pencil correction just here. Translated by the sense.
intervention by a supernatural cause.

The defenders of the hypothesis find it by no means difficult to meet such misgivings. To prevent the orbits
from becoming parabolas, no divine
% --- p. 177 ---
intervention was in truth needed. Among infinitely many possible cases there is for the parabola only a
single
% The mark stands after „gunstigt“ and BEFORE the colon, as printed.
favourable one\footnote{% printed mark: *) --- the only note on p. 177.
If the radius vector be denoted by r, the semi-major axis by a, the intensity by $\mu$ and the velocity by
v, then the formula common to all three conic sections is:
% The sign in the parenthesis is a true minus (not Nielsen's ÷); read on the KB copy.
\[ v^2 = \mu\left(\frac{2}{r} - \frac{1}{a}\right). \]
Since now a is positive, negative or $\infty$ according as the conic section is an ellipse, parabola or
hyperbola, it is seen from this that it must depend upon the velocity (the initial velocity) which of the
three conic sections is to come forth. For the ellipse $v^2 < \mu\,\frac{2}{r}$, for the hyperbola
$v^2 > \mu\,\frac{2}{r}$, for the parabola $v^2 = \mu\,\frac{2}{r}$. Of the infinitely many possible cases
only a single one will accordingly be favourable to the parabola, while the number of favourable cases is
infinitely great both for hyperbola and for ellipse.}:
the parabola is, on the presuppositions contained in the hypothesis itself, the very least probable form;
but that the very least probable does not happen is quite natural. As regards the hyperbola, this curve
has, when one considers the matter in all generality, certainly a probability in its favour, just as much
as the ellipse; but when the closer circumstances are brought into account, then it is in the first place a
question whether the force of rotation in the rotating mass has ever been great enough to give an initial
velocity answering to the hyperbola. But even if the possibility of this be conceded, what then? Then the
consequence will in any case only be that some planets must in primeval times have left the system, and
that we have only those left which got a smaller initial velocity.

If one goes into detail with the hypothesis's application, then there are indeed not a few facts --- all
motions from west to east; the axial rotations answering to the sun's
% --- p. 178 ---
rotation and to the directions in the orbits; small inclinations; slight eccentricity; the outer planets'
more rapid rotation and slower orbital
velocity\footnote{% printed mark: *) --- the first of two notes on p. 178.
Spectral analysis too seems to speak in the hypothesis's favour. We will not, however, omit to adduce the
following, though, as it appears to us, somewhat weak objections by a distinguished French geognost: Mais
l'hypothèse émise par Laplace, il ne la présentait lui-même qu'avec „defiance“ . . . ses conjectures
n'expliquent point la forme elliptique des orbites planétaires, l'inclinaison de leur axe; elles paraissent
en outre démenties par le mouvement rétrograde des satellites d'Uranus
% The points are set spaced, as printed: three after „defiance“, four after „d'Uranus“. The italic
% begins at „Enfin“ --- confirmed on the KB copy; all before it is ordinary antiqua.
. . . . \textit{Enfin la découverte de l'analyse spectrale, qui sera l'éternel titre de gloire de M. M.
Kirchhoff et Bunsen, autorise à croire que la composition chimique du soleil diffère
% „très-notablement“: the hyphen falls at the line break, so it cannot be decided whether it is
% the compositor's division mark or a real hyphen. The older French spelling with the hyphen is
% followed.
très-notablement de celle des planètes de son système, car cet astre ne contient, du moins dans ses couches
extérieures, ni silice, ni étain, ni plomb, ni mercure, ni argent, ni or (!).} La Terre par Élisée Reclus.
I. Paris 1870. pag.\ 19.}
and so forth --- which in broad outline speak for its validity; but there are nevertheless also several
phenomena which will not readily accommodate themselves to
% The errata leaf corrects p. 178, line 5 from the top: „dem. --- læs den.“ The printing has
% „dem.“ --- read on the KB copy, because the working copy carries a pencil correction here.
% Translated by the sense.
it. Everything does not fall out so neatly as Kant assumes.

% The quotation runs over the turn of the leaf and is closed on p. 179 („Men“). It is not a
% misprint and gets no note.
``The telescopic planets Vesta, Juno, Ceres, Pallas (and the new
% The errata leaf corrects p. 178, line 9 from the top: „hældende --- læs heldende“. The printing
% has „hældende“ --- read on the KB copy, because the working copy carries a pencil correction
% here.
Astraea), with their mutually intertwined, strongly inclined and more eccentric orbits, have,'' says
Humboldt%
\renewcommand{\thefootnote}{**)}%
\footnote{% printed mark: **) --- the second of two notes on p. 178.
Kosmos.\ I.\ S.\ 71 flgd.}%
\renewcommand{\thefootnote}{*)}%
, ``been considered as a zone dividing our planetary system into two spatial sections, indeed as a kind of
middle group within it. On this view the \emph{inner} planetary group (Mercury, Venus, the earth and Mars)
presents, in comparison with the \emph{outer} (Jupiter, Saturn and Uranus), several striking contrasts. The
inner planets, nearer the sun, are of middling size, denser, fairly alike and slowly rotating (in about a
24-hour period), less flattened, and, with one exception, namely the earth, altogether moonless. The
\emph{outer} planets, further from the sun, are
% --- p. 179 ---
considerably larger, five times less dense, have twice as rapid an axial rotation and are therefore strongly
flattened as well as moon-rich, in the ratio of 17 to 1, if it is confirmed that Uranus really has
% The mark stands AFTER the full stop, as printed.
6 satellites.\footnote{% printed mark: *) --- the first of two notes on p. 179.
In the most recent time only 4 moons have been expressly observed.}
But'' --- he adds --- ``this general consideration of certain properties characteristic of whole groups does
not admit of being applied with equal right to the single planets in each of the groups; nor to any relation
between the single circling heavenly bodies'
\emph{distance from the central body}, or to their absolute size, their density, period of rotation,
eccentricity, inclination of orbits and
% THE QUOTATION IS NOT CLOSED. The citation opens with „denne almindelige Betragtning“ and gets no
% closing “ after „Axe**).“ --- both copies show it, and the errata leaf does not mention it.
% Reproduced as printed; the page's quote balance therefore stands at +1.
% The mark stands BEFORE the full stop, unlike line 4 from the top of p. 179.
axis%
\renewcommand{\thefootnote}{**)}%
\footnote{% printed mark: **) --- the second of two notes on p. 179.
As an example that we know no mechanical law of nature which, ``like the beautiful law binding the squares of
the periods of revolution to the third powers of the major axes of the orbits, makes the planetary bodies'
six elements just named, or the form of their orbits, dependent either upon one another or upon the mean
distances,'' it is adduced that Mars, further from the sun, is smaller than both the earth and Venus; that
Saturn is smaller than Jupiter and yet much larger than Uranus; that Uranus itself seems to be denser than
Saturn; that we, as regards the inner group, find on both sides of the earth both Venus and Mars less dense
than the earth itself. The time of rotation (about the axis) does indeed on the whole decrease with the
distance from the sun; but it is nevertheless greater with Mars than with the earth, greater with Saturn than
with Jupiter. Saturn stands midway between Jupiter, whose axis of rotation is almost perpendicular, and
Uranus, whose axis by contrast almost coincides with the orbit; and so forth.}%
\renewcommand{\thefootnote}{*)}%
.

% MISPRINT: „antage ig“ for „antagelig“ --- the l is missing, and a clean word space stands in its
% place. Both copies show it, so it is the compositor's fault (a dropped sort); the errata leaf
% does not mention it. Translated by the sense.
The hypothesis may on the whole be acceptable, even if it does not succeed in accounting completely for all
the details and removing all misgivings. What is peculiar in the co-working of conditions and forces under
which those first formations went on loses itself in the darkness of primeval time, inaccessible to special
investigations. The hypothesis has probability in its favour so long as no facts can be pointed out that
stand in manifest conflict with its assumption. But suppose even that irrefutable objections
% --- p. 180 ---
led at length to a straightforward rejection of the hypothesis; would it then follow that natural science
were compelled to resort to supernatural presuppositions and make a loan from theology? By no means! If the
hypothesis cannot be defended, then one of two things happens: either the question of the solar system's
formation is held in abeyance as unanswerable, or a new and more satisfactory hypothesis forms itself. In
both cases science will with equally unshakeable certainty hold fast its fundamental proposition, the
irrefusable fundamental proposition that everything in nature must go on naturally. This is what is
decisive; the manner of explanation is not what is decisive. Even if the solar system's stepwise
coming-to-be, by transitions from state to state out of a primal state, could be established by a theory as
trustworthy as the doctrine of gravitation, a theory of a series of alterations of state would still be no
explanation of the origin of the whole; but such a theory would assuredly, by clarifying our view of the
inner concatenation of the natural causes, furnish a great real contribution to the determination of
\emph{the concept: nature}.

% RUN-IN HEAD WITH THE WRONG LETTER. The errata leaf corrects p. 180, line 13 from the foot:
% „c) Systemets --- læs b) Systemets“; the Indhold and the sequence (a) on p. 173) require b) too.
% The printing has „c)“, read on the KB copy. Reproduced as printed, not renumbered.
\runin{c) The system's self-formation: the immanent Idea.} While physics with its hypothesis holds
exclusively to phenomena and conditioned causes, the speculative doctrine of immanence goes the other way,
taking decisive aim at the unconditioned: \emph{the Idea}. The Hegelian deduction of the solar system is
in this respect remarkable.

Nature is more than a swarm of facts; it is the Idea itself in the outwardly existent, or in the form of
being-other. It lies in the divine Idea to open itself up, in order to be subjectivity and Spirit. From this
standpoint the solar system must be apprehended; for its ``absolute mechanics'' is the form in which the
Idea of nature has realized itself. The abstract forms, space and time, the self-conceiving Idea has
% --- p. 181 ---
filled with matter, whose essence is gravity; for gravity consists precisely in this, that matter, seeking
its point of rest outside itself, moves from place to place and by place and motion cancels the
contradiction of time and space. The motion is threefold: α) the mechanical, imparted from without, is
uniform; β) the half-conditioned, half-free motion, namely the motion of fall, shows that a body's
separation from its gravity is still accidental, while the motion nevertheless already belongs to gravity
itself; γ) the unconditionally free motion, ``the great mechanics of the heavens,'' is seen in gravitation.

``\emph{Gravitation} is the true and determinate concept of material corporeality, realized into Idea.
Universal corporeality divides itself essentially into particular bodies and joins itself together, as the
phenomenal being of motion, into the moment of singularity or subjectivity, which is herewith immediately a
system of several bodies.''\footnote{% printed mark: *)
Vorlesungen über die Naturphilosophie. Berlin 1842. S.\ 94 flgd.} --- There can be no well-grounded question
whence these bodies come; they are real moments of the Idea's being: they exist, for in them the Idea has
existence. The independent bodies of which the solar system consists refer themselves essentially to one
another, are heavy, but hold themselves in this relation and set their unity in an other outside themselves.
Thus the plurality is no longer indeterminate, as with the stars, but the difference is posited; and the
determinateness rests solely on \emph{absolutely universal centrality} and on \emph{particular centrality}.
From these two determinations follow the forms of motion in which matter's concept is filled out. The motion
falls to the relative central bodies, which are in themselves the universal determinateness of place: and
yet their place is at the same time indeterminate, insofar as it has its centre in an other; this
indeterminateness must then likewise have being, while the place determinate in
% --- p. 182 ---
and for itself is only one. It is therefore indifferent to the particular central bodies at what place they
are; and this comes into view in their seeking their centre, i.e., leaving their place and setting
themselves at another place. The third is this: taken immediately, these bodies might be removed equally far
from their centre; were they so, then they would not be removed from one another. If they moreover all moved
in the same orbit, then they would not be different from one another at all: they would then be one and the
same, the one only a repetition of the other, and their difference of course only an empty word. The fourth
is this, that as they alter their place at different distances from one another, they return into themselves
through a curve; for only thereby do they present their independence against the central body --- as also
their unity with the centre by moving about it in the same curve. As independent against the central body
they keep to their place and no longer fall into the centre. The double-sidedness of the relation is that
the particular bodies at once posit themselves a central body and are posited by the central body: the
centre is meaningless without periphery, and the periphery without the centre.

On this foundation the Keplerian laws are to be deduced from the pure Idea. As regards first the orbit's
form, the circle is only an orbit for plainly uniform motion; but in free motion space and time must hold
according to their difference, and hereby the abstract, simply uniform velocity is cancelled. Thinkable, as
one says, it may well be that a uniformly increasing and decreasing motion can happen in a circle. But this
thinkability or possibility is only an indeterminate, superficial representation. The circle is the line
returning into itself in which all the radii are equal; that is, it is completely determined by the radius,
or: the radius is the whole determinateness.
% --- p. 183 ---
But in free motion, where determinations of space and time in their \emph{difference} enter into a
qualitative relation to one another, this relation must in the spatial itself necessarily come forward as a
\emph{difference} within the same, which expressly demands \emph{two} determinations. The demand is that
different arcs be traversed in the same time. If the arcs are to be different, then the radii must be so
too: from the accelerating motion follows difference in the radii. Thereby the form of the orbit returning
into itself becomes essentially an \emph{ellipse}; --- the first Keplerian law.

The abstract determinateness which is the circle's distinguishing mark shows itself also in this, that the
arc or angle enclosed by two radii is a magnitude \emph{independent} of them, an empirical one. But in the
motion determined by the concept, the distance from the centre and the arc traversed in a certain time must
be comprehended in \emph{one} determination, must make up \emph{one whole} (the concept's moments are not
accidental to one another); then there comes forth a spatial determination of two dimensions, \emph{the
sector}. The arc is in this way essentially a function of the radius vector, and, as unequal in equal times,
carries the inequality of the radii with it. That the spatial determination becomes, by means of time, a
determination of area hangs together with what holds of fall: that the same determinateness must occur the
one time as time in the root, the other time as space in
% THE ERRATA LEAF corrects p. 183, line 9 from the foot: „Qvadratet*) --- læs Qvadratet*).“ --- a
% missing full stop after the footnote mark. The KB copy has NO full stop after „*)“. The working
% scan by contrast shows a small mark on the baseline after the parenthesis --- the pencil has
% again got there before us, just as at p. 160. Reproduced AS PRINTED, without the stop.
% The note extends over the turn of the leaf: it begins at the foot of p. 183 and ends at the foot
% of p. 184 (without the mark repeated); per the house rule it is set here in its entirety.
the square\footnote{% printed mark: *)
Running time is negative unity, while space spreads itself out as something existent. Space is measured by
time in such a way that the units of space give the numerator, the chosen unit of time the denominator (time
in the root). Now as the units of space, under the accelerated velocity of the motion of fall, reflect
themselves in the unit of time, there comes forth a number which is the unit of time's own product, and so
the square of time, a measure for space. This is the most intelligible paraphrase we are able to give of the
enigmatic words of which Hegel (Anfr.\ Skr.\ p.\ 88) says
% THE ERRATA LEAF corrects p. 184, line 2 from the foot: „ans --- læs \emph{aus}“. The KB copy
% unambiguously prints „ans“ (n, not u) in the German sentence; the working copy has over the same
% letter a light pencil mark making n into u --- again the pencil before us. Reproduced AS PRINTED:
% „ans“. No colon after „siger“, and the German is ordinary antiqua, not italic.
Diess ist der Beweis des Gesetzes des Falls ans dem Begriffe der Sache.} Only that
% --- p. 184 ---
here the quadratic element in space is, by the line of motion's return into itself, limited to the sector.
These are, as one sees, the general principles on which the second Keplerian law --- that in equal times
equal sectors are cut off --- speculatively rests.

This law, however, concerns only the arc's relation to the radius vector; and time is moreover the abstract
unity in which the different sectors are compared, because it is what determines as unity. But the further
relation is a relation of time, not as unity but as quantum in general, as period of revolution, to the
orbit's size or --- what is the same --- to the distance from the centre. As root and square we saw time and
space stand to one another in \emph{fall}, the half-free motion, which is indeed from the one side determined
by the concept, but from the other side nevertheless determined outwardly. But in absolute motion, in the
sphere \emph{of the free measures}, every determinateness gets its totality. As root, time is a merely
empirical magnitude, and in a qualitative respect only abstract unity. As a moment in the developed totality,
by contrast, it is at the same time a unity determined thereby, a totality for itself, produces itself and
refers itself therein to itself. As what is in itself dimensionless, time in its production comes only to a
formal identity with itself, to the \emph{square}; space, by contrast, as the positive outside-one-another,
comes to the concept's dimension, to the \emph{cube}. This is the third Keplerian law, the relation of the
cubes of the distances to the squares of the times; --- ``a law which is so great precisely because it is so
simple and presents the matter's own reason.''

% The word-break runs over the turn of the leaf: „Be-“ on p. 184, „stemmelse“ on p. 185.
As regards the bodies themselves, they have, for their
% --- p. 185 ---
determination of their different nature, the principal sides that make up the moments of their concept.
One of them is the universal centre for the abstract relation to itself. Against this extreme stand the
\emph{immediate}, centreless \emph{singulars}, existing outside themselves, likewise in the form of
independent corporeality. The particular bodies stand under the determination of being-outside-self as well
as of being-in-self; they are centres for themselves and yet refer themselves to that first as to their
essential unity. The central body presents the abstractly rotatory motion. The centre is to be a point; but
being a body, it is at the same time extended and consists of what seeks. This dependent matter, which the
central body has in itself, requires that it must rotate about itself. For the dependent points, removed
besides from the centre, have no fixed determinate place --- they are only falling matter and thus
determined only in one direction. The remaining determinateness is wanting; every point must therefore
occupy all the places it can occupy. What is determined in and for itself is only the centre; the remaining
outside-one-another is indifferent, for hereby only the place's distance is determined, not the place
itself. This accidentality of the determination comes to existence in matter's altering its place, and is
expressed by the sun's rotating-in-itself about its centre. This axial motion is no change of
place\footnote{% printed mark: *)
That, according to ``the mechanics of the heavens,'' there can hardly be any axial rotation without a
translatory motion, may easily be overlooked when one looks only to the ``absolute Idea.''}; for all the
points keep the same place with respect to one another. The whole is a reposing motion. --- The
\emph{dependent} bodies, which at the same time have a free existence and do not make up connected parts of
the extension of a body endowed with a centre, also have rotation, but not about themselves, for they have
no centre. They rotate accordingly about a centre that
% --- p. 186 ---
belongs to another individual body, from which they have been expelled. Their place is in general this or
that; and this accidentality of the determinate place they then express by their rotation. But their motion
is a sluggish and rigid motion about the central body, since they constantly remain in the same position
with respect to it, as for example the moon in its relation to the earth. The dependent heavenly bodies form
that side of the concept which is called particularity; therein lies that with their difference they fall
apart from one another, for in nature particularity exists as twoness, not as in Spirit in the form of
\emph{one}. First, then, the moment is posited that the reposing motion becomes a restless motion, a sphere
of excursion or striving out beyond immediate being toward a beyond. This moment of being-outside-self is
itself the substance's moment as a mass, a sphere; for every moment here obtains its own being. This sphere
is the \emph{cometary}, and expresses a vortex, a constant standing on the point of dissolving and
dispersing into the infinite and empty. The extreme of excursion consists necessarily in approaching the
central body's subjectivity once and then yielding to the repulsion. But this being-other has immediately
its opposite in itself, an opposite which is: to seek its centre, its rest --- the cancelled future, the
past as moment; this is the \emph{lunar} sphere, which points not to an excursion from immediate being but
to what has come to be, or to being-for-self, to the self. The cometary sphere is therefore referred only to
the immediately axis-turning body, the lunar by contrast to the new centre reflected into itself, to the
planet. --- The \emph{planetary} sphere is finally the one that expresses relation both to itself and to
another; it is axis-turning motion and yet has its centre outside itself. The planet has its centre in
itself; but it is only a relative one; it does not have its absolute centre in itself, and is so far also
% --- p. 187 ---
dependent. The planet has both determinations in itself and presents them both as change of place. As
independent it shows itself in this, that its parts themselves alter place with respect to the position they
have to the straight line connecting the absolute and the relative centre; this grounds the planet's
rotatory motion. The planet is, as the third, the syllogism in which we have the whole. The solar, planetary,
lunar, cometary nature! this fourfoldness of heavenly bodies forms the completed system of rational
corporeality. The planetary nature is the totality, the unity of the oppositions, while the others present
only its isolated moments; the planetary nature is the most perfect, also with respect to the motion, which
here alone comes into consideration.

The fundamental weakness in this deduction is conspicuous. By one-sidedly fastening attention upon the Idea
and ideal causality, the thinker comes at every decisive turning point to overlook the actual conditions and
mechanical causes. Here one argues from grounds; that is the old Aristotelian manner. But as little as causes
in nature admit of being known without grounds, just as little do grounds in nature admit of being known
without causes: \textit{vere scire est per causas scire} --- that is a fruitful word of Baco's.

None the less everyone who understands how to read will get an impression that a magnificent endeavour with a
great thought stirs and breathes in this deduction, an endeavour aiming at cancelling the patchwork, carrying
the phenomenal determinations back to the Idea's unity and beholding the universe in the light of reason. That
nature is Idea means, in other words, that it is not a piecing together of countless fragments, but one
infinite essence connected in itself with one ground; that all single facts, all particular phenomena and
laws, are side-determinations
% --- p. 188 ---
of the one fundamental essence, real moments in the one content-rich concept.

% Run-in head, read on the page itself (KB copy): „c)“ in ordinary antiqua, „Natur og Aand.“
% SEMIBOLD --- the same setting as the other a)/b)/c) heads. This edition does not reproduce the
% difference between semibold and letterspaced.
\runin{c) Nature and Spirit.} Physics and the philosophy of immanence are agreed that nature is a whole
connected in itself, and the world-system an active revelation of eternal reason; but from neither side is
there offered us any clear or penetrating account of what is really to be understood by the word
\emph{nature}.

Physics simply refers the question away from itself; it treats of everything that is in nature, but what
nature itself is it does not treat of. When it goes back to the primal nebula, it understands this in the
sense that the primal nebula was originally in nature, that there was nature in the primal nebula; but that
nature is in its essence a primal nebula, or that the primal nebula is nature itself, it by no means asserts.
The philosophy of immanence takes the question up and answers it resolutely with the declaration that nature
is the absolute, the Idea itself in sensible outward being; but the answer is illusory.

The mysticism of the philosophy of immanence shows itself in its making, without more ado, the absolute into
Idea and the Idea into the absolute. The Idea for itself as Idea is by no means the absolute. The pure Idea
is the hollowest, most powerless, most ambiguous being, when it is to be held fast in its own absoluteness,
separated from every corresponding medium. \emph{No Idea without a medium!}

On the deduction given above it looks almost as though it were the planets themselves, in concrete
corporeality, that out of logical interest in a system of three syllogisms agreed to meet --- not at the same
place but, for difference's sake, each at its own place --- and came to an understanding about setting
themselves a common centre, the sun; but since simple planets must at bottom be supposed too stupid to see the
concept's logical point, it cannot be
% --- p. 189 ---
the planets simply, but the planet-Idea, the Idea existing as planet in planets, that conceives the whole.
Matter is indeed conceiving too, but in unclarity, in obscurity; the clear conceiving, the pure thinking, is
in space and time. It is then not existing matter with properties and forces, but a dialectical process
between speculative space and reflecting time, that has made the Keplerian laws necessary. In order that two
such conceiving parties as time and space might in the concept's name each get its own, the difference between
them had absolutely to come forth, and the planetary orbits to become ellipses; time had to become root, so
that the quadratic element in space could be limited to a sector. But time would not be content with being
root; it strove further forward, it wanted to be squared, wanted by the Idea's way to rise to being square.
Only when time attained ``the formal identity'' of being square, while space advanced to the concept's
dimension, to cube, did the solar system come into order: the solar, cometary, lunar, planetary nature merged
into the Idea.

As essential content the Idea is common to nature and Spirit; but, it must be expressly noted, the Idea stands
in another relation to Spirit than to nature; for Spirit it is a derivative, for it is determined and formed by
an omnipotently conceiving will; for nature it is an original, for it determines and forms the natural whole.
With this standpoint we now return to the cosmogonic hypothesis, to the primal nebula.

% Greek sub-head, run-in and LETTERSPACED; read on the KB copy under strong magnification: α
% (italic alpha), not „a“ or „&“, which is what the OCR gives.
\subrunin{α) Origin and beginning.} The primal nebula with which the beginning is made must itself have a
beginning; for if it be assumed that the matter spread through space exists from eternity, then we have, as
shown above, in this beginningless thing an existent without coming-to-be, a nebular phenomenon without
essence, an existent without ground, a dreamt presupposition, a natureless, unnatural first. But with a
natureless, unnatural first no formation of nature can begin. According to
% --- p. 190 ---
Kant, the whole infinite world-space is from the beginning to be filled with a matter whose properties and
forces lie at the ground of all alterations. ``This fundamental matter, an immediate consequence of the divine
existence, must,'' he says, ``at once be so rich, so complete, that the development of its combinations in the
course of eternity can spread itself over a plan which encloses in itself everything that can be, which admits
of no measure --- in short, which is infinite.''\footnote{% printed mark: *)
Allgem.\ Naturgeschicht.\ u.\ Theor.\ des Himmels. 7tes Hauptstück. S.\ 319.} The
theological-supernaturalistic representation of a single miracle of creation is the scientific makeshift Kant
still%
\renewcommand{\thefootnote}{**)}%
\footnote{% printed mark: **)
1755. Kritik der reinen Vernunft appeared only in 1781.}%
\renewcommand{\thefootnote}{*)}
has in reserve.

Now it might well seem as though a miracle lying an eternity back in time, a single general miracle which
moreover happened with necessity and did the whole thing at one stroke, might after all be tolerated in science
--- just for a beginning; and yet the contradiction is easy to discover. Existing matter cannot by an immediate
word of power spring out of self-power; matter belongs to nature as nature and must have a natural beginning.
Matter's coming-to-be is object-power's first real possibility, and the issuing of power is determined solely by
the real concept of the atoms. The beginning happens as the concept is realized. Of what character the matter
thus coming to be will otherwise be, what properties its atoms and masses will assume from the beginning,
depends solely upon the peculiarity of the concepts of power that determine from within; but their peculiarity
depends in turn upon the peculiarly determinate plan, the content-rich disposition, the Idea.

The Idea is accordingly indwelling in nature; the development is in its beginning as in its continuation a
conditionally necessary
% --- p. 191 ---
realization of determinate possibilities, of dispositions determining from within. Nature does not give itself
its possibilities, but realizes the given possibilities; it does not impart dispositions to itself, but develops
the imparted dispositions: \emph{all development of nature is, essentially regarded, a development of imparted
dispositions}. Hereby we are carried with scientific consistency from the concept \emph{beginning} to the concept
\emph{origin}.

The beginning is natural, but the origin supernatural. Nature's imparted dispositions presuppose an imparting of
dispositions; this imparting does not stem from anything whatever in nature; it is to be sought solely in what
--- not outwardly but inwardly --- is above nature, and so in the supernatural itself, in Spirit. The Idea of
nature's development could not possibly be determining for a self-developing natural whole if such a determining
Idea did not originally have an inner determinateness in itself, if it did not in its essence contain a system of
strict, inviolable fundamental determinations. Whence, then, has the Idea these determinations? A nature-Idea can
as little determine itself as a nature-thought can think, or a nature-concept conceive itself. The
all-thinking, all-conceiving power-essence that determines all Ideas is Spirit as Spirit, the omnipotent
self-being. If it is then true that object-power does not stand finitely outside self-power; if it is true that
all determinations of object in the doubling of power, through real concepts and dispositions, point to sole
power's absolutely overreaching self-determination: then it is also true that this world is created, that the act
of creation is Spirit's omnipotent act of will; then it is true that the whole of actual nature has a natural
beginning, not \emph{although}, but precisely \emph{because} it has a supernatural origin.

% Greek sub-head, run-in and letterspaced; the letter is read on the KB copy: β, not „B“, „ß“ or
% „P“, which is what the OCR gives.
\subrunin{β) Fundamental will and formative drive.} Although physics is on the whole inclined to presuppose an
eternally
% --- p. 192 ---
% MISPRINT: the Danish has „i Universitet samme Ælde“, where the sense requires „i Universet“. Both
% copies print „Universitet“ (checked on the KB copy), and the book's own errata leaf does not touch
% p. 192. Translated by the sense.
existing matter, it is nevertheless far from wanting to ascribe the same age to the many worlds in the universe;
it assumes on the contrary that the different world-systems begin their formation at highly different points of
time. ``Just as in our forests,'' says Humboldt, ``we see the same manner of formation simultaneously at all
stages of vegetation, and from this consideration, from this coexistence, receive a visible impression of the
onward-striding development of life, so we know also the different stages of star-formation in the great sea of
worlds. The process of condensation, which Anaximenes and the whole Ionian school taught, seems here to go on as
it were before our eyes.''

To explain to himself the naturally necessary connexion between the many world-formations setting in at different
times, Kant assumes that the infinite whole must from the beginning have a centre of condensations and motions
determined by the Creator's wisdom, a first material central point, and that formation then follows upon formation
as the ordering motion spreads itself and gets the better of the confusion in the original chaos. There is in this
way of seeing a double crookedness. In order to preserve the natural series of causes, Kant overlooks that to
every peculiar world-formation there must answer a peculiar disposition, and so a peculiar act of creation. He
further overlooks that an act of creation, scientifically understood, by no means interrupts but on the contrary
makes possible a law-determined connexion among the world-causes themselves: the many worlds in celestial space
make up a system of systems because all the world-Ideas make up moments of one and the same Idea of the universe,
because all creating acts of will proceed from one divine will, an eternal fundamental will.

That the arrangement of the solar system cannot immediately derive from God's hand, but must be a product of
mechanically working causes, Kant will prove in the following way. All the planets move in the same direction
about the sun;
% --- p. 193 ---
this circumstance can be explained only from the mechanism out of which the motions have arisen. If God with his
own hand had set the planets in motion, his choice --- since there is here not the least motive why he should
bind himself to one single determination --- would have displayed more freedom in all manner of deviations and
differences. Why do the planetary orbits fall almost in one and the same plane, and yet not quite? Why do they
approach being circles and are yet slightly eccentric? In all such things we recognize nature's usual procedure;
it strives after something determinate but does not quite reach it; it must employ various intermediate causes,
but the intermediate causes make it bend a little aside, so that it never hits the matter to a nicety. If it were
true, as a philosopher has said, that God constantly practises geometry, and if this could be clearly seen along
all the ways of the laws of nature, then quite another perfection and geometrical exactness would be traced in
the immediate works of the almighty word.

The kernel of truth in the Kantian reasoning is that actual formations of nature are due to working causes in
nature, and that it is particularly the mechanical causes which in their reciprocal actions take up the outward
cases with all their accidentalities. It is precisely by the accidental that Kant will recognize the natural;
against this there is nothing to object. But the manner in which he in his ``theory of the heavens'' sets the
fulfilment of the general laws against the plan itself in the legislating will betrays a one-sided reflexion. With
a human plan and its execution it does indeed sometimes go so that supervening circumstances, incalculable
accidents, unforeseeable hindrances can effect such thoroughgoing alterations in the execution that the realized
plan does not at all, or at any rate only in part, answer to the master-builder's end. But a divine plan cannot
possibly be exposed to such unforeseeable hindrances.
% --- p. 194 ---
There is in the solar system's disposition, in the primal nebula, an inner manifold of dispositions to
dispositions, of possibilities, of conditions for conditions, of causes for causes. The contraction is cause of
condensation and rotation, the rotation again a cause of the parts' heaping together at the equator, and so forth.
The cause of the earth's formation and its motion about the sun, the cause of the moon's formation and its motion
about the earth, all such causes with their conditions, which develop themselves time after time, are from the
beginning enclosed in the ground-cause of the disposition, so that the will which has willed the disposition and
has fully known what it itself willed has in each and all willed the whole development. As the many outward causes
with the many conditions, inward and outward alike, gradually enter into activity, the supervening circumstances
do indeed co-operate, but nothing, nothing at all, is left to mere accidentality: the laws lay claim to the cases
they meet, so that all of them without exception must punctually serve the execution of the plan.

The creating will uses the cases, wills the accidental, because it in earnest wills the natural; it wills a
naturally necessary formation of the parts because with Spirit's freedom it wills a naturally necessary formation
of the whole. The whole's development through millions upon millions of ages is in all its details a natural
fulfilment of the creating fundamental will. Between \emph{the fundamental will} within and the manifold of
mechanical causes without lies \emph{the formative drive}, spiritual and yet natural, an essential link between
nature and Spirit.

% Greek sub-head, run-in and letterspaced; the letter is read on the KB copy: γ, not „y“ or „7“.
\subrunin{γ) Teleology and mechanism.} The old dispute about the relation between the mechanical causes and the
so-called causes of purpose must now in principle admit of being resolved. What is the purpose of the solar system?
A natural development of everything enclosed in its original disposition:
% --- p. 195 ---
the development of the disposition is the attainment of the end. The development happens in virtue of mechanical
causes; but when the development is the goal and the mechanical causes the means, then there can in the matter
itself be no conflict between the mechanical and the teleological. Only when the consideration is entangled in
finite representations can the conflict arise. Either one drags ``the purposes'' down into the finite and offends
by arbitrariness, or one holds the real causes fast in the parcelling out and offends by necessity.

To drag the purposes down into the finite is to bring out certain arrangements beneficial and salutary for life,
especially for men, and to apprehend them so singly and so separated from the connexion of the whole that it looks
as though the Creator had arbitrarily had the attainment of such results in view. It is a teleology so egoistic
that it varies and becomes different according to the standpoint on which the beholder is placed. Why does the full
moon stand higher in the heaven in winter than in summer? On the Northerners' teleology it must be the Creator's
purpose that it shall light up the long winter nights. A pity that there are also long nights without a full moon!
Why has the earth's axis got an inclination to its orbit? From the standpoint of an earth-teleology this must of
course have happened with the all-wise purpose that the separation into hot, temperate and cold zones, with the
alternations and differences thereby arising, should furnish an inexhaustible wealth of conditions for the
development of organic life on the earth. Unfortunately those who live under the line, and those who dwell about
the poles, must in spite of the all-wise fatherly purpose feel themselves stepchildren.

To hold the causes fast in the parcelling out is to make conditioned necessity everything. To this answers the
altogether perverse, though among several physicists not
% --- p. 196 ---
uncommon, representation that a natural phenomenon cannot possibly express a certain relative end when it demonstrably
proceeds as a law-determined effect of mechanical causes. If the teleological consideration at the stage of finitude
is one-sidedly subjective, then the merely mechanical one is assuredly just as one-sidedly objective. As an
investigating science, physics has meanwhile --- however finite its mechanism may be --- the unmistakable advantage
that its investigations are fruitful and its groundings scientific. Finite teleology, by contrast, whether it appeals
to inner or to outer purposiveness, is and remains unscientific; of it holds Baco's word: \textit{Causarum finalium
inqvisitio sterilis est, et tanqvam virgo Deo consecrata nihil parit.}

From the standpoint of infinity, of principle, that is, of Spirit, the matter stands otherwise. The development of
nature is first completed when it in each and all satisfies Spirit. It does indeed satisfy Spirit at every stage, but
surely only insofar as every stage marks a step toward the completion. The question whether the solar system has its
goal solely in its mechanical form, or is destined to furnish the foundation for the attainment of other and higher
ends of nature, has scientific significance and coincides with the question of principle: what forms the development
of nature requires if the whole conceptual content answering to the disposition is to be thoroughly, all-sidedly and
completely realized. Herewith the transition from the abstractly mechanical to the elemental, more concrete nature is
motivated.

% THE CLOSE OF PART ONE. The text on p. 196 does not run to the foot: it ends about three fifths down,
% whereupon follow blank setting and then a short centred rule about three quarters down the leaf; the
% rest of the page is blank. The rule is set here DOUBLE and so heavier than the usual division rules;
% that difference is not reproduced --- \streg serves throughout.
\streg

% ============================================================
%  Part Two. Elemental Nature.
% ============================================================
% Danish: \deel{Anden Deel.}{Den elementære Natur.}
\deel{Part Two.}{Elemental Nature.}
% A short centred rule stands between the Deel head and the Kapitel head.
\streg

% Danish: \kapitel{Første Kapitel.}{Elementerne.}   (printed pp. 196--249)
\kapitel{Chapter One.}{The Elements.}

% Danish: \afdeling{A.}{Elementernes Oprindelse.}   (printed pp. 197--214)
\afdeling{A.}{The Origin of the Elements.}

On the usual way of representing it must seem almost contradictory to speak of the elements' origin: an
element is precisely that which is without origin. All things can be resolved into their elements, but the
elements themselves cannot be resolved; they cannot pass away, and therefore cannot arise either: they are
eternal. By examining more closely the significance of \emph{matter and element}, of \emph{matter and form},
and then testing matter's essential share in \emph{the substantiality of things}, we shall however obtain a
new confirmation of what the doctrine of mechanical nature has already shown: that the elements, in order to
have a natural existence, must have a supernatural origin --- in virtue of Spirit.

% Run-in head, read on the page itself (KB copy): „Stof og Element.“ is set in SEMIBOLD antiqua,
% not letterspaced; „a)“ before it stands in ordinary type. The house macro \runin renders the
% whole head alike.
\runin{a) Matter and element.} Greek philosophy begins, as philosophy of nature, in all immediacy with a
doctrine of matter. What is characteristic of this beginning is
% --- p. 198 ---
that matter and principle go together into one; the opposition of the real and the ideal, of nature and
Spirit, does not yet exist at this stage of consciousness: it is pure immanence in speculative innocence. The
contradictions hidden in the immanence then gradually emerge, now under this form, now under that; the
represented phenomena and the thought concepts come into conflict with one another; the problems become
constantly more difficult, and the discussions more entangled.

The general thought that bears the Ionians' representation of a primal matter is the fundamental thought of
\emph{transformation}. How the matter was more closely represented in itself --- as water, as air, as fire,
and so forth --- is inessential; the transformation of matter itself is the main thing. The energy with which
the thought of transformation is developed in Heraclitus already makes it clear that nothing material has
eternal subsistence, that matter in its immediacy is a contradiction. Fire is for him not an existing matter
by whose condensation or rarefaction other existing matters might come forth; that at any rate is only a
sensible way of representing. What is decisive is that every thing is just as much a non-existent as an
existent. The visible fire that can be extinguished has in its being a non-being; for how else should it be
able to be extinguished? That a fire is extinguished means that it passes over from being to not being; and
so with everything existing. But fire itself, which in the ceaseless becoming, streaming and changing of the
things of sense is the true universal fundamental essence --- this fire is not kindled and is not
extinguished; it is a fire of thought whose unity is difference and whose difference is unity; it is the
world-producer, the everliving fire which in itself
% COMPOSITOR'S FAULT in the Greek, reproduced as printed: „ἀποσπεννύμενον“ stands with π where
% Heraclitus's word has β („ἀποσβεννύμενον“). Read on the KB copy under high magnification; both
% copies show it. Read „ἀποσβεννύμενον“. The face's ϰ and ϱ forms are reproduced per the house
% rule with κ and ρ.
kindles measures and extinguishes measures (ἁπτόμενον μέτρα καὶ ἀποσπεννύμενον μέτρα). That in this way ---
since every part of the whole is not merely opposed to other parts but is something in itself opposed --- there
must show itself strife, discord and tension everywhere in nature, goes without saying. The like becomes
unlike; the
% --- p. 199 ---
unlike joins itself together. Strife is the father and lord of all things
% The Greek read on the KB copy. „βασιλεὺς“ has a grave, not an acute --- the printing sets a grave
% before a closing parenthesis; see also „τὸ κενὸν“ below, „τὸ στερεὸν“ p. 200 and „μορφὴ“ p. 207.
(πόλεμος πάντων μὲν πατήρ ἐστι πάντων δὲ βασιλεὺς); life is a mixed draught; the world's strings are strung,
like those of the bow and the lyre.

Against Heraclitus's pervasive dissolution of everything existent stood the Eleatics' sharp denial of all
becoming, arising and passing away. If becoming were to admit of being thought, then it would have to be either
the like coming from the like, or the unlike coming from its unlike. The first is impossible, for the like
\emph{is} --- and does not become --- the like. The second is also impossible. If something were to come forth
from its unlike, it would have to come forth from an other in which it was not, and so from non-being, from
nothing; but from nothing comes nothing.

Between the sharply opposed views Empedocles attempted a mediation. On the one side he makes concessions to the
Eleatics and asserts with Parmenides that arising and passing away --- that is, a qualitative alteration in the
original matter --- is unthinkable; on the other side he will nevertheless not close his eyes to the fact that
things in the world, like the whole world, are subject to a constant change. Nothing then remains for him but to
refer the phenomena of alternation to motions in space, to combinations and separations of unoriginated,
imperishable and qualitatively unalterable substances or \emph{elements}. Of such elements there must necessarily
be several of unlike kind, if the manifold things with their different properties are to be derived from them.
Thus arose the doctrine of the four elements --- \emph{earth, water, air and fire} --- a foundation for the view
of nature of antiquity and the Middle Ages which the newer physics has dissolved.

In another way the representation element was apprehended in atomism. To explain actual becoming, Leucippus and
Democritus started from the opposition between the full (τὸ πλῆρες) and the empty (τὸ κενὸν), between
% --- p. 200 ---
the atom and space. Earth, water, air and fire consist of atoms; the atoms themselves are accordingly the real
elements. In every atom matter is independent and solid (τὸ στερεὸν). The atoms are without becoming; there is
in them nothing non-existent: they are indivisible. No particular properties can be ascribed to them; they are
the immediately, positively existent. The only thing by which they are distinguished among themselves is their
figure, size, weight, arrangement and position --- nothing but quantitative differences. Where a combination of
atoms takes place, the empty is always present as interval. Out of the atoms things are composed. A thing arises
when a complex of atoms forms itself; it passes away when the combination of atoms dissolves; it alters when the
atoms' position and arrangement change; it grows when new atoms enter into the combination; it decreases when
atoms are separated out without being replaced. As regards their constituents, all bodies are alike; as regards
their properties, they may by contrast be highly unlike.

Without dwelling on details, which concern the individual ways of representing more than the fundamental thought
as a whole, and belong particularly under the history of philosophy\footnote{% printed mark: *) --- the only note
% on the leaf. The title stands in ordinary antiqua, not italic.
Cf.\ Zeller's clear and excellent presentation. Die Philosophie der Griechen. 1ster Theil. Tübingen. 1856.}, we
can already see from what has been adduced how matters stand with immanence in these first attempts at a
philosophy of nature. In Heraclitus the thought is pervasive and unfalsified. The everliving fire which is the
reason of things is also the reason in the soul. The essence in the transformations is a being of sense only
insofar as it is a being of thought: the thinking element in nature and the thinking element in man is the same
universal. While Parmenides, in order to keep the being of thought pure as the differenceless one, lets the world
of sense for itself furnish an empty, deceptive spectacle of the
% --- p. 201 ---
non-existent, Heraclitus by contrast reconciles the non-existent with the existent and purifies them both in the
fire of thinking, for the sensible furnishes only fuel for an infinite burning. In the transformation the
material comes everywhere into view with the immediacy of a being, as earth, as air, as finite things; but that
the finite things are only appearances of the infinite being of thought is seen from this, that non-being, the
negative, is everywhere present. The existent is --- in \emph{becoming} --- a non-existent: that is immediate
immanence.

It is otherwise when matter gets absolute independence and becomes element. The element is no being of thought;
it is precisely thought's hard resistance, a barrier impenetrable for thought. And yet thinking does not give up
without more ado; it grasps the element in representation; the representation is determined by sensation;
sensation by impressions of sense; impressions of sense by the immediately qualitative; on this presupposition
the original can be laid into elemental qualities, and Empedocles show the necessity of four elements. To what
degree representation hereby gets the upper hand of thinking comes to light immediately from this, that the
existent and the non-existent are held fast alongside one another in such a way that both of them, each in its
own manner, get an immediate being. It is an advance in the atomists that they would not begin with sensible
qualities. They would conceive the qualities by deducing them from the atoms' combinations; they exerted
themselves to the utmost to make the atomistic elements of matter into pure moments of thought. The atom is not
something straightforwardly sensible; it is a categorical expression for what thought expressly demands, an
absolute existent, an existent whose only quality is to exclude all non-being by having the empty, the
non-existent, as its limit. But with all this the being of thought is nevertheless originally denied, and thought
has become a stranger to itself. In the place of transformation there now comes
% --- p. 202 ---
the category of \emph{alteration} to play the principal part: position, arrangement, separation, combination
denote nothing but alterations in the outward. With such outwardness immanence is in danger. Everything, the
non-sensible too, becomes sensible; and yet something thinkable must be presupposed. Democritus is as little able
as Empedocles to derive thought from the element or to show any essential unity-in-difference of sensation and
thinking. If the being of thought is to be maintained and immanence preserved, it must happen by another way.

% Run-in head, read on the page itself (KB copy): semibold antiqua, „b)“ in ordinary type.
\runin{b) Matter and form.} The simplest expedient for a reconciliation between elements and the being of thought
must, from a popular physical standpoint, seem to be the one indicated by Anaxagoras. Anaxagoras assumes, as
regards the material, a chaos of unlike elements which are, each in itself, determined as uniform in quality
% The printing has no comma after the parenthesis; reproduced as printed (both copies).
% „(Homøomerier)“ stands in apposition to „Gulddele“ etc.
(homoiomeries) parts of gold, parts of flesh, parts of bone and so forth, infinitely divisible, in infinite
mixture; as regards the ideal, he presupposes an ordering understanding, Spirit (νοῦς), all-knowing, almighty, the
finest, purest, absolutely simple, indivisible being, which in all beings is and remains like itself. On the
foundation of this dualism immanence can be preserved so far as it can be presupposed that Spirit is present in
everything, moves and penetrates everything; that the mechanical and teleological causes complete one another in
such a way that the actual formations go on naturally by the parts' separation and combination, while the
all-ordering, all-moving Spirit organizes the whole and builds a world of it answering to the ends of reason and
wisdom. But on closer consideration one will soon see that a dualism of this nature must dissolve into
contradictions. The very assumption that outside Spirit, independently of it, there subsists from eternity a
manifold of absolutely independent fundamental parts proves plainly enough that Spirit is not almighty, but
powerless.
% --- p. 203 ---

It is the same fundamental contradiction that still haunts half-philosophical brains when it is asserted in one
and the same breath both that matter in its elements is eternal, without origin, and that there is nevertheless a
governing reason and a rational order in nature. Reason is a being of thought. If reason is to govern, it must
unconditionally have the elements in its power; but that is possible only on the presupposition that the elements
themselves are absolutely dependent upon reason. If one assumes that the elements' eternal existence is reason's
own eternal reality, that reason not merely exists in the elements but that the independent elements are precisely
the existing reason: then one has either outright denied that what we call reason is a being of thought, or brought
so unreasonable a dualism into reason itself that the rational ceases.

% COMPOSITOR'S FAULT, reproduced as printed: the Danish has „hvcri“ for „hvori“. The letter is an
% unambiguous c --- the bowl is open at the right --- on both copies. Read „hvori“.
The Anaxagorean dualism must indeed from its side have helped to strengthen the consciousness of the
spiritual being in which the knowledge of nature has its truth; but the explanation was, as Socrates complains, in
actuality mechanical, and the dualism unthinkable. One had then to choose between one of two things: either to hold
to the independent elements and give up the being of thought, or to hold fast the being of thought and apprehend
matter as a kind of non-existent, a formless stuff, and deduce the elements by means of form.

With Plato's doctrine of Ideas the being of thought came, in the Greek sense, to full recognition. Without concepts
full of essence there is neither any true being nor true knowing. The concepts, which at once express the truth of
being and of knowing, are self-valid forms of self-valid content, enclose all actuality and subsist by themselves
as Ideas. The sensible alters, the phenomena change, natural things arise and pass away; but what subsists in and
for itself, the universal, the self-like in the manifold of comparisons
% --- p. 204 ---
(ἓν ἐπὶ πολλῶν) is eternal, unchangeable, the object of knowing, the existent itself, \emph{the Idea}.

The Idea's unity is indeed also a plurality; the existent has in itself manifold determinations of non-being; but
the differences hereby arising all lie in being's essence and agree with one another. The existent cannot possibly
be the differenceless one. If we say: the one is, then we make a difference between unity and being; we have
accordingly two: the one and the existent. If we say: the existent is a whole, then it lies in the concept whole
that it must have parts: the whole is a plurality of parts. If anyone will assert that the existent has no parts
and consequently is not a whole either, then he makes the existent impossible. How should it be able to be? In
itself! But then it must be different from itself; for there is after all a difference between that in which
something is and the something that is therein. From these dialectical indications it will be evident how matters
stand in Plato with the manifold of Ideas. Every Idea is a determinateness of being because it is a determinateness
of thought, a determinate thought, and conversely. The Ideas must accordingly originally have been considered from
two sides, namely: 1) as beings of thought, whose being is a necessary thinking presupposing an eternal thinker; and
next: 2) as objectively existent, universal substances in the fundamentally existent. But an even development of the
two principal sides Plato has not been able to give. The Ideas' relation to the Godhead, the subjective side of the
being of thought, is touched upon only in the form of representation under mythical clothing, while the objective
side, the Ideas' real being-for-self, is developed through the very finest dialectic.

Now if the true world is a world of Ideas, whence comes the world of sense? A something must necessarily be
presupposed that can make a co-being of the sensible, the material, the becoming intelligible. But what is that
% --- p. 205 ---
something? It would have to be matter, stuff. But a positive expression for matter and stuff is not to be found in
Plato. Matter cannot be thought; for what is in truth thought is the Idea. Matter, the stuff formless in itself,
cannot be sensed either; for what is sensed is always only a determinate, formed stuff, not the universal, formless
foundation of everything material (only a τοιοῦτον, not τόδε).
% COMPOSITOR'S FAULT in the Greek, reproduced as printed: „ἀναιθησίας“ stands without the sigma, where
% the Timaeus has „ἀναισθησίας“. Read on the KB copy under high magnification. Read „ἀναισθησίας“.
% The face's ϑ form is reproduced per the house rule with θ.
The formless stuff, the mass (ἐκμαγεῖον), is accordingly to be apprehended by means of a certain spurious thinking
(μετ’ ἀναιθησίας ἁπτὸν λογισμῷ τινι νόθῳ). Stuff in and for itself is a non-existent, which nevertheless in a
sensible sense must become a co-existent. If stuff is to be more closely determined, it must happen by its
opposition to the Idea. While the Idea is the eternally existent, stuff is the always becoming, the unlimited, which
gets its limitation, its form, its stamp from the Idea. By a combination of the Ideas with the material there come
forth the sensible things, the finite yet eternal nature, the whole outward world existing in space. There is then
in nature not a mixture of two sorts of being, the Ideas' and stuff's, for the stuff in things is essenceless; it is
the Ideas that give things reality: the whole of actuality is in the Idea.

That stuff, the non-existent, can, by means of the Idea, acquire significance in the sensible existence subject to
becoming and alteration, rests then upon its being formed, upon its merging into the forms of Idea necessary for
outward being --- that is to say, into fundamental geometrical forms. From this standpoint Plato's peculiar
deduction of the four elements will be intelligible.

As a corporeal world, it is said in the Timaeus, the world had to be visible and tangible. Visible it could not be
without fire; tangible not without earth, the ground of everything solid. Between fire and earth there had to be, as
a middle, a third that could bind them together. Since now the fairest binding-together is
% The word-break runs over the turn of the leaf: „Propor-“ on p. 205, „tionen“ on p. 206.
a determinate ratio, propor
% --- p. 206 ---
% THE ERRATA LEAF corrects p. 206, line 3 from the top: „tilstrækkelig; --- læs tilstrækkeligt;“.
% The KB copy prints „være tilstrækkelig; men nu,“ --- with a semicolon and without the -t --- and
% so it stands here. The working copy is PENCIL-CORRECTED at just this place.
tion, this third must stand in a ratio to both. Had we merely to do with surfaces, then one intermediate term
would be sufficient; but now, since bodies are in question, there must necessarily be two. We obtain then four
elements which among themselves form a continuous proportion, so that fire stands to air as air to water, and
air to water as water to earth.

How matters stand more closely with this proportion in an arithmetical sense is in the present connexion
without interest; characteristic, by contrast, is that the elements are formed without existing stuff out of
geometrical figures. All surfaces, it is said, consist of triangles; all triangles are derived from the
right-angled, which are partly isosceles, partly scalene. Among the right-angled and scalene triangles the
fairest is the one whose smaller leg is half as great as the hypotenuse. Of six such triangles an equilateral
one arises; but of four isosceles triangles a square. Of the square the cube, the hexahedron, is formed; of
equilateral triangles the tetrahedron, the octahedron and the icosahedron. Fire's fundamental parts are
tetrahedra, air's octahedra, water's icosahedra and earth's hexahedra.

As stuff shrinks to a non-existent, and the elements are formed of pure limiting forms in space, atomism is
dissolved: only stuffless discreteness, pure mathematical punctuality, remains. One may now indeed say that
immanence is maintained, that sensible things are immanent in the Ideas, since the Ideas themselves are
things' own reality. But whence the sensible comes; how it comes about that a non-existent can look like
stuff, and bring it about that the eternal world of Ideas becomes an outward, sensible world tangible in
becoming and alterations --- that is from a Platonic standpoint inconceivable. The emphasis falls on matter
and form.
% --- p. 207 ---

The relation between matter and form (ὕλη and μορφὴ) plays a principal part in Aristotle, in the philosophy of
nature as in the metaphysics. Matter is the foundation of all becoming; for becoming consists precisely in a
matter's assuming a determinate form, passing over from being in one form to being in another. Matter and form
stand as possibility to actuality, as the passive to the active, in a particular sense as accidentality and
mechanical necessity to purposiveness. Only form as form is conceivable; formless matter is without concept
and so far a non-existent. None the less a certain positive being must be ascribed to this non-existent, for
matter resists form: the accidental and inessential in natural things stems from matter's resistance to form.
The Platonic deduction of the elements is, according to Aristotle, still more mistaken than atomism. In the
atoms there is at least a formed matter, and so a reality; but Plato's elements are hollow figures, as much at
variance with mathematics as with physics. To compose bodies out of surfaces leads consistently to assuming
indivisible lines, indeed to resolving magnitudes into points. As pure mathematical figures the elements
cannot possibly fill world-space, and yet Plato will not concede any empty space. From a physical standpoint
the attempt is altogether unreasonable. For how can the corporeal, which after all has weight, consist of
surfaces which have no weight? Whence, on this presupposition, should the single elements' specific gravity or
lightness derive? --- The assumption of solid, indivisible atoms is
nevertheless also false. On the absolute indivisibility it might happen that the atoms' figure did not fit the
space in which they were to be. Nor does one see any ground why there should be only invisibly small and not
also large atoms. If the atoms are moved by something else, then they are subjected to an action, and the
apathy ascribed to them is cancelled;
% --- p. 208 ---
if they move themselves: then either the moving element in them is different from the moved, and they are then
not indivisible; or, if the moving is one with the moved, then there are in one and the same atom conflicting
properties.

Far from placing things' arising and passing away in a heaping together and dispersing of indivisible atoms,
whose quantitative differences are insufficient to explain the corporeal qualities, Aristotle assumes a
pervasive alteration of matter effected by form, a real becoming in essential transitions. When water freezes,
or when ice melts, it does not happen by a mere alteration in the parts' position, but by a re-forming of the
parts themselves, in the course of which the properties change; when bones and flesh are formed in the body
out of the nourishment, it does not happen by single parts being drawn out like bricks from a wall, but a new
matter forms itself. It is the Heraclitean doctrine of transformation, developed into an articulated concept.

The conditions for the qualitative alteration lie, as with every motion, in the relation between the possible
and the actual, more closely determined as a relation between the suffering and the working. The working and
the suffering must be partly of like kind, partly opposed. The alteration consists in the working making the
suffering like itself, a transformation which can again happen in several ways; for the mechanical causes in
matter are only intermediate causes; the emphasis lies on motion. Nature itself is motion's inner ground;
every motion has in itself a goal by which its direction and limit are determined: nature works teleologically.
It is the peculiarity of motion downward and upward that forms the foundation for the separation of the four
elements. Yet account must be taken not merely of the opposition between the heavy and the light; the bodies'
felt properties also come into consideration. The elemental properties are warmth and cold, dryness and
moisture, the first two
% --- p. 209 ---
% The four names of the elements are letterspaced, not the whole phrases; checked on the KB copy.
active, the last two passive. By placing these properties together in pairs we obtain, with the deduction of
two impossible combinations, four possible ones: warmth and dryness in \emph{fire}, warmth and moisture in
\emph{air}, cold and moisture in \emph{water}, cold and dryness in \emph{earth}; and so again the four elements.

The elements pass over into one another; for everything becomes from an opposite to an opposite. The more
complete the opposition, the more difficult the transition; the less complete, the easier. When water's cold is
cancelled, air arises; for water and air have moisture in common and are incompletely opposed. Water and fire,
by contrast, are completely opposed; the transition must therefore happen mediately, in that water first passes
over into air and then, when the moisture is cancelled, over into fire.

Without laying weight on the fact that the construal of the phenomena conflicts at essential points with
present-day physics, we shall now, from the standpoint of the philosophy of nature, examine in particular how
far Aristotle has by his explanation of the elements' origin really been able to reconcile form with matter,
maintain the being of thought and carry immanence through.

% Run-in head, read on the page itself (KB copy): semibold antiqua, „c)“ in ordinary type.
\runin{c) Matter and substance.} The atomists hold fast the being of matter in unconditioned independence and
exclude the being of thought; Plato maintains the being of thought by reducing matter to the semblance of a
non-existent. Aristotle makes the being of thought into form and the complement of sensibility into a positive
substratum, a matter that lets itself be formed. He acknowledges with a certain evenness both the ideal and the
real; in him we should then expect to find what the newer speculation seeks so eagerly: the ideal-real principle
that can ground the material world's immanent unity with objective reason.
% --- p. 210 ---
But precisely in Aristotle we learn how unreasonable and self-contradictory the whole speculative endeavour is
when immanence lacks its transcendent presupposition, and the unity of nature and Idea is to be sought
immediately and at first hand --- not in Spirit, but in nature. There is in corporeal things a peculiar union
of something sensible with something thinkable; upon this union rests every body's peculiar subsistence or
substantiality. If we abstract from the thinkable, that is, from the being of thought, then we have only a
fleeting, dreamlike appearance, but no substance, no actual body. If we abstract from the sensible, the
material, then we have a general concept of body, but the existing substance, the actually determinate body, is
gone. So it is incontestably certain that matter as matter must have just as decisive a significance for a
body's substantiality as form. This fact puts Aristotle in an embarrassment --- an embarrassment which has by no
means yet been overcome, either by Schelling or Hegel or Feuerbach.

The universal essence --- the genus, the species --- has, according to Aristotle, no subsistence for itself; it
subsists only in and through the single. Substance in the fullest sense, the true, actual essence, which is
subject but never predicate, is a single being. Since now form as form is the universal, a predicate which only
in an abstract, improper sense can become subject, it follows that form can only be one side of substance, not
substance itself. If the ground, the \emph{real} ground, of the single being lies not in form but in matter; if
it is matter that gives things their particular subsistence by giving them individuality, then what is properly
substantial must lie in matter. But that Aristotle will not concede. If matter were substance, then the
substantial would fall absolutely outside knowing, for all knowing goes upon form. Substance must then be
conceived in such a way that matter is a single being only according to
% --- p. 211 ---
its possibility, but that the single being's realization is due to form. And yet this assumption too must be
rejected. That matter has the possibility of becoming a corporeal single being by means of form must necessarily
be thought in such a way that the same matter which is the single being's possibility is also form's possibility.
Matter accordingly does not get form as an other, a fully finished actuality supervening; on given conditions it
develops its form, its actuality, out of its own possibility: matter lets itself be formed because it is its
nature to have form in itself, to form itself. It is a pure fancy that there should be a general substratum
(ὑποκείμενον) falling outside thought because it fell outside form! The formless matter is --- as may indeed be
recognized from this, that in its formlessness it is apprehended as something general --- itself a thought, an
abstract, abortive thought. Formless matter is a designation for something untrue, for a non-existent and yet
existent formable thing waiting for form. Formless matter is no possibility; it is on the contrary an
impossibility, for it is a contradiction. The contradiction does not concern matter at all, but thought. Matter
formless in itself is just as contradictory a thought as form formless in itself, or matter matterless in itself.

% „til det Tænkelige eller det Ikketænkelige.“ stands with a full stop, not a question mark, although
% the question continues; reproduced as printed.
The manner in which Aristotle will deduce the elements, and in general everything he teaches about the properties
of bodies, shows plainly how contradictory it is to make form without more ado into the thinkable and the
material into the unthinkable. Where do such elemental properties as warmth and cold, dryness and moisture
belong? to the thinkable or the unthinkable. Insofar as they are apprehended as forms, they must fall
exclusively under knowing, and so under the thinkable. But according to the deduction they answer to feeling, and
so to the material, the unthinkable. How should dryness let itself be transformed into a thought? It would have
to be a dry thought. But
% --- p. 212 ---
thoughts are in actuality as little either dry or moist as they are, in the natural sense, either warm or cold.

Matter in all indeterminateness is something non-sensible; it is determinateness of form that makes matter
sensible. In consequence of the Aristotelian conception one may then with just as much right assert that form is
sensible and matter thinkable, as the reverse. One can let matter be a moment in form's concept and make form the
whole actuality; but one can also let form be a moment in matter's concept and place the whole actuality in
matter. This swinging between a one-sided idealism which in form holds to the concept, and a one-sided realism
which seeks reality in matter, has not merely shown itself in antiquity but --- a necessary consequence of the
ambiguity of the doctrine of immanence --- has repeated itself in a quite characteristic manner in more recent
times. Hegel, who makes matter and stuff into categories as logical as form and substance, culminates in a
conceptual idealism; Feuerbach, who out of offence at the natureless hollowness in conceptual idealism goes to
the opposite extreme, sets form as a moment in matter's self-individualization and ends with a sensible
this-here, an unsayable.\footnote{% printed mark: *) --- the only note on the leaf. The mark stands in
% the text after the full stop, separated by a space: „et Uudsigeligt. *)“.
The echoers assure us that Hegel is one-sided because he overlooks the real, and Feuerbach one-sided because he
overlooks the universally ideal; so one may become all-sided by setting the ideal immanent in the real and the
real immanent in the ideal; so one has the ideal and the real in pure immanence --- that is to say, in immanent
talk.}

Unclarity about what matter, element and form of matter are must necessarily coincide with a corresponding
unclarity about what nature is. That Aristotle, for all his
% --- p. 213 ---
penetrating eye for experience, his superior acuteness and admirable endurance in comprehensive thinking about
nature, does not know what nature is --- is not at one with himself about how he is really to determine nature's
essence --- is demonstrable from his own system. The transition from the metaphysics to the philosophy of nature
he marks as a transition from the doctrine of the unmoved to the doctrine of the moved. In the circle of the
motions and everything moved there present themselves immediately the corporeal single beings, consisting of
matter and form. What now is matter? It is not known and cannot be known, for matter as matter is no object of
knowing. What then is form? That is precisely what is known; and yet, when we look more closely, it is not known
either. Form, it is said, is actuality. The perfect form is the perfect actuality, the highest energy: that is the
Godhead. This form-being is in the absolute sense a being of thought; it is the eternal thinking that thinks only
itself, the blessed reason that conceives itself, an actuality without any possibilities. But the forms of nature
--- what sort of essentialities are they? Are they the Godhead's thoughts? By no means! To assume this would be in
equal degree at variance with the Godhead's essence and with the forms' nature. With the Godhead's essence! For
the perfection in the divine thinking rests precisely on its thinking the perfect alone. To think the perfect it
must exclude the thought of everything else; for nothing is perfect but the Godhead itself. With the forms'
nature! For if the forms of nature were the Godhead's thoughts, then they would have a being in an other, namely in
the thinking Godhead, and would be derived from an other, namely from the divine thinking. On this presupposition
form could not possibly be things' substance. Substance is precisely that which neither is in an other nor can be
derived from an
other, but unconditionally is in itself. What a splitting of form's
% --- p. 214 ---
concept! The highest form is not merely thinkable by us but thinkable in itself: its thinking is being and its
being thinking. The forms of nature, by contrast, are thinkable by us but in themselves unthinkable, for in the
form of thought they are not substances.

Determinateness of form is then in itself unthinkable too. Natural things' properties and determinations lie in
the moving and determining forms. Matter is determined by the form determinate in itself; but form itself is not
determined; form is determinate. Has form then formed and determined itself in advance? By no means! Were the
forms of nature to have formed themselves in advance before they could begin to form matter, then they must in
this self-forming have passed through a becoming; they must already have had the possible, the passive, the
material in them before they came into contact with matter; there must then again have been a form finished in
advance which could give the forms of matter their form. Determinateness of form immediately presupposes itself:
the whole system of forms working in the realm of nature comes to be taken as a given, inexplicable fact. This fact
then has again the questionable feature that it suffers from insoluble contradictions. Insofar as the forms are
apprehended in knowing, they are apprehended for themselves in abstracted universality and cannot be substances;
insofar as they go together with matter, they are burdened by the material and undergo an essential alteration;
they are determined in determining; they become passive by being active, just as matter, by means of its
resistance, becomes active by being passive. By continuing the analysis one can find contradiction upon
contradiction. The secret is that the immediate immanence of matter and form rests upon a fundamental
contradiction; but that a fundamental contradiction, wedded to dialectical reflexion, breeds strongly is a
well-known matter.

% The sentence runs over the turn of the leaf and is completed on p. 215.
The yield of the old doctrine of immanence is then the following: matter is not immediately to be apprehended as
element;

% --- p. 215 ---
% Division A runs six lines onto this page; the sentence continues from p. 214 without a
% paragraph break.
it is in itself a dependent thing, which first acquires quality and qualitative subsistence when it is determined
by a form answering to knowing. The elements have an origin; the demand is that their origin shall be shown, that
the elements shall be deduced. The demand in itself is true; but it is not yet fulfilled.

% Short centred rule in the middle of p. 215, between division A and division B.
\streg

% Danish: \afdeling{B.}{Elementer og Naturgrund.}   (printed pp. 215--232)
\afdeling{B.}{Elements and Natural Ground.}

The hard opposition of element and being of thought has its original ground in the opposition of immediate
object-power and overreaching self-power. Absolutely overreaching self-power conceives the content of the ground of
nature as a totality of moments in one being of thought; upon this rests \emph{the ideality of the elements}. In
object-power the reflected moments come to existence as the immediate issuing of the fundamental differences, as
independent members in outward actuality; in this consists \emph{the reality of the elements}. In order to be
grounded in itself, nature must be grounded by Spirit: the infinite conversion of ideality into reality, and
conversely, goes on in \emph{the ground-medium}.

% Run-in head, set in semibold on the same line as the paragraph it opens. Read on the page itself.
\runin{a) The ideality of the elements.} The being of thought, determined solely by human knowing, cannot possibly
be apprehended as the real essence at work in the elements of nature. One of two things: either consciousness must
halt before a barrier of impenetrable reality, bow to outwardly mighty nature and ascribe to its elements and
forces an eternal, absolute
% --- p. 216 ---
independence; or thought must, in speculative forgetfulness of its powerlessness, declare matter --- the
unthinkable ground of all outward being --- to be a non-existent, posit \emph{its} essence as the essence of things,
and make itself
% „sit“ stands letterspaced in the setting, also on the KB copy, although the word is short and
% stands at the end of the line.
into a godhead. If the being of thought is to master the elements, then it must be a superhuman essence, resting in
divine knowing, grounded by an omnipotent thinking.

The Ideas, it is said in Plato, are not things of thought
% Greek, set in italic Greek type: νοήματα (acute over eta), read on the KB copy.
(\textit{νοήματα}) which merely have existence in the soul. Right! But when it is thence derived that the Ideas
cannot be the Godhead's thoughts, the inference is false. The soul's thoughts, apprehended for themselves, are
thoughts of powerlessness, abstract things of thought which have the realities outside them; but the Godhead's
thoughts are thoughts of omnipotence with inner, essential reality. That the forms of Idea are necessary forms of
thought the Platonic dialectic proves; but forms of thought are thinking's own forms, forms determined by thinking
itself, forms which only a thinking self can bring forth. If the Ideas' eternal being is not grounded in an eternal
thinking, then there is no true knowing, and the Platonic dialectic is an empty play on words.

What does it really mean that thinking is indeed dependent upon the Idea, but that the Idea, what exists in and for
itself, is independent of thinking? It means expressly that the Idea is in itself an unthinkable essence. An
existent that precedes all thinking is an existent whose determinateness, far from resting upon any determining
thought, on the contrary falls outside every determination of thought. For such an existent to become an object
for thought, it must be represented; it cannot melt together with thought, but must be held fast in existing
independence as thought's resistance, its impenetrable other. Thought may now quite well in each and all direct
itself by this object, invisible in representation, and assure itself that it exactly apprehends all
% --- p. 217 ---
the determinations contained in the existent; but what has happened? The eternal determinations of being,
independent of thinking, have been falsified; they have been transformed into something which they nevertheless
are not in themselves, namely into nothing but determinations of thought. What is in thought is accordingly not the
Idea itself --- the determinations of Idea are not determinations of thought --- but an imitation, an intellectual
copy of an original existing outside; the Ideas' actual being falls outside their thought being: that is a hard
dualism of thinking and being.

To form the world and bring the best possible order into the disorderly, the highest master-builder had, says the
Timaeus, constantly to look to \emph{the archetypes}, the Ideas, things' eternal \emph{models}; hereby, it is
claimed, the Ideas' substantiality is incontestably maintained. Right! But when one will thence derive that the
archetypes to which the Godhead itself looks cannot possibly be products of the Godhead's own thinking, the
inference is false. The long, unclear dispute between the Ideas and the divine understanding must surely at length
admit of being brought to an end.

So long as human consciousness, be it in intoxicating speculation or fasting criticism, arrogates to itself an
absolute knowing of the absolute, the mystical representation of ``ein unvordenkliches Sein'' will constantly crop
% The German phrase stands in quotation marks but in ordinary antiqua: neither italicized nor
% letterspaced (checked on the KB copy).
up and bring disturbance into the much-praised unity of thinking and being. In the psychologically limited self,
real being necessarily precedes every thinking, however speculative, turn it how one will; only in the absolute self
is the finite succession in time cancelled. Divine being and divine thinking presuppose one another to infinity:
being follows from thinking, and thinking from being, in Spirit's eternal substantial essence. To make being the
absolutely first and let a thinking follow after is, with respect to the concept of the Godhead, just as thoughtless
as to make thinking the first and derive being from it.
% --- p. 218 ---

To prove that there is in eternal being something which divine thinking cannot alter, and by which it must
accordingly direct itself, such general propositions have been adduced as: that two is less than three, that a
triangle must have three angles, that the whole must have parts, that an effect must have a cause, and so forth.
But who does not at once see that such so-called ``eternal truths'' are necessary determinations of being only
because they are necessary determinations of thought? If God could ever grow weary of thinking, then he would at
the same time be weary of being.

On the Aristotelian definition that substance is what cannot be in an other, the divine thoughts certainly cannot be
things' substances; but from the definition an either--or follows. Either there are no corporeal substances at all
--- which Aristotle cannot concede, since he most emphatically insists that the corporeal single beings are the real
substances. Or the corporeal substances must be absolute --- which he cannot concede either, since it stands in
express conflict with his doctrine of the elements' transitions, of corporeal formations, combinations and
dissolutions. Were the corporeal substance a single being in the sense that its being in itself made its being in
another impossible, then such a single being would have to be an eternal, indissoluble atom: the corporeal
substances would then have to be raised above all alterations. That things undergo alterations, perish and arise
through a transformation of matter is a proof that their substantiality is not absolute but relative. The
substances' relativity makes it incontestable that their subsistence in themselves is a subsistence in another.

From a general concept of divine thoughts to a sensible representation of such determinate substances of matter as
carbon, sulphur, gold and the like there is, outwardly regarded,
% --- p. 219 ---
undeniably a great leap; but when the matter is thought through, it will nevertheless show itself that the
categories substance of matter and substantial being of thought at bottom presuppose one another.

The substantial being of thought is objectified in a thinking that avails. In the infinite thinking the mutually
different determinations of thought must shine through one another: every particular concept is a mirror of thought
reflecting the whole system of concepts in which the content full of essence is formed out into Idea. What now gives
this being of thought substantiality? The thinking, knowing, all-availing self-power! That thought is power, and
thinking an infinite availing, means that thinking is working and has in itself all the moments of actuality as
reflexes in the overreaching energy of the infinite reflexion. But when all the moments of actuality without
exception are enclosed in the substantial being of thought, it follows that no matter, no real moment, can come
from without, but that all reality, all actuality, be it ever so material and outwardly tangible, must come from
within, from the being of thought, the essential ground of the real world's, outward nature's, elements and
substances of matter.

We have seen how the being of thought in and by itself presupposes the substances of matter; we shall now see how
the substances of matter in their reality presuppose the being of thought.

How far chemistry's so-called fundamental substances really are simple substances and mark the outermost limit of
bodies' possible resolution is a question of its own; at science's present stage they represent the elements, and
that is enough for us.

The elemental substances of matter have first the mark of ideality that their unaltered subsistence comes into view
in nothing but alterations. When oxygen and hydrogen unite and form water, they alter one another mutually; they
give up their independence; their particular
% --- p. 220 ---
properties are transformed and disappear in new properties: they perish as the water comes to existence. The
substances of matter's being in themselves is a being in another; their immediate independence is a one-sided and
transient state; to know their essence one must know their combinations. What they are in themselves, they are for
one another mutually; they play a new part in every new combination, they become incessantly an other through
another; in the stream of alterations there remains no trace or residue of absolutely immediate independence: for
their modifications and alterations there are no outward, finite limits. And yet there are limits, essential,
law-determined limits; for the elemental substances have in all this being-other each its own particular substantial
essence; the one cannot be transformed into the other, oxygen not into hydrogen, hydrogen not into carbon. But that
the elemental beings of matter keep themselves unaltered precisely by undergoing alterations proves that their
fundamental being is an essence reflected into itself, a \emph{substantial} being of thought.

The next mark of ideality is that the elemental substances of matter are, as regards qualities and properties, both
independent of one another and yet essentially determined for one another mutually. If we cast a glance at the
somewhat more than sixty fundamental substances, metals and metalloids, at present known, it is indeed seen that
some among them stand nearer to one another, others further apart; but it also shows itself that every single one
has its properties, its peculiar essence, for itself. There is --- if we except the few which, although they are
composite, nevertheless play a part as fundamental substances --- no trace of the one elemental substance having
borrowed anything from the other. From this standpoint the fundamental substances are accordingly independent of one
another; their differences look as if they had come forth by a pure accident. Who would see, looking at carbon, that
its peculiar existence in any way rests upon
% --- p. 221 ---
that of hydrogen, or the reverse? But the hydrocarbons teach us that there must be an inner, essential and quite
peculiar kinship between the two substances; they fit one another --- indeed, it looks for all the world as though
they had been fitted to one another. But how is that to be understood?

It is a fact that the properties of the one substance, although they immediately belong to the substance in question,
acquire significance only through the properties that just as immediately belong to another substance. This fact
shows that the properties of the one fundamental substance must originally, independently of outward coincidence, be
reflected in the properties of the others; in the formation of the one there is expressly a regard to the formation
of the others: here is among them all a harmony predetermined in divine knowing
% The Latin phrase stands in italic in the setting (antiqua italic), not letterspaced; the æ ligature
% is printed in „præstabilita“.
(\textit{harmonia præstabilita}). But this again proves that the elemental qualities are moments of one and the same
fundamental essence, the substantial being of thought.

Finally --- and this is ideality's fully developed mark --- the whole sum of elemental substances has the character of
a system of concepts carried through in all directions and answering originally to the content of the Idea of nature.
The real concepts of the substances of matter present to consideration two principal sides, a finite and an infinite.
From the side of finitude they are, as concepts of existence, bound, each of them, to its peculiar matter; from the
side of infinity they are in the Idea raised above parcelling out into an eternally clear and transparent unity.

The inner real connexion between the existing substances we learn to know in chemistry from their affinities. The
ground of the difference of the affinities, upon which the whole manifold of combinations and separations essentially
rests, is a real ground, the ground as being of matter. But that this real ground must be ideally articulated, and so
originally determined by a rational ground, can be expressly
% --- p. 222 ---
shown. Affinity is in its essence double-sided: chlorine has affinity to gold because gold has affinity to chlorine,
and conversely. Oxygen has affinity to hydrogen because hydrogen has affinity to oxygen, and conversely. And why,
then, has oxygen no immediate affinity to fluorine? Quite straightforwardly, because fluorine has no immediate
affinity to oxygen. But in this outward referring from the one to the other there is no grounding at all. Affinity's
formal, self-cancelling because---because
% The dash in „Fordi---Fordi“ stands without spaces, as printed.
demands a deeper-lying motive, a rational ground decisive for the essence.

Affinity presupposes that the concepts of matter in question must necessarily be side-concepts of one and the same
comprehending total concept. Oxygen's and hydrogen's mutual affinity, in virtue of which they unite and form water, is
the expression in existence of the natural truth that the real concepts of oxygen and hydrogen are, as side-concepts of
water's real concept, comprehended in one another. What holds of two concepts of matter holds of them all: they are
all concepts of the same conceiving. From the side of infinity the concepts of matter are concepts of Idea; in each
single one of these concepts of Idea the Idea's whole content is present in a particular manner. The real ground is
grounded in the rational, the being of matter in the being of thought; that is the ideality of the elements.

% Run-in head, semibold, on the same line as the paragraph; read on the page itself. No rule before it.
\runin{b) The reality of the elements.} Matter is as little an existent in and for itself as the empty, abstract form.
Matter and form are side-determinations of real power, object-power, which is originally opposed to self-power's
ideality. In the ideal form of the concept of power the element is a reflected moment; in object-power's real form the
moment is a formed matter, an existing element. Nothing is added to the moment from without in order that it may
become an element; nor is anything drawn from the element in order that it may become a moment: moment and element
are, as regards content,
% --- p. 223 ---
essentially the same. And yet they are opposed to one another --- indeed so essentially opposed that the one is
something thinkable, the other something in and for itself unthinkable.

Only a being of thought lets itself be thought; but an existing element --- coal, gold, sulphur --- is no being of
thought; as an existing object it is thought's actual resistance. For powerless thought the resistance is a hindrance,
an outward reality which by sensation can be apprehended in image and representation, but not conceived. For
omnipotent thought the object is the substantial satisfaction, the other with which thought is sated, the object in
which the objectifying self recognizes its own power. That there is ideal power in the divine thinking is categorically
expressed by there being in the elements a real power, an actual object-power. Spirit does not think the elements; it
creates them; it masters them in form-giving concepts: to \emph{conceive} is to \emph{master}.

From concepts of element to existing elements human thinking can show no finite transition. It can be seen in science
that the elements express concepts, and that the concepts are conceived by Spirit. If the concepts be taken away, then
the elements are impossible; if the elements be taken away, then the concepts are unactual: the conversion of concept
into element is an issuing of power. But the transition itself, the actual act of creation, cannot be deduced from the
general presupposition lying in the doubling of power. The doubling of power expresses the eternal fundamental relation
between the ideal being of the concepts and the real being of the elements in general; but the bringing forth
of these, just \emph{these} elements with \emph{these} qualities, of \emph{these} substances with \emph{these} forces
in \emph{these} properties, is for every particular world-system a self-conditioning, unconditioned act of will, an act
of creation which in the unity of the disposition determines the particular dispositions.
% --- p. 224 ---
The peculiar character of the fundamental substances must be found out by the way of experience; but the science
of experience for itself does not have the task of uniting the representation of the elements' manifoldness with
a scientific knowledge
% THE ERRATA LEAF: „S. 224, Lin. 5 f. o.: mechaniske, læs: chemiske.“ The word is reproduced as
% printed. The KB copy prints „mechaniske“ clean and uncorrected; the working copy has the word
% struck through in pencil with „chemiske“ written in the margin by the earlier owner --- the
% correction had already reached this place in the working scan and is not type.
of the unity of the ground of nature. Mechanical analyses have, especially in this century by means of the
voltaic pile, led to the discovery of a great quantity of previously unknown fundamental substances. The question
is now this: will a future separation of substances hitherto regarded as fundamental substances end with an
increase or a decrease in the number of the fundamental substances? By the isolation of potassium, sodium,
norium, thorium, erbium, terbium and so forth, the number has undeniably become greater and greater. But there
are then also indications that the turning point might one day arrive at which the results turned about, and the
movement went from manifoldness toward unity. It has shown itself that saltpetre, for example, far from being what
the Middle Ages assumed, a modified element, is a salt (nitrate of potash)
% Berzelius's oxygen dots: the N carries five dots, two over three.
and nitric acid a combination of nitrogen and oxygen ($\ddot{\dddot{\mathrm{N}}}$); it has shown itself that
cyanogen, though it behaves as a fundamental substance toward the other fundamental substances, is a combination
% The chemical formulas are printed upright with lowered figures and spaces between the signs:
% „C2 N“, „N H4“; so reproduced.
of carbon and nitrogen (C$_2$ N), and ammonium, of which the same holds, a combination of nitrogen and hydrogen
(N H$_4$). There are now those who assume that by eliminating one alleged fundamental substance after another one
will at last be able to resolve them all into a very few, indeed possibly into a single one. Upon so magnificent
a reduction there would then have to follow a corresponding deduction. Just as one can already now transform
oxygen into ozone --- indeed let the transformed oxygen come forward both as ozone and as antozone --- so one
would then, when the ideal was reached, be able to let the many qualitatively different substances arise and form
themselves out of one and the same common fundamental substance.

However highly we prize the analytical art, to which the honour of so many great and fruitful discoveries belongs,
% --- p. 225 ---
we may nevertheless never forget what investigating science constantly recalls: that merely empirical analyses,
even if they be continued with ever so wonderful success, cannot possibly by themselves lead to a true insight
into the essence of the ground of nature. If one stops at the representation of a manifold of eternally existing
fundamental substances, then one denies the unity of the ground, indeed the ground itself. If the fundamental
substances are eternal, then they exist without being grounded, and then atomistic separations and combinations
step everywhere, in all formations, into the place of grounding; things are then pieced together but not grounded.
Without an original ground nature becomes only an exceedingly fine mechanical patchwork, a natureless nature, a
contradiction. If one ends in the representation of a common, differenceless, general fundamental substance and
puts this existing matter in the ground's place, the contradiction is still more conspicuous. If substances of
different kinds were to come forth from a first matter in which there were no differences, then the differences
would have to arise out of nothing. If in the presupposed fundamental substance there really were a quantity of
different constituents, then it would in actuality not be one but several substances; but when there are no
differences at all in it, then there is also nothing from which determinations of difference can be derived, either
immediately or mediately. The differences belong just as originally and necessarily to the presupposition as the
unity. The content-rich unity of the determinations of difference, the real ground of the substances of unlike
kind, is no fundamental substance, no object of experience. If the ground of nature is not to be known, then we
remain in the contradiction; if the ground of nature is to be known and the elemental differences grounded, then
the limit of empirical knowing must be overstepped.

Another side-representation accompanying the science of experience will likewise obscure the apprehension of the
fundamental relation between the original and the derived. The simple substances, the elemental realities that
cannot be resolved, show themselves to the observer
% --- p. 226 ---
as the constant, the original, while the composite substances represent the inconstant and derived. A double salt
has then from this standpoint less independence than the particular salts into which it can be separated; these can
subsist without it, but it not without them. For the same reason the salts have less independence than their acids
and bases, and these in turn less independence than the elements themselves. This way of seeing
% The printed mark: *) --- it stands AFTER the full stop, as reproduced.
may within the circle of the realities themselves be perfectly warranted.\footnote{%
Cf.\ Forelæsninger over „Philosophisk Propædeutik“ from the university year 1860--61, p.\ 138.} But for the
knowledge of the elements' relation to the ground of nature it is misleading. Knowledge of the ground reduces the
relative difference between simple and composite substances to a derivative and holds fast the elemental cycle. In
the cycle the ceaseless alternation of arising and passing away holds as much for the simple substances as for the
substances of matter composed of them. That the composite arises means that the elements perish; that the composite
perishes means that the elements arise. Every substance is a derivative; nothing is at bottom original but the
ground itself.

The superficial side-representation of the elements' unconditioned priority, which does not concern research,
confuses concept and conceiving. It constantly comes to look as though the first, fixed, unalterable concepts of
nature lay exclusively in the elements, while the particular, more content-rich concepts uttered in the formations
of composite substances were only derived, accidental concepts of mixture dependent on outward circumstances. That is
a fundamental error. The elements are, according to their concepts, expressly laid out toward particular formations.
Only from the particular formations can it be seen what they really signify as elements. The concept of water,
% --- p. 227 ---
the concept of quartz, the concept of feldspar are just as original concepts of nature as the concepts oxygen,
hydrogen, silicon, aluminium, potassium. It is precisely the more all-sided, more concrete concepts that contain the
more one-sided in themselves; not the reverse.

Upon insight into the relation between more abstract and more concrete concepts rests all essential insight into the
teleological. The science of experience is within its rights when it dismisses loose reasonings about lower and
higher, about means and ends in nature's housekeeping; but concepts such as teleology of nature and reason in nature
it cannot treat. Reason in nature and teleology rest upon dispositions in nature, dispositions in nature upon the
ground of nature; but the question which the science of experience must by its character leave unanswered is
precisely the fundamental question, the question of the ground.

% Run-in head, semibold, on the same line as the paragraph; it stands about the middle of the leaf.
\runin{c) The ground-medium.} Power's fundamental doubling necessarily brings with it a doubling of the ground. In
the self-medium the ground is being of thought, ground of reason; in the object-medium it is being of matter, real
ground. The ground of reason has indeed its truth, the confirmation of its ideality, in the real ground, and this in
turn the confirmation of its reality in the former; but the former is reflected in inwardness, the latter merges into
existential immediacy: they form extremes. The ground-medium, the immanent intermediate being in which the opposition
of ideality and reality ceaselessly arises and just as ceaselessly disappears, is the middle of the extremes.

Outward, material actuality is an extreme of positive immediacy, an antilogical reality that excludes all negation,
all reflexion, and with that all ideality. Every element expresses its concept in so realistic a manner that the
concept disappears in the expression. If we say: this piece of coal has a limit, a limit is negation, therefore it
has in its being a negation --- then we confuse what belongs to thought with what belongs to the object. The piece of
coal's
% --- p. 228 ---
limit is the positive in its surface; the parts lying in the surface have just as immediate an existence as those
found in the interior. The limit as negation belongs to thought: the surface is the body's ceasing, its non-being ---
that is a judgment of the reflecting consciousness. The body's non-being does not concern the existing body: the
element's negation falls outside the element.

The antilogical immediacy of the actual is found indicated in age-old problems. The Eleatics denied becoming because
the existent must exclude the non-existent. The Megarians went further and denied possibility, because the possible
is in actuality a contradiction, an existent which nevertheless is not. To
\emph{be able}, they said, is to \emph{work}; what nature can do, it does; what it does not do, it cannot do
% Greek, set in italic Greek type; polytonic, read on the KB copy. „φασιν“ stands without an accent,
% as printed.
(\textit{ὅταν ἐνεργῇ μόνον δύνασθαι, ὅταν δὲ μὴ ἐνεργῇ οὐ δύνασθαι, φασιν}): only the actual \emph{is}; the possible
is not at all. Aristotle explains becoming and possibility by overlooking the point and breaking the edge off the
problem. To explain becoming he refers to possibility; to explain possibility he refers to matter --- the
inexplicable! What is matter? It is possibility. What is possibility? It is matter. In more recent times Hegel has
given a dialectical impetus toward the clearing up of real possibility. The immanent negation is the secret of all
possibilities; that is the matter to a nicety. A pity that the immanence of the negation has, as regards natural
things, itself remained a mystical secret.

For the elements to be able to work, there must be in every actual element countless possibilities. In an isolated
state existing carbon must indeed express its concept; and yet cannot in actuality express its whole concept. Where
carbon comes forward as diamond, only one side of the concept is realized, for the possibility is that it can also
come forward as graphite; where it in actuality comes forward
% --- p. 229 ---
as graphite, the possibility is that it can also come forward as in anthracite. In carbon's combination with
hydrogen the possibility of carbonic acid is excluded. That the whole concept cannot possibly be realized at once and
in one place in one and the same actuality rests upon existence's being exclusive, upon the one form of actuality
necessarily excluding the other. According to its concept pure carbon is \emph{both} graphite \emph{and} diamond, but
in actuality it is \emph{either} graphite \emph{or} diamond. It is accordingly not, as Aristotle assumes, matter, but
the concept, that is the possibility; when the form of the concept has become a form of matter, we have the actuality.
Every element is always something other in actuality than it is in possibility; the actuality is each time only
\emph{this}, only one, but the possibilities are manifold. The concept's manifold possibilities have a negative being
in the one positive actuality with which an element exists. This negative being of the possibilities is the actual's
immanent negation, its ideality, its inner.

But according to what we have seen above, the actual in nature has no inner; the natural thing is, in antilogical
immediacy, what it is, neither more nor less. So there is dispute about the relation between nature's inner and its
outer; the dispute goes on both in verse and in prose. Against Haller's strict
% The two German verse citations stand in ordinary antiqua, without quotation marks and without
% italic; checked on the KB copy.
separation of the inner and the outer --- Ins Inn're der Natur dringt kein erschaffner Geist --- Goethe pledges his
intuitive genius for the inner's immediate unity with the outer: Natur hat weder Kern noch Schaale, Alles ist sie mit
einem male. Who now is right, Haller or Goethe? In reality there is everywhere a unity of the inner and the outer,
for reality is that the inner merges into its outer. But reality is finite; against reality's finitude stands
ideality's infinity, and so there is after all a decisive opposition between the inner and the outer.
% --- p. 230 ---

Ideality, the inner, answers to the being of thought, the ground of reason; in the ground of reason the concepts of
element are concepts of Idea: there is in every existence in nature an infinite surplus of ideality. Reality in its
immediacy answers to the ground of matter, the real ground; for existence's real ground merges, with the given
conditions, immediately into the grounded existence.

For intuition and representation it then easily comes to look as though ground of reason and real ground furnished two
foundations lying one upon the other, so that the real ground went together with actuality while the ground of reason
was to be regarded as a kind of invisible container of the possibilities. The realization of a possibility would then
come to be apprehended thus: that the possibility came by an inflow from the inner and, by intervening in the outer
conditions, got power to alter what subsists, take actuality from what is given in advance and thereby give itself
actuality. In this way every new possibility would accordingly come as a conqueror from the depth of the idealities,
ally itself with a circle of conditions and tear actuality to itself. An unnatural machinery! One represents the
elements as swimming in a sea of possibilities; and then there is between possibility and actuality no substantial
unity. One might then just as well represent the elemental realities as being nothing but occasional causes, loose
counters which an omnipotent will set up and moved at pleasure.

All activity in nature rests upon a substantial unity of possibility and actuality, but this unity rests in turn upon
a pervasive opposition of ideal and real power. In the way of representation everyone knows well enough what is
generally meant by an activity in nature. So one assumes that working and activity is something that needs no
explanation, and explains without more ado the transition from possibility to actuality by activity,
% --- p. 231 ---
pointing particularly to causes, forces, conditions as the active element. One forgets that precisely that by which
one explains is in high degree in need of an explanation. What use is it to appeal to ``the active element'' and say:
the causes \emph{work}, the forces \emph{work}, the conditions \emph{work}, when one is not able to give an account of
what really lies in the concept: to \emph{work}?

Reality for itself, the antilogically determined actuality, cannot work, for it lacks possibility. Two pieces of wood,
for example, cannot be rubbed against one another without the one having, as regards the impression, to furnish
possibility for the other's actuality; in being mutually both possible and actual for one another, they work, and heat
is brought forth. Oxygen and hydrogen can be mixed without giving up their immediate actuality; no new production is
then effected. But if an electric spark passes through the mixture, then the possibility is unlocked: the elements
work and bring forth water. Possibility for itself, pure ideality, cannot work either, for it lacks actuality. The
possibility lying in the concept of water cannot realize itself; it must have the oxygen and the hydrogen for
actuality in order to be able to work. Working is under all forms an infinite conversion of ideal into real, of
possibility into actuality, and conversely.

This conversion presupposes that the ideal is power. If ideality were not power, it could not possibly have the moment
of actuality in itself and work out its own reality. All working is, seen from possibility's side, a categorical
expression of the ideal's mastering the real. But the conversion also presupposes that the real is from its side power
and masters the ideal. If immediate reality did not have ideality, possibility, as a hidden moment, an indwelling
negation, in itself, then it would be powerless over against possibility; it could then not take part in the activity,
it could not work at all. From actuality's
% --- p. 232 ---
side all working is a categorical expression of the real's mastering its ideal.

% Greek and Latin both stand in italic, separated by an ordinary comma; read on the KB copy.
Possibility is an availing (\textit{δύναμις}, \textit{potentia}); its essence is an ideal power. Actuality,
antilogically determined, is a being's independence, is resistance, object; its essence is a real power. Actuality is
an issuing of power; possibility is ideality's surplus. Water is according to possibility ice, and yet here is a
qualitative difference. The transition from water to ice is an abrupt, decisive conversion; the degrees of cold with
their preparatory intermediate determinations lose themselves in the disjunction: either water or ice. So long as the
water is actual, the ice is only possible; the new form of state, in which the ice is actual and the water possible,
rests upon an issuing of power. The issuing of power may otherwise be of highly different kinds; but it is always
recognizable by the qualitative conversion. Real possibility requires conditions. The conditions can enter one by one;
so long as a single condition is wanting, there is here only a possibility and a number of conditions; but with the
entrance of the last condition everything is altered; the possibility is no longer possibility, the conditions not
conditions; the whole is a new actuality --- it is an issuing of real power. The conditions' transformation into
moments in the ground is the essence of the issuing. The real power masters the conditions, for it transforms them;
and yet this is only the one side of the matter. The opposite is that real power is dependent upon the conditions,
inasmuch as it is dependent upon possibility and upon ideal power. Upon the restless, yet at bottom quiet, unity of
real and ideal power rests all working and activity. The elements combine, and the composite substances of matter come
forth by an issuing of real power --- on the ground of power's ideality. The composite is separated, mighty ideality
unlocks its possibilities, and the elements return --- in power's real issuing.

% --- p. 233 ---
% p. 233 opens with a NEW paragraph, so division B closes here with two whole paragraphs in the
% upper third of the leaf.
The ground-medium, the universal essence in which the ground of reason is one with the real ground, reflects in
its inner what is parcelled out in the outer, and unites in the outer what is different in the inner. The
ground-medium expresses the unity-in-difference of nature and Spirit; in the ground-medium the natural is
spiritual and the spiritual natural. That an existing element can carry all its possibilities, its immanent
negations, in immediate existence is possible only because the real power, whose issuing determines its
restricted actuality, is in the ground-medium one with the inner ideal power that carries ideality's surplus.

So here, then, is an immanence, an original identity of being of thought and being of matter. Yes! But the unity
is grounded upon an opposition, the immanence upon a transcendence; the ground of nature itself is a derived
ground. Nature is grounded by Spirit.

% Short centred rule (single) between division B and the following division, about the middle of the leaf.
\streg

% Danish: \afdeling{B.}{Elementer og Legemer.}   (printed pp. 233--249)
% DIVISION HEAD WITH THE WRONG LETTER --- and the head does not open the page, but stands about the
% middle of p. 233, after the rule. The errata leaf corrects p. 233, line 13 from the foot:
% „B. --- læs C.“; the sequence (A. on p. 197, B. on p. 215) requires C. too. The printing has „B.“,
% read on the KB copy. The working copy is PENCIL-CORRECTED here: beside the printed „B.“ stands a
% hand-drawn „C“ and the B is struck through --- the seventh place where the pencil hand has applied
% Nielsen's errata; cf. the exactly corresponding case at the run-in head on p. 180 (c) --> b)).
% Reproduced as printed, not renumbered.
\afdeling{B.}{Elements and Bodies.}

The natural difference between element and body rests upon an immediate separation of independent and dependent, a
movable determination of limit which comes particularly into view when the concept \emph{corporeality} is
investigated. That the elemental concepts of independence are concepts of power makes itself known in facts by this,
that the determinations of concept, although their inner essential connexion must necessarily be presupposed, do not
admit of being derived from one another in a rational manner; the factual connexion is known \emph{from the corporeal
properties}. Bodies' independence
% --- p. 234 ---
is a material substantiality, a selfless independence, which in elemental nature has its measure and its standard in
\emph{the individuality of bodies}.

% Run-in head, read on the page itself (KB copy) and not in the Indhold: set in semibold antiqua, on
% the same line as the paragraph.
\runin{a) Corporeality.} From the foregoing it is seen that the word element can be taken in very different senses;
now it is the fundamental atoms, now the masses of the simple substances, now again the universally corporeal forms
of state and activities that are meant. Every element is a body, every smallest part of matter fills a space; but not
every body is an element. In determining the element one has originally started from the premiss that it must be the
absolutely simple, the independent in itself; but simplicity and independence do not straightforwardly answer to one
another. The atom, the abstractly
% MISPRINT: „Eukelte“ for „Enkelte“ --- an inverted n. The letter is read on the KB copy, where the u
% stands unambiguously, and calibrated against the correctly set „Enkelte“ three lines above. The
% errata leaf is silent. Translated by the sense.
simple, is a dependent thing; the atoms subsist in the mass and are known only from the relations of weight and
measure in which the masses combine.%
\footnote{% printed mark: *) --- it stands AFTER the full stop.
On chemical proportions, equivalents and atoms cf.\ Forelæsninger over Phil.\ Propæd.\ 1860--61, p.\ 200 ff.}
Even though the fundamental substances are in themselves chemically inseparable, they are by no means on that account
absolutely independent; a glance at the elemental cycle shows that the simple has at bottom no greater independence
than the composite. Hegel protests against chemistry's exclusive right to determine what an element is; in him the
limit between independent and dependent is \emph{logically} fluid. The peculiar way in which he has dialectically
refreshed the doctrine of the four old elements is too characteristic to be passed over here.

Matter, it is said, has individuality; it has shown in the solar system that it can determine and form itself out of
its own essence. What gravity is in the mechanical, light is for the elemental qualities. Light is the first
qualified matter, reflected being-in-itself, nature's immediate, universal self. Light's outward existence
% --- p. 235 ---
is \emph{the star}, as a moment in a totality \emph{the sun}. The sun is the existing light presupposed by the
planets themselves. Light itself is no element, but the principle for the elements' development. For light has not
merely darkness; as qualified matter in the form of ideality it has at the same time a plurality of qualified matters
in the form of reality for its opposite. It is these realities, the light-being's immanent negations and moreover
matters for themselves in outward being, that come forward as four elements, answering to the four natures uttered in
the solar system: air to the solar, fire to the lunar, water to the cometary and earth to the planetary nature.

Air is the element of universality, not in the form of selfhood but immediately as being-other; air, light degraded to
being for another, is heavy; light itself is without weight. Passive toward light, transparent, mechanically elastic,
a creeping, corroding power, air is everywhere on the point of dissolving and volatilizing the individual.

The elements of particularity, of \emph{opposition}, are fire and water.

In the moon immediate being-for-self has stiffened into a sulphurous being-for-self, a fire-principle without fire ---
the reader will perhaps best be able to think of it as a frozen fire --- this stiffness (Starrheit) now passes over,
in the region of the elements, into a restlessness existing for itself, into burning fire. Air is at bottom fire, as
shows itself when it is compressed. In fire, then, air is expressly posited as what it is in the concept, as negative
universality or self-referring negation. Fire is materialized time or selfhood (light identical with heat), the simply
restless and consuming --- a consuming of another which at the same time consumes itself and so passes over into
neutrality.

The cometary element, the second member of the opposition, the neutral, is water. As the opposition sunk together in
itself,
% --- p. 236 ---
water has no singularity existing for itself, no fixity or determinateness in itself. There is in this element a
pervasive equilibrium of possible states; every mechanical determinateness posited in it is dissolved. The limit of
form is received from another and sought in another (adhesion). Water is fluid, like air, but not elastically fluid,
so as to expand itself on all sides. It is more earthly than air, seeks a centre of gravity, stands nearest to the
individual and strives toward it, because it is essentially the concrete neutrality which is not yet \emph{posited} as
concrete.

The planetary, concrete element, earth, is first the still indeterminate \emph{earthy} in general; more closely
determined, it is the totality which takes the other elements up into itself and comprehends the differences into
individual unity. It is the earth's destiny to appropriate the universal astral powers which as heavenly bodies have a
semblance of independence, to digest them, bring them under the dominion of an individuality in which these giant
limbs are reduced to moments. The first great act of digestion is the elemental process.

The physical elements, it is added, are not individualized in themselves but shapeless; they are moments of
corporeality, but not subsisting bodies. The inorganically individual is for the concept the most obstinate, because
matter is here specified into individuality; and yet the individuality is immediate, and as quality immediately
identical with the universal.

This, now, is Hegel's central thought; it is the last and most complete development of the doctrine of immanence that
we have before us here. Of the Hegelian deduction holds what was brought out above on the occasion of the Aristotelian.
The question is not in the first place how far the explanation given satisfies physics --- here, with the exception of
a few
% --- p. 237 ---
small errors, there is really as good as nothing that concerns physics --- the fundamental question is whether such a
method in the philosophy of nature has validity, and what worth can be accorded it in science.

% Schelling's words are set in ITALIC (antiqua italic against the surrounding antiqua), read on the KB
% copy. No letterspacing.
\textit{Ueber die Natur philosophiren heisst die Natur schaffen} --- these words of Schelling, which the editor of
Hegel's philosophy of nature has chosen as a motto%
\footnote{% printed mark: *) --- it stands BEFORE the comma.
G.\ W.\ F.\ Hegels Vorlesungen über die Naturphilosophie als der Encyclopädie der philosophischen Wissenschaften im
Grundrisse zweiter Theil. Herausgegeben von D.\ C.\ L.\ Michelet. Berlin 1842.}, are well suited to set the immanent
standpoint in a blazing light. There is something oracular, something high-priestly, something loftily
dogmatic-episcopal in the speculative self-certainty with which nature's secrets are here revealed. It looks
undeniably as though Hegel had himself created the four elements in order to introduce them, if not into nature, then
at any rate into the philosophy of nature.

It is in truth not a want either of scientific earnestness or of logical depth, but a misunderstanding rooted in the
doctrine of immanence concerning the objective concept's self-conceiving, that lets itself be seen in this abstruse
creation of elements. The apprehension of the relation of the concept of corporeality to existing corporeality is
mistaken from the bottom up. Here is no ground-medium for rational-real necessity, for here are no qualitatively
determinate extremes; here is, in other words, no trustworthy starting point, either in the a priori or in the
empirical. Not in the a priori; for concepts such as air and fire, water and earth do not admit of being derived with
logical necessity from pure thinking. The elements are not beings of thought but beings of matter. So determinations
borrowed from sensation and from ordinary everyday experience are smuggled in. Not in the empirical; for
% --- p. 238 ---
it is after all supposed to look as though the smuggled-in determinations of sense were not due to experience, but were
most strictly derived from the immanent concept. The logical concept repeats itself with subjective necessity in three
principal forms: the universal, the particular and the singular; ergo it is to repeat itself with objective necessity in
four elements. And why? Because nature, which cannot hold the concept fast, must let the moments of particularity fall
apart from one another into two members and therefore needs two elements, fire and water, to express the particular. One
gets the impression that the last addition derives not from any higher conceiving but from a plain everyday experience.
With the application of the four natures in the solar system it looks still more questionable; here there is nothing
whatever, sensible or thinkable, to hold to. That the moon is sulphur is about as scientific as that it is a green
cheese. What a leap from the fire-principle in the moon to a burning fire on the earth! Worst of all, however, is the
comet-water. Suppose the comets were waterless --- and that might well be the case --- then there is no water-principle
at all. But when there is no water-principle in the comets, it must go either against ``the immanent concept,'' which
then by the new theory of comets suffers a blow it never recovers from; or, if ``the concept'' is absolutely to be in the
right, it must go against poor nature --- for how shall nature be able to hold its water when there is no water-principle?

Hegel maintains that in such deductions it does not really turn upon the character of the outward, material existence;
the main thing is that one can everywhere in the corporeal world see the concept. Paracelsus, he remarks, has said that
all earthly bodies consist of four elements --- of mercury, sulphur, salt and the virgin earth --- just as one also had
four cardinal virtues. Such things must not be understood according to the letter; the higher meaning is that
% --- p. 239 ---
real corporeality has four moments. --- One learns from this, then, that on these presuppositions it does not help to
make much of realism by appealing to the unity of the real with the ideal. When the ideal demands no more than that real
corporeality shall have four moments, then one may just as well seek the moments in Paracelsus and Jacob Böhme as in
Schelling and Hegel.

% LETTERSPACING BROKEN AT THE LINE BREAK. On the page „E x i s t e n s-“ stands letterspaced as the last
% word in the line, while „grændse;“ at the top of the next line is ordinary antiqua. Reproduced as
% printed: only the first element is set as emphasis.
The limit of the independent and the dependent in the world of bodies is a movable but nevertheless law-determined
\emph{limit of existence}; the moments in the concept of corporeality are actual facts.

% Run-in head, read on the page itself (KB copy): semibold antiqua.
\runin{b) Corporeal properties.} The relation between the concept of corporeality and the existing body is, as we have
already seen from several sides, a relation of power. In the ground-medium the unity of concept and actuality is
reflected as working and activity; but in the object-medium the concept is so immediately one with its existence that
the object stands in the concept's place. Of what nature a body is must be known from its origin; but the body's
character, the whole peculiarity with which it exists, must be known from its properties. The body with its properties is
an immediate given; the given wills, in antilogical immediacy, to be seen and sensed from all sides in all ways just as
it is given. This necessity of the object involuntarily makes an impression of power; but just as involuntarily
object-power sets reflexion in motion again. It is impossible for the observer to sense without thinking: this is a body,
this body is hard, is yellow, and so forth. So the body becomes logical subject (substratum) and the properties
predicates; for experience happens only in thinking and judging. It now turns on what direction the reflexion takes. The
reflecting consciousness can, as the history of science teaches,
% --- p. 240 ---
leap over the ground-medium and throw itself upon one of power's extremes, either upon self-power in all arbitrariness or
upon object-power in all accidentality. Scholastic theological metaphysics is an example of the first one-sidedness,
empirical scepticism of the last.

% The Latin is set in italic, read on the KB copy; the book here writes „distingvendum“ with the qv form.
The body's existence is a word of power, a revelation of self-power's will; that is scholasticism's theological
presupposition. How this presupposition was in its time applied is seen from the Catholic doctrine of transubstantiation.
Bread and wine are transformed into body and blood and yet retain their sensibly corporeal properties. How now is this
possible? Here a distinction must be drawn --- \textit{distingvendum est} --- between substance and accidents. The
arbitrary self-power which has created substance out of nothing and imparted certain accidents to it can at pleasure let
two substances of unlike kind share one and the same sum of accidents, or else re-create the substance but let its
accidents remain unaltered. When a mouse eats the Most Reverend, what then does the mouse get? Bread or breadness
(\textit{panis} or \textit{paneitas})? God knows, answers Petrus Lombardus. An apt answer. In a natureless nature the bond
between substance and accidents is so arbitrary that omnipotence itself immediately looses and binds it, according as the
priest chants.

% The fraction is printed as a true fraction (1 over 20000) in mid-line.
Empirical scepticism takes its starting point in the immediately given; the given's necessity is hidden in accidentality:
object-power is antilogical in this and this and this. Knowing rests upon experience, but what is experienced is each time
a single thing, and every single thing something accidental. We suppose we know that there is a substance called gold;
but that this is a fancy can be seen from the way in which gold's properties are put together. Gold, it is said, is a
yellow metal which can crystallize in octahedra; it is so malleable that it lets itself be rolled out into leaves of
$\frac{1}{20000}$ of a line in thickness, and so forth --- nothing but empirical marks.
% --- p. 241 ---
Were it to be shown that gold was a substance, and the substance gold the necessary unity of these properties (accidents),
then it would have to be possible to show how each single property followed from the others and from the combination they
form, just as with a triangle, for example, it can be shown that each single side depends upon the two others and upon the
angle they enclose. But this is so far from being the case that on the contrary each single property is experienced by
itself without the least indication of what it really has to do with the others. To be yellow, to be malleable, to
crystallize in octahedra are in actuality three things so unlike in kind that no one readily thinks of representing them
in indissoluble unity. That gold has of all the metals the least affinity to oxygen is adduced in chemistry as a
universally valid proposition. But how does one prove this proposition? Grounded upon experience the proposition cannot be;
for one can experience a \emph{number of single} cases, but not \emph{all} cases, still less \emph{a universally valid}
case. One appeals to induction; that is to say: one puts some cases for all cases; one might just as well put $1 = 2$,
$2 = 3$, $3 = \infty$.

To avoid arbitrary determinations from self-power and disturbing accidentalities in object-power, the doctrine of
immanence, the philosophy of the absolute with its absolute knowing, has briskly struck a line through both extremes; it
will in virtue of its ideal-real principle articulate the content of rational necessity; but the articulation founders on
the arbitrarily accidental determinateness of the corporeal properties.

The doctrine of immanence stands in conflict of principle with physics. Physics is indeed raised above the looseness of
scepticism, for it maintains reason's right and knows the necessity in ``the unalterable laws.'' But rational necessity is
in nature real, and real necessity inseparable from
% --- p. 242 ---
accidentality. In every corporeal property the necessary has its immediate being in the accidental. The accidental
is that gold is yellow, that it is malleable, that it has a stronger affinity to chlorine than to oxygen. But this
accidental is precisely existence's expression for an articulated necessity, a law-determined connexion: physics
relies on induction because it relies on the law; it knows a substance in gold because it knows a factual necessity
in the combination of the properties. That the necessity is factual, untransparent, physics utters by declaring
that we know only things' properties, not their essence.

% QUOTATION MARKS OVER THE TURN OF THE LEAF: the Hegel citation opens here and is closed only on
% p. 243 at „Hemmelighedsfuldt.“ The page is therefore +1 in the quote balance and p. 243
% correspondingly -1; it is a quotation over the leaf, not a misprint.
Against such an utterance the doctrine of immanence supposes it can set an explanation of this sort: that the thing's
essence reflects itself in its properties, so that a knowledge of the properties must necessarily go together into
one with a knowledge of the essence. Such an explanation is good enough in a logic; but how does it look when it is
to be carried through in a physics? Let us, for example, choose the property of a body that it is \emph{magnetic}, in
order to see what the philosophy of immanence avails to elucidate about this property. Magnetism's essence is,
according to Hegel, altogether transparent in its appearing (Erscheinung). What is essential is that the concept's
different moments are not set identical in an outward manner, but that they set themselves identical. The brittle
point opens itself up to separation in the concept; thus we have the poles. On the physical line, which has the
difference of form in it, they are the two living ends, each of which is posited so that it is only in relation to
its other and has no significance when the other is not. Both are one determination. ``This appearing,'' he adds,
``is \emph{the whole theory} of magnetism. The physicists say that it is not yet known what magnetism is, whether it
is a current, and so forth. All this belongs to the sort of metaphysics that is not
% --- p. 243 ---
acknowledged by the concept. Magnetism is nothing mysterious.''

After such an utterance one should now expect that the essence cleared up in the concept would have to spread an
almost divine light over the otherwise so enigmatic difference between magnetic and non-magnetic bodies; but the
expectation is disappointed. ``Upon what bodies magnetism comes to display itself is,'' it is said, ``perfectly
indifferent to philosophy. Chiefly it is found in iron, but also in nickel and cobalt.'' We learn from this
declaration that actual bodies' actual properties stand in an altogether outward, accidental and indifferent relation
to corporeality's essence; and yet the body's essence is immanent in its properties! The mark of accidentality is
inseparable from the fact; it is facticity's immediate determinateness, its own antilogical moment. The doctrine of
immanence does not dissolve accidentality, because it loftily overlooks the accidental; it does not crack
actuality's hard nut, because it throws the facts away as useless shells.

That the philosophy of immanence cannot come to terms with object-power has in a deeper sense its ground in its
having excluded self-power; it cannot digest the accidental because it cannot appropriate what is arbitrary in the
concepts of nature; instead of nature's categories, fixed from within, it slips its own self-chosen concepts
underneath, and becomes falsely arbitrary: its idealism is as hollow as its realism.

We are to go to school with nature, says the physicist; it is not we who are to teach nature something, but nature
that is to teach us something. The physicist accordingly acknowledges that in natural things he has to do with
determinate, expressly prescribed thoughts and concepts. And in this he is right. The corporeal properties show that
the concepts of nature
% --- p. 244 ---
have their actual expression in a categorical language of power. Gold is a metal, iron is a metal; but from the
universal concept metal it is impossible to derive the particular metals' particular properties. Just as every metal
has in outward being a stamp of accidentality, so it has in essence and concept the stamp of an omnipotent
arbitrariness. Why is gold not magnetic just as much as iron? Indeed, why is there in the earth both gold and iron?
To such questions no physicist will answer. Real concepts such as gold and iron are, in their categorically ideal
determinateness, concepts of will. There is in every determination of thought a moment of arbitrariness, just as there
is in every property a moment of accidentality.
% ANOTHER LETTERSPACING BROKEN AT THE LINE BREAK, as on p. 239: „r a t i o n e l l e“ stands
% letterspaced at the end of the line, while „Nødvendighed,“ at the top of the next line is ordinary
% antiqua. „r e e l l e“ and „N a t u r n ø d v e n d i g h e d“ are letterspaced in their entirety.
The moment of arbitrariness stands to \emph{rational} necessity as the moment of accidentality to the \emph{real};
the unity of both is \emph{natural necessity}, a unity articulated in the ground-medium.

% Run-in head, read on the page itself (KB copy): semibold antiqua.
\runin{c) Corporeal individuality.} That every body's peculiarity must display itself in its properties goes without
saying; he who fully knows all a body's outward (physical) as well as inward (chemical) properties must in truth be
said to know the body's essence. But from a human standpoint the knowledge will, whether one otherwise holds to the
real or to the ideal principal side, always rest upon an approximation.

The science of experience is realistic; its separation between hard and soft, tough and brittle, elastic and inelastic
bodies is grounded upon the representation of the mutual relation of the invisible parts: the forms of state are of
essential significance for bodies' individuality. And yet the form of state can alter without the body's thereby losing
its substantial peculiarity. Ice, water and vapour may indeed be said to be different things, since with each mode of
aggregation there follows an alteration of properties; but the
% --- p. 245 ---
substance, the chemical character, is for water-vapour the same as for ice and water. The physical properties can
alter, and the chemical remain unaltered. Alterations of individuality can go on without alteration of substance; but
an alteration of substance cannot go on without altering the individuality.

Upon what, then, does an alteration of substance really rest? This question gets a different answer in physics from
the one it gets in philosophy.

It lies in the nature of the case that the science of experience with its exact analyses must lay particular weight
upon material discreteness. ``In our times,'' says H.\ C.\ Ørsted%
\footnote{% printed mark: *) --- it stands BEFORE the comma.
Naturl.\ mech.\ Deel. 1859. S.\ 14.}, ``the investigations of the laws by which substances of unlike kind enter into
those sensibly perfect unions which we call chemical have led to the necessary presupposition that the substances in
these unions do not completely penetrate one another, as Kant and his adherents assume, but that the act of union
stops at certain fundamental parts of bodies, which hang together with great intimacy without penetrating one another.
Hereby we are led to a doctrine of secret fundamental parts in bodies which has so much likeness to the old atomism
that occasion was hereby taken to give the chemical fundamental parts here in question the name of atoms.''

In the formation of new substances it depends of course upon the constituents' character, quantity and position. A new
body arises not merely by corresponding fundamental parts uniting in a certain proportion of quantity; bodies with a
new individuality can also come forth from the same constituents in the same quantity. With these so-called
\emph{isomeric} bodies a difference must then, as regards the manner of composition, be made between \emph{empirical}
and \emph{rational} formulas.
% --- p. 246 ---
% CHEMICAL FORMULAS. The element letters are printed as upright antiqua with lowered figures; they are
% set here as ordinary text with the figures in math mode, so that the letters keep the upright form the
% printing has. All series of figures were read under magnification on the KB copy --- the OCR is no
% witness here. NB. The printing writes „12 Vægtdele Brint“ (12 parts by weight of hydrogen) between
% „6 Atomer Kulstof“ and „4 Atomer Ilt“; both copies have it so, and the errata leaf is silent.
Let a given product be, for example, C$_6$ H$_{12}$ O$_4$, and so, on the empirical formula, consisting of 6 atoms of
carbon, 12 parts by weight of hydrogen, 4 atoms of oxygen; what substance is it now? There is in this matter a
possibility of two substances, a doubleness which finds its explanation in the empirical formula's admitting of being
separated into two rational formulas: C$_4$ H$_6$ O$_3$ + C$_2$ H$_6$ O, in which case the substance is acetate of
methyl oxide, and C$_2$ H$_2$ O$_3$ + C$_4$ H$_{10}$ O, in which case the substance is formate of ethyl oxide. The
substantial difference rests upon a rearrangement of the parts, for which reason such isomeric formations have also
particularly received the name of \emph{metameric}. In a corresponding way the alteration of substance is in certain
cases derived from the combinations containing the same substances in the same relative but not in the same absolute
quantity, so that \emph{polymeric} formations come forth. Thus C$_8$ H$_{16}$ O$_4$ gives the same percentage
composition as C$_4$ H$_8$ O$_2$, but since the quantity of atoms in each of the elements is different, the difference
between the two products may be substantial.

The atoms' different manner of position can of course as little be an object of immediate observation as their size,
figure or weight; but what substantial significance must be ascribed to the manner in which the constituents are
thought to be grouped can be taught us by the influence alone which the crystal form exercises upon the determination
of the physical properties. The crystal forms are determined chiefly by the axes, the central line about which the
whole formation is regularly laid out. That the axis really furnishes the essential point of support for the position
--- which must accordingly be different according as it happens parallel with the axis or perpendicular to it ---
confirms itself from many sides. The determinate directions along which the crystal cleaves stand in a relation to the
axis; the same holds of the passage of light which, since it rests upon the vibrations of the parts of the aether, is
precisely dependent upon
% --- p. 247 ---
the position of the atoms. Parallel with the axis tourmaline, for example, is almost opaque, whereas seen in a
direction perpendicular to the axis it is transparent; dichroite, which is blue in the direction of the principal
axis, is brownish yellow perpendicular to it. Something similar holds of heat. Crystals of the cubic system (the
spheroidal system), which has three kinds of principal axes, conduct heat equally well in all directions; whereas
calcspar and quartz, which belong to the rhombohedral system, which has only one axis, show a conductivity that is
greatest parallel with the principal axis and least perpendicular to it. --- From all this it is seen that the science
of experience must by its character on the whole go in an atomistic direction.

Over against physics, which derives the composite bodies' substantial individuality from the peculiar grouping of
unaltered fundamental parts, the speculative philosophy of immanence goes to the idealistic extreme of making the
substances of matter into immediate self-beings and ascribing individual subjectivity to the elemental bodies. So
Schelling in intuitive mysticism; so Hegel in overstrained logic.

% QUOTATION MARKS OVER TWO TURNS OF THE LEAF: the great Hegel citation opens here on p. 247 and is
% closed only on p. 249 („i sig selv.“). p. 247 is therefore +1 and p. 249 -1 in the quote balance;
% p. 248 has neither opening nor closing. „Cohæsion“ is letterspaced in BOTH halves across the line
% break, unlike pp. 239 and 244.
``The corporeal properties, predicates of the self-individualizing matter, are real moments of an indwelling activity
of form striving toward subjectivity. The movement of form springing from the principle of subjectivity runs through
two stages: from gravity to sound, and from heat to the electric light. Gravity is qualified into \emph{specific
gravity}; its articulating moments --- punctuality, linearity and superficiality --- are particularly posited in
qualified \emph{cohesion}. Cohesion's ideality consists in the material parts, existing outside one another, being set
in one another; the ceaseless alternation of the parts' being in one and being outside one another is the vibrating
and trembling that lets itself be heard in \emph{sound}. Sound for itself is individuality's self. The tone touches our
inmost part because it is itself an inner, a
% --- p. 248 ---
subjective. To individuality there belong matter and form; sound is the total form making itself known in time --- the
whole individuality, which is nothing else than that this soul is now set in one with the material. The next is that
the provisional, still imperfect individuality is dissolved in \emph{heat}. By heat matter is restored to its original
formlessness, its fluidity; it is its abstract homogeneity celebrating a triumph over the first specific
determinatenesses. The advance consists in matter's continuity, existing only in essence, being expressly posited as
the negation of negation, as activity, as existing dissolving. Out of this dissolving the individualizing activity of
form begins afresh in a deeper sense. The immediate shape, formless in itself, is in its one extreme punctuality, in
its other extreme the spherical fluidity --- shape as inner shapelessness. The point must now pass over into the line,
and form set itself on the line in polar extremes with an indifferent middle. This syllogism makes up the
shape-forming principle of the giving of form --- \emph{magnetism}. As the magnetic activity deposits its product, the
shape is determined as \emph{crystal}. The individual bodies are now developed into physical totalities, selfish
subjects whose particular predicates show themselves as physical properties in relation to the universal elements with
which they enter into real processes. In relation to light, the selfhood identical with itself, the corporeal
individualities utter their peculiar self-being in opacity, transparency and the many darkenings that give them
\emph{colours}; in relation to air the individuality is \emph{smell}, in relation to water \emph{taste}. The
development of form of the still immediately independent individualities ends with an actual conflict, a physical
tension of their particularity, their real selfishness, which is equalized in the electric light. It began with the
bodies having light outside them; it ends with
% --- p. 249 ---
% Two notes on this leaf, printed *) and **). The mark for the first stands AFTER the closing
% quotation mark.
every corporeal individuality having its light, its own subjective essence, in itself.''%
\footnote{% printed mark: *) --- the first of two notes on p. 249.
Hegels Vorlesungen über die Naturphil. S.\ 129--360.}
--- The arbitrariness in such a deduction is conspicuous.

Only by the doctrine of the ground-medium can the philosophy of nature be brought into essentially true agreement with
the science of experience without giving up its independence as a science of its own. In the ground-medium we have the
principle for the full acknowledgment of experience's facts; for object-power is presupposed, and the concepts' issuing
in power happens in masses and fundamental parts, so that these must naturally play the principal part in the
composition and dissolution of the substances of matter, the corporeal individualities. In the ground-medium we have
the principle for an actual articulation of the concepts of body; for self-power is presupposed: the articulating
thinking goes together into one with an articulating will. The system of the concepts of power, in its infinite,
peculiar connexion, is a system of will in which the concepts' transitions punctually furnish the transformations that
in actuality ground the corporeal alterations.%
\renewcommand{\thefootnote}{**)}%
\footnote{% printed mark: **) --- it stands AFTER the full stop.
On „Forandringsproblemet i Philosophien“ cf.\ Forelæsn.\ over „Phil.\ Propædeut.“ 1860--61, pp.\ 143--177.}%
\renewcommand{\thefootnote}{*)}

What elements and bodies are for one another mutually, in the ideal as in the real sense, must be known from their
mutual reciprocal actions, their law-determined activities; the law-determined activities are functions.

% THE CHAPTER'S CLOSE. The text on p. 249 stops about three fifths down the leaf; beneath it stands
% blank space and then a short centred rule, PRINTED DOUBLE, about two thirds down --- as at pp. 61 and
% 125. Per the house rule it is set as a single \streg.
\streg

% Danish: \kapitel{Andet Kapitel.}{Elementære Functioner.}   (printed pp. 249--317)
\kapitel{Chapter Two.}{Elemental Functions.}

% Chapter intro: prose stands between the chapter head and the A. head.
% --- p. 250 ---
% The chapter head OPENS p. 250 --- the page begins with the folio, beneath it the head, and beneath
% that at once the chapter's first paragraph. There is NO rule above the head (unlike pp. 1 and 197,
% where a rule separates \deel from \kapitel: here no Deel head precedes). Nothing of p. 249 runs over
% onto p. 250. The division head A. does NOT stand at the top of the leaf but about three quarters
% down: first the chapter's introductory paragraph runs to its end, then a short centred rule, then
% „A.“ on its own line.
In the elemental world of bodies the existing bodies are to be apprehended: partly in reflected being as moments for
the corporeal whole, and so subject to the \emph{universal}; partly as parts, subsisting members which under
many-sided conflicts hold themselves against one another in the \emph{particular}; partly, finally, as substantial
products, peculiar productions of the ground of nature at work in \emph{the individual}. Insofar as bodies stand,
according to their particular character, in a relation to a qualitatively determinate, universally corporeal
essentiality, the activities flowing from this show themselves as \emph{dynamic} functions, answering to \emph{light
and heat}. Insofar as what is proper to the bodies themselves maintains itself under mutually straining oppositions,
the reciprocal action passing through the dynamic universal essence will come into view in \emph{articulating}
functions, answering to \emph{magnetism and electricity}. Insofar, finally, as the whole activity of form is so
pervasive that the members become moments and the moments members of individual formations, we have the
\emph{body-forming} functions, answering to \emph{chemism and chemical processes}.

% Short centred rule between the chapter's introductory paragraph and division A, in the middle of
% p. 250.
\streg

% Danish: \afdeling{A.}{Dynamiske Functioner: Lys og Varme.}   (printed pp. 250--273)
\afdeling{A.}{Dynamic Functions: Light and Heat.}

The usual representation of the dynamic --- that substances are transformed into forces, penetrate one another and the
like --- is unclear and without significance in science. The dynamic has the mechanical for its foundation and consists
in
% --- p. 251 ---
the active parts' finite discreteness approaching the infinite to such a degree that it makes an impression of
continuity.%
\footnote{% printed mark: *)
Cf.\ Grundideernes Logik I, S.\ 354--69.}
When light and heat are referred to vibrations in the aether, the activity is mechanical; but what is peculiar here is
that the parts of the aether and their motions melt together in representation in such a way that the single terms of
the relation are not held fast for themselves but merge as moments in the activity. The closer difference is determined
by the qualitative. For gravitation's abstractly mechanical relations of weight all matter is of one kind; here there
are only the heavy masses and the empty space. The dynamic in light and heat is by contrast conditioned by a
qualitative difference in matter itself, an essential opposition between ponderable and imponderable matter. To call
the aether an ideal matter is however superficial, for the aether fills its space just as much as the weighable masses
do. When it is said that the aether ``penetrates transparent bodies,'' then the dynamic in this penetration is a reflex
of the mechanical. Even if every atom of a body is ``surrounded by a sphere of aether,'' it is nevertheless unthinkable
that the atom of aether should penetrate the body's atom and punctually occupy the very same space as it. That the
aether is without weight means only that on account of its extremely slight density it is unweighable; but it by no
means proves that it is in itself without gravity. If its density be assumed to be 845,000 million times less than
air's, so that a cubic mile of aether would contain no more mass than
% The fraction stands in the setting as a slanted fraction: 1 and then 1/3. The footnote mark **) stands
% BEFORE the comma, as printed.
$1\tfrac{1}{3}$ pounds%
\renewcommand{\thefootnote}{**)}%
\footnote{% printed mark: **)
C.\ Holten. Lysets Naturlære. Kbhvn.\ 1861. S.\ 167.}%
\renewcommand{\thefootnote}{*)}%
, then it is intelligible from this that the aether may very well be unweighable without being withdrawn from the law of
gravity. The ideal moment in the aether is that it immediately represents universal matter, a hypothetical foundation
for pervasive dynamic activities.
% --- p. 252 ---
But the opposition between the ideal and the real does not lie in matter itself; it rests upon the peculiarity of the
functions. The dynamic ideality is in \emph{light} what the corresponding reality is in \emph{heat}; the reciprocal
determination of the ideal and the real is articulated in \emph{the relation between light and heat}.

% Run-in head, read on the page itself (KB copy): set in semibold, not letterspaced.
\runin{a) Light.} To know what light is, we must begin by investigating the active element upon which its being rests.
Here, then, two activities of unlike kind must be presupposed, one in the form of otherness and one in the form of
selfhood, a making-visible and a seeing, a material and an ideal activity: light itself is the reflected unity of the
two activities of unlike kind. This can from the standpoint of natural science be uttered in the following two
propositions: 1) Light rests upon a trembling in the aether; when the aether is at rest, there is no light. 2) Light
rests upon an activity of sight; where there is no seeing, there is no light. --- Light's essence is accordingly at
once both objective and subjective. But although the two activities cannot possibly be separated from one another
without light's disappearing, there is nevertheless nothing to prevent each of them from being considered for itself
--- only that in considering the one, one must constantly be conscious of understanding the other along with it. We
treat first the objective essence and then the subjective, in order to reach a developed reflexion of light's real-ideal
nature.

% The Greek letter is read on the KB copy: α. The head is run-in and letterspaced.
\subrunin{α) Light's objective essence.} Seen from the side of nature, light must necessarily have a material
foundation answering to its activity. By apprehending the cause insofar as it reflects itself in its effects, natural
science seeks to make clear from all sides the activity conditioned by the foundation. What is not seen reflects itself
in what is seen: light's objective essence can be determined only by its objectification. What is peculiar about the
hypothesis of vibration
% --- p. 253 ---
is that the sensible element understood along with it disappears for observation; for what is observed is not the
vibrating parts of the aether, it is phenomena which as effects are referred to vibrations as to invisible causes;
what is fleeting in the representation transforms itself into a reflex of the motion: concepts and laws enter into
reflexions of
% p. 253 carries two notes, *) and **). Note **) runs on to p. 254, where its end stands at the TOP of the
% note block, before p. 254's own note *). The book does not repeat the mark on the continuation. Per the
% house rule the whole note stands here at its mark.
the cause.%
\footnote{% printed mark: *)
It lies near to compare water-waves, air-waves and aether-waves with one another; but they are nevertheless very
different. Cauchy has --- in his work Memoire sur la dispersion de la lumière --- treated the aether-waves with
atomism for a foundation and derived correct formulas for the dispersion of colours. From this Fechner has concluded:
ergo atomism is
% MISPRINT which the errata leaf does NOT correct: „Fcrmler“ for „Formler“ --- a wrong sort (c for o).
% Both copies and both OCR witnesses show it, so it is the compositor's fault. Translated by the sense;
% for comparison the correct „Formler“ stands three lines above.
proved. But the same formulas can be derived with the doctrine of continuity for a foundation; one might then with
the same (in)justice conclude: ergo continuity is proved. What can be shown on both presuppositions is that the
mechanical here passes over into the dynamic.
% „Memoire sur la dispersion de la lumière“ stands in ordinary antiqua, NOT in italic (checked on the
% KB copy); hence no \textit. The accent on „lumière“ is a grave.
}
The cause of the vibration divides itself originally into two causes, a moving and a living force, which are ceaselessly
converted into one another, ceaselessly consume and bring forth one another. This working is a pervasive reflecting;
that the cause works dynamically means: it is reflected in its mechanical moments.%
\renewcommand{\thefootnote}{**)}%
\footnote{% printed mark: **)
% The note begins on p. 253 and ends at the top of the note block on p. 254; it is gathered here at its
% mark. ± is a plus over a stroke, without dots, read on the KB copy; minus is by contrast an ordinary
% long minus. „v = o“ and „y = o“ are set with the letter o, not with the numeral nought; so reproduced.
That a part of the aether is in vibration means that it moves to and fro perpendicularly over a given straight line;
in the line itself it has its position of rest. If the distance from the position of rest be denoted by $\pm y$
answering to the time t, then the maximum for y can be denoted by $+a$ on the one side and by $-a$ on the other side of
the line. If the aether-part's mass be called $m$, $m\varphi$ the moving and $mv^{2}$ the living force, then the relation
between $\varphi$ and $v^{2}$ is the following:
$\varphi = -ky = \dfrac{d^{2}y}{dt^{2}}$;
$\dfrac{dy^{2}}{dt^{2}} = v^{2} = k(a^{2}-y^{2})$.
From this it is seen that when $y = a$, then $\varphi$ is at its maximum but $v = o$, that is, when the moving force is
greatest, the living force has disappeared. When $y = o$, and so the aether-part is in the position of rest, then
$v^{2} = ka^{2}$ but $\varphi = o$, that is, when the living force is greatest, the moving force has disappeared.}%
\renewcommand{\thefootnote}{*)}

% --- p. 254 ---
If the opposition of light and darkness were merely immediately qualitative --- either light or darkness --- then the
concept of vibration in the form of universality would be satisfactory. But the very fact that the illumination is
weakened with the spreading presupposes a law and demands a grounding. The grounding follows naturally from the
concept's development, for the concept of vibration presupposes that the waves spread themselves as spherical surfaces
about the luminous point, and that the aether is perfectly elastic. In two concentric spherical surfaces the sums of the
living forces will accordingly be equal. Since now the spherical surfaces' mass grows with the distance, and the
spherical surfaces stand in direct ratio to the squared radii, the living forces --- and so the strength of the light
--- must stand in inverse ratio to the squares of the radii, that is, of the distances.%
\footnote{% printed mark: *)
% The whole enumeration m, n, r, v and their primed terms is set in math, so that the series does not
% change face on the way. In the formula $r_{1}$ stands with a lowered figure one, while the running text
% writes $r'$ with a prime; both copies show it, and the same happens on p. 256. Reproduced as printed.
For a spherical surface with mass $nm$ let the radius be $r$ and the velocity of vibration $v$; for a spherical surface
with mass $n'm$ the radius $r'$ and the velocity of vibration $v'$; one will then have $nmv^{2} = n'mv'^{2}$. Since now
$\dfrac{n'}{n} = \dfrac{r_{1}^{2}}{r^{2}}$, it follows that
$\dfrac{mv^{2}}{mv'^{2}} = \dfrac{r_{1}^{2}}{r^{2}}$.}
But now the opposition between strength of light and colour? For the grounding of this opposition too the concept of
vibration contains the necessary conditions, namely such particular moments as the amplitude, the distance from the
position of rest in the running time, and the time of vibration.%
\renewcommand{\thefootnote}{**)}%
\footnote{% printed mark: **)
% The note begins on p. 254 and runs on at the top of the note block on p. 255; gathered here.
% „arc (sin = y/a)“ is Nielsen's own way of writing arcsine and is kept. All „o“ in the equations are
% the letter o, not the numeral nought.
From the equation $\dfrac{dy^{2}}{dt^{2}} = k(a^{2}-y^{2})$ one obtains
$\sqrt{k}\,dt = \dfrac{dy}{\sqrt{a^{2}-y^{2}}}$ and from this again, by integration,
$\sqrt{k}\,t + C = \operatorname{arc}\left(\sin = \dfrac{y}{a}\right)$, and so
$y = a\sin(\sqrt{k}\,t + C)$. To $t = o$ answers $y = o$,
% The setting has no equals sign before the last o: „og C --- i første Qvadrant --- o.“ stands so in both
% copies. Reproduced as printed.
consequently $o = a\sin C$, and C --- in the first quadrant --- o. One accordingly has
$\sqrt{k}\,t = \operatorname{arc}\left(\sin = \dfrac{y}{a}\right)$.
If the time of vibration be denoted by T, then $t = \tfrac{1}{4}T$ answers to $y = +a$, and so
$\sqrt{k}\,\tfrac{1}{4}T = \operatorname{arc}(\sin = 1) = \dfrac{\pi}{2}$,
$\sqrt{k} = \dfrac{2\pi}{T}$, and the formula of vibration
$y = a\sin\dfrac{2\pi t}{T}$, in which a is the amplitude.}%
\renewcommand{\thefootnote}{*)}
% --- p. 255 ---
The strength of the light rests upon the amplitude, the colour upon the time of vibration. From this it is seen that the
amplitude can be different and the time of vibration the same, that is, the strength of light can be different and the
colour the same; conversely the amplitude can be the same and the time of vibration different, that is, the strength of
light can be the same and the colour different.%
\footnote{% printed mark: *)
% The Greek letter is read on the KB copy: β. ± is a plus over a stroke, without dots. The period between
% 2πa/T and cos is Nielsen's multiplication sign and is kept. The citation opens „Da det ...“, closes
% after „lige stærke“, re-opens after the interpolation „(eensfarvede)“ and closes after „af dem.“ ---
% the book's usual habit; the quotation marks balance.
If the velocity of the light-waves be denoted by V, the time of vibration by T, the wavelength by $\beta$ and the number
of vibrations by N, then one gets $T = \dfrac{\beta}{V}$ and $N = \dfrac{V}{\beta}$, from which it is seen that the time
of vibration is in direct and the number of vibrations in inverse ratio to the wavelength. The red rays have the
greatest, the violet the smallest time of vibration; when 500 billion vibrations in the second give red, then 733
billion give violet. The strength of the light, by contrast, has for its measure the square of the velocity with which
the aether-part passes through the position of rest. From the formula $y = a\sin\dfrac{2\pi t}{T}$ there follows, namely,
$\dfrac{dy}{dt} = \dfrac{2\pi a}{T}\,.\,\cos\dfrac{2\pi t}{T}
 = \dfrac{2\pi}{T}\sqrt{a^{2}-y^{2}} = v$. In the position of rest $y = o$, the velocity
$\dfrac{dy}{dt} = v = \pm\dfrac{2\pi a}{T}$, and consequently
$v^{2} = \dfrac{4\pi^{2}a^{2}}{T^{2}}$. ``Since, however, it does not turn upon an absolute determination, one may just
as well use $a^{2}$ as the measure of the strength of light, so that two equally strong'' (monochromatic) ``rays of
light which meet under conditions as favourable as possible will give an illumination 4 times as strong as one of
them.'' C.\ Holten. Lysets Naturlære. S.\ 173.}
Finally,
% --- p. 256 ---
the concept has in it the possibility of taking up all the particular conditions from which the laws for the
aether-waves' forward-advancing motions
% The footnote mark stands here without preceding punctuation, in mid-sentence.
%
% MISPRINT, corrected by the book's own errata leaf: p. 256, line 12 from the foot: „vil --- læs ville“.
% Per the house rule the text is reproduced AS PRINTED. IMPORTANT: the working scan is pencil-corrected by
% an earlier owner, and the pencil HAS reached this place --- a written-in „le“ stands after „vil“, which
% the Google text layer reads along and renders as „ville“. The KB copy, which is not corrected, clearly
% prints „vil“. The reading has therefore been taken on the KB copy.
%
% In the text m', x' stand with primes, in the formulas $a_{1}$, $x_{1}$, $t_{1}$ with lowered ones --- as
% printed, see p. 254. „y = 0“ is here by contrast a NOUGHT (not the letter o), checked on the KB copy.
%
\footnote{% printed mark: *)
If a wave propagate itself from the point $m$, in which the vibration is denoted by
$y = a\sin\dfrac{2\pi t}{T}$, to $m'$, the way being $x$ and the time employed upon it $\theta$, then one gets for the
vibrations of the aether-part $m'$ the formula $y = a\sin\dfrac{2\pi(t-\theta)}{T}$.
If two waves are simultaneously to run through the same line, but from different starting points, so that the one
traverses the way $x$, the other $x'$, before they arrive at $m'$, then --- on the presupposition that the wavelength
$\beta$ is for both the same, but the amplitude for the one $a$, for the other $a_{1}$ --- the vibrations of $m'$ will
make up the sum of the vibrations of both and be expressed by the formula:
\[
y = a\sin 2\pi\left(\frac{t}{T}-\frac{x}{\beta}\right)
  + a_{1}\sin 2\pi\left(\frac{t}{T}-\frac{x_{1}}{\beta}\right).
\]
If $\dfrac{t}{T}-\dfrac{x}{\beta}$ be denoted by $\dfrac{t_{1}}{T}$, $x'-x$ by $e$, and $a = a_{1}$ be assumed, then the
calculation shows that
$y = 2a\cos\dfrac{\pi e}{\beta}
 \sin\left(\dfrac{2\pi t_{1}}{T}-\dfrac{\pi e}{\beta}\right)$;
from this it is seen that for an odd number of half wavelengths one gets $y = 0$, so that light with light gives
darkness; for an even number, by contrast, $y = 2a\sin\dfrac{2\pi t_{1}}{T}$, so that light with light gives
strengthened light.}
and their mutual actions upon one another admit of being derived. It can be exactly calculated how two light-waves which propagate themselves simultaneously along the same line, but from
different starting points, will meet, in which cases they will strengthen and in which cases they will cancel one
another; from this it can be explained that light with light can give not merely strengthened light but also darkness.

From these indications it must surely be possible to judge in what sense the concept of vibration answers objectively to
% --- p. 257 ---
the concept of light and determines light's cause in nature in such a way that the whole system of particular causes of
particular phenomena of
% The footnote mark stands without preceding punctuation, before „o. s. v.“
interference, reflexion, refraction, polarization%
\footnote{% printed mark: *)
% MISPRINT, corrected by the book's own errata leaf: p. 257, line 13 from the foot: „Reflection --- læs
% Reflexion“. Per the house rule the Danish is reproduced AS PRINTED. THE READING WAS TAKEN ON THE KB COPY
% and is POSITIVELY settled: after „Refle“ the note has TWO sorts --- a round, open c and then a t with a
% plain ascender and cross-stroke --- while the text's „Reflexion“ on the same page has ONE narrow sort
% without an ascender. The working scan is NO USE here: the pencil has also reached this place and written
% an x over the „ct“.
Reflection (tilbagekastning) and refraction (brydning) always happen according to the same laws as regards direction,
but not as regards intensity. Two parallel mirrors A and B with an inclination of a little over $33^{\circ}$, so that
the angles of incidence and emergence at both become about $57^{\circ}$, show that the rays cast back from A are either
wholly or only partly or not at all cast back by B, according as the angle between the mirrors' planes is $0^{\circ}$,
$90^{\circ}$, $180^{\circ}$ or something in between. Hereby light's polarization is determined.
% The letterspacing in the note is checked on the KB copy: „Almindeligviis“ is letterspaced, „hedder det,“
% is NOT, then everything is letterspaced down to and including „plansat“, while „(polariseret)“ again
% stands in ordinary antiqua.
\emph{Usually}, it is said, \emph{the vibrations perpendicular to the ray go on in all directions, but under certain
conditions the vibrations are bound to a certain plane, and the light is then said to be polarized}.}
and so forth, with all the modifications and shadings belonging under them, admits of being developed from one
fundamental cause. In the reflected unity of the phenomena of light with the natural concept of cause we know light's
objective essence.

% The Greek letter is read on the KB copy: β. The head is run-in and letterspaced.
\subrunin{β) Light's subjective essence.} Without sight no light. Apprehended from this side, light accordingly becomes
not a function of the vibrating aether, not a function of otherness but one of selfhood, a function of the seeing self.
This, considered for itself, is light's subjective essence, its ideality. This ideality, though it is demonstrably
subjective, the speculative philosophy of nature has none the less transformed into something objective in nature, and
thereby given a new proof of what ambiguities of abstraction the doctrine of immanence must resort to when it seriously
attempts to engage with actuality. Light, says Hegel, is matter's pure selfhood, over against which matter itself, even
according to its immediate being, is the selfless, darkness. Light is the infinite oscillating in itself, the pure
reflexion into itself, the same as what Spirit's I is
% --- p. 258 ---
in a higher form. Space, light and Spirit so far denote a stepwise development of form of the same essence. ``I am the
infinite space, self-consciousness's infinite likeness with itself, the abstraction of empty self-consciousness, the
pure identity of me with me. I am only the identity in the relation of myself as subject to myself as object. With this
identity light is parallel; of this it is a faithful image. It is only not I because it is not darkened and broken in
and by itself, but is only the abstract shining.''

By zealously insisting on light's objective ideality the doctrine of immanence has supposed that it was really fighting
for the right of spirituality, and has on this account raised a vehement --- one may well say fanatical --- polemic
against physics. It is especially Newton who has to suffer for it; Hegel will show in his experiments insipidity,
incorrectness, indeed dishonesty. The coarsest representations are ascribed to the physicists: that they speak of
bundles of rays, of four-cornered rays, that they --- \textit{horribile dictu} --- regard light as something material
and composite and let every single ray have its single colour in the white light, just as every single horn has its
single tone in a Russian horn band, and so forth. With all this, however, what is really the main thing in physics'
exact investigations is completely misunderstood, and that to great damage --- not for physics, but for philosophy.

Physics does not teach us that the vibrating aether is in and for itself light; it makes the eye the mediating organ
between the active element in nature and the activity of sight, and apprehends the cause of light as a law-determined
reciprocal action of two causes of unlike kind, an objective one in the aether and a subjective one in the soul. That
physics will not engage in a particular explanation of light's ideal essence is no proof that it denies light's
ideality; it is only a proof that it neither is nor will be philosophy. Physics holds
% --- p. 259 ---
fast its standpoint, knows its problems and knows what it has to do as an investigating science. It cannot, as regards
light's side in nature, content itself with general turns of speech about ``oscillating,'' ``reflecting,''
``manifesting''; it develops the doctrine of vibration into an exact system of functions and laws. In expressly
presupposing the activity of sight, physics teaches neither that the strength of light straightforwardly \emph{is} a
square of the amplitude, nor that the time of vibration \emph{is} the colour, so that 500 billion vibrations in the
second are red and 733 billion violet; but it teaches that the different sensations of light which the eye receives are
effects of different corresponding causes, and shows in this respect exactly that the strength of light is a function of
the amplitude, and the colour a function of the time of vibration. Physics does not teach that the rays of light are
material lines shot out in bundles from the luminous body; it teaches on the contrary that they are mathematical lines,
radii, that is, lines of direction upon which the wave-surfaces stand perpendicular. Expressions such as cones of light,
cylinders of light, bundles of rays and the like are partly geometrical determinations of form, partly figurative
designations aiming at securing the representation and supporting the intuition. Optics divides and composes the rays of
colour, not in a material sense, but in a manner corresponding to the way mechanics divides and composes
% The footnote mark stands AFTER the full stop. One footnote on the leaf.
forces.\footnote{This can among other things be seen from the fact that Cauchy, for example, introduces imaginary waves
into the analysis. Cf.\ Holten, Anf.\ Skrift, S.\ 280.} Physics cannot, as regards the wave theory, possibly assume that
white light is a sum of colours existing in it; for when all colours melt together in the white, then they are precisely
by this melting together cancelled; the vibrations for white light must then form a common wave in which the particular
colour-waves are taken up and contained. It assumes, by contrast --- and with
% --- p. 260 ---
full right --- that white light is the actuality in which all the colours have their possibility, that the different
colours come forth when the light is broken and the common wave divides itself into particular waves with particular
times of vibration.

Had the speculative philosophy of nature seen itself able to give as sufficient an account of the subjective functions
in the activity of sight as physics has given of the objective ones in the activity of nature, then it would have stood
otherwise with the apprehension of light's ideality than it now does.

To place light's ideality in ``matter's pure selfhood'' is meaningless, for matter has no selfhood at all; it is
precisely matter's character to be the selfless other. There is as little any immediately and straightforwardly luminous
matter as there is any seeing, self-intuiting matter. What is luminous in matter is an activity, but an activity in the
form of otherness; therefore this luminous element can shine neither for itself nor for any matter whatever. This is the
contradiction: that the luminous element in matter cannot shine unless it is taken up into an activity of sight, and so
into a self-activity; only as the self-activity sets in does nature, so to speak, get the benefit of the activity in
another, for now what is luminous can shine, now matter can be illuminated.

That the activity of sight, though conditioned and determined by mechanically dynamic causes, makes what is not in
itself light into something luminous, makes functions of another into functions of its own, of selfhood, is in the form
of ideality a categorical expression of light's subjective essence.

% The Greek letter is read on the KB copy: γ. The head is run-in and letterspaced; the hyphen in
% „reel-ideelle“ stands in the printing.
\subrunin{γ) Light's real-ideal nature.} From the objective side light is a function of the vibrating aether; in this
consists its reality. From the subjective side it is a function of the seeing self; in this consists its ideality. Light
is accordingly an indissoluble unity, reflected in its essence, of real and ideal, of objective and subjective. None the
less
% --- p. 261 ---
this indissoluble unity is at the stage of finitude nevertheless dissoluble; the objective activity of sight falls in
actuality outside the subjective. Light shines, after all, whether we see it or not; the sun of day would shine over the
earth even if all the eyes of animals and men were closed. So strong is the fundamental sensation of the self-validity
in light's shining, its changing clarity and the colours of things, that we are tempted at every moment to transfer its
ideality into outward nature: only when by conferring with the speculative philosophy of nature we have properly come to
our senses does the temptation pass.

The contradiction that light in nature has no ideality, and that light's nature is nevertheless both real and ideal ---
this apparent contradiction, existing in the semblance as semblance --- cannot be resolved from a merely finite
standpoint; its resolution is in the fundamental relation of nature and Spirit. If Spirit, the Spirit almighty in
itself, which grounds and bears the whole of nature, were merely a thinking and not at the same time an intuiting,
seeing being, then the concept of natural light would in actuality be as contradictory as the concept of a round square.
But Spirit is intuiting, seeing, objectifying self-power; in self-power's objectification light has an objective natural
ideality independent of all finite seeing.

When, however, it is particularly brought out that natural light --- the light here spoken of --- is not a spiritual but
a sensible light, an objection arises. Sensible light is, after all, only insofar as it can be sensed and seen with
sensible eyes. If Spirit by its seeing and intuiting were to give light sensible ideality, then this seeing and
intuiting would have to be sensible, not spiritual; Spirit as Spirit would then have to have sensible eyes. We must, as
regards objections of that kind, dismiss all loose representations of a loose imagination. Here it is not asked what and
how many eyes Spirit has, or how a Spirit looks with sensible eyes; here it is asked how it is possible
% --- p. 262 ---
that light can have an ideality independent of finite seeing; and the answer runs: it is possible only in virtue of an
infinite seeing; possible only because the whole of nature has its ideality from Spirit.

That nature's ideality and its sensible reality are not finitely put together but are grounded in an original unity can
be known from the ground-medium. In the ground-medium the causes of unlike kind whose co-working is the distinguishing
% The footnote mark stands BEFORE the comma. One footnote on the leaf.
mark of all aesthetophysical functions\footnote{Cf.\ Grundideernes Logik I, S.\ 369--87.} have their common source. From
the unity in the ground-medium the separation issues in the object-medium; upon this rests all the parcelling out of
finitude. Outwardly regarded, the eye and light stand in an accidental relation; they can meet accidentally and be
separated accidentally, for the eye can be without light, and light without an eye. But in nature and essence eye and
light are mutually presupposing concepts, real concepts of the same conceiving: light's concept is originally reflected
in the eye's concept, just as the eye's concept --- as its structure sufficiently shows --- is reflected in light's
concept: upon the double-sided reflexion of the concept of light in the aether and the concept of light in the eye rests
light's real-ideal nature.

% Run-in head, checked on the KB copy: „Varmen.“ is set in semibold, „b)“ in ordinary antiqua.
\runin{b) Heat.} Of heat too it holds that the objective is determined subjectively; feeling is for heat what sight is
for light. None the less a difference shows itself here; for in the doctrine of heat the outwardly immediate objectivity
plays a far more prominent part than in the doctrine of light. In all phenomena of light the effect of light must
unbrokenly be determined by sight; but the phenomena of heat cannot be determined in that way by feeling. Feeling is a
most imperfect measure of heat; so the objective side --- the alterations which heat effects in bodies, the expansion
particularly --- must be brought out from the beginning. The expansion
% --- p. 263 ---
lies at the foundation of the construction of the thermometer. By the thermometer the subjective in the apprehension is
altered into something objective for sight, the objective sense steps into feeling's place: heat, though invisible,
becomes visible. But although heat is real in another way than light, it will nevertheless be possible, in good
agreement with physics, to show that it too has its ideality, and that it is precisely this ideality that grounds its
real activity.

% The Greek letter is read on the KB copy: α. The head is run-in and letterspaced.
\subrunin{α) Heat's ideal essence.} Physics speaks of ``the unknown cause of what we call heat.'' How matters at bottom
stand with the representation of known effects of an unknown cause has already been elucidated in the doctrine of force.
As regards heat we have only to show how the concept of cause originally stands to the concept of the elemental
function. In the form of fact, and so outside the concept, the function is pieced together out of finite presuppositions.
Heat's cause, what warms, is represented as a something for itself, an unknown matter, while the effect --- the
alteration that goes on with the warmed body --- is observed for itself and drawn in under what is known. The warming,
the function itself, then becomes from this standpoint both inexplicable and explicable: inexplicable insofar as it has
an unknown cause; explicable insofar as the visible effects admit of being observed. It is otherwise when the function is
apprehended according to its concept, and so from the side of essence. That heat essentially consists in an activity no
one will surely deny. One may now provisionally form whatever representations of the material element in the cause one
will; it is nevertheless incontestable that what warms is a cause which must expressly be determined and known from its
effects. Were heat without more ado a property of a matter of heat, and a body's warming to consist solely in the warm
matter's penetrating the body in smaller or greater quantity, then there would be no activity of heat at all: the warming
% --- p. 264 ---
itself would be a pure fancy. The matter of heat does not become warmed, for it is warm. Nor does the body become warmed,
for the heat in the body does not really concern the body; it is on this presupposition not the body but the matter of
heat in which the heat is, and the matter of heat is warm. If one will not stop at the empty, meaningless representation
that the warm is a matter of heat and the matter of heat the warm, then one must necessarily concede that heat is
activity, and that what is warmed takes part in the activity just as much as what warms. That there is something which
warms rests upon there being something which receives heat, which lets itself be warmed. But to receive heat, to let
oneself be warmed, is just as surely conditioned by an activity as is warming. The working and the suffering do not fall
apart as two things: the working is suffering, and the suffering working. Here is no unknown cause in outward opposition
to a known effect, for the cause is reflected in its effect and the effect in its cause: the elemental function is a
function of dynamic reciprocal action. Heat as cause, heat as effect, and the warming answering to the reciprocal action
are side-concepts of one and the same concept of nature, heat's real concept. This real concept is heat's ideality.

Considered for itself as a universal cause reflected in its effects, heat is not an extensive but an intensive magnitude.
It rises and falls in continuous transitions, and there is as little any absolute limit to its falling as to its rising.
The whole of heat is in every degree of heat, and yet heat is different from itself in its different degrees. It has
according to its ideality all conditions enclosed in its real concept; it grows and decreases in itself, by itself, out
of itself: there is nothing from without that comes either to give or to take away heat.

% --- p. 265 ---
But, one objects, the bodies which are warmed do come from without! That is, as regards the essence, a great
misunderstanding. Just as the cause must reflect itself in its effects, so in turn the many effects reflected upon the
bodies must reflect themselves in the cause. What the bodies are for themselves in other respects, and whence they come,
does not concern heat; only what they show themselves to be according to their share in the activity of heat concerns
heat. The corporeal forms of state have no subsistence outside or alongside universal heat; they are on the contrary
forms of transition for a determining activity of heat, phenomena of heat's essence, functions of determinate degrees of
heat. That the forms of state stand in heat's service and make up moments of heat's concept can be seen from the
transitions. If a body passes over from solid to fluid, or from fluid to gaseous state, then heat is bound; if the
transition is the reverse, then heat is set free. From this it is seen that heat comes to existence in the form of
corporeal states, that every form of state expressly represents a certain determinate capital of heat. That heat is bound
means that the capital is increased, because the state entering is itself the measure of the increased capital; that heat
is set free means that a deduction is made from the capital, because the state entering is precisely the expression in
existence of the diminished capital. If specific heat be made the measure, the same repeats itself: bodies are measured
with one another in being measured with heat; but heat is, by means of the bodies, measured with itself, for the measuring
is done with units of heat. From this standpoint it shows itself that the otherwise different concepts of body are all
comprehended in heat's universal real concept; that is heat's ideal essence.

% The Greek letter is read on the KB copy: β. The head is run-in and letterspaced.
\subrunin{β) Heat's real essence.} From the standpoint of ideality, working and activity merge into reflecting and
% --- p. 266 ---
conceiving; but the reflexion of power demands at the same time the reverse: that the actual conceiving be expressed in a
real manner in working and activity. For reflexion and concept the active members are reflexes, moments,
side-determinations; for the sensible, outward-going activity, the actuality visible in facts, the real concept's
essential moments are existing, active members. Thus heat gets its reality in virtue of its real concept. Heat as
universal cause must then, like light, have its material foundation and in it a particular subsistence for itself. Over
against this the different bodies then come forward, each with its material independence, and come from without under the
heat communicating itself in different measure, and communicate heat to one another in turn. The warming, which rests
precisely upon the real opposition of heat and body being converted into a real unity, hereby gets the character of
material activity. That an actual unity of opposition cannot possibly be apprehended as straightforwardly existent is
easy to see. The immediately existent unity of opposition is, namely, a contradiction, since it is composed of two
mutually excluding determinations. If the opposition be posited as existent, then the unity is excluded; if the unity be
posited as existent, then the opposition is excluded: the resolution of the contradiction is activity. The activity, the
dynamic function of warming, consists in a ceaseless conversion, a reciprocal action in the course of which the unity
proceeds out of the opposition --- so that heat becomes corporeal and bodies warm --- and the opposition again out of the
unity, so that there is nevertheless an actual difference between heat and bodies. Hereby heat's real essence is
determined in general.

If this actual essence is now to be more closely investigated; if the investigation is to go into detail; if it is the
activity's inner material \emph{how} that is the object of the investigation --- then physics has surely ground to speak
of the ``unknown'' in heat's cause. The unknown, so
% --- p. 267 ---
understood, is to be discovered by the way of observation, through an approximation. It lies near to let heat have the
same foundation as light; so physics then assumes that heat, in analogy with light, must rest upon vibrations in the
aether. The real concept to which the philosophy of nature has to hold is accordingly the concept \emph{analogy}.

As a concept of fact, analogy everywhere shows in the likeness a recognizable unlikeness, but in the unlikeness again a
recognizable likeness; for the likeness must in the foundation be what overreaches. As regards radiant heat, the likeness
to light stands in the foreground; as regards the conduction of heat and a number of other phenomena, it seems by
contrast as though the unlikeness were on the point of veiling the likeness. It is a well-known matter that heat can
spread itself in rays like light; that the rays of heat are refracted, and refracted differently, like the rays of light;
that they are cast back according to the same laws: the angle of emergence $=$ the angle of incidence, and so forth. What
the difference between transparent and opaque bodies is for light, the difference between diathermanous and athermanous
bodies is --- as Melloni has shown --- for heat. In all this the likeness is striking. But heat is not light; so the
unlikeness also comes into view, among other things in this, that dark bodies --- bodies which do not radiate light ---
can radiate heat. So analogy takes the unlikeness up into the likeness, and the explanation is that the mere rays of heat,
even though they can, as Tyndall has recently shown, pass through the eye's aqueous medium, nevertheless cannot act upon
the optic nerve, because the waves are too broad and the vibrations too slow.

The likeness's confirmation through the unlikeness is seen in a striking manner, among other things, from the fact that a
dark body --- a piece of iron, for example --- which during the heating sends out ever more and more refrangible rays,
becomes luminous as
% --- p. 268 ---
it becomes glowing, first red-hot and at last white-hot.

It is otherwise when heat propagates itself not by radiation but by conduction, passing by immediate contact from the one
body to the other and spreading itself through the whole mass. Here the unlikeness becomes so great that the analogy at
first glance seems to disappear; for to the highly characteristic difference between good and bad conductors of heat
there answers no such difference between good and bad conductors of light. But the unlikeness bends again under the
likeness, and the explanation is that the waves of heat --- the vibrations of the aether which bring forth heat in the
body --- at the same time set the body's own small parts in vibration. The peculiar reciprocal action between the
vibrating parts of the aether and the body's vibrating small parts finds a confirmation in the doctrine of specific heat
and atomic heat: the analogy is re-established.

The attempts made from several sides (Redtenbacher and Clausius) to form a quite clear representation of the relation
between the vibrations of the aether and those of the small parts --- and so of the inner how of the reciprocal action
--- have no scientific significance as hovering hypotheses, but would have significance if there could thereby one day be
reached a presupposition that really grounded the laws of heat; for thereby a deep and secure insight into heat's real
essence would be opened to us.

% The Greek letter is read on the KB copy: γ. The head is run-in and letterspaced; the hyphen in
% „ideel-reelle“ stands in the printing.
\subrunin{γ) Heat's ideal-real nature.} How intimate the natural connexion is between heat's ideal and real essence can
be seen from the laws. Heat penetrates and expands all bodies; this fundamental law is in equal degree decisive for the
nature of bodies and for the nature of heat. What is realistic is that heat and bodies have each their own nature with
which they enter into reciprocal action; what is idealistic is that the two natures are at bottom one nature, and that it
is precisely the unity of the fundamental nature that makes the reciprocal action, and with it
% --- p. 269 ---
the fulfilment of the laws, possible. Both conceptions find equal confirmation in the doctrine of cohesion.

On the realistic conception the body's fundamental parts would, if heat did not work, only draw themselves together, for
which reason every body becomes firm and stiffens in the cold; but the aether surrounding the fundamental parts works by
its vibrations repulsively, removes the trembling parts from one another and effects an expansion. Cohesion accordingly
rests upon a double nature at work in bodies: that of the parts of the aether and that of the fundamental parts. The
realistically mechanical way of seeing, however, dissolves of itself into the idealistically dynamic. It is a fact that
it is by no means always increased heat that causes a body's expansion, but that a diminution of heat can effect the
same. Pure water at $4^{\circ}$ expands on being warmed, but it also expands on being cooled. Now it would be
inconsistent to assert that the cause of the water's expansion above $4^{\circ}$ lay in the vibrating parts of the
aether, while the cause of its expansion below $4^{\circ}$ lay by contrast in the vibrating parts of the water. But when
the parts of the aether and the corporeal fundamental parts have, in the dynamic reciprocal action, an equally essential
share in the expansion, this means in other words that the nature of heat and the nature of water work as one fundamental
nature. The unity of the two natures has its exact expression in the law; the law is essence, for it is concept.

If we consider the particular laws, the realistic separation returns: the double nature comes forward, but always on the
foundation of unity, that is, of ideality. The law for the expansion of solid bodies is in its universality to be
expressed thus: the coefficient of expansion of volume is three times as great as the coefficient of expansion of length.
In this universal the unity of essence, and so the ideality, is uttered. The coefficient of expansion of length --- the
figure indicating how great a part of its length at $0^{\circ}$ a body expands
% --- p. 270 ---
under an increase of heat from $0^{\circ}$ to $100^{\circ}$ --- is however different for different bodies: platinum,
glass, iron, copper and so forth. Does this difference now fall outside or inside heat's nature? In reality it falls
outside, for the coefficients rest upon the peculiarity of the bodies and are so far no concern of heat's. And yet in
ideality it falls inside; for the actual expansion must, according to the law, always presuppose an actual coefficient,
and the one actual coefficient is in the determination of the activity of heat just as essential as the other: the
difference of the coefficients is a necessary condition for particular utterances of heat's nature. --- The law for
melting has its universal expression in the two determinations: the unalterability of the melting point, and the binding
of heat. What is realistic is seen at once from the position of the melting points at different heights --- when ice
% The sign before 39 is ÷, Nielsen's minus (cf. p. 46), read on the KB copy. It stands here in running
% prose, not in a formula.
melts at $0^{\circ}$, then hammered English iron melts at $1600^{\circ}$ and quicksilver at $\div\,39^{\circ}$ --- but the
different binding of heat proves in turn what is idealistic, namely that the corporeal differences concern heat as much as
they concern bodies; for the bound heat is just as much the ground of the altered state as the alteration of state is the
ground of the binding of heat: the nature of the bound heat is the state, and the nature of the state the bound heat. ---
When the degree of heat in a body whose weight is m and whose specific heat is c is increased or diminished by t degrees,
the body must receive or give off a quantity of heat which can be expressed by the product mct. This law, which utters in
the product that weight, specific heat and degree of heat are not merely factors of a mathematical multiplication but real
moments in an actual concept of nature, then also shows --- as on the whole do all the laws of heat --- how intimately the
nature of heat and the nature of bodies make up, through all real differences, one indivisible fundamental nature.

Heat without bodies and bodies without heat are
% --- p. 271 ---
unnatural things. The subjective moment in heat's essence does not indeed come into view so immediately as in light; but
the real concept, the pervasive concept of power which makes the actual differences into moments for the active unity, is
for heat just as subjective-objective as for light. Heat has, like light, its real conditions and sources in nature,
because it has, like light, its ideal conditions and sources, its real concept, in self-power's infinite conceiving, in
Spirit.

% Run-in head, checked on the KB copy: „Lys og Varme.“ is set in semibold, „c)“ in ordinary antiqua.
\runin{c) Light and heat.} The dynamic functions of light and heat presuppose, as we have seen, a universal material
foundation upon which the activity of light becomes real and the activity of heat ideal --- yet so that ideality is what
overreaches in light, reality in heat. There form themselves accordingly, upon the foundation, two systems of functions
analogous to one another and yet of unlike kind.

That light with all its reality is idealistic is seen from its relation to sight. The whole outward world is a world of
images of light because it is a world of images of sight. We suppose we see the actual, material objects themselves, but
that is --- as optics sufficiently proves --- a fancy. The space with objects that we see in the room has as little
immediate actuality as the space with objects that we see in the mirror. There is indeed an actual difference between the
mirrored space and the space in the room, but this difference is not objectified by mere sight. In light for sight
everything is semblance, reflected semblance, mutual semblance; the difference of the colours rests upon a difference of
semblance: it is a pervasive illusion. But this illusion is no deception; it is on the contrary an ideality of things,
whereby they show themselves as phenomena in visible forms and relations and therein display their essence.

But what now is the world of bodies in relation to heat?
% --- p. 272 ---
Heat governs the forms of state: all solid bodies can melt. One has indeed hitherto, it is said, not
% The quotation marks stand as printed; the closing mark stands BEFORE the full stop (checked on the KB
% copy). The citation is closed on the page; the balance remains nil.
availed to melt coal, ``even though several physicists have claimed to have found traces of melting on the edges of the
diamonds they have experimented with. By analogy one must conclude that there is no perfectly unmeltable body, but that
all would melt if only one could bring forth sufficiently high degrees of heat.'' By analogy one must then also conclude
that all bodies can evaporate at sufficiently high degrees of heat. With respect to heat the whole corporeal world is,
according to its material essence, to be apprehended as a world of vapours; the difference between these vapours is
again, when heat is thought to overstep all known limits and to rise to a mythical heat, vanishing. Everything is a
vapour, and matter a uniform, space-filling being; that is heat's judgment upon the corporeal world. Although there is
nothing illusory in the activity of heat, the ideality is nevertheless recognizable in this, that bodies' immediate
independence perishes in heat: heat is accordingly the ground of inwardness upon which their outward coming-forward,
their peculiar being, generally rests.

The two systems of functions, light's and heat's, proceed from one fundamental system; in this fundamental system all the
real concepts of the world of bodies are enclosed; the whole of nature is of one casting, a total concept grounded and
borne by infinite conceiving.

But that bodies, in virtue of the comprehending total concept, nevertheless also have an independence of their own over
against light and heat can be seen from their properties. Colour, for example, is just as much a property of certain
bodies as of light. The body that casts back the red rays is red; the one that casts back the blue is blue. The relations,
in this respect very entangled, rest as much upon the
% --- p. 273 ---
% The footnote mark stands AFTER the full stop. One footnote on the leaf.
inner structure of bodies as upon the refrangibility of the rays of light.\footnote{Cf.\ Holten, Lysets Naturl., S.\ 190,
308 and elsewhere.} As regards heat, bodies show by the difference between their melting points, binding of heat,
specific heat, diathermancy, conduction of heat and so forth, that it is they which, by their different properties, bring
something peculiar of their own into the universal activity of heat.

By being set in relation to light and heat, bodies are at the same time set in relation to one another. If rays of heat
which have passed through a glass plate be let fall upon a plate of alum, they are almost entirely absorbed, whereas a
plate of alum lets through very nearly all those rays of heat which have previously passed through a plate of citric acid.
Rays of light which have passed through a green glass can, as is well known, easily pass through another green-coloured
glass, but are absorbed when they fall upon a red glass.

Bodies' peculiarity is an individual that serves the universal and is mastered by the general in offering it resistance.
Were bodies so dependent that light could really transform them into semblance, then pure, unbroken light without shadows
and without colours would be as good as pure darkness. As phenomena of light, bodies are phenomena of their own essence in
being phenomena of light's essence. If all bodies had one and the same melting point, so that heat could transform them
all into vapour at the same time, then universal heat would be as uniform and as insignificant as the universal vapour in
the universal space. It may then also be said that bodies, in being phenomena of heat's essence, are at bottom phenomena
of their own essence.

% p. 273 ends mid-word: the line ends „Eiendom-“, and p. 274 begins with „melighed“.
If particular weight now be laid upon bodies' own
% --- p. 274 ---
peculiarity and upon the self-subsistence reflected in the universal, then the dynamic universal that dwells in the bodily
individualities takes on new and peculiar forms, and the functions alter their character in consequence of what is
particular.

% Division A ends four lines down p. 274; then a short centred rule and the head B.
\streg

% Danish: \afdeling{B.}{Leddelende Functioner: Magnetisme og Elektricitet.}   (printed pp. 274--293)
\afdeling{B.}{Articulating Functions: Magnetism and Electricity.}

In relation to light and heat, bodies do indeed display a multitude of different properties, whereby they make good their
mutual inequality and the corresponding self-subsistence; but the opposition of the individual and the universal is
nevertheless equalised, in the merely dynamic functions, in such a way that the universal determinations retain a
one-sided preponderance. Only when the dynamic functions pass over into becoming articulating functions, functions of
polar oppositions, as in magnetism and electricity, do the bodily peculiarities begin to come forward in full light; but,
remarkably enough, it is then also that what is proper to the bodily essence begins in earnest to show itself as an
impenetrable riddle.

The universal in magnetism and electricity is neither so at one with the essence of ponderable matter that it can be
grasped as a property in being, after the manner of gravity; nor has it so uniform a persistent being for itself that it
can be referred to a uniformly represented substratum after the manner of light and heat. On the contrary! When one tries
to hold fast the essence of electricity in particular as something universal for itself, be it in the aether or in some
other
% --- p. 275 ---
% The Danish breaks mid-word here: p. 274 ends „ube-“, p. 275 begins „kjendt“.
unknown fluid, it goes together with the bodies concerned; and when one then will hold it fast in the several bodies, it
frees itself and takes on the stamp of the universal.

The enigmatic is increased still further by the unrest and tension that hereby arises between representation and thought.
The physicist is not content to observe the phenomena and abstract the laws; he feels a need to form for himself a
representation of what the inwardly operative really is. But the inwardly operative is essence, and an essence does not
let itself be represented. What is magnetism? A peculiar force in certain bodies! This explanation satisfies neither
representation nor thought. Representation is not in a position to hold fast what is peculiar in the force, for no force
lets itself be represented straight off. Nor can what is peculiar in the magnetic force be held fast by thought; thought
demands something universally valid and necessary, but this force is precisely so peculiar that no one has any inkling
whence it comes, why just iron ore and iron, in preference to so many other bodies, should be endowed with the magnetic
force.

Just as untrustworthy are the representations that have been formed of electricity with thought's support. Dufay assumed
that there were two kinds of electricity or electric fluids, whereas Franklin maintained that there was only one kind;
and, as it is said, ``the last century's dispute between the Franklinists or Unitarians and the Dualists cannot yet be
regarded as finally decided.''

Not even by keeping to the laws can the observer secure for himself any clear representation of the common operative
cause. By analogy with magnetism one assumes that the electricity e at one point must, at another
% The formula is read on the KB copy: $\mu$ stands BEFORE the fraction, not in the denominator.
point at the distance r, produce a distributing force $K = \mu \dfrac{e}{r^{2}}$. By experiment it can further be shown
that two small bodies with the
% --- p. 276 ---
quantities of electricity e and e' must at the distance r repel one another
% Here the print sets an apostrophe (prime) on e', whereas pp. 279, 280 and 287 use a raised
% figure one (m$^1$, e$^1$). Both forms are read on the KB copy and are reproduced as they stand.
with a force $F = \mu \dfrac{ee'}{r^{2}}$. If now $\mu$ be assumed identical in both expressions, one would suppose that
the equation for the moving force F and the equation for the distributing force K must point to one and the same
fundamental force, one and the same common cause; but an acute physicist does not let himself be dazzled by a merely
formal
% printed mark: *) --- it stands AFTER the full stop.
likeness.\footnote{%
``It has been the view hitherto prevailing that electric
% The row of dots stands as four spaced points, as printed.
distribution . . . . is due to the electricities' mutual attractions and repulsions, and that these latter pass over into
the corresponding mechanical forces between the electrified bodies themselves, when the electricity is bound to the
bodies by surrounding insulators. The first assumption is an analogy drawn from our acquaintance with bodies' mutual
attractions and repulsions; but this is hardly fit to be carried through any further, partly because the distributing
force at a point appears independently of the electricity present at that point, partly also because the effects of this
force are not an acceleration of the electricity's motion at that point, but must be measured by the quantity of
electricity flowing through. But even if this representation be nevertheless held fast, these forces cannot, because the
electricity meets a resistance, be regarded without more ado as the cause of a mechanical pressure upon the body in a
definite direction, for this ``resistance'' is nothing else than a greater specific conductive resistance''. L. Lorenz.
Lectures on Natural Science. Copenhagen 1870. pp.\ 192--193.}

Over against everything enigmatic and uncertain in these hypotheses and representations, the philosophy of nature has to
lay down two wholly certain fundamental
% The letterspacing runs from „alle“ to „Opfyldelse.“ on p. 277. The figures „1)“ and „2)“
% are NOT letterspaced and therefore stand outside \emph.
propositions: 1) \emph{all bodies, however peculiar they may be in properties and combinations, are realised
power-concepts, so that every utterance of force is a natural utterance of concept}; 2) \emph{the power-concepts'
expression in existence makes itself specially known in the laws, for the
% --- p. 277 ---
% The emphasis runs across the page break: the Danish divides „Al-“ / „menes“.
universal's going together into unity with the individual is to be seen only in the fulfilment of the laws.} The relation
between physics and the philosophy of nature will then, with regard to these two fundamental propositions, be determined
thus: physics starts from the facts and discovers the laws; the philosophy of nature starts from the discovered laws and
develops the concepts of nature.

The articulating functions thus acquire natural-philosophical significance in that they lead, step by step, to an ever
deeper insight into the individualisation of the bodily fundamental relations. It will then be possible to see how the
polar opposition shows itself first more outwardly and immediately in \emph{magnetism}, next more inwardly reflected in
\emph{electricity}, and at last unfolds its doubly articulating reciprocal action in \emph{electromagnetism}.

% Run-in head, read on the page itself (KB copy): set in bold roman, not letterspaced.
\runin{a) Magnetism.} An attraction according to the law of gravity presupposes at least two members, two masses or parts
of mass; nevertheless there is in this no proper articulation; the one mass is as the other, and every difference in
properties, even in magnitude, is irrelevant to the essence of gravity taken universally. It is otherwise with the
attraction between two magnets; each of them has its poles, the unlike poles attract, the like repel one another;
attraction is conditioned by repulsion, and conversely; thus pole works against pole, member against member: these are
member-forming functions. In the poles the one mass is not merely outside the other,
% --- p. 278 ---
but the otherness is accentuated twice over, the one pole is qualitatively opposed to the other; the south pole has its
determinate other in the north pole, and just for that reason the attraction is mutual; the south pole, on the other
hand, has only a formal other in another south pole;
% MISPRINT: the print has „Eensformighedon“ for „Eensformigheden“ (o for e) --- a compositor's
% fault shown by both copies; the errata leaf does not correct it. Rendered here by the sense.
uniformity is at variance with the principle of articulation, and the repulsion is mutual. How originally the opposed
members presuppose one another is seen in the fact that one can indeed procure magnets with many poles, but not a magnet
that has only one pole. The polar opposition is grounded in qualified unity; the unity is a unity of opposition expressly
determined by articulation.

Seen immediately, it looks as though the magnetic force sat in each of the magnet's poles, while in the so-called
% „Ligevægtsrum“ is letterspaced; checked against the KB copy.
\emph{space of equilibrium} there were no such force; but already Gilbert has shown that the articulation goes all the
way through. By breaking a bar magnet across the middle one does not get two half magnets, each with its pole, but two
whole ones, each with two poles. The part is a complete whole, and the division can be continued down to the smallest
particulars: the articulated whole consists of nothing but articulated parts (small magnets, elementary magnets).

The thoroughgoing articulation opens us an insight into the connexion of the whole; for the reason why the operative
poles form themselves toward the end-faces lies in this, that the small magnets in the space of equilibrium equalise and
cancel one another's effects, whereby there arises a peculiar opposition, conditioned by the circumstances, between bound
and free
% printed mark: *) --- it stands AFTER the full stop.
magnetism.\footnote{%
``If we conceive that in a steel bar all the small parts had become equally strongly magnetic under the influence of the
same magnetising force, with their poles turned alike and in the bar's lengthwise direction, then the poles of the free
magnetism would lie in the bar's end-faces, since the everywhere equally great magnetic forces within these would have to
cancel one another's outward effects entirely. This magnetic state cannot, however, subsist, on account of the elementary
magnets' distributing effects upon one another, whereby there arise alterations in the inner magnetism, which consist
partly in a weakening, and this must then be greatest in the bar's longitudinal axis, partly in a strengthening, which
becomes greatest at the middle of the bar.''
% The print writes „Anf. Skrift.“ here but „Anfr. Skr.“ in the note to p. 280; the two forms
% are reproduced as printed, rendered as „Work cited“ and „Wk. cit.“ respectively.
Lorenz. Work cited. pp.\ 144--45.}

% --- p. 279 ---
If one will now say that the magnetic functions do indeed presuppose an articulation but are not themselves articulating,
then we must for the present content ourselves with referring to what physics calls ``the distributing force''. A piece of
iron becomes, by contact with a magnetic pole, itself a magnet that can carry small pieces of iron. The distributing force
is doubtless as yet only a superficial representation, in so far as it immediately presupposes what it really ought to
explain, namely the polarity of the elementary magnets; none the less it does cast a peculiar light upon its own
presupposition and discloses a secret of cohesion.

For here there shows itself a peculiar relation between the mobility of the small parts and the fixity of the mass, a
relation more closely determined by the
% „Coercitivkraften“ is letterspaced; checked against the KB copy.
\emph{coercive force}. Steel is magnetised with more difficulty than soft iron, but then holds the faster to its
magnetism; the reason is that the coercive force is greater in steel than in iron. Changes of temperature, shocks,
vibrations and the like weaken the coercive force, in that they alter the cohesion.

The distribution that a magnet produces, especially in soft iron, always works back strengtheningly upon the magnet's
inner magnetism, but now strengtheningly, now weakeningly upon its free magnetism --- again a modification of articulating
activity.

Magnetism's mechanical effects recall in certain respects those of gravity. Two magnetic poles, in which the quantity of
free magnetism is m and
% Raised figure one, not a prime: „m og m$^1$“, read on the KB copy.
$m^{1}$, will at a distance r act upon one another with a force K that can be determined by the equation
$K = \dfrac{mm^{1}}{r^{2}}$; and yet, as indicated above, the magnetic attraction is qualitatively different from that of
gravity; a steel needle, for example, does not become heavier because it is magnetised. If the influence of the repelling
pole be reckoned in, then the whole effect of a magnetic bar will, on
% --- p. 280 ---
% printed mark: *) --- it stands AFTER the full stop.
particularly given conditions, prove to stand in inverse ratio to the third power of the distance.\footnote{%
Hansteen. Untersuchungen über den Magnetismus der Erde. Christiania 1819. --- On Gauss's formula
$K = \dfrac{mm^{1}}{r^{p}}$ cf. Lorenz. Wk.\ cit. pp.\ 152--56. A.\ Wüllner. Lehrbuch der Experimentalphysik. 2ter Band.
Leipzig 1866. Wirkung magnetischer Massen auf einander aus der Ferne. pp.\ 540--58; Versuche von Gauss. pp.\ 558--64.}

In that the magnetic needle points, as it is said, toward the north, it points at the same time to a magnetism that does
not let itself be explained by immediately existing small magnets, namely to terrestrial magnetism. It is a leap from the
particular over into the total. Totality-magnetism points from many sides to articulating activity of diverse kinds, a
widely ramified system of articulating functions. The old notion of magnetic mountains in the neighbourhood of the earth's
poles has long since been given up; the earth is magnetic because it is unceasingly being magnetised. The periodic
changes, both the daily and the secular, to which the needle is subject, partly in respect of declination, partly of
inclination, are, like the changes in the strength (intensity) of terrestrial magnetism, so many testimonies that what is
constant in the magnetising activity is united in a peculiar way with what is variable. What is immediate in the phenomena
has, in the essence itself, a thoroughgoing reflexion of power for its presupposition. Provisionally it looks as though
the whole system of magnetic forces turned within the earth from east to west.

A closer elucidation of the possible ground of the magnetic phenomena can, in relation to the present standpoint of
physics, first be given when magnetism is brought into
% --- p. 281 ---
connexion with another and still more deeply intervening system of articulating functions.

% Run-in head, read on the page itself (KB copy): set in bold roman, not letterspaced.
\runin{b) Electricity.} The analogy between magnetism and electricity is striking on account of the polar opposition; but
as regards the inwardly operative, the likeness must rather seem to belong to the phenomenal than to the essence. For
magnetism is, as ``force'' and ``property'', bound to the magnetic body, whereas electricity ``can pass freely over from
the one body to the other'', and in this passage shows itself analogous to heat. As regards conduction there is so
remarkable a connexion between electricity and heat, that for the metals as well as for their alloys the same ratios have
been found between their conductive power for heat as between their electric conductive power. If heat had poles, it
would lie near at hand to conceive electricity as a polarised heat; but since no fact points to there being a positive
heat-pole that attracted a negative one and repelled another positive one, since by bound and liberated heat something
quite other is understood than by bound and free electricity, since heat, in other words, is only quantitatively, not
qualitatively different in itself, it is seen from this that there must be a fundamental difference between electricity
and heat.

It is characteristic to see how heat and electricity, however inwardly they may be united, yet in an essential sense each
play their own part in the world of bodies. Out from its invisible, universal ground heat moves through the manifoldly
different bodies, dissolves the differences, so far as it is able, into the uniform, effaces the bodily individualities
and sets its highest in transforming the whole into vapour. In opposition to this, electricity must be said to be laid
out precisely upon asserting the differences by force of something common, upon articulating what is proper
% --- p. 282 ---
and making the bodily individualities good. Heat can indeed be converted into electricity, and electricity into heat; but
this by no means signifies that electricity and heat are one and the same; it is, on the contrary, an intimation that the
difference is preserved in the unity. To the development of the unity of difference there corresponds the activity of
form. When the common is posited, form demands the particular; and when the particular is posited, it demands the common
again. Thus heat and electricity work, each in its own way, in form's service. And yet it is only the extremes of form's
movement, not the substantial activity of form itself, that we as yet have before us. Electricity has indeed for itself a
universally pervading essentiality after the manner of heat; for that glass-electricity is called positive and
resin-electricity negative by no means says that the electricities in the one body should be of another nature than in
the other; on the contrary! the same two different kinds of electricity are found in every body, save that in the body's
so-called natural state they cancel one another: the negative electricity is in all bodies like itself as negative, the
positive as positive. The preluding activity of form shows itself in this, that the difference between the universal and
the individual is immediate just as much as the unity is. From the standpoint of the universal it must indeed look as
though electricity were a certain common something that ran through the bodies, spread itself upon the surfaces according
to their different form, passed from the one body over to the other, and so forth; but if we look more closely, all this
points to activities and changes in the bodies concerned themselves. Electricity does not reduce bodies to mere members of
transition, in order everywhere to set itself in peculiar relations to itself alone; it is just as much the bodies that,
by means of electricity, set themselves in a peculiar relation to one another.

How deeply electricity is implicated in the peculiarity of the different
% --- p. 283 ---
bodies can already be seen from the manner in which it is awakened. It is awakened by friction; but friction is precisely
what makes the bodily differences recognisable; it is a significant expression of the difference between the glass and
the resin, that the two bodies, on being rubbed with the same stuff, wool or silk, become, the first positive, the second
negative. It is awakened by distribution.
% The sentence lacks a comma after „Elektricitet“ and after the parenthesis
% „(Influenselektricitet af første Art)“; both copies show this. Rendered as printed, the
% English omitting the commas at the same two places.
If a body A be charged for example with positive electricity it will, across an insulating interval, act upon another body
B in such a way that by distribution there arises in B negative electricity on those parts of the surface that are nearest
to A (influence-electricity of the first kind) but in the surface furthest from A, on the contrary, positive electricity
(influence-electricity of the second kind). To this there correspond again the stronger or weaker attraction, for since
the distribution is greatest in the good conductors, the attraction to these must also be strongest. The peculiar
difference, determined by electricity, between good, half-good and bad conductors,
% „Isolatorer“ is letterspaced; checked against the KB copy.
\emph{insulators}, which is so significant for articulation, heat will not merely alter but at higher degrees even
efface: one knows no body that at white heat remains an insulator. But when heat works not inhibitingly but
conditioningly, the articulation is so pronounced that even the mere contact of two different bodies is enough to awaken
the electricity. How qualitatively peculiar the opposition is can be seen from a series such as: zinc, lead, iron,
copper, silver, gold, platinum, in which every preceding member becomes positively electric by contact with one that
follows. The electric difference between zinc and copper and the electric difference between copper and platinum are
together equal to the electric difference between zinc and platinum; or: when one lays a copper plate upon a zinc plate
and upon this a platinum plate, the electric tension in the end plates is just as great as if one had laid the platinum
plate and the zinc plate immediately upon one another.
% --- p. 284 ---
% The whole of p. 284 stands in ordinary roman. The JBIG2 working scan makes the
% type look bold and the Hegel citation look letterspaced; the KB copy shows light,
% ordinary type throughout --- no \emph.
The peculiarity, tension and prickliness that bodies display in the polar opposition has led Hegel to inculcate with real
emphasis what idealism, even in the very newest time, is still so glad to persuade itself of, namely that the elementary
bodies are at bottom self-beings. ``Reflexion,'' he says, ``is accustomed to conceive the body-individual as something
dead, which only comes into outward mechanical contact or enters into a chemical relation. --- We, on the contrary,
conceive the electrical tension as the body's own selfishness, which is physical totality and asserts itself in contact
with another. It is its own wrath, the body's own effervescence, that we see; there is no one present but itself, least
of all a foreign matter. Its youthful courage breaks out, it rears up on its hind legs; its physical nature gathers
itself against entering into dealings with another, and that as an abstract ideality of light. Not only do we compare
bodies, but they compare themselves and therein assert themselves as physical; it is a beginning of the organic,'' and so
forth. --- It is the eternal mania of the philosophy of immanence: if one cannot make Spirit into matter, then one must
absolutely be able to make matter into Spirit; if one cannot explain the conscious will from mechanical attraction, then
one must explain mechanical attraction from an unconscious will; if one cannot derive the self from the selfless, then
there is nothing else for it but to make the elementary bodies, the selfless individuals, into self-beings. The attempt
is equally ingenious, whether it be ``the physical totalities'' or the naked atoms, the pure ``force-points'', that are
made into self-beings.

But if there is to be earnest with selfhood, then there must also be earnest with otherness; if Spirit is to be Spirit,
then one had better let matter be matter and take the physical bodies for what in fact they are, lifeless, selfless
realities. The peculiar union of the universally bodily with what is
% --- p. 285 ---
% The Danish breaks mid-word here: p. 284 ends „Sær-“, p. 285 begins „egne“.
proper to bodies, which electricity displays, gives clearly to be known that both the bodies themselves and the
fundamental parts of which they consist must have a selfless independence. For it is precisely the mark of the selfless
that all utterances of a given body's peculiarity can be exactly determined and measured by influences proceeding from
other bodies, so that action and reaction, determined by an exact relation between inertia and force, correspond to one
another down to the very smallest particulars. The mechanical laws are in the strictest sense laws of otherness; but the
various exact measurements of electricity's distributing force, its density, tension, current-strength and so forth put
it beyond all doubt that the effects that are electric in the special sense are in every respect subject to the laws of
otherness, just as much as the
% printed mark: *) --- it stands BEFORE the full stop, as printed.
mechanical\footnote{%
In t seconds the quantity of electricity that has passed through the conductor at the current-strength s will be
$e = st$; that is a law of otherness; Ohm's law is a law of otherness, and so on.}.

By discharging a Leyden jar through the wire windings of a multiplier one gets an electric current, which however ceases
instantly; by the voltaic pile (and the galvanic batteries) one gets a perennial current. Volta compared his pile with an
extraordinarily large but weakly charged battery of Leyden jars that charged itself again instantly when it had been
discharged: the electricity in the Leyden jar and the electricity in the pile is therefore essentially the same
% printed mark: **) --- it stands BEFORE the full stop, as printed.
% The note runs over the leaf: it begins at the foot of p. 285 and is continued at the
% foot of p. 286. By the house rule it stands whole at its mark.
electricity. None the less the mode of operation is different%
\renewcommand{\thefootnote}{**)}%
\footnote{%
``When a Leyden jar is discharged, it is not only everywhere the same quantity of electricity that passes through the
conductor; the quantity of electricity is also the same whether the conducting wire be
% The row of dots stands as four close points, as printed --- unlike the spaced points
% in the note to p. 276.
short or long .... It stands otherwise with the electricity that a galvanic element produces; as a rule one needs only
to close it with a tangent galvanometer in order to observe the effect upon the magnetic needle of the far more
considerable quantity of electricity that such an element sends through the conducting wire each second; but if one
inserts a long metal wire into the circuit, the galvanometer shows that the current-strength has thereby been diminished
to a noticeable degree.'' L. Lorenz. Lectures. pp.\ 213--14.}%
\renewcommand{\thefootnote}{*)}%
.

% --- p. 286 ---
The different modes of operation do not lead straight off to an understanding of the inner how, but to the discovery of
different laws and thereby to modes of representation which, with the laws' help, must be brought into mutual agreement,
reduced to fixed universal presuppositions and by degrees made clear in such a unity of connexion that they yield real
moments for an ever richer development of electricity's peculiarly articulated concept of nature: reduction and deduction
condition one another mutually.

Of the current and of what flows one can doubtless form various representations; for although it is given that in the
% The sign before E is Nielsen's MINUS (÷), read on the KB copy; cf. the house rule from
% batch 4 (p. 46).
conducting wire there is $+$ E at the pile's positive and $\div$ E at its negative pole; although it may be assumed that
the current must be a double current, namely a motion of the opposite electricities in opposite directions: still the
% MISPRINT: after „Hvorledes“ the print sets a full stop where the sense requires a comma.
% Both copies show the full stop, so it is the compositor's fault; the errata leaf does not
% correct it. Reproduced here, the English taking the same full stop for a comma.
inner how. the unceasing conversion of positive into negative and of negative into positive, of repulsion into attraction
and attraction into repulsion throughout the whole motion, is thereby in no way made
% printed mark: *) --- it stands BEFORE the full stop, as printed.
% The note runs over the leaf: it begins at the foot of p. 286 and is continued at the foot
% of p. 287. By the house rule it stands whole at its mark, and the page mark for p. 287
% therefore stands after the whole note, where the text of p. 287 begins.
clear\footnote{%
How Weber has sought to explain the electrodynamic forces by starting from the mutual repulsions and attractions of
electricities at rest is an intricate investigation belonging to the special science. Only one point shall be
particularly brought out. Against Helmholtz's objection that in Weber's formula there is a contradiction between the law
of electric force and the law of the conservation of force, Weber has urged in his reply that the more closely determined
principle of the conservation of energy finds application to two elements only when the potential of both is of like
form. For two electric elements at the distance r, with a constant c, the potential is:
% The formulae are read on the KB copy. The raised figure one in $ee^1$ and $mm^1$ is a
% figure one, not a prime. The minus within the parenthesis is an ordinary minus sign, not
% Nielsen's ÷.
\[
V = \frac{ee^{1}}{r}\left(\frac{1}{cc}\cdot\frac{dr^{2}}{dt^{2}} - 1\right);
\]
for two ponderable small parts it is, on the other hand:
\[
V = \frac{mm^{1}}{r}.
\]
If one now disregards the sign, one needs to ascribe to c not an infinitely great but only a very great value in order to
obtain the same potential for the ponderable as for the imponderable elements. (For if the law of attraction be
designated by $f(r) = \dfrac{\mu}{r^{2}}$, then, with respect to the
% MISPRINT: „Polentialet“ for „Potentialet“ --- a wrong sort (l for t). The errata leaf
% corrects it: p. 287, line 15 from the foot: „Polentialet --- read Potentialet“. The KB
% copy prints „Polentialet“; in the working scan the pencil-bearing hand has drawn a
% cross-stroke through the l and made it a t. Rendered here by the sense.
potential, $\int f(r)dr = \varphi(r) = - \dfrac{\mu}{r}$; cf. C. Ramus. Analytic Mechanics. pp.\ 53--54; pp.\ 59--60).
The point is that Newton's law of gravitation, so far as it is confirmed by facts, should be able to come under Weber's
law. Cf. Zöllner, Ueber die Natur der Cometen. Vorrede. p.\ LXIV ff.; the piece: Ueber die allgemeinen Eigenschaften der
Materie. pp.\ 313--41.

If a law for electricity has now really been found, or can be found, that takes Newton's law up into itself, then there
will doubtless also be opened to the philosophy of nature a deeper insight into the inner connexion of the laws and
thereby of the real concepts; but the fundamental relation between nature and
% The sign after „Gjenstandsmagt“ is read as a comma: the sense requires it (three members
% in the enumeration), tesseract reads a comma, and the space after it is an ordinary word
% space, not a sentence space. The rejected reading is a full stop.
Spirit, between object-power, real concept and comprehending self-power, will be and will remain unchanged the same.}.
The only decisive yield for the philosophy of
% --- p. 287 ---
nature is: that matter shows itself articulated down to the smallest parts; that the thoroughgoing articulation rests
upon a punctually determined co-operation, fixed by the real concept, of the bodily individual with something universally
bodily; that it is this co-operation which, in the form of articulation, secures to the elementary bodies, the selfless
beings, an independence that in all relations, qualitative as well as quantitative, is unindependent.

% Run-in head, read on the page itself (KB copy): set in bold roman, not letterspaced.
% It falls on p. 287, not p. 288.
\runin{c) Electromagnetism.} The relation between the two kinds of articulating functions, the magnetic and the
% --- p. 288 ---
% The leaf opens flush (without indentation) and mid-word: the preceding leaf ends on
% „elek-“, and the sentence from p. 287 is continued here. No new paragraph. The head
% „c) Elektromagnetismen“ therefore stands NOT on this leaf but on p. 287.
% The working copy has a vertical pencil stroke in the left margin down the whole leaf;
% it is not text.
electric, is again determined by an articulation. It is the first impulse toward a further determination of this
articulating reciprocal relation that gives H. C. Ørsted's discovery its world-historical significance.

``The likeness that in several respects obtains between electricity and magnetism had already long led physicists to
surmise that these two forces must stand in a closer relation to one another. Yet all attempts to demonstrate such a
relation were in vain, until Ørsted in 1820 showed that an electric current in the neighbourhood of a magnetic needle
brings the latter out of its position. For if one leads an electric current from south to north beneath a magnetic
needle, its north end will make a deflection toward the east, and after some oscillations it will come to rest in a
position that approaches the more nearly to being perpendicular to the magnetic meridian, and therefore to the current's
direction, the stronger the current is. If the magnetic needle be brought beneath the current, the north end will make a
deflection toward the west. By a counterweight one can bring a dipping needle to lie horizontal; if such a one be placed
on the eastern side of the current, the north end will give a deflection downward; if it be placed on the western side,
the north end will make a deflection upward. If the current's direction be reversed, the deflections go in the opposite
% The footnote mark stands AFTER the closing quotation mark: „Nordenden.“*)
direction, and the south end is always driven in the direction opposite to the north end.''%
\footnote{% printed mark: *) --- the only note on p. 288.
C. Holten. The Doctrine of Nature's Universal Laws. Copenhagen 1870. pp.\ 318--19.}

It now became, as discovery followed upon discovery, gradually clear that ``the electromagnetic force'' was in its
phenomenon a single member belonging to a widely embracing chain. Thus Ampère showed that in the relation between two
current-conductors there was ``an electrodynamic force'' which uttered itself now in a repulsion, now in an attraction,
from which it followed that two parallel current-conductors must attract one another when the currents go in the same
direction, but repel
% --- p. 289 ---
one another when they go in opposite directions. He went further and demonstrated that two current-conducting spirals
attract and repel one another like two magnets. He went further still, in seeking to prove that magnetism was really a
peculiar phenomenon of electricity.

A screw-thread of copper through which an electric current is led acts in altogether the same way as a magnet of equal
size. One can now regard a bar magnet as a bundle of linear bar magnets, which however, on account of their mutual
influences, are not quite rectilinear but somewhat bowed toward the magnet's axis; these linear magnets are then again to
be regarded as a series of elementary magnets of unequal strength. According to Ampère's explanation every elementary
magnet will then behave as a circular electric current whose plane is perpendicular to the elementary magnets' axis, and
the linear magnet will represent a series of such currents. Such a system of current-conductors Ampère called a solenoid,
% LETTERSPACING, checked against the KB copy: it begins at „enhver“ (not at „at“) and runs
% unbroken to and including „Solenoider.“, across the word division „Sole-/noider“ in the
% setting. It is the only letterspacing on pp. 288--291.
and the explanation is: that \emph{every magnet can be regarded as a system of solenoids.}

If it could now merely be assumed that electric currents could maintain themselves, and that even in good conductors,
without either giving off or receiving any work, then magnetism would at the same time be deduced from electricity; but
the assumption is unreasonable. The investigation must then be carried on from another starting point.

By the help of strong electromagnets Faraday has instituted experiments upon the magnetic behaviour of different bodies
and found that all bodies are in greater or lesser degree acted upon by magnetism, yet with this difference, that some
are attracted, others --- ``the diamagnetic'' --- are on the contrary repelled by a strong magnetic pole. ``The magnetic
distribution in the diamagnetic bodies grows, at weaker magnetising forces, proportionally with these; but if these
exceed a certain
% --- p. 290 ---
limit, the distribution approaches a certain maximum. This is therefore the same as takes place with the magnetic bodies;
but these latter reach this limit more quickly, and it may therefore happen that bodies which contain a mixture of
magnetic and diamagnetic substances, for example charcoal or ferruginous bismuth, show themselves magnetic at weaker
magnetising forces but diamagnetic at stronger ones, since the effect of the diamagnetic substance now gets the
preponderance. Among the most strongly diamagnetic bodies belong bismuth, phosphorus and antimony. Water and quicksilver
show themselves diamagnetic, whereas solutions of salts containing magnetic metals are as a rule magnetic. Oxygen shows
itself to be magnetic; most other gases are on the contrary diamagnetic. All bodies are, according to Becquerel's
assumption, more or less magnetic, and if a body shows itself diamagnetic, that is only because it is surrounded by a
more magnetic medium. If one thus encloses a dilute solution of ferrous sulphate in a glass tube, then this tube is
magnetic in air, but diamagnetic when it is held immersed in a concentrated solution of the same salt, because this
solution is more strongly magnetic. --- Since, however, a body in a vacuum can also behave
% SETTING FAULT, checked against the KB copy: the second „l“ in „tillige“ stands raised
% above the line of type (the sort sits too high in the setting). Both copies show it, so it
% is the compositor's fault and not the copy's. Rendered as ordinary text; the fault cannot
% be set.
diamagnetically, one is accordingly compelled also to assume that a magnetic distribution can take place in a
% The footnote mark stands AFTER the closing quotation mark and with a space: „Rum.“ *)
vacuum.'' %
\footnote{% printed mark: *) --- the only note on p. 290.
% MISPRINT in the page number, checked against the KB copy: between „2“ and „7“ stands a
% space of one sort's width with a small dot on the line --- that is, a sort that has not
% come to print. Both copies show the same, so the fault is the compositor's. The sense is
% plainly „277--79“ (cf. the note to p. 292: „Lorenz. Anf. Skr. S. 270.“), but the text is
% reproduced AS PRINTED.
% NB. Here „Lorenz“ stands WITHOUT a full stop after the name; p. 292 has „Lorenz.“ with
% one. Both are reproduced as printed.
Lorenz Work cited. pp.\ 2.7--79.}
This last circumstance sets a question mark to the explanation instead of a full stop. If Ampère's theory be now set
beside the great result of Faraday's investigations, then it is seen that magnetism has ceased to be anything particular:
the magnetic articulation is in its kind just as thoroughgoing as the electric.

% --- p. 291 ---
The ground both of the magnetic and of the electric polarity is essentially to be sought in the original relation between
what is real in object-power and what is ideal in self-power. That the articulation cannot be explained from object-power,
that is, from the material existence of things alone, is easy to see. If we begin with the elementary magnets, then every
such magnetic fundamental part must be divided into three parts, of which the one is a north pole, the second a south
pole and the third a neutral middle; but if we look more closely, each of these parts falls again into three parts. The
part that forms the north pole is, when it is separated from the two others, itself the whole elementary magnet with
north pole, south pole and neutral middle, and so on. The contradiction is that every part of the magnetic whole is to be
a part existing in space, that every such part can again be divided, and that the divided part is itself a whole. The
same contradiction repeats itself with electricity. In its natural state every part of the body is to have both
electricities in indifferent unity. When the two electricities are one, then the body is not electric, which is to say:
then electricity is not; when electricity is, then the neutral unity is not. But when actual electricity is only in the
twoness, how can the possible be in the unity? Electricity's differenceless unity is precisely its impossibility! These
contradictions can be solved only on the presupposition that the twoness's unity is in the real concept, and hence in
power's ideality, and that the unity's twoness, precisely on that ground, comes to expression under given conditions in
material existence, and hence in power's reality. This does not indeed make the mode of operation itself clear, but it
casts light on the essence of what is operative.

By referring what is in the phenomenon to what is operative in the essence, Ampère's theory is well fitted to illuminate
the interior of the articulated totality and to lead to an explanation of terrestrial magnetism. ``As is well known, a
magnetic needle is acted upon by the earth
% --- p. 292 ---
as by a distant magnet, and quite in the same way the earth behaves also toward electric current-conductors. Thus a
movable plane current-circuit sets itself perpendicular to the magnetic meridian, in that it behaves, with respect to the
earth as well, as a magnetic plate whose north-magnetic side will turn toward the north; in the lowest part of the
conductor the current will accordingly go from east to west. In the same way the current-conducting movable spiral will
set itself with its axis in the magnetic meridian, and one can even get a movable part of a current-conductor to rotate
by terrestrial magnetism alone. All these effects we can now ascribe with far greater probability to electric currents in
the earth's interior than to really present magnetism. One has in fact already observed, in the part of the earth's
surface accessible to us, now weak, now stronger electric currents, and it is therefore probable that further down in the
earth's interior too there are electric currents, whose combined effects are what we observe on the
% The footnote mark stands AFTER the closing quotation mark and with a space: „Overfladen.“ *)
surface.'' %
\footnote{% printed mark: *) --- the only note on p. 292.
Lorenz. Work cited. p.\ 270.}

% THE SIGNS: the sentence sets „+ E og --- E“. The minus is here an ordinary horizontal
% stroke of the em-dash's length, NOT Nielsen's division-minus „÷“ (cf. pp. 46 and 233).
% Checked against the KB copy. Reproduced as printed, that is, with ---.
Here, in these facts, is a peculiar interplay of forces, a remarkable alternation of the manifold and yet uniform. North
pole and south pole, + E and --- E in endless repetitions: it looks uniform. But what a manifoldness of qualifying
activities and qualitatively determined differences comes to light upon this comparatively uniform ground! They are
naturally necessary utterances of punctually determined presuppositions, regardless outbreaks of inner regard; the real
ground is without regard, for the rational is nothing but regard.

% LETTERSPACING, checked against the KB copy: „telluriske“ (p. 292) and „kosmiske“ (p. 293)
% are letterspaced; the antithesis runs across the turn of the leaf.
That the functions of magnetic and electric polarity, with their mutual polar crossings, are not merely \emph{telluric},
holding only for the terrestrial whole, but
% --- p. 293 ---
also \emph{cosmic}, holding for the whole solar system; that one can speak of a magnetic storm just as well as of an
electric one; that the northern
% THREE NOTES on this leaf, printed *), **) and ***); the last two are set with local
% \renewcommand, restored immediately after. The first two marks stand BEFORE the comma:
% „Nordlyset*),“ and „Zodiakallyset**),“; the third stands within the sentence with no
% preceding punctuation: „Cometerne***) ere Phænomener“.
lights%
\footnote{% printed mark: *) --- first of three notes on p. 293.
% The note ends WITHOUT a full stop after the page number: „Love. S. 285“. Checked against
% the KB copy; both copies show the same, so it is the compositor's omission. Reproduced
% as printed.
That there must be a peculiar relation between magnetism and the northern lights is already indicated by this, that the
luminous arc which it most often forms stands perpendicular to the magnetic meridian, and that the northern lights act
upon magnetic needles even where they cannot be seen. Strong oscillations of the magnetic needle have been observed in
Paris at the same time as the northern lights were seen in Christiania or in northern Scotland. C. Holten. Nature's
Universal Laws. p. 285}, the zodiacal light%
\renewcommand{\thefootnote}{**)}%
\footnote{% printed mark: **) --- second of three notes on p. 293.
% MISPRINT, checked against the KB copy: the print has „Phoshorescensen“ (without the
% second p). Both copies show the same, so it is the compositor's fault. Rendered here by
% the sense.
On ``Corona, Zodiacal Light, Northern Lights and the Phosphorescence of the dark ground of the heavens in its relation to
the Moon'' see Zöllner, Work cited. pp.\ 127--30.}%
\renewcommand{\thefootnote}{***)}%
, indeed even the relation of attraction between the sun and the comets%
\footnote{% printed mark: ***) --- third of three notes on p. 293.
% „Anfr. Skr.“ with an inserted r. FORMERLY ALLEGED A MISPRINT; THE CLAIM IS WITHDRAWN.
% The book uses both forms indifferently --- „Anfr.“ occurs ten times, „Anf.“ seven --- so
% the form is no compositor's fault but one of the book's two spellings. The English keeps
% the distinction: „Anf. Skr.“ = Work cited, „Anfr. Skr.“ = Wk. cit.
On Olbers as founder of the electric theory of comets, cf. Zöllner p.\ 146 ff. On ``the difference between gravitating
and electric action at a distance upon ponderable masses'' Zöllner lays down the following: ``If a body stands at the
same time under the influence of gravitation and of another's free electricity, then with increasing mass gravitation,
but with decreasing mass electricity, will have the preponderance as moving force. For this reason the cometary nuclei,
as liquid masses, stand under the influence of gravitation, the developed vapours, as aggregates of very small parts of
mass, under the influence of the sun's free electricity.'' Wk.\ cit. pp.\ 117--120.}%
\renewcommand{\thefootnote}{*)}
are phenomena of electric or electromagnetic origin --- all this may well in general be assumed, although physics is not
yet in a position to demonstrate the causal connexions in detail. What is decisive for the real concept is that the
thoroughgoing articulation yields an essential presupposition for all forming, body-shaping activity of nature.

% CLOSE OF THE DIVISION. The text on p. 293 breaks off about a third of the way down the
% leaf; below it stands blank space and then a short centred rule (single-printed), about
% the middle of the leaf --- as at p. 146. Below that follow the note rule and the three
% notes, and NOTHING more. Division B therefore ends clean at the foot of this leaf: the
% head „C. Legemdannende Functioner: Chemisme.“ stands NOT here but begins on p. 294.
\streg

% Danish: \afdeling{C.}{Legemdannende Functioner: Chemisme.}   (printed pp. 294--317)
\afdeling{C.}{Body-Forming Functions: Chemism.}

In all the utterances of articulation and polar opposition hitherto considered, the bodily members hold themselves
immediately, and so far still mechanically, outside one another: tension is converted into living force, and living force
into tension, heat into electricity, and electricity into work; but under all this the bodies concerned nevertheless
assert their relative independence. Only in chemism does the articulating activity become so thoroughgoing that the
members surrender their independence in their affinity: bodily existence is dissolved in order to be brought forth, and
brought forth in order to be dissolved again.
% Letterspaced on the KB copy: „Elektrolysen“, „det elektrochemiske System“ and „elementære
% Legemdannelser“ (the last broken across the line-end „Legem-“ / „dannelser“, letterspaced
% in both halves). „et stort chemisk Laboratorium“ stands in ordinary roman and is NOT
% letterspaced, though the working scan makes it look bold.
Phenomena of the polar functions' dissolving activity are seen in \emph{electrolysis}; the reciprocal determination of
dissolution and production, bound to peculiar relations of measure, upon which nature's body-forming capacity essentially
rests, is made intuitable in \emph{the electrochemical system}; the result is that nature, fruitful in real productions
of form, is at bottom a great chemical laboratory, a self-working workshop for \emph{elementary formations of bodies}.

% Run-in head, printed in bold on the same line as the paragraph's first word. Checked
% against the printed leaf, not taken from the table of contents.
\runin{a) Electrolysis.} In and by articulation the natural cause has acquired a quite peculiar character, different from
the mechanical and the merely dynamic. But articulation itself takes on in turn a differing stamp according as it is a
magnetic, an electric or a chemical activity in which it gives itself to be known. The connexion between magnetism and
electricity is so intimate that a physicist could be tempted to take magnetic polarity for a peculiarly conditioned
expression of the electric, although the facts, hitherto at least, speak for a peculiar mutually corresponding
co-operation in which the magnetic, so far as the comprehensive whole of nature (terrestrial magnetism) is concerned, is
dependent
% --- p. 295 ---
upon the electric. The investigation now goes further, and the question becomes:
% The whole question is letterspaced on the KB copy, from „hvorledes“ to „Chemismen?“,
% across three lines.
\emph{how does galvanism, particularly as streaming electricity, stand to chemism?}

From the ideal point of view --- and every natural cause's reality is of course determined from within by its real
concept, its ideality --- different conceptions are possible. It is possible that electricity is one thing and chemical
affinity another, as gravity is one thing and cohesion another. Affinity is indeed conditioned by a qualitative
opposition after the manner of electricity; but the unlike electricities are found in every substance, whereas affinity's
unlike members are separately existing substances. There is a free and a bound electricity, but there is not in the same
sense a free and a bound affinity; one can speak of an electric current, but not of a streaming affinity. Electric
attraction can take place at a distance, chemical attraction only by immediate contact; under the equalisation of
electric tension the bodies concerned hold themselves as existing members, under the equalisation of the chemical
opposition the bodily members become moments for the formation of a new body. It is therefore possible that the chemical
real difference has, in the system of power-concepts, another principle than the electric. But it is also possible that
it is the same principle, only at a different stage of development, so that precisely what is preformed in
electromagnetism becomes realised in chemism and receives its form of limitation in crystallisation. How the matter
stands herewith cannot be settled by a priori speculation, but only by the categorical dictate of the facts.

What, in a voltaic pile, is the originally operative cause? Is it the electric current that is the cause of the chemical
activity; or is it conversely the chemical activity that is the cause of the current; or is there here
% --- p. 296 ---
perhaps a peculiar and equal co-operation of both factors? According to Volta's assumption the mutual contact of the
different metals is the sole source of the pile's electricity; Volta overlooked the chemical effects and regarded the
contact. Since, however, the pile's chemical energy had shown itself in striking results, Faraday went to the opposite
extreme and declared the chemical action that the fluids exercised upon the metals to be the only real cause of the
electric current. Fechner has established the correctness of Volta's fundamental experiment, but Faraday has shown that
galvanic currents can arise even without contact between different metals. Faraday's theory has brought with it a new
terminology. The conductors at whose surfaces the chemical separation takes place are called
% Five letterspaced words in succession, all checked against the KB copy: Elektroder,
% Anoden, Kathoden, Elektrolyt, Elektrolyse. „Kathoden“ was not caught by spacing.py.
\emph{electrodes}; the one that is in connexion with the chain's positive pole is the \emph{anode}, the way upward; the
one that is in connexion with the negative pole is the \emph{cathode}, the way downward; the body that lets itself be
chemically separated by the current is an \emph{electrolyte}, and the process itself an \emph{electrolysis}.

As water is, so are all compound substances separated, when they --- be it noted --- conduct electricity and are at least
in part in a molten or dissolved state. As regards what is immediately observable, there is no lack of exact
determinations.%
\footnote{% printed mark: *) --- first of two notes on p. 296. The mark stands AFTER the
% full stop: „Bestemmelser.*)“.
% „Anfr. Skr.“ with an inserted r stands thus here too, in both copies. The book uses both
% forms, „Anf. Skr.“ and „Anfr. Skr.“ (five places each); it is therefore a wavering in the
% abbreviation and no misprint. English keeps the distinction as elsewhere.
Lorenz. Wk.\ cit. pp.\ 223--234.}
The electrolytic law says that the weights of the substances separated by the same quantity of electricity stand to one
another as those substances' equivalent numbers.%
\renewcommand{\thefootnote}{**)}%
\footnote{% printed mark: **) --- second of two notes on p. 296. The mark stands AFTER the
% full stop: „Æqvivalenttal.**)“.
% THE NOTE RUNS OVER THE LEAF: it begins at the foot of p. 296 and ends at the foot of
% p. 297; the mark is not repeated on the continuation. By the house rule the whole note is
% set here at its mark. „Naturlærens chemiske Deel“ is letterspaced on the KB copy (a book
% title). The figures 32,3 --- 12,48 --- 403,32 read on the KB copy; tesseract read 403,82.
For every part by weight of hydrogen set free between the pole plates there is dissolved in each ``cell'' 32.3 parts by
weight of zinc. But now the weights of the chemical equivalents of hydrogen and zinc stand to one another as 12.48 to
403.32, or as 1 to 32.3. For every equivalent of hydrogen developed in the cell in which the electrolysis takes place,
there must accordingly be dissolved in each of the chain's cells 1 equivalent of zinc. \emph{The Chemical Part of the
Science of Nature}, translated by C. L. Petersen. Copenhagen 1856. p.\ 81.}%
\renewcommand{\thefootnote}{*)}
% --- p. 297 ---
% The leaf begins WITHIN the paragraph from p. 296: the first line stands without
% indentation.
How far can these investigations now open us an insight into the fundamental relation between galvanism and chemism?

It goes without saying that physics must from every side seek to give an account of bodies' molecular states and the
possible changes these undergo. The separation of water, whose constituents can each be collected and measured in a
voltameter, Grotthuss has explained in the following ingenious way. When hydrogen is combined with oxygen into water, the
oxygen atoms become negatively and the hydrogen atoms positively electric through the intimate contact between the
smallest parts; but on account of the uniform distribution of the small parts of oxygen and hydrogen the combination
shows no free electricity. There now lie --- one may represent it to oneself --- in a given line a certain number of
water-parts between the two poles of a galvanic chain. Each water-part consists of one oxygen and one (double) hydrogen
atom. The positive pole acts upon the first water-part in such a way that the oxygen atom, which is negative, is
attracted, the hydrogen atom, which is positive, on the contrary repelled. But the first water-part acts in turn upon the
second, this upon the third, and so on. When now the attraction that the positive pole exercises upon the oxygen atom in
the first water-part is sufficiently great, this is torn away from its hydrogen atom; but the separated hydrogen atom
combines immediately thereafter with the oxygen in the second water-part; the hydrogen atom thereby set free draws in
turn the oxygen from the third water-part, and so on. Throughout the whole stretch there takes place a continuing
separation and composition of the water; only at the poles themselves can its constituents become free.

The separation of water by a galvanic current in a voltameter makes the voltameter itself a galvanic element.
% --- p. 298 ---
If the current be broken and the voltameter be connected immediately thereafter with a multiplier, one gets a short-lived
deflection from a current in the opposite direction. That the current's direction has turned about rests on this, that
the chemical activity has turned about; it is not a separation that now takes place, but a combination, a reunion of the
water's constituents, oxygen and hydrogen, separated by the first current. The electric activity that the substances
separated at the electrodes exercise is called polarisation, and the electrodes themselves are said to be polarised.

There can then be no doubt that galvanism and chemism stand to one another alternately as cause and effect. When
therefore physics, although it finds itself in steady progress with the discovery of the laws, and although its
investigations are far from completed, has already attained a manifold knowledge of the peculiar coming-forth of the
phenomena, continues to assure us that it knows nothing, or as good as nothing, of the real essence either of magnetism
or of electricity or of chemism: then a surmise is awakened that behind this alleged not-knowing there lies concealed a
misinterpretation of the nature and essence of knowledge.

% The working scan shows a raised dot after „Indhold“ and another in mid-line on p. 302
% („Forbinden. · Og hvad“). Neither is on the KB copy: it is marginal dirt in the working
% copy, not type.
Every scientific knowledge of physical content is formed by a peculiar reciprocal action of sensation, representation and
thought: to sensation there corresponds the fact, to thought the law, to representation the hypothesis. As
representation, on the subjective side, stands mediating between sensation and thought, so the hypothesis, on the
objective side, stands mediating between fact and law: the fact is what is sensed, the law what is thought, the
hypothesis what is represented. Representation starts from what is sensed, prepares the resolution of the particulars
into the universal and remodels the whole for thought's benefit; it slips presuppositions underneath, that thought may
draw consequences; in short, it unites itself with reflection and forms hypotheses. If one takes the word hypothesis in
the most comprehensive
% --- p. 299 ---
% The Danish breaks mid-word here: p. 298 ends „Be-“, p. 299 begins „tydning“.
sense, then one can say with truth that all physics' knowledge is in the last instance a knowledge of hypotheses. The
mark of the hypothetical is then not the uncertain, but what is presupposed by representation: an unsensed fact set up
behind what is sensed, and yet after the likeness of the sensible.

From this point of view we can understand that physicists hold the essence of light and heat to be more unknown than that
of gravity, and the essence of magnetism and electricity in turn more unknown than that of light and heat. The cognition
of the law of gravity rests upon only one universal presupposition (hypothesis), whose consequential consequences confirm
themselves everywhere in the observation of the phenomena; the doctrine of light and heat, on the contrary, has a highly
intricate hypothesis enclosing manifold presuppositions: the aether and its waves are for consciousness not immediately
actual but merely possible, represented facts. Even if the hypothesis gradually gains in probability, even if the
probability rise to a maximum that excludes all uncertainty, the hypothesis does not thereby change its nature, does not
cease to be hypothesis. The highest that chemical physics can strive for is the formation of a hypothesis that explains
the phenomena of magnetism, electricity and chemism in a way as many-sided and satisfying as the hypothesis of light at
present explains the phenomena of light. But even if such a hypothesis could be set up, there is nevertheless a
qualitative leap from the very likeliest hypothesis to a real insight into the essence of the matter: unsensed sensible
facts, facts behind facts, can only be represented; but the representations must be resolved into concepts --- and that
is the law of all knowledge --- the concepts must be thought.
% SECTION a) ENDS HERE. The text breaks off about three quarters of the way down the leaf;
% below it stands blank space and NO rule. Checked against the KB copy. Section b) begins at
% the head of p. 300 as a run-in head.
% --- p. 300 ---
% THE TABLE OF CONTENTS AND THE HEAD DISAGREE, and by the house rule the disagreement is
% recorded rather than smoothed away:
%     Contents : „b) Det elektro-chemiske System“  (with hyphen)
%     Head     : „b) Det elektrochemiske System.“  (without hyphen)
% Settled on the KB colour copy under strong magnification of the word itself, and gone over
% on the working scan: both copies set „elektrochemiske“ as one word, without hyphen and
% without space. The word is bold and unbroken --- it stands in mid-line, so there is no
% question of a word division. The book elsewhere writes the word as one throughout
% (pp. 294, 303, 305); it is the table of contents alone that has the hyphen.
\runin{b) The Electrochemical System.} The thinking in which the concepts of nature are clarified can as little dispense
with representation as representation can dispense with the observation of facts. But in the concept, that is, in the
logical essence of the matter, all representations determined by the facts must be resolved. This resolution is no
thoughtless levelling that lets everything run out into one; it is on the contrary an ideal articulation, under which
representation's elements are reflected and transformed into moments of content for the concept in question. As regards
the essence in existence, the matter in fact, on the other hand, the concept's content must be separated out into
representations, and what is represented exactly determined by what exists. To the subjective conversion of
representation into concept and concept into representation there corresponds the objective conversion of ground into
existence and existence into ground, of possibility into actuality and actuality into possibility, a conversion which
nature accomplishes in a peculiar way in chemism and the chemical processes.

% „Væsen“ in „Tingenes Væsen“ is letterspaced on the KB copy; spacing.py found nothing on
% this leaf.
It is impossible to understand the objective conversion without being clear about the corresponding subjective one.
Existence can be represented, but the ground cannot be represented, it must be thought; actuality can be represented, but
possibility cannot be represented, it must be thought. What thought holds fast in this sense is something more than
formal determinations of connexion, combinations, relations; it is essence as essence, the essence of nature in its
concept, an ideal essence of power, a substantial essence of thought. If one confuses what can merely be represented with
what alone can be thought, if one supposes that by ever more ingenious combinations of bodies' invisible parts, by
continued inventions of atoms, molecules and molecular states, one can work one's way forward to a penetrating cognition
of the \emph{essence} of things, then one deceives oneself; for it is as impossible to represent an essence as it is
impossible to see a tone or hear a colour.
% --- p. 301 ---
The relation between the represented and the essential interior will show itself more clearly by degrees, when we mark
how it has gone in the history of science with the interpretation of the chemical phenomena.

% The Latin is italicised in both copies: annihilatio rerum, creatio, and again annihilatio
% and creatio. The word „Paracelsisk,“ with its comma checked on the KB copy; tesseract read
% „Paracelsiskø“. The leaf bears a vertical pencil stroke in the inner margin (the same hand
% as elsewhere in the working copy). It is not text.
To represent to oneself facts behind facts, fundamental parts so small that not even the strongest microscope is able to
make them visible, is by no means inadmissible. It is in such good agreement with the nature of things that the
objections which speculative philosophy of nature has wished to raise against it betray a questionable leaning to the
mystical. The mystical shows itself when thought, by a leap from possibility to actuality, skips the necessary
intermediate determinations. Paracelsus assumed that the copper which came forth from ``the blue vitriol'' in which one
had laid ``metallic iron'' came forth --- not, as science now understands it, by a separating out from the vitriol
(sulphate of copper oxide), of which it constitutes a component --- but by a mode of formation that stood in essential
consonance with the theological explanation of a \textit{creatio} answering to an \textit{annihilatio rerum}. There is in
Schelling, and even in Hegel, much that is Paracelsian, save that in the latter negation and the negation of negation
step into the place of \textit{annihilatio} and \textit{creatio}.

What is precarious in the hypothesis-forming representation is that its imaginary facts can so easily be false, and the
false presuppositions settle fast as dogmas. A striking example of this is the hypothesis of phlogiston. The changes that
all combustible bodies, the organic as well as the inorganic, undergo in fire, G. E. Stahl gathered under a common point
of view, in that he derived combustion from a substance dwelling in all bodies, a universal combustibility-substance
which he called phlogiston. Thus was formed the phlogistic theory,
% --- p. 302 ---
% The Danish breaks mid-word here: p. 301 ends „phlogisti-“, p. 302 begins „ske“.
accepted from the latter half of the seventeenth to toward the close of the eighteenth century. The representation of
phlogiston stands in close connexion with the mediaeval representation of the metals as compound substances. How the
metallic calx, a designation that answers roughly to the metallic oxides of our own day, stands to phlogiston according
to the theory can easily be made clear. The more phlogiston a body contains, the more combustible it is, and the less ash
there is after the combustion. Carbon, which leaves behind only very little ash, is then especially rich in phlogiston.
If now rusted, that is, burnt iron, the calx of iron, be heated with charcoal, the carbon will give up its phlogiston to
the calx; there will thus come forth calx of iron with phlogiston, and dephlogisticated charcoal; the calx of iron will
then, in its union with phlogiston, constitute the reduced metal, the pure iron, a compound substance.

Investigating science, which cannot dispense with a represented substratum, understands meanwhile with admirable
shrewdness how to procure itself means for sifting the representations and testing their objective validity. This showed
itself in a striking way when Lavoisier came forward with the balance and called the phlogistians to account.

If one weighs the iron before and after its combustion, one will convince oneself that the burnt mass weighs more than
the unburnt. The combustion can therefore not take place by something going away; it must on the contrary rest on
something being added; it consists not in a separating but in a combining. And what then is phlogiston? An arbitrary
representation, irreconcilable with the nature of the matter! If phlogiston is to be the combustible, what then is the
fire-sustaining? How can combustion rest upon a separating out of phlogiston, when the body increases in weight by
combustion? When the adherents of the phlogistic theory had to answer these questions, they entangled themselves in
contradictions and at last laid bare their weakness by resorting to such an assertion as
% --- p. 303 ---
this, that phlogiston must be a substance that made bodies lighter.%
\footnote{% printed mark: *) --- the only note on p. 303. The mark stands AFTER the full
% stop: „lettere.*)“.
% „lige lange. Kaldes“: the full stop after „lange“ is so faintly printed on the KB copy
% that it is only guessed at as a pale dot, and the working scan's thresholding has lost it
% altogether. The compositor's wide sentence space is there, and the sense requires the full
% stop; it is reproduced. The rejected reading is that the full stop is wanting altogether.
% The formula is printed centred on its own line with the fraction under the radical and a
% full stop after. The printed letters are upright; by the house rule (pp. 39--47) this is
% not reproduced, and the ordinary mathematical italic stands.
The phlogistic hypothesis has even been refuted by an exact mechanical proof. Tobias Meyer showed that pendulums of iron
and of oxide of iron have the same time of oscillation, whereas the iron, according to the phlogistic theory a
combination of calx of iron and phlogiston, ought to have less weight and greater mass than the calx of iron, and hence a
greater time of oscillation; on the presupposition, of course, that the pendulums are of equal length. If the time of
oscillation be called T, the moment of inertia I, the distance of the centre of gravity from the point of suspension a,
and the weight P, there results:
\[ T = \pi \sqrt{\frac{I}{aP}}. \]
For a pendulum with greater mass I, which contains the sum of all the parts of mass, becomes greater; for a pendulum with
greater mass and less weight $\frac{I}{aP}$ consequently becomes greater, and consequently T also greater.}

By combining quantitative with qualitative analysis chemistry can establish that matter maintains itself under all
changes, that no creation and no annihilation takes place, that the weight of every chemical product is equal to the sum
of the constituents' weights, and that every increase therefore signifies a combination, every decrease a separation.
Within the circle of these fundamental determinations new representations may now form themselves for the grounding of
the phenomena, new hypotheses be laid down; but they will all be subject to science's correctives, they must all pass
through criticism's ordeal by fire.

% „Støchiologien“ with ø, checked on the KB copy; the Google layer read „Stochiologien“.
% Likewise Qv- throughout here (Qvalitative, Qvantitative), where the Google layer gives Qu-.
The electrochemical theory, indicated by Ørsted, particularly developed by Davy and Berzelius, sets out in stoichiology
to let both the qualitative and the quantitative come to full validity, the quantitative in the chemical proportions, the
qualitative in the electric articulation. The mode of representation answering to this is in brief the following.
% --- p. 304 ---
As zinc and copper in contact with one another acquire opposite electricities, so, according to the electrochemical
theory, the atoms of two elementary substances acquire opposite electricities. That the atoms are not in and for
themselves electric, but first become so by mutual contact, is a presupposition of representation which explains the fact
that one and the same substance can, according as the combination's answering member is, be now positive, now negative.
Thus sulphur in combination with oxygen forms the electropositive, but in combination with hydrogen the electronegative
member. By the mutual contact of the fundamental parts the electricity becomes bound, so that the body is everywhere in
its inner composition equally both positive and negative; but the articulation repeats itself in ever more concrete
forms; for when the body produced by the combination is brought, under the requisite conditions, into chemical relation
with another body, the neutrality ceases.
% Letterspaced on the KB copy: Salte, Syren, Basen. „ikke“ and „Ilt“, which spacing.py
% flagged above, are on the contrary ordinary roman on the KB copy --- rejected.
% „o. s. fremdls..“ with TWO full stops (the abbreviation's and the sentence's) stands thus
% in both copies. Reproduced as printed.
As the elementary substances unite in pairs into binary combinations --- combinations of the first order --- so the
binary combinations unite in turn in pairs into combinations of the second order and form \emph{salts}, in which the one
of the combinations, \emph{the acid}, is the electronegative, the other, \emph{the base}, the electropositive member;
salts are combined in turn into double salts and so forth.. In the series of the elementary substances oxygen stands
highest as the most negative, potassium lowest as the most positive substance; from this it can then be explained that an
acid's character is in general the more strongly stamped the nearer its constituents lie to the negative, and a base's
the more strongly the nearer its constituents lie to the positive end of the series: sulphuric acid is the strongest of
all acids, potash the strongest of all bases. It is likewise in consequence of the theory, and a fine proof of the
reciprocal determination of the quantitative and the qualitative, that a body which combines with oxygen in several
proportions takes on now acid, now basic properties, according as the oxygen, the electronegative member, is more or
% --- p. 305 ---
% BERZELIUS'S OXYGEN DOTS. Six formulae on this leaf, all read on the KB colour copy at
% 6--8 times magnification (the working scan shows only indeterminate blots, and both OCR
% witnesses lose the dots):
%     manganic acid       Mn with THREE dots over the M    (= MnO3)
%     permanganic acid    Mn2 with SEVEN dots over the M, three over four (= Mn2O7)
%     protoxide of Mn     Mn with ONE dot over the M       (= MnO)
%     sesquioxide of Mn   Mn2 with THREE dots over the M   (= Mn2O3)
%     alumina             Al with THREE dots over the A    (= Al2O3)
%     magnesia            Mg with ONE dot over the M       (= MgO)
% Set after the precedent of p. 224 ($\ddot{\dddot{\mathrm{N}}}$, two over three, the
% outermost sign being the UPPER row) and after p. 246, where the formulae are set with
% upright letters (\mathrm) and lowered figures in mathematics. The fourfold dot accent is
% wider than the M and pushes the n away; \! sets it right. That is the setting's affair,
% not the leaf's.
% Note the setting's punctuation: the first two formulae stand WITHOUT a comma, the last two
% WITH. Reproduced as printed.
less predominant. Thus $\dddot{\mathrm{M}}\mathrm{n}$ is manganic acid and
$\dddot{\ddddot{\mathrm{M}}}\!\mathrm{n}_2$ permanganic acid, whereas $\dot{\mathrm{M}}\mathrm{n}$, protoxide of
manganese, and $\dddot{\mathrm{M}}\mathrm{n}_2$, sesquioxide of manganese, are both bases, though of differing
character. The quality of articulation is quantitative, its quantity qualitative: they are articulated relations of
measure. It is finally intelligible from this that one and the same binary combination can in one composition play the
part of the base, in another that of the acid. In silicate of alumina the alumina is base; in alumina-magnesia,
$\dddot{\mathrm{A}}\mathrm{l}$ $\dot{\mathrm{M}}\mathrm{g}$, the alumina is acid: the negative is negative only by its
opposition to the positive, the acid only acid by its opposition to the base, and conversely.

Even if it be granted that the electrochemical theory at its present stage of development cannot sufficiently explain all
chemical phenomena, that there is doubt about the place of individual elementary substances in the series, that chlorine
for example shows itself now more positive, now again more negative than oxygen, and so forth, it must nevertheless also
be granted that the system of presuppositions which the hypothesis-forming representation has here set up has a quite
other scientific significance than the representation of phlogiston. But suppose even that the electrochemical theory
were to develop into so comprehensive a many-sidedness that it at last took the other theories, the theory of types%
\footnote{% printed mark: *) --- first of two notes on p. 305. The mark stands BEFORE the
% comma: „Typetheorien*),“.
% „organiske“ is letterspaced in the note (checked on the KB copy); „der i de“, which
% spacing.py also flagged, is merely loose justification.
By J. B. Dumas, who assumes that in the \emph{organic} combinations there are certain types in which every element can be
replaced by another element. The particular case in which the exchange of an equal number of equivalents takes place he
calls metalepsy.}, the thermochemical theory%
\renewcommand{\thefootnote}{**)}%
\footnote{% printed mark: **) --- second of two notes on p. 305. The mark stands BEFORE the
% comma: „Theorie**),“.
J. Thomsen. Contribution to a Thermochemical System. Transactions of the Royal Danish Academy of Sciences. 5th Series.
3rd Volume. Copenhagen 1853.}%
\renewcommand{\thefootnote}{*)}%
, spectral analysis and so forth up into itself as moments, and thus took on the stamp of a fully valid system of nature:
there would yet
% --- p. 306 ---
show itself, for the hypothesis-forming representation, a limit, precisely that limit which qualitatively separates fact
and essence.

% The Liebig citation: the quotation marks are opened and closed three times, once about each
% piece around the interpolations („siger Liebig *)“, „tilføier han“). Here, then, each piece
% is closed --- the book's usual habit of reopening and closing only at the end is NOT found
% at this place; the balance comes out even. The mark of omission is printed as three spaced
% full stops. The mark *) stands with a space after the name and BEFORE the comma.
In the exposition of the chemical forces' activity, the body-forming functions, the doctrine of atoms must necessarily
play an essential part, from whatever point of view one otherwise regards it. An exact chemical analysis without any
application whatever of atomistics may well be taken for an impossibility. But atomistics falls especially under
representation. ``There must,'' says Liebig %
\footnote{% printed mark: *) --- the only note on p. 306.
Chemical Letters. 5th Letter.}, ``be a cause which makes the elements' coming together in other proportions impossible,
which sets up an insuperable hindrance to their diminution or increase. The fixed proportions are utterances of this
cause . . . the cause itself cannot be observed with the senses, and can become an object only for speculation, for the
spiritual faculty of representation. In setting out,'' he adds, ``to develop the view which has for the present become
the prevailing one concerning the cause of the chemical proportions, I must beg that it not be forgotten that its
falsity or truth has not in the least degree anything to do with the law itself; this last remains, as an expression of
experience, always true and does not alter, however the representation of its ground may shift about''. Herewith the
hypothesis-forming speculation of representation is characterised in a telling way; but herewith too the limit of
empirical knowledge is clearly and definitely indicated.

The philosophy of nature has now the task of making clear the relation between the represented atoms and the chemical
real concepts; from this point of view the atoms are to be conceived as substantial points of function. That the atoms
are in the chemical sense points of function means: 1) that their peculiar existing is determined by their peculiar
acting; 2) that by their activity they both bring forth and themselves undergo a series of law-determined changes; 3)
that under all changes they hold themselves unchanged, not in material immediacy, but in essential form as reflected
constants.
% --- p. 307 ---
% The Danish breaks mid-word here: p. 306 ends „eiendom-“, p. 307 begins „melige“.

From the first determination it follows that the real concepts let themselves be articulated by chemical analyses.
Object-power's ideality, the power in the bodily essence, has its real expression in the atoms, oxygen's essence in atoms
of oxygen, hydrogen's in atoms of hydrogen, water's in atoms of water. According to immediate object-power either the
total concept (water's concept) must exist and the existence of the side concepts (oxygen's and hydrogen's) be cancelled,
or conversely; this is power's doubly-sided expression, but in both cases the expression will rest upon the functioning
atoms.

% The citation opens with „ on this leaf and is NOT closed here: it runs on to p. 308.
% A citation that merely runs over the turn of a leaf is by the house rule no fault, but the
% quotation balance for this piece thereby becomes +1.
From the second determination it follows that the body formed by a chemical combination is a product of the atoms'
arrangement only in so far as the arrangement is itself a conditioned effect of the atoms' nature. ``With the number of
the atoms that are united into a group, the direction of the attractions is multiplied; wherefore the strength of the
attractions also decreases in the same degree as the manifoldness of the directions increases. A part of common salt, one
of the smallest parts of cinnabar, presents a group of not more than two atoms; an atom of sugar, on the other hand,
contains thirty-six, the smallest part of olive oil several hundred single atoms. In common salt the attraction utters
itself in only one, in the atom of sugar in thirty-six different directions. Without anything being added or taken away,
we can conceive the atom of sugar's thirty-six single atoms as ordered in a thousand different ways; with every change in
the position of a single one of these atoms the compound atom ceases to be an atom of sugar, for the properties peculiar
to it vary with the manner in which its atoms are
% --- p. 308 ---
% P. 308 begins WITHOUT indentation: the first line continues the sentence from p. 307
% (checked on the image). No new paragraph here.
% The citation was opened on p. 307 and is closed here; by the house rule a citation that
% merely runs over the turn of a leaf is no fault and gets no note. The quote balance on this
% page is therefore minus 1.
arranged'' (Liebig). But the manner in which the atoms are arranged must surely, under given conditions, have its ground
in their mutual attraction, in their mutual disposition, in their nature and essence. In representation the cause can be
separated from the conditions, but not in the concept. Sugar's concept and olive oil's concept are in the essential sense
total concepts; it is in these total concepts that the ideal-real causality of the arrangement must originally be sought.
The compound atoms are not aggregates --- an aggregate of atoms is no atom --- they are concrete points of function,
newly formed atoms, for which the single atoms are substantial moments in power's expression.

% The letterspacing of „Forvandling“ and „Forandring“ is checked on the KB copy: both words
% are letterspaced, the words following („i Grunden“, „i Existensen“) are not.
From the third determination it follows that the body-forming functions become form-giving in that they effect a relative
\emph{transformation} in the ground, while they effect a corresponding \emph{alteration} in existence.

If one denies the transformation, then there is no essential formation of body; no new body comes forth by the
combination. Water, for example, is then, strictly speaking, only an illusory, not a substantial unity of body; the
phenomenon of water as an immediate thing of sense is a semblance without essence: water as water is not there in
actuality, but only in our imagination. It is, according to the atomistic mode of representation, not the water but the
atom-groups of oxygen and hydrogen that is the thing. The elementary atoms are the things of things. Instead of a ground
of things the atomist sees everywhere things of things, things within things, things behind things, the real things
behind the unreal. Since now the real things are unsensed things of sense, the hypothesis-forming speculation of
representation comes at last, out of sheer interest in sensation, into conflict with the senses.

If one overlooks the alteration conditioned by the arrangement, then one overlooks the articulation and loses oneself in
a foggy speculative formalism. Oxygen and hydrogen then do not become, by
% --- p. 309 ---
mutual alteration, the fundamental parts of water, discrete moments in actual parts of water, but abstract side-forms of
the concrete water-form; then one mediates form by form in formal infinity, but to an essential insight into any actual,
natural formation of body one does not attain.

% The head c) is printed in bold and run into the section's own first line; checked on the KB
% copy, where it clearly stands bold against the surrounding roman. This edition does not
% reproduce the distinction bold/letterspaced; \runin covers both.
\runin{c) Elementary Formations of Bodies.} No formation of a body is in the physical sense to be conceived as a single
effect of a single cause; every body comes forth, be it as educt or as product, from the co-operation of different
causes. As it goes with the formations, so also with the transformations. In the reciprocal action of cohesion and heat
by which the forms of state are determined, the mechanical, the dynamic, the magnetic, the electric and the
electrochemical can take part in the most various ways. So the productions are highly unlike. What a swarm of simple and
compound substances, of chemical formations, mechanical aggregates, amorphous lumps, crystalline masses and clear
crystals!

% spacing.py flagged „enkelte“ and „alle“ here; both are merely loose justification and are
% plain on the KB copy. Rejected. „Formvirksomheden“, on the other hand, is letterspaced,
% checked on the KB copy.
Seen under the point of view of ideal causality, the elementary formations of bodies are, without regard to their
apparently greater or lesser originality relatively to one another, all to be referred to a formative drive prevailing in
nature, and this formative drive in turn to the fundamental will that utters itself in the system of real concepts. From
the point of view of the physically real causes, on the other hand, the formations are to be conceived partly singly,
partly in groups, and without exception to be referred to the forces operating in given substances under given
circumstances. For a union of two such opposed points of view a third must necessarily be sought; this third is
unambiguously designated by \emph{the activity of form}. Form is the natural middle category between the determining real
concepts and the corresponding material causes.

% The Greek sub-head is LETTERSPACED and run-in. The letter is read from the page (KB copy):
% α, not a. „ere“ (and „vorde“ on the next page) are likewise letterspaced, checked on the KB
% copy. „Begreberne“ itself is plain.
\subrunin{α) Real Concepts and Forms.} Concepts \emph{are},
% --- p. 310 ---
but forms \emph{become}. The concepts oxygen and hydrogen do not come to be as consequences of the formation of oxygen
and hydrogen; but oxygen and hydrogen are formed in that the concepts are \emph{realised}, and the concepts are realised
in that the substances are \emph{formed}. Water's concept does not result from the formation of water, but the formation
of water results from the concept by means of the form. There are manifold forms of mixture, but there is, strictly
speaking, not a single concept of mixture. Concepts can in new combinations unfold their consequences into new concepts,
but they cannot be mixed. The quartz-form, the breccia-form, the gravel-form and the like are evident forms of mixture;
but the concepts themselves, by which such formations are originally determined, are none the less to be conceived as
fully valid categories of nature. The essence of mixture can indeed also reflect itself in a concept; but that proves
just how far the concept of mixture is from being a concept of mixed kind. To mixture belongs accidentality, and
accidentality too has its concept; but precisely that reflexion in which the concept of accidentality is demanded makes
the accidental itself a necessary moment: the reflexion of essence cancels accidentality by positing it, and posits it as
cancelled.

Only with existence does accidentality come to be, and with accidentality mixture. The form of nature, the middle
category between concept and existence, has in itself this doubly-sidedness, that it is present both in the logical being
of the concepts and in the real becoming of things; from the activity of form it will be understood that the accidentality
cancelled in the essence can, in the elementary formations, step forth in outwardly co-operating accidents, and that the
real concepts, which in themselves permit of no mixture, are realised in mixed forms.

A naive running together of concepts and forms has only too often led the philosophy of nature astray. How great the
blunders are that abstract speculation has committed in letting the activity of form fall metaphysically together with
the form of
% --- p. 311 ---
% The Danish breaks mid-word here: p. 310 ends „Le-“, p. 311 begins „gemsformen“.
body, and this in turn with the a priori determining concept of body, is well known. Plato's doctrine of ideas in
particular has not seldom had to serve as a point of support for a conception, as poor in spirit as it is mystically
obscure, of the divine archetypal forms' implantation in matter. (Cudworth). The plastic art of nature in the crystals
has served as proof that crystallisation itself was subject to the magic of the types. From this point of view one would
then have to represent to oneself that the whole manifoldness of stereometric figures which, according to crystallography's
now prevailing systematics, occurs in the regular, tetragonal, hexagonal, rhombic, monoclinometric and triclinometric
systems, has been from eternity in the region of the ideas; that it is these self-subsistent typal forms
% The Greek is read from the KB copy at high magnification; in the print the word is divided
% at a line-end (παραδεί-γματα) and is here set whole.
(τύποι, παραδείγματα) which --- for the sake of the plastic art --- everywhere press forward and, where it can be done,
stamp their images in matter, much as a signet ring stamps its inscription in the wax or in the melted sealing-wax.
Eagerly as Aristotle combats Plato's doctrine of the eternal archetypes of things, he is by no means himself free from
harbouring an obscure accessory representation of the metaphysical pre-existence of the forms operative in matter.
% The working scan shows a comma after „Betænkelighed“; the KB copy has no comma there. The
% blot therefore belongs to the one copy and is not text. Both OCR witnesses read it as a
% comma.
The schoolmen accordingly found no scruple in making use of the Aristotelian forms of nature as divine potentates that
ordered matter according to their own demand. Thus matter was there only for form's sake and had necessarily to conform
itself to the
form.\footnote{% printed mark: *) --- the note stands wholly on p. 311.
% The Latin is printed in upright roman, not italicised (checked on the KB copy). The
% ellipsis stands in the print as THREE spaced points.
Non habent esse formæ propter materiam, sed magis materiæ propter formas .\,.\,. Neque igitur formæ ideo sunt diversæ, ut
competant materiis diversis, sed materiæ ideo sunt diversæ, ut competant diversis formis.\par
Unde fit, ut qvædam formæ reqvirant materiam simplicem, qvædam vero materiam compositam; et secundum diversas formas
diversam partium compositionem oportet esse congruentem ad speciem formæ et operationem ipsius. Thomas Aqvinas.}
% --- p. 312 ---
% The Hegel citation: the compositor reopens „ after the interpolation „siger han“ and closes
% in the usual way; here it comes out even, for „Formens Virksomhed“ is a closed citation of
% its own. The quote balance on the page is 0. The letterspacings are all checked on the KB
% copy. The ellipses stand in the print as THREE spaced points; the first is moreover preceded
% by the sentence's own full stop.
The same one-sidedness prevails even in Hegel, only with this difference, that the dialectic of the concept is here so
thoroughly carried through that the blunder must really strike the eye. ``Form's activity,'' he says, ``is none other
than the concept's in general: to posit \emph{the identical as different and the different as identical}. .\,.\,. To the
extremes of subjectivity, which holds itself at one point, and of the fluid, which is only as continuum but in itself
completely undetermined, magnetism yields \emph{the middle}, form's abstract liberation, which in the crystal becomes a
material product, as it already shows itself for example in \emph{the ice-needle}. The activity that has passed over into
its product is the shape, and determined as \emph{crystal}. In this totality the different magnetic poles are reduced to
neutrality, and the linearity of the place-determining activity is realised into the whole body's face and surface
.\,.\,.
% α) and β) here are Hegel's own enumerating letters, not heads; in the print they stand by
% chance at a line-end, but the text runs on unbroken.
Form, the one, works all this, in that it α) bounding the sphere, crystallises the body outwardly, and β) forming out the
punctuality, crystallises its inner continuity in a thoroughgoing manner in \emph{the cleavage of the laminae}, that is,
in \emph{the nucleus-shape}.''

One sees how reflexes of material determinations of form --- bounding surfaces, axes, fracture surfaces, inner and outer
crystallisation --- have hovered before the thinker in the development of the real concept of form; but one sees also how
entirely the operative causes with their conditions are dissolved and vanished in the a priori formal analysis of the
concept. The consequence of this is that elementary formations of bodies in the most outwardly concrete forms of mixture
must all let themselves be derived with the same logical necessity from the forming concept: in granite, the earth's real
kernel, there is then revealed a ``granite principle'' with quartz, mica and felspar, and therein is recognised ``the
terrestrial trinity''.

% The sub-head is LETTERSPACED and run-in; the letter β read from the KB copy (OCR gave
% „B“/„ß“).
\subrunin{β) Material Causes and Material Forms.} That the physicist has from every side directed his attention
% --- p. 313 ---
to the actual forms of bodies goes without saying; but to conceive form as a substantial essentiality for itself, or to
develop the universal real concept of the movements of form, does not lie in physics' way at all. The crystal form is a
physical result of a body's passage from the fluid to the solid state. How little there can in this connexion be any talk
of a peculiar principle of crystallisation is seen from this, that at the passage an amorphous or merely crystalline
formation can come forth just as well as a formed-out crystal; in this respect the conditions rule. In order to bring
bodies to crystallise one can employ different procedures --- solution in water for example, melting and cooling,
evaporation and condensation of vapours --- according to the constitution of the bodies concerned; but the influence of
the conditions is on the whole known by the rule that the crystals become the larger and the finer the more slowly and
quietly the crystallisation takes place. This
% spacing.py gave „Det“ (2.00) and „men“ (1.29) here. On the working scan the passage is
% underlined in pencil by a former owner; the KB copy shows plain roman all the way.
% Rejected.
fact points in turn to a peculiar connexion of conditions and operative causes. It is the figure and magnitude of the
mutually attracting atoms upon which the formation of the crystal essentially depends. Thus the material cause answers
exactly to the condition, for the more slowly and quietly the crystallisation takes place, the better can the atoms get
time to arrange themselves in full agreement with their inner mutual attraction.

% „Hauy“ stands without the diaeresis in the print (Haüy); reproduced as printed.
% „molécules intégrantes“ is italicised; checked on the KB copy.
By setting physical causes in the place of the metaphysical forms, investigating science has already attained great
results. Hauy, crystallography's real founder, started at the beginning of this century from the presupposition that to
every peculiar crystal form there must answer a peculiar chemical composition. His theory of the primitive forms
determined by \textit{molécules intégrantes}, and of the lawfulness with which the layers of the molecules altered, a
lawfulness from which the derived forms naturally let themselves be explained, actually led him to predict from the
crystal
% --- p. 314 ---
% The Danish breaks mid-word here: p. 313 ends „Krystal-“, p. 314 begins „formen“.
form what analysis afterwards confirmed, a difference in the chemical composition of certain minerals; thus he
distinguished heavy spar from strontian earth, dioptase from emerald, epidote from actinolite, and so forth. His point of
view was indeed one-sided, but the one-sidedness was in a scientific form. It is an honour to the founder that his theory
let itself be exactly corrected; science went a step forward when Mitscherlich proved both that different substances
crystallise in the same form --- isomorphism --- and that one and the same substance crystallises in different forms ---
dimorphism, polymorphism.\footnote{% printed mark: *) --- the mark stands AFTER the full stop:
% „Polymorphie.*)“. The note stands wholly on p. 314.
% Both OCR witnesses read „here“ for „høre“ (the familiar ø/e confusion); the KB copy shows
% „høre“ plainly.
The fundamental forms for polymorphous substances belong as a rule to different systems; yet they can also belong to the
same system. Titanic acid crystallises dimorphously in the tetragonal system as rutile and anatase; but anatase's
pyramids cannot be referred to rutile's, or conversely.} If polymorphism can bring it about that from one substance (for
example titanic acid) several substances are formed (rutile, anatase and brookite), then isomorphism can conversely bring
it about that different substances alternately form one and the same substance (garnet for example), the substantiality
being essentially determined by a typical formula. These are fruitful glimpses into the atomic constitution of bodies
that have opened up through this analysis of the activity of form. And yet, with all this, form as form has lost its
essence's originality, its priority as that which determines from within.

% „Naturformen“ is letterspaced only in its FIRST half: the print gives „N a t u rformen“ ---
% N, a, t, u, r with spaces, „formen“ closed up. Checked on the KB copy at threefold
% magnification; the break does not fall at any line-end, so the letterspacing is really given
% up in mid-word. Whether it is the compositor's negligence or Nielsen's emphasis on the
% member „Natur-“ cannot be settled from the page; by the house rule only what is actually
% letterspaced is marked. Not on the errata leaf.
Here there now shows itself this singular thing: metaphysics, which makes an essence of form, cannot explain the activity
of form; physics, which gives so clear an account of the activity of form, cannot recognise form's a priori essence! The
matter is, according to what has gone before, easy to clear up. As middle category between concept and existence, form
--- the \emph{nature}-form, be it noted --- must belong equally to both. In its unity with the concept form is ideal
causality, its power
% --- p. 315 ---
an ideal power, in which material causes and conditions enter as reflected real determinations, cancelled in the reflexion
of power. In its unity with existence form's real causality is dispersed and divided into a multitude of material causes
and supervening conditions; form's ideal, original power disappears for the observer and loses itself in a parcelling out
of determinations of power, under whose reciprocal action the actual, bodily form can first come forth as the result of
outwardly conditioned activity of form. Both ways of seeing, the metaphysical and the physical, have their common ground
in the fundamental medium, in the reciprocal determination of ideal and real power: nature's elementary activity of form
is a presentation, articulated in matter, of Spirit's infinitely forming comprehension.

% The sub-head is LETTERSPACED and run-in; the letter γ read from the KB copy (OCR gave
% „y“/„7“).
% THE WHOLE of this page is underlined in pencil on the working scan in seven strokes. On the
% KB copy all of this is plain roman. spacing.py's three single calls are rejected for the
% same reason. The pencil is not text.
\subrunin{γ) The Reciprocal Action of the Elementary Forms.} It is the formed substances, the existing bodies, that
unceasingly transform one another. Every body has according to its essence an inward form, and so far the form-giving
comes from within; but since bodily existence is outward and must necessarily have an outward form, bodies must find
themselves in outward relations, and the form-giving must come from without. What is characteristic of elementary
form-giving is a many-sided reciprocal determination of a from-within with a from-without, a reciprocal determination
that answers in every respect to natural otherness. In mechanical form-giving --- bounding forms by rubbing, cutting,
breaking up, crushing and so forth --- the inward is a moment for the outward; but that outwardness none the less has an
inward for its substratum is evident. If a body is to take on a form by being cut by another, it must be of a peculiar
constitution. The aptitude to receive the outer form rests upon cohesion, the inner form, and this in
% tesseract read „Varmen,“; the KB copy shows a full stop (the familiar full stop/comma
% confusion). Likewise „udvortes, og nødvendig“ above, where the KB copy has no comma.
turn upon a peculiar relation to heat. In chemical form-giving the
% --- p. 316 ---
% tesseract read „Udvortes, Moment“; the KB copy has no comma. This page too bears pencil
% strokes in the working copy; spacing.py's three calls after them are rejected on the KB copy.
outward is a moment for the inward. The siliceous substance (silicon) and oxygen are mutually transformed into silicic
acid; the mutual transformation is a consequence of their indwelling form, their interior. None the less an outward is
altogether necessary for the development of this interior; the two elementary substances, each existing for itself, must
by outward circumstances be brought together under such conditions that they can enter into combination. In the formation
of the mass of which a granite is formed, the mechanical with its outward and the chemical with its inward have had an
equal share. When it then seems as if the formation of quartz (crystallised silicic acid) were more original than the
formation of granite, then a distinction must surely be made between existence and concept. Seen in the concept, there is
exactly the same kind of reciprocal determination of an inward and an outward in the formation of quartz as in that of
granite. From the side of existence, on the other hand, it is quite true that there can be quartz without granite, but
not granite without quartz. That the more compound formations take on the stamp of the derived, while the uncompounded
and less compounded look like the original ones, rests upon time and parcelling out, not upon concept and essence. The
physically abstract can be derived from the concrete just as well as conversely. The coming forth of alumina, for
example, may just as well be owing to a dissolution of felspar as the formation of felspar to a formation of alumina;
minerals can come forth by the melting of mountain masses just as well as mountain masses by the heaping together of
minerals.

There is in the essence the same relation between the real concepts and the material activity of form as between the laws
and their fulfilment. In the formation of a sand-heap the laws are just as punctually fulfilled as in the formation of
quartz or felspar. However wildly it may have whirled, and however much that is accidental there may be in a sand-heap's
outer form and inner composition, no physicist will
% --- p. 317 ---
say that for every new accidental sand-heap a new law, accidental according to the circumstances, must have been formed;
but he will say that it is the same eternal old laws that at every new formation take up the accidents they meet, master
them as conditions and transform the accidentality into a moment in conditioned necessity. So too with the real concepts
of form; they do not come to be, occasionally and according to circumstances, for they are in self-power's eternal
comprehension; but they become realised in the course of the ages, now here, now there, according as the development of
form takes place. The existential difference between the real concept sand-heap, the real concept quartz and the real
concept granite reflects itself in existence as a difference of essence between accidental and necessary forms of
formation; but whence does this come? Plainly from this, that in otherness's elementary world there must be a relative,
actual difference between the more constant, the substratum, which expresses the necessary, and the inconstant, the
single and particular, the incessantly varying, which presents the accidental. And yet necessity is accidental, and
accidentality necessary: it is precisely a mark of elementary otherness that necessary and accidental are abruptly
separated and abruptly fall together. In the form of existence the necessary is antilogical just as well as the
accidental, the relatively constant just as well as the variable. How the elementary functions look from this point of
view must be seen more closely from the reciprocal determinations between the formed parts and the comprehending whole.

% Close of the chapter. The text stops about the middle of the leaf; below it stands ONE
% single centred rule, and the rest of the page is blank. The rule is SINGLE, not double ---
% checked on the KB copy under magnification; the double rules at pp. 61, 125 and 249 look
% different. Nothing of the Third Chapter begins on this leaf; the chapter head stands on
% p. 318.
\streg

% Danish: \kapitel{Tredie Kapitel.}{Det elementære Hele: Jordlegemet.}   (printed pp. 318--374)
\kapitel{Chapter Three.}{The Elemental Whole: the Body of the Earth.}

% Danish: \afdeling{A.}{Jordlegemets Sammensætning.}   (printed pp. 318--345)
\afdeling{A.}{The Composition of the Body of the Earth.}

The earthly whole, the earth-concept's totalised expression in power, encloses a compass, given in facts, of elements,
physical causes, functions and elementary formations of bodies. Unfolding its rich content of masses and forms,
object-power shows itself at once as power in the manifold and as unity-power; but from both sides as power in another.
% The working scan seems to show „Eénhedsmagt“ with an accent; on the KB copy it stands
% plainly „Eenhedsmagt“ without one. The accent is therefore the pencil (or a speck) in the
% working copy, not type.
The unity-power that holds the body of the earth together has the existence of the parts for its actual other, comes by
this otherness outside the unity, becomes to itself an other, a manifoldness-power that lets the whole subsist in
outwardly independent parts. The air, the sea and the earth's crust with its fire-fluid interior are particular parts of
mass of the great body. The whole is a composite; the composition, a unity in the form of otherness, is conditioned by
the mutual opposition of the parts. But the parts' opposition is determined only by forms of state, elementary forms of
otherness; the opposition is formal, not substantial: the whole body of the earth could, under certain conditions of
pressure and temperature, be by turns a mass of air, a fire-fluid mass, a mass of water and a solid mass of earth; then
the whole is as each part, and each part again as the whole. The masses opposed by the forms of state pass over into one
another: the parts are totalised. Through the totalising of the parts the unity comes to light, and object-power shows
itself from the opposite side as otherness-power. The other of the manifold is the unity itself, the mighty unity that
lets the parts subsist in the wholeness and become moments for the articulation of
% --- p. 319 ---
% The letterspacing on this leaf is checked word by word on the KB copy. The working scan has
% here a long pencil stroke in the margin and pencil underlinings beneath „Cohæsionen“, „Ilt“,
% „Qvælstof“, „Kulsyre“, „Vanddampe“, „Luftmassen“ and „cohæsionsfri“; these words are
% ORDINARY ROMAN on the KB copy and therefore stand outside \emph.
the thoroughgoing whole. It is in the reciprocal action of \emph{the comprehending parts and the articulated whole that
earthly elemental nature} is to be known.
% The print's letterspacing runs to the end of „Natur“ and leaves „erkjendes“ outside it;
% English word order cannot mirror the break exactly, and the span is set to the same phrase.

% The head a) is printed in bold and stands run-in in the paragraph.
\runin{a) Comprehending Parts.} With the body of the earth the concept of elementary corporeality becomes completely
developed, not in the form of the concept but in the form of existence. The concept of corporeality has three principal
sides, each of which presents corporeality whole, namely: \emph{matter}, materialised possibility; the \emph{body}
determined by fixed limits, materialised actuality; and the unity of possibility and actuality materialised as
\emph{mass}. Relative as this formal distinction is, it will nevertheless show itself decisive in the determination of
the body of the earth's comprehending parts: \emph{the atmosphere, the earth's crust and the sea}.

% The sub-head is printed LETTERSPACED and stands run-in; the Greek letter is read on the KB
% copy (OCR renders it as „a“).
\subrunin{α) The Atmosphere.} That the atmospheric air, regarded as mass, is an actual body just as well as the solid
earth, no one will draw into doubt. None the less a clod of earth, a stone and the like make quite another palpable
impression of corporeality than pure air does. If this difference rested on a subjective estimate alone, it would be
without significance in science; but it has in itself an objective ground, for it depends on cohesion. It is cohesion
that gives a body its mass-like unity and determines its independence: the weaker the cohesion, the less independent the
body. The parts of the air --- oxygen, nitrogen, carbonic acid, aqueous vapours --- find themselves in mechanical
mixture, and the mass of air is as good as free of cohesion. By the parts' being outside one another the mass is
discrete, but by their interpenetration it is nevertheless a continuum. The form of air, the tensely fluid, discretely
continuing form of otherness, makes air a materialised possibility, a unity of space and time equalising for existence,
uniformly extended like space, fugitive as an element of time. A body disappears when
% --- p. 320 ---
it is dissolved into air; it disappears not merely for sight and feeling%
\footnote{% printed mark: *) --- it stands BEFORE the comma: „Følelse*), men“
Chlorine, bromine and iodine are indeed visible as gases; but that does not concern the airy in general. One must on the
whole remember that the marks of character adduced here acquire their relative significance only by particular opposition
to the other forms of state.}%
, but it disappears in actuality itself by losing its own determining form. Air's form is a formless form, the hovering
form of the being of matter, of materialised possibility, the form as well for the bodily existence that has been
cancelled as for that which has not yet come to be.

As the being of matter existing in the universal, but none the less antilogical, form of possibilities, air finds itself
in peculiar fundamental relations to the universally operative essentialities: gravity, light, heat and electricity.
These fundamental relations are in turn specified by particular reciprocal action with the earth's other parts, mainland
and sea. Under heat's influence air takes up vapours and thereby sucks in moisture; then the differences come to light:
the formless matter clothes itself in form and makes its existence recognisable in phenomena and forces.

Pure, still air is in abstract generality a transparent and therefore invisible medium; but gravity, heat and moisture
unite to impart to the layers such a difference that the rays of light passing through are refracted and bent --- the
observing astronomer must calculate the refraction. By the cooling of the vapours the being of matter announces its
hidden corporeality: the invisible becomes visible in \emph{mists and clouds}. The formative drive stirs as in dreams,
the ideal causality of the activity of form, nature's indwelling fantasy, is set in motion, and the variable reflects
itself antilogically in the constant. The clouds, these tossed-off pictures of matter's transformation, possibility's
unripe
% --- p. 321 ---
foetuses, fugitive productions of improvised corporeality, change their shape with every moment; and yet there is in
their appearance, on account of the repetition of the conditions, a certain constancy:
% The three Danish popular names answer to cirrus, cumulus and stratus; they are rendered
% here by their sense, as Nielsen does not use the Latin terms.
\emph{feather-clouds, globe-clouds and trailing clouds} are typical forms for alterations and transitions.

Air big with vapour is not perfectly transparent. Under the states of transition the aqueous vapours let only the red and
yellow rays pass through them; so there are produced such phenomena as the red of morning and of evening. From an optical
point of view the region of clouds is on the whole to be regarded as the region of illusory formations, of shifting
appearances, rich in refracting media, rich in colours, images of light and mirages; \emph{the rainbow, mock suns, the
Fata Morgana} and the like are phenomena, but not bodies.

That materialised possibility is itself a material actuality is known by its activity. The activity, mechanical in
motion, is mechanically dynamic in polar tension. By indicating the atmosphere's pressure the barometer indicates the
air's unequal warming, as the hygrometer indicates its unequal moisture. The unequal warming cancels the equilibrium and
produces currents of air; this is what meteorology explains in detail, \emph{the origin of the winds}. The air which on
account of the strong warming in the earth's equatorial regions rises into the air lifts itself above the colder masses
of air to both sides toward the poles, while the air below streams from the poles toward the equator. As the earth turns
about its axis, masses of air are carried from the higher latitudes to the equator with a lesser velocity of rotation
over lands that move with greater rapidity from west to east, and thereby acquire, in relation to the moving surface of
% „Passatvinden“ is letterspaced on both sides of the word division („P a s s a t-“ /
% „v i n d e n“); the whole word is set letterspaced.
the earth, a motion from east to west; thus there comes forth \emph{the trade wind} with the regularly shifting
\emph{monsoons}. The electrically articulating activity that goes through all ponderable matter goes on peculiar
conditions through
% --- p. 322 ---
the layers of air and sets them in peculiar activity. Currents in higher and lower layers cross one another; the masses
of cloud set themselves against one another as positive and negative; the forces enter into tension, the discharge takes
place in \emph{the lightning} and is accompanied by the rolling of thunder.

In their general outlines the phenomena fall together with the fulfilment of the laws in such a way that the constant is
predominant; but under the incessant variation of the conditions the constant passes over into the variable: the
manifoldness of accidents in rules and exceptions makes the boundary of the law-bound possibilities a movable one.

For intuitive representation calm, transparent air is an incorporeal corporeality, a pure resistanceless passivity;
but the forces slumber in the possibility; when the storm rises and the hurricane rages, when the lightnings crackle and
the thunder rolls: then the passive has become active; then it is recognised that there are forces in possibility's
womb, that the air, ``the bearer of sound'', has a voice to raise in earthly affairs.

% The sub-head is printed LETTERSPACED, run-in; the Greek β is read on the KB copy (OCR
% renders it as „2“, „B“ or „ß“).
\subrunin{β) The Earth's Crust.} The solid form, the form of objectivity in materialised immediacy, the completely
determining form for bodily actuality, is as opposed as possible to the form of air. Regarded generally, the earth's
crust constitutes a solid continuum; but fixity itself conditions a breaking, whereby discreteness is actualised; the one
continuum becomes innumerable separate continua, a body of bodies, a whole in nothing but fragments. It is the hard
otherness, outwardness posited in composition: unity-power and the power in the manifold are equally strong.

From this presupposition there follows a peculiar relation between the constant and the variable.

The variable is a constant; for the constituent parts are, each in its place, resting results of a completed activity of
form: the great body that has come to be has its coming-to-be behind it.
% --- p. 323 ---
The widely spread combinations --- combinations of silica, of clay, of lime, of talc, of iron and so forth --- are found
as existing; their becoming belongs to the past, to primeval time. In the infinite manifoldness of all this present,
abiding, relatively constant, time is at once confirmed and cancelled. The time of the comings-to-be vanished in the
unrest of the formations, in the manifold movements of the elementary transitions; now time stands still, contemplates
its work and rests in its results. Such is the impression of the constant.

But the variable comes to its right just as emphatically; it comes first and immediately to light in space. Thus geognosy
speaks of the occurrence of minerals, stones, strata of earth, mountain masses; and
% MISPRINT (setting, not scanning): in „naar“ the FIRST a is a damaged sort --- the bowl is
% wanting, so that the letter stands as an open angle. Both copies (the working scan and the
% KB colour scan) show the same, so the fault is the compositor's. The reading „naar“ is
% beyond doubt from the context and is rendered here without alteration. The errata leaf does
% not mention it.
when it means the same, it is nevertheless always something else. The quartz that is \emph{here} is in actuality another
than the one that is \emph{there}; even if the composition be the same, the existence cannot be: the same mass of quartz
cannot possibly exist at the same time in two, let alone in a thousand different places. Nor is the same in the chemical
combinations typical for minerals everywhere the same. Felspar is indeed always felspar; it has, like all minerals, its
determinate chemical formula; but that the felspars vary, that there is a real difference between an orthoclase (potash
felspar), an albite (soda felspar), a labradorite (lime felspar) and so forth, is well known.

The variable shows itself still more prominently in mechanically typical relations of mixture. Lava is indeed always
lava; but there is nevertheless a considerable difference between the kinds of lava. ``Etna's lava,'' says the geognost,
``consists especially of a mixture of augite; in Vesuvius's, on the other hand, the felspar is predominant, and the
augite plays a more subordinate part; in some of the older lava streams from this volcano (on Monte Somma) we find even a
great quantity of leucite
% The quotation opened here („bestaaer især …“) is not closed until p. 324 („… under
% særegne Omstændigheder.“); the citation therefore runs over the leaf, which by the house
% rule is no fault.
% --- p. 324 ---
instead of the labradorite. The minerals named are, however, the predominant ones in the lavas of all volcanoes, though
in different proportions of quantity, and according as one of them is preponderant the lava is termed
% „(Doleritlava)“ and „(Trachytlava)“ are NOT letterspaced; only the names themselves.
\emph{augite lava} (dolerite lava), \emph{felspar lava} (trachyte lava) or \emph{leucite lava}. It is true that now and
then still other minerals%
\footnote{% printed mark: *) --- first of THREE notes on p. 324
Lava is on the whole, as may particularly be concluded from Bunsen's analyses, no determinate species of rock.}
appear as subordinate constituents of the lavas, thus especially magnetic ironstone, which indeed is never quite wanting,
mica, olivine, hornblende and so forth; but these may be regarded as unessential or as having arisen under peculiar
circumstances.''%
\renewcommand{\thefootnote}{**)}%
\footnote{% printed mark: **) --- it stands AFTER the closing quotation mark
B. Cotta. Geological Pictures, translated by C. Fogh. Copenhagen 1859. p.\ 127.}%
\renewcommand{\thefootnote}{*)}

Still more outward is the variation in the manner in which the one species of stone is said to pass over into the other;
\emph{granite into syenite}, this in turn into \emph{diorite}, ``whose dense and schistose varieties are with difficulty
distinguished from those of \emph{diabase}, in which hornblende is replaced by augite and the felspar is either
labradorite or oligoclase, wherefore these dense, schistose, porphyritic, porous or amygdaloidal varieties are often
gathered under the common name of \emph{greenstones}.'' ``These greenstones,'' it is further said%
\renewcommand{\thefootnote}{***)}%
\footnote{% printed mark: ***) --- third of three notes on p. 324
% A THIRD abbreviated form: „Anfrt. Skr.“ The book elsewhere has „Anf. Skr.“ and „Anfr. Skr.“;
% this one occurs here. Reproduced as printed, with the English distinguishing the three as
% „Work cited“, „Wk. cit.“ and „Wk. cit'd.“
Wk.\ cit'd. p.\ 137.}%
\renewcommand{\thefootnote}{*)}%
, ``resemble again very much the melaphyres, \emph{the dolerites and basalts}, in which pyroxene'' (that is, augite)
``forms the chief mass in combination with labradorite, olivine and magnetic ironstone, or with nepheline; but the
dolerites and basalts appear already as lavas.'' Such transitions, which all bear the stamp of antilogical accidentality,
can be determined subjectively and yet scientifically by geognostic comparison; but so hard is the otherness, so
dominant is manifoldness-power in the elementary outwardness, that the composition of the one
% --- p. 325 ---
% The Danish breaks mid-word here: p. 324 ends „Sammen-“, p. 325 begins „sætningen“.
species of stone cannot possibly be derived from the relation between the constituents of the other. That the varying
formations have, each according to its conditions in its place, a fully valid cause, a sufficient ground, cannot be
doubted; but the real ground has gone immediately together with its existence; the confusion of the parcelling out rules
in the local causes, and the necessity is antilogical: a systematics of the formations in the logical sense there can be,
but a systematics of the stones there cannot be.%
\footnote{% printed mark: *) --- it stands AFTER the full stop
Only from a general point of view can it be brought out that the species of rock richest in silicic acid (granite and
syenite) are on the whole the older, and the more basic ones (dolerite and basalt) the younger.}

Has unity-power then given itself up for lost? By no means. It has everywhere, in the unbroken systems of composition,
its material middle instances, peculiar central points for completed processes of formation, great groups of mass-like
productions; in these productions the materialised universal has been co-operative, partly in passive, partly in active
form. It is water, the passive-active solvent, and fire, the warm furtherer of elementary activity, that have each in its
own way functioned in the service of the universal and thereby of the unity of power. There is, as regards constituents
and composition, great difference between the stratified
% The series „Overgangskalksteen … Travertino“ is ordinary roman on the KB copy; spacing.py's
% result for „Kridt,“ is rejected.
species of stone --- transition limestone, Zechstein, shell limestone, Jura limestone, chalk, travertine and so forth ---
but they all have nevertheless a distinctive stamp of common origin, in so far as they are \emph{neptunian}, formed,
that is, by solution and deposition in the waters. In opposition to these the massive stones with their corresponding
porphyries and lavas have the distinctive stamp of having passed through the process of fire, partly as older,
\emph{plutonic}, partly as younger, \emph{volcanic} formations. As witness to a many-sided reciprocal action both between
the process of water and that of fire
% --- p. 326 ---
and between the masses themselves we have the \emph{metamorphic} species of stone.%
\footnote{% printed mark: *) --- first of two notes on p. 326; the mark stands after the full
% stop: „Steenarter. *)“
The greatest difficulty in the chemical investigations of metamorphism lies in the limitation of the means, and
especially in the circumstance that a quantity of substances which, on account of the extraordinarily small quantities in
which they occur, can hardly be demonstrated by the sharpest reagents, none the less exercise a mighty effect in the
course of the ages and sometimes appear in mighty unfoldings of mass. Carl Vogt. Grundriss der Geologie. Braunschweig.
1860. pp.\ 462--63.}
``Exact observations comprehending great mountain tracts have,'' says A. Humboldt, ``proved that eruptive masses of stone
do not appear as any wild, unbound operating force. In the regions of the world most remote from one another one often
sees granite, basalt or diorite exercise their transforming influence upon clay slate and dense strata of lime, as well
as upon the grains of quartz in sandstone, and that in a manner alike down to the single utterances of force.''%
\renewcommand{\thefootnote}{**)}%
\footnote{% printed mark: **) --- it stands AFTER the closing quotation mark
Kosmos I. p.\ 224.}%
\renewcommand{\thefootnote}{*)}

A natural tension of unity-power and manifoldness-power, grounded in the strong elementary otherness, displays itself in
a tense opposition between the earth's active --- fire-fluid? --- interior and the solid, resting exterior. One who could
survey at one glance the whole stiffened surface would surely discover that it is always, now at this place, now at that,
in a trembling; these local tremblings then indicate what volcanic eruptions and earthquakes, with all the omens of
subterranean storms, confirm in the strongest way, that even the firmest parts with their composition are in the power of
the wholeness, subject to time and its vicissitudes, and that the earth's crust is indeed a solid and comprehending part,
but still only a part of the whole.

% The sub-head is printed LETTERSPACED, run-in; the Greek γ is read on the KB copy (OCR
% renders it as „y“ or „7“).
\subrunin{γ) The Sea.} The drop-fluid element has in its form of state this peculiarity, that the matter-like possibility
and
% --- p. 327 ---
the bodily actuality neutralise one another into existing mass. Air and earth form extreme members of opposition,
immediate extremes; but the sea lays itself mediating between the two: it is especially the peculiarity of the form of
mass, as elementary form of otherness, that allots to it the mediating role. In the cohesionless mass of air the boundary
is so movable that what is becomes nothing but becoming, the actual nothing but possibility. The mass of earth, the mass
of metal, of mineral, of stone, has so strong a cohesion that each piece for itself is a determinate body with fixed
boundaries and existing properties. In water as water, on the other hand, corporeality has too much independence to be
fugitive matter and too little independence to hold itself objectively fast as a body bounded in itself; so its
characteristic form becomes a particular existing form of mass, undulating between air and earth, between matter and
body.

Water evaporates and passes over into air; it freezes into ice and becomes as a part of the earth's crust. It is water
that in the masses of cloud gives air its phenomena, and it is water that in masses of snow and ice mingles with the
strata of earth, loosens their connexion, tears them along with it, thaws and washes them away in dissolved masses. In
mud, the thickly flowing mass, water can assimilate itself to earth, as in vapour it can assimilate itself to air; but
according to its peculiar cohesion and specific gravity it holds itself between the two, resting upon the earth, akin to
the air.

The sea, with its depths and currents, its relations of temperature, light and electricity, is, so to speak, a fluid
stratum of earth, an embodied counter-image to the air, by whose winds it is moved, whose play of colour it borrows,
whose shifting phenomena it mirrors. There are currents in the sea just as well as in the air, and there is on the great
scale a certain correspondence between currents of air and currents of the sea. --- In the ship's journal
% --- p. 328 ---
on his third voyage Christopher Columbus says%
\footnote{% printed mark: *) --- first of two notes on p. 328
A. Humboldt. Kosmos I. p.\ 275.}%
: ``I take it for settled that the waters of the sea move from east to west, as the stars of heaven do
% The Spanish parenthesis is set in ITALIC (roman italic) against the upright setting round
% about; checked on the KB copy.
(\textit{las aguas van con los cielos}).'' The equatorial current must in particular be regarded as a consequence of the
trade winds and of the advancing tide.

The earth's crust can bulge and burst; but it has no ebb and flood: the cohesion is too strong in the body; the air can
sink and rise, but not in ebb and flood: the cohesion is too weak in the fugitive matter. The tides, ebb and flood, are
changes particularly characteristic of the mass as mass. The mass, which receives its form from its surroundings, is too
little coherent to hold its bounding fast and withstand influences from without, and yet coherent enough to swell, gather
and form itself under attractions from the moon and the sun.

Although the ocean covers about three quarters of the earth's surface, it might nevertheless, as regards the composition
of the whole, seem as if the fluid element were more restricted than the solid and the aeriform; for while the air
surrounds the whole exterior of the earth, and the earth's crust its whole interior, the sea must content itself with a
bed assigned to it within narrower bounds. This restriction is, however, completely cancelled by the series of
transformations --- from water to vapour to clouds to rain and so forth --- that is so characteristic of water's
circulation.

Sea water is, as is well known, not a uniform unity of mass, merely a chemical combination of oxygen and hydrogen, but a
concrete mass containing salt, with chlorine, sodium, magnesium and so forth.%
\renewcommand{\thefootnote}{**)}%
\footnote{% printed mark: **) --- it stands after „o. s. v.“
G. Forchhammer. On the Constituents of Sea Water. Copenhagen 1859.}%
\renewcommand{\thefootnote}{*)}
The quantity of salt, which at all depths is about the same, exercises the most considerable influence upon the sea's
physical properties. Salt water has its greatest
% --- p. 329 ---
% THE MINUS SIGN: here it is printed as a LONG STROKE (the same length as the em-dash
% elsewhere on the leaf), not as ÷. On p. 267 Nielsen sets ÷ in the same sense (÷ 39°); the
% signs are not harmonised but each reproduced as printed. The degree sign is set after the
% practice of p. 267: $^{\circ}$. „0°“ is a real zero, read on the KB copy.
% The print uses the Danish decimal comma (2,67); the English gives the decimal point, as
% elsewhere in this translation.
density at --- 2.67$^{\circ}$ C, its freezing point at --- 2.55$^{\circ}$, whereas fresh water has its greatest density
at 4.5$^{\circ}$ C and its freezing point at 0$^{\circ}$ --- an essential difference with regard to upper and lower
layers during freezing. That the form of mass rules, that the constituents, however different they may be in other
respects, flow away into the dissolving, neutralising unity of mass, is seen from this, that the differences which
physics demonstrates in all the intermediate layers between the surface and the deep essentially concern heat and cold,
density, pressure, motion and rest, and thus yield, not particular determinations of objectivity in particular bodies,
but nothing but predicates of mass, constant-variable fundamental relations, material universal determinations in
undulating transitions.

% The working scan seems to show „dét er en Selvfølge“ with an accent; the KB copy has
% plainly „det“. The accent is the pencil in the working copy (cf. „Eenhedsmagt“ p. 318), not
% type.
It goes without saying that all the parts of the body of the earth must themselves be bodies, the air and the sea just as
well as the solid earth; it goes without saying that every body must unite in itself both possibility and actuality; but
in that the fundamental difference utters itself in elementary forms of otherness, corporeality unfolds its three
principal sides in three different relations of possibility and actuality. Thus we have the three circles of
corporeality: in the one circle actuality is dissolved into its moment of possibility, that is, in \emph{the atmosphere};
in the second possibility is stiffened in its moment of actuality, that is, in \emph{the earth's crust}; in the third an
equilibrium of possibility and actuality is set in undulating masses for a mediating transition between air and earth;
that is the circle of the waters, the tensionless yet fluid element, \emph{the sea}.

% The head b) is printed in bold and stands run-in in the paragraph.
\runin{b) The Articulated Whole.} On account of the developed difference between the comprehending parts, the earthly
whole has attained, not merely an articulation of the substances of matter in the abstract dynamic sense, as by magnetism
and electricity, but, in consequence of the formation of bodies particularly conditioned by chemism, a concrete and
actual articulation, so that the parts,
% --- p. 330 ---
unindependent as they are in and for themselves, appear \emph{each in its own circle} with elementary independence. How
the articulation is categorically carried through --- partly by the composition's own relations, in \emph{the opposition
of land, sea and air}; partly by a union of telluric and cosmic relations, in \emph{the zones of the earth and the
climates}; partly, finally, by those individualisations which yield the object of the description of the earth
(geography) and whose manifold particulars are all enclosed by the relation between \emph{the ocean and the great
continents} --- shall be touched upon in what now follows.

% The sub-head is printed LETTERSPACED, run-in; the Greek α is read on the KB copy.
\subrunin{α) Land, Sea and Air.} If from the beginning there had been perfect equilibrium between unity-power and the
power in the manifold, the somewhat flattened globe of the earth would, according to what has gone before, simply have
had three concentric envelopes: the solid crust, the drop-fluid sea and the sea of air. But unity-power, which makes the
parts unindependent, and the power in the manifold, which gives them independence, find themselves in mutual tension. So
the equilibrium is unquiet. The power in the whole shows itself as a movable power of connexion; it goes with heat and
cold, in great and small, through fire and water, and drives the elementary activity of form forward in all directions to
a many-sided development of articulating consequences.

The unity of the body of the earth's interior and its exterior is, without regard to the differing constituents, an
actual unity of opposition realised in actual forms of state. Since the emphasis falls on the opposition, the solid crust
shrunken by cooling is too narrow to hold the warm interior, the hot vapours with the fire-fluid masses, and whatever
% spacing.py gives a result for „der“ in „her og der“; on the KB copy both „her“ and „der“ are
% ordinary roman here (unlike p. 323, where they really are letterspaced). Rejected.
else there may be, within fixed, immovable bounds. So the exterior must for all its hardness give way here and there; the
earth's crust bulges and is broken; something raises itself while something else sinks: mountain peaks, islands, groups
of islands come to light above the mirror of the sea. But it is precisely through the manifoldness
% --- p. 331 ---
that unity's way to wholeness goes: the separation of land and sea is an articulation of the whole.

In the formal respect the articulation is a bounding, and the bounding mutual: the mainland is the sea's, and the sea in
turn the mainland's boundary. In reality bodily independence still holds good over against the mass, in so far as it is
originally the mainland that sets a boundary to the sea, and not conversely. And yet, remarkably enough, the qualitative
difference between body and mass is lost in the mutual bounding: the body is mass and the mass body, for the articulating
power is equally divided. The sea has power to alter the bounds of the lands, as conversely the lands have power to alter
those of the sea. So one speaks of masses of land just as well as of masses of water; what is common is that the masses
are formed in being bounded. For the form of bounding there is indeed an inward cause, but the inward is, as regards the
articulation, everywhere a moment for outward division. Elevations, fractures, subsidences hold good just as well for
that part of the earth's crust which constitutes the sea bottom, and hence for reciprocal actions of the solid and the
fluid in the deep, as for configurations of the mainland's surface: there are valleys and cliffs \emph{beneath} the sea
just as well as \emph{above} it.

Mountain tracts and valley regions appear not as separate bodies but as great forms of mass. As a whole the mountain is
the materialised universal in which the difference of the constituents disappears, and hence mass. As part of the whole
mass of land the mountain has particular existence only in the form, that is, in that bounding by which the mass is
there: the mountain is a formed mass just as well as the wave and the cloud. The inner difference between independent and
unindependent, which rests essentially upon the form of state, has under the articulation drawn itself into a mass-like
independence and has thereby come to expression in the form of bounding: the mountain with its fixed and the cloud
% --- p. 332 ---
% The Danish breaks mid-word at the foot of p. 331: the word is „Constante“, p. 332 beginning
% with „stante“. Checked on the KB copy. No paragraph break, no rule and no sheet signature at
% the turn.
with its fugitive outlines stand to one another as the constant to the variable; but the constant form of bounding is a
form of mass just as well as the variable one.

% The letterspacing runs over all four mountain names and stops before „o. s. v.“;
% „Alpebjerge“ is divided over two lines in the setting and letterspaced in both halves.
% Checked on the KB copy.
Mountains have, like clouds, typical forms: \emph{massif mountains, chain mountains, plateau mountains, alpine
mountains} and so forth. The causes of formation are for mountains as for clouds inward-outward (with the difference, of
course, that lies in the opposition of the fixed and the fugitive); for repetition of something common, common causes
with common conditions, is in the fixed as in the fugitive the only way in which the universal can be expressed in
elementary existence. According to the figurative form of the bounding, the mountain masses are earth-fast, petrified
clouds, the cloud masses airy, hovering mountains.

As regards the fugitive masses, the form of bounding is movable, and the change rests upon transformations in which they
disappear in order to come to light again; it is otherwise with the fixed masses. The more fixed the masses are, the more
resistance the form of bounding offers to influence coming from without, the more does the transformation take on the
stamp of destruction and disturbance. Air in uproar threatens destruction; but it is something else that is threatened;
air itself is not exposed to destruction. How should that element whose essence is sheer change be able to suffer by
being changed? The sea in uproar threatens destruction, but destruction does not threaten the sea. The mighty sea can
swell and roar, can seethe in breakers and lay itself to rest in cool transparency: the sea in all transitions is always
the same. It is the mainland, what is immediately independent, that is threatened with destruction from all sides, from
the earth's interior, from air and sea. It is by these destructions that the earthly reciprocal action between wholeness
and parts displays itself in striking relations of power. The parts can be destroyed; therein is seen the impotence of
manifoldness-power; but the whole cannot be destroyed: unity-power comes to the manifold's
% --- p. 333 ---
relief on account of the wholeness. The destruction of the parts is the elementary activity of form by which new parts
are produced; the more different these parts are, the more richly they group themselves into characteristically opposed
whole masses, the more prominent is the articulation in the whole.

% The Greek letter is read on the page itself (KB copy): β, not „B“, „ß“ or „P“, as the OCR
% gives. The head is run-in and letterspaced throughout its whole extent.
\subrunin{β) Zones of the Earth and Climates.} As a whole for itself the body of the earth has a compass of telluric
relations answering to its parts, relations of reciprocal action determined exclusively by its peculiar composition. As
part of a greater whole, as planet, member in the solar system, the earth finds itself in cosmic relations, and is
dependent upon the sun not merely mechanically with regard to gravity, but also dynamically with particular regard to
light and heat. For an outward consideration the two systems of relation can doubtless be held apart and so far
distinguished \textit{in abstracto}; but in actuality they are neither parcelled out nor pieced together; they are in the
most various ways fused into concrete, operative unities.

In so far as the telluric relations are taken up into the cosmic, the articulation is designated by \emph{the zones of
the earth} (zones); in so far as the cosmic are taken up into the telluric, it is carried through in the
constant-variable differences of \emph{the climates}.
% „Jordbælter“ and „Klimaternes“ are letterspaced; „(Zoner)“ is not. Checked on the KB copy.
% „Led“, which spacing.py flagged, stands in ordinary setting and is rejected.

% The footnote mark *) stands after the full stop: „uden at give. *)“.
In the relation of gravitation there is reciprocal action between the earth and the sun, they attract one another
mutually; but as regards light and heat the reciprocal action ceases: the earth receives without giving.\footnote{One can
indeed say that the sun also gets its share of the heat that the earth radiates into the space of the world; but so far
as the sun is concerned this may be regarded as a vanishing magnitude.} The foundation of the receiving is its double
motion, the daily about its axis and the yearly in its orbit. That the earth's orbit, the ecliptic, cuts the equator
% --- p. 334 ---
at an angle of $23^{\circ}$ $28'$ is of rich consequence for the distribution of light and heat. The torrid zone, the
belt between the tropics on both sides of the equator; the temperate zones, the belts between the tropics and the polar
circles; and finally the frigid zones about both poles as far as the polar circles, are unequally placed.

At the equator day and night have each their twelve hours; with removal from the equator the day's length alters with the
season; about the poles day and night have each their season. By day the earth's surface is warmed by the sun's rays,
whose effect is as the sine of the angle at which they fall; after sunset the earth radiates heat toward the space of
heaven without the loss being made good: the relation between warming and cooling depends upon the relation between the
length of day and of night. Meteorology has derived from this the following propositions: 1) that heat must decrease from
the equator toward the poles; 2) that heat in the neighbourhood of the equator must throughout the whole year be fairly
% MISPRINT: „temmelg“ for „temmelig“. It stands thus in both copies (checked on the KB copy)
% and is not corrected on the errata leaf. Rendered here by the sense.
evenly distributed; 3) that the seasons stand forth the more distinctly the further one removes from the equator, since
the difference between summer and winter heat grows ever greater, so that the summer even in the neighbourhood of the
polar circle can still be very warm.

The heat-absorbing and heat-radiating earth stands in an active relation to everything that it receives. The telluric
conditions co-operate with the cosmic and, by taking these up into their own circle, form those particular relations of
temperature which, in concretions of the constant with the variable, give the climates. ``If the whole surface of the
earth were covered with water, or if it were formed of solid land, everywhere flat, which was everywhere of the same
constitution and had at all places the same capacity to receive and give off rays of heat, then a place's relations of
heat would depend only upon its geographical latitude, and all places on the same
% --- p. 335 ---
degree of latitude would have the same climate. But as it is, the effect that the sun's rays can produce is altered by
manifold causes, and a place's climate depends not only upon the direction of the sun's rays but also upon the
circumstances under which they act; it depends upon the form of the land and of the sea, upon the direction and height of
the mountain tracts, upon the direction of the prevailing winds and so forth.''\footnote{The Chemical Part of the Science
of Nature, after Joh. Müller's Grundriss der Physik und Meteorologie. Translated by C. L. Petersen. Copenhagen 1856.
p.\ 205.}
% The citation opened on p. 334 with „ is closed here; the quotation mark therefore runs over
% the leaf, which by the house rule is no fault. The footnote mark *) stands AFTER the closing
% quotation mark.

How the matter stands in detail with land and sea climate, with the mean heat of the year and of the seasons, with the
curvatures of the isotherms; why isotheres and isochimenes wind like serpents across the circles of latitude, and so
forth --- all this belongs under that new circle of empirical observations which Alexander Humboldt has introduced and
which the science of nature restlessly carries on. Philosophy must restrict itself to an indication of established
fundamental relations.

The functionary of which the power of wholeness particularly avails itself in all elementary articulations is heat. Heat
works for unity in that it works for the uniform, the vaporous. So it meets resistance in the different constitution of
different bodies: the manifold and many-formed stands outwardly up against unity. But heat is prepared for the
resistance; it has itself introduced it. Through degrees and gradations the intensive unity of heat enters into the
manifoldness of the many-formed, takes up the bodily differences and constitutes itself into \emph{temperature}. The
elemental whole in all its parts, the oppositions of land, sea and air, the differences of the climates, the vertical
ones in the heights of mountains and the horizontal ones in the tracts of the earth --- all have their subsistence in and
by temperature. Without a thoroughgoing unity-power conditioning all
% --- p. 336 ---
activities of form, the composite whole would be only a finite piecework; but heat, this inwardly intensive and yet so
widely extended unity, cancels the piecework, melts the finite into the dynamic infinity of gradations, alterations,
transitions: heat in its temperature is an articulating functionary.

% The Greek letter is read on the page itself (KB copy): γ, not „y“ or „7“, as the OCR gives.
% Run-in and letterspaced head.
\subrunin{γ) The Ocean and the Great Continents.} The body of the earth, this great whole divided into masses of water
and of land, is a realised power-concept; if such a power-concept be straightforwardly conceived as a logical universal
essence and its self-development presented in the light of absolute necessity, then the blunder is evident.

``It is,'' says Hegel, ``the concept's business to conceive as necessarily determined what appears to sensible
consciousness as accidental.'' On this presupposition he supposes he can deduce ``the continents'' from the concept of
the body of the earth. --- ``In general the striving directed toward a more determinate shape opposed to the sphere goes
out upon the pyramidal, in that within the sphere it forms a base, a breadth, which tapers toward the other side; and
thence comes the land's falling apart (das Zerfallen des Landes) toward the south. But the restless, revolving current
hollows out this figure by degrees from the west inward toward the east, drives and presses the solid as it were toward
the east, and lets the figure swell toward the eastern side like a bent bow, so that it is curved and rounded inward on
the west. On the whole the land is torn asunder into two parts, the old and the new world . . . . The new world
% The omission is set with FOUR spaced points, read on the KB copy; reproduced as printed.
% „(das Zerfallen des Landes)“ stands in ordinary roman, not in italic.
presents on the whole the undeveloped twoness --- a northern and a southern part after the manner of the magnet; the old,
on the contrary, the complete division into three parts, of which the one, Africa, the solid metal, the lunar, stiffens
with heat, where man is dully lost in himself --- the mute Spirit that does not step into consciousness; the second,
Asia, is the bacchantically cometary extravagance, the wildly
% p. 336 ends „den vildt kun af sig“ and p. 337 begins „fostersvangre Midte“; the English
% divides at the nearest point that its word order allows.
% --- p. 337 ---
big with young of itself alone, the formless begetting which yet could not become master of its own middle; but the
third, Europe, forms consciousness, the rational part of the earth, the equilibrium of streams and valleys and mountains,
whose middle is Germany. The continents are accordingly not accidental, divided merely for convenience' sake; there is an
essential difference between them.''

There is in this conception something of a grand kind; it passes like a breath of spirit over the ocean and the great
continents; but it is nevertheless evident that the deduction of ``what appears to sensible consciousness as accidental''
does not here reach its goal; the immanent deduction does not attain to ``the rational part of the earth'', but loses
itself in a ``bacchantically cometary extravagance'' and cannot ``become master of its own middle''.

To set out upon a deduction according to which the globe's present shape --- the five continents, Asia, Europe, Africa,
America, Australia, together with the division of the world-ocean into the Atlantic, the Indian and the Pacific with
their intermediate seas, gulfs, indentations and so forth --- should proceed with logically absolute necessity from the
concept of the globe, is to set oneself a superhuman and just therefore unscientific task. The mutual bounding of the
ocean and the great continents is an outward object for intuition, observation, measurement, description and delineation;
but it is in truth no object for immanent logical deduction. In the articulated whole the concept of the earth is ---
in accordance with its present stage of development --- doubtless completely realised; but precisely by the realisation
the concept has gone under as concept, its logical determinations have in all continents, on all boundaries, at all
coasts, in all valleys and on all mountains become antilogical. It is not the concept of the earth with its indwelling
sides of essence, but the body of the earth with its great limbs, that we immediately have before us. Here, in this
factually
% The left margin of p. 337 bears a long pencil stroke (a reader's hand), which is not text.
% --- p. 338 ---
given whole, there is not a certain sum of accidents falling outside the concept alongside a certain number of
determinations of necessity falling together with the concept; but \emph{everything is necessary because everything is
accidental}: antilogical necessity is the stamp, with those ineffaceable features that object-power has set upon the
composite whole, upon the unity of connexion realised in the articulation.
% The letterspacing runs from „Alt“ to and including „tilfældigt“; „den“ is not letterspaced.
% Checked on the KB copy. „Sum“, which spacing.py flagged, stands in ordinary setting and is
% rejected.

It is not a logical a priori deduction, but an actual recognition and demonstration of the realised concept of
corporeality, that matters with such a unity of connexion. The logical in the concepts of earth, air and sea glimmers
everywhere as a reflex of what is antilogical in representation; only so can the development of the concept prepare the
understanding of the actualised concept. Physics has, as regards the constitution of the body of the earth in the whole
as in the particular, still many observations to make and many problems to solve; that is its interest in infinity in
the midst of the finite: to know and know and know --- the constant in the variable, the universal in the particular, the
necessity of the laws in the many and different cases; but however great the advances made, one never gets the better of
the accidental, the antilogical point of the fact. The philosophy of nature follows in physics' track and controls the
development of the real concepts; but however many-sidedly thoroughgoing the connexion may be, it sees nevertheless
everywhere the accident and accidentality peeping forth from every single member: it is in the unity of composition not
an \emph{absolute} but a finite, \emph{conditioned} and therefore \emph{accidental} necessity that, with object-power's
immediacy, governs the earthly whole.
% Only „absolut“, „betinget“ and „tilfældig“ are letterspaced; „endelig“ and „altsaa“ are not.
% Checked on the KB copy.

% The head is printed in bold (not letterspaced like the Greek sub-heads); „c)“ itself stands
% in ordinary roman. The same as at c) Aarsagsmomenternes Eenhed, p. 79. The difference of cut
% is not reproduced; the house macro \runin serves for both.
\runin{c) Earthly Nature.} In so far as the body of the earth is an outward object for observing consciousness, we have
% --- p. 339 ---
in its composition the many things of nature but not the natural unity of things --- the \emph{geognostic} standpoint. If
consciousness takes up impressions of objects and phenomena into reproductive fantasy, then the idea of nature comes to
light in a double gleam of objective and subjective --- the \emph{aesthetic} standpoint. Only in the scientifically
cognised reciprocal determination of ideal and real causality can the idea realised in the elementary formations step
forth distinctly as \emph{nature} --- the \emph{teleological} standpoint.
% In all three cases the letterspacing catches only the adjective (and „Natur“), not „det“ and
% not „Standpunkt“. Checked on the KB copy.

% The Greek letter is read on the page itself (KB copy): α, not „a“, as the OCR gives. Run-in
% and letterspaced head.
\subrunin{α) The Geognostic Standpoint.} If we gather all those parts of the science of nature that belong to the
investigation of the shape and composition of the body of the earth under the common name of \emph{geognosy}, then
geognosy, with astronomy, physical geography, chemical physics, chemistry, mineralogy and lithology for auxiliary
sciences, is so widely ramified and comprehensive a science that it may well be said to span the whole of earthly
elemental nature. Can this science, then, teach us in the proper sense what nature is? No! It investigates the things of
nature, but not nature. It teaches, to be sure, that every peculiar body has its peculiar nature, and that the activity
under which bodies are formed, grouped, altered, is a thoroughgoing activity of nature; but nature itself in the unlike
natures, the many differing activities, stands before the beholder's thought as an unopened, enigmatic essence; the only
thing known with certainty about this unknown essence is that it works according to laws. ``We do not know what force is,
nor what light, what heat, what magnetism, what electricity is in itself''; those are the investigators' own words; but
those who acknowledge such a not-knowing must surely also, for consistency's sake, grant that they do not know what
nature in itself is.

It is the singular thing, the at once attractive and repellent thing, about the consideration of the body of the earth,
that this
% --- p. 340 ---
whole of nature, so rich in manifold formations, stands so near to us and yet so far, is so knowable and yet so
unknowable. It looks as if the whole outward nature lay open to the observer's gaze and gave itself straightforwardly as
it is in its facts, as if there were nothing at all behind; and yet there is something behind, something hidden, a secret
that always draws itself further back the deeper the investigator penetrates. In so far as geognostic knowledge is clear
about itself and its essential limit, there is in this no disappointment at all. The things of nature are all without
distinction objects for scientific cognition. There may be in the uppermost of the region of air, in the innermost of the
earth, in the bosom of the mountains and in the abysses of the sea manifold things of which no one, be he physicist or
philosopher, has yet dreamed; but the investigator knows that what has hitherto been unknown, however surprisingly new it
may be at first glance, must nevertheless always in one way or another let itself be ordered into series of the old, the
already known. Only when the boundary of experience is overstepped, when one expects of the advance of geognostic science
something which by its nature it neither can nor will yield, can disappointments set in and illusion begin its play. For
even if all the strata of the earth in all five continents were examined, turned over and rummaged through; nay, if the
most improbable thing of all were to happen, that a way opened into the earth's middle and the forces from the secret
workshops of darkness stepped forth into the light of day: we should nevertheless not come a hair's breadth further with
nature as nature. This doubly-sidedness --- that nature in itself is an inner unity of essence which indivisibly
penetrates everything, and yet in actuality a divisible essence, completely present in the connexion of the parts --- is
a riddle which geognosy, however much the masses of knowledge be increased and however strongly insight grow in empirical
compass, is unable to solve. Investigating consciousness continues to have existing nature
% --- p. 341 ---
outside itself as a completed object, an independent other, something known and akin, and yet a stranger of which it dare
not say: it is flesh of my flesh and bone of my bones --- it is power in another, elementary otherness.

% The Greek letter is read on the page itself (KB copy): β, not „6“ or „ß“, as the OCR gives.
% Run-in and letterspaced head.
\subrunin{β) The Aesthetic Standpoint.} The secret bond of kinship between the beholder of nature and the nature beheld
is first of all discovered by fantasy; but the discovery is coloured by the gleam of light in which consciousness at
different stages of development conceives its own essence.
% The footnote mark *) stands between „Schiller“ and the comma.
``When one,'' says Schiller\footnote{% MISPRINT in the note: „Wercke“ for „Werke“. Read on the
% KB copy, where the „c“ stands plainly; the errata leaf does not touch the place. Reproduced
% as printed.
Schillers sämmtliche Wercke. Zwölfter Band. Stuttgart und Tübingen 1838. pp.\ 187--88.}, ``recalls the beautiful nature
that surrounded the ancient Greeks; when one considers how familiarly this people could live with free nature under its
happy sky, how very near its mode of representation, its mode of feeling and its manners lay to that simple nature of
which its poetic works are so faithful a mirror, then one must be startled to remark that one finds so few traces of the
sentimental interest with which we moderns can cling to scenes and characters of nature. ``A. Humboldt does not seem
rightly to have understood Schiller's thought when, with regard to the passage adduced, he instructs us that these
individually excellent utterances can by no means be applied to all the peoples of antiquity, that ``a deep feeling for
nature utters itself in the oldest poems of the Hebrews and the Indians: hence among peoples of very different, Semitic
and Indo-Germanic, descent.''%
\renewcommand{\thefootnote}{**)}%
\footnote{% printed mark: **) --- second of two notes on p. 341.
Kosmos. II. p.\ 9.}%
\renewcommand{\thefootnote}{*)}
% The footnote mark **) stands AFTER the closing quotation mark.
% QUOTATION RECKONING, p. 341: five „ against two “. The house habit (cf. pp. 10, 11, 16)
% repeats „ at every interruption of the citation and closes only once: both „erindrer sig …“
% and the interpolated „A. Humboldt …“ stand open and are closed together at „Afstamning.“
% The last citation, „enten er Digteren Natur …“, runs over the leaf and is closed on p. 342.
% Reproduced as printed.
Schiller's fundamental thought is briefly and tellingly uttered thus: ``either the poet \emph{is} nature or he will
\emph{seek} it. The former gives us the naive, the latter the sentimental poet . . . . So long as man is still pure ---
not, be it understood, raw --- nature, he works as an undivided sensible unity and as a harmonising whole. Sense
% Both „er“ and „søge“ are letterspaced --- spacing.py does not catch the two-letter „er“.
% Both checked on the KB copy. The omission is set with FOUR spaced points.
% The leaf bears two notes, printed *) and **); the second is set with a local \renewcommand,
% cf. p. 12. Neither runs over the leaf.
% --- p. 342 ---
and reason, receptive and self-active faculty, have not yet parted in their office, still less do they stand in
contradiction with one another . . . . If man has entered into the state of culture, and art has laid its hand upon him,
then that \emph{sensible} harmony in him is cancelled, and can now utter itself only as \emph{moral} unity, that is, as
striving after unity. The agreement between his feeling and his thinking which in the first state \emph{actually}
obtained now exists merely \emph{ideally}; it is no longer in him but outside him, as a thought that has first to be
realised, no longer as the fact of his life.''\footnote{Work cited. pp.\ 196--98.}
% The footnote mark *) stands AFTER the closing quotation mark. Here the citation opened on
% p. 341 is closed; the quotation reckoning for the two leaves together is +2, owing to the
% two repeated „ on p. 341. The omission above is set with FOUR spaced points, as on pp. 336
% and 341.

If the thought uttered by Schiller be given a particular direction toward the doctrine of consciousness, then it will be
understood how in principle the matter stands with the Hindu's, the Greek's and the Hebrew's consideration of nature.

For the Hindu there is no essential difference between the conscious and the unconscious; subjective nature in the soul
and objective nature in the world of things are one and the same essence. The opposition between the self and the
selfless has not yet come to be, or, in so far as it is immediately felt, the self as the finite is reduced to being the
untrue, the evil, and the true, blessedness-bringing infinity is sought in the selfless. To sink oneself in oneself is to
lose oneself in the unbounded; all the particular appearances of the infinite One are intuited in fantastic symbolism as
changing forms of the universal life, dreamlike revelations of what stirs obscurely in the ground of the soul; no wonder
that at this stage of consciousness there can utter itself a sympathy with the enigmatic, the unconscious, which has a
deceptive likeness to the ambiguous pantheism of modern romanticism.

For the Greek the main thing is the essence of man in the harmoniously totalised form of contemplable individuality. If
% --- p. 343 ---
it is now man himself who is to be regarded as the nature that is in itself completed, filled with content within its
bounding, individualised, then it is also intelligible from this that the delineations of nature occurring in the Greek
poetic works, however full of fidelity, life and grace they may be, yield only the background or at most the middle
ground in those pictures whose foreground is formed by human actions, sufferings and passions.

For the Hebrew the opposition of nature and Spirit is both posited and cancelled. Nature is created; only to the
Creator's praise can the Hebrew be poetically inspired for nature. The Hebrew consideration of nature in antiquity and
the physical one of the present day form extremes so hard that at first glance it looks as if they were irreconcilable.
The Hebrew sees in the prominent events of nature nothing but miracles, visible signs of Spirit's unconditioned
self-power. The almighty will is the cause of all causes: Spirit commands, nature obeys; Spirit does everything, nature
nothing. For present-day physics the causes of nature are one and all, Spirit is gone; here are laws, but no lawgiver;
here is unity of connexion, but no unity of self; here is necessity, but no freedom: nature does everything, Spirit
nothing.

It will, however, according to what has gone before, be possible to see that the two extremes are in the
natural-philosophical sense complements of one another. In conceiving self-power as the cause of causes it must surely be
expressly understood that the almighty will goes through natural intermediate causes: the real causality is
indispensable to the ideal. In bringing out the causes of nature it must be expressly understood that the endless series
of causes all point back to the fundamental cause, and hence to the cause of causes, self-power: the ideal causality is
indispensable to the real. Not thought alone, but fantasy too makes it evident from many sides that nature is nature only
by means of Spirit.
% --- p. 344 ---
% The Greek letter is read on the page itself (KB copy): γ, not „y“ or „7“. Run-in and
% letterspaced head.
\subrunin{γ) The Teleological Standpoint.} The aesthetic consideration of nature is by no means so irrelevant to science
as, seen in passing, it might seem. Fantasy big with thought, however capriciously it may go to work in the particular,
is essentially bound to the fundamental law of objectivation; and the law says: \emph{there is nothing objective but what
is objectified.}
% The letterspacing runs over the whole sentence, and across the line-end as well
% („som objectiveres.“). Checked on the KB copy.

Selfhood and otherness, subjectivity and objectivity, the conscious and the unconscious are to all eternity correlative.
Take the self-being away, and there is no essence in another; for an essence that posits and represents itself by itself
is, on account of this its self-positing and self-representing, not an other-being but a self-being. Take subjectivity
away and then see what becomes of the objective. An object that is in a position to objectify itself for itself and for
other objects is more than object, it is a self-being, a subjectivity. Take consciousness away and consider what becomes
of the unconscious! It is one of the fundamental errors of the present day that the objective should be older than the
subjective, and the natural world rest upon a substratum of eternally selfless, unconscious being. How so great an error
has been able to enter upon science's way and remain standing as a dark, ingrown prejudice that shades the truth can be
briefly shown.

That the investigator must use his consciousness in every investigation goes without saying; but two things especially
contribute to dazzle him. In the first place he must see to it that the object can in each and all be presented in its
own light, so that it is neither distorted nor coloured by subjectivity; thus the opposition of consciousness and object
is sharpened to the utmost. In the second place the investigator has to do only with the objects of experience in their
finite relation to human consciousness. Since now it cannot occur to any rational person to make the
% --- p. 345 ---
actual existence of objects dependent upon human consciousness and to declare that there are no visible things when we
do not see them, none audible when we do not hear them, none tangible when we do not feel them, there is thereby
strengthened the in itself fundamentally false representation that the sensible in things can be without any sensation,
the thinkable without any thinking. At this standpoint nature is and remains a riddle.

Elementary nature is the selfless essence that is grounded in eternal selfhood, a content rich in thought but unconscious
of an infinitely self-thinking consciousness. Nature as nature cannot explain itself; it is immediate in its facts, and
the fact of facts is in turn the blind, ungrounded facticity. Only when it is fully recognised that nature, bearing,
forming, fruitful in formations, in all unconsciousness faithfully, that is, immediately, renders those plans, thoughts
and determinations of will which an infinitely mighty and rich consciousness has laid in its constitution, is it
recognised for what it in truth is: the teleologically real other-being whose end is set by a seeing self.

From an exclusively geognostic point of view the reflexion of the earthly whole in all its parts is outward, empty and
formal: as the whole is, so must the parts be, and conversely; further one does not come. From a teleological point of
view the development of the whole becomes the end that is set, and the parts means for the end. On the side of the facts,
the things of nature, nothing is hereby altered; the essential alteration concerns only the conception, the
understanding, the grounding: the teleological conception must go as one with the geological.

% Division A ends HERE, at the foot of p. 345, with a short centred rule; the head
% „B. Jordklodens historiske Udvikling.“ stands first on p. 346 and belongs to the next batch.
% The rule is checked on the KB copy. (The same shape as at p. 146; cf. by contrast pp. 87, 106.)
\streg

% Danish: \afdeling{B.}{Jordklodens historiske Udvikling.}   (printed pp. 346--361)
\afdeling{B.}{The Historical Development of the Earth.}

To have a history is, essentially understood, to have a determination of one's own, an end. When geology speaks of the
history of the globe, the representation of a progressive formation, a movement from lower to higher stages of
development, involuntarily presses itself forward. But the distinction between a lower and a higher would surely be
without all significance if there were no standard by which it could be determined. It is a representation of the end
that yields the standard. The body of the earth, the elemental whole, has the end in itself only in so far as it has it
outside itself, namely in the organisms; from this point of view the history of the globe falls in a natural way together
with the history of the organisms. With regard to this we may distinguish three comprehending sections: a
\emph{primeval period}, when the body of the earth took on its first rudimentary shape and prepared itself for a
production of organic beings; \emph{intermediate periods}, in which the body of the earth's particular transformations
keep step with corresponding formations of plant and animal forms; finally a \emph{present-day} or \emph{concluding
period}, in so far as the organic end may be said to be completely realised with man's appearance on the earth.

% Run-in head, printed in bold on the same line as the paragraph; the letterspacing in
% Urperiode, Mellemperioder, Nutids- and Slutningsperiode above is on the other hand genuine
% letterspacing. Both checked on the KB copy.
\runin{a) The Primeval Period.} Of the very first state of the globe physics can naturally give us only obscure,
indeterminate representations. In a natural-philosophical interpretation of the phenomena the emphasis must then fall
again upon the general opposition between real and ideal causality, and thereby upon the opposition between the physical
and the teleological conception; the task is to show how the first of these conceptions is completed in the
last.\footnote{% printed mark: *)
On grounding reflexion. Cf. The Logic of the Fundamental Ideas I. pp.\ 92--105.}

% --- p. 347 ---
So long as the beholder looks exclusively at the real causality, there can --- it lies in the nature of the case ---
show itself nothing but physical causes and corresponding effects: ends and means are nowhere to be discovered. The
succeeding formations are results of the preceding, but only accidentally necessary results, not intended advances, not
realised purposes.

According to the cosmogonic hypothesis the moon must have been formed out of the mass of the earth in the same way as
this out of the mass of the sun. From the moon's formation as represented result the consideration goes back to the
representation of the necessary presuppositions. The earth pregnant with the moon must have had another volume than it
now has; for while the globe's radius is at present only 860 miles, it must at that time have been about 60,000 miles.
% The figures are in Danish miles (the „Miil“ of the print, about 7.5 km), not English ones.
Before the moon was born, the mother earth must then also have turned rather slowly about its axis: the day, which is now
only 24 hours, must then have amounted to 27 days. Various things may further be adduced in explanation of the earth's
% The printed mark stands BEFORE the comma: „Fladtrykning*),“. Rendered thus.
flattening\footnote{% printed mark: *) --- first of two notes on p. 347.
% „Plateaus's“ with both an s and an apostrophe-s, so printed in both copies.
Plateaus's experiment with the drop of oil.},
condensation, increased
\renewcommand{\thefootnote}{**)}%
velocity of rotation\footnote{% printed mark: **) --- second of two notes on p. 347.
Cf. for example Zimmermann. Die Wunder der Urwelt. Berlin 1855. pp.\ 26--33.}%
\renewcommand{\thefootnote}{*)}
{} and so forth; but with all this the consideration moves in a circle of represented phenomena without essential
grounding. The whole grounding comes to adjusting represented presuppositions to a merely represented result; for the
moon's formation out of the earth is after all, for the beholder, only a represented and not an actual event. But even if
we grant that the event has such probability in its favour that it may be regarded as fact, is any advance really
thereby designated? By no means. If
% --- p. 348 ---
in this connexion there is talk of an advance of formation, then it is the beholder who is slipping in his own
subjective representation. The only thing that from the standpoint of real causality can be seen in all events, all
formations and all changes is the fulfilment of ``the unalterable laws''; but that with regard to the fulfilment of the
laws no advance takes place is on physics' own side incontestable: in the earth as rotating sphere of vapour the laws
were just as completely fulfilled \emph{before} as \emph{after} the moon's formation.

If we consider the changes which according to the geognostic hypothesis must have taken place with the body of the earth
in the millions of millennia that the primeval period has swallowed up, then the representation of gradual formations has
undeniably a quantity of indications to support it. It has probability in its favour that the body of the earth, after
having been vaporous, may have passed over into being drop-fluid and at last into being fire-fluid. At this last stage
geology links itself to geognosy. ``The earth was accordingly once a fire-fluid mass, whereby it is nevertheless not
denied that it may already then have been surrounded by an aeriform atmosphere; but this must under all circumstances
have had an essentially different composition from the present one, since many substances that are now solid or fluid
can, at the higher temperature, have been present only in aeriform state . . . Since, however, the space of the world
possesses a heat that is indeterminable but in any case very slight --- doubtless below
% The sign before 50 is Nielsen's MINUS (÷), read on the KB copy; the sign before 2000 is an
% ordinary plus. The degree marks stand as a raised o on the page.
$\div\,50^{\circ}$ --- the fire-fluid globe, whose heat, to judge by the substances known from its surface, certainly
exceeded $+\,2000^{\circ}$, had constantly to radiate heat into the space of the world and thus to become colder and
colder. A solid crust began to form on its surface, like a covering of ice on water; but since ebb and flood, proceeding
from the moon's and the sun's attraction, called forth constant movements in the
% --- p. 349 ---
completely fluid mass of the earth, this formation of a crust could not go on without continual and great disturbances.
Many a time, no doubt, the thin covering was burst asunder; the loosened flakes sought again and again to unite, but for
a long time were unable to form any lasting crust, until the earth's constantly advancing cooling at last created a
covering so mighty, consisting of stone flakes baked together, that this could at least no longer be burst into pieces,
though it might well crack, now at one place, now at another.''

``A solid crust now encloses the fire-fluid kernel and is itself in turn surrounded by an aeriform atmosphere; but water
there is not yet. It could first form when the cooling of the earth's surface had advanced so far that its heat exceeded
the boiling point at the then presumably much higher pressure of air.''\footnote{% printed mark: *)
Cotta. Wk.\ cit. pp.\ 11--13.}
% The note abbreviates „Anfr. Skr.“; on p. 356 it is „Anf. Skr.“ The book uses both forms;
% neither is corrected.

By following the description in detail and observing in particular how the earth's crust, as a first rough draft, is
constantly transformed and enriched, in that it is broken through by eruptive masses --- granites, porphyries,
greenstones and so forth --- which differ in composition and specific gravity according as they come from a different
depth; by seeing how the waters, under their destructive influence upon the constituents of the surface, work together
upon further formations; how the masses dissolved in the waters and thereafter deposited increase the crust's thickness
here and there, by this increase hinder the radiation of heat and bring it about that a part of the crust is melted again
from below, so that by the heating there are formed crystalline schists, gneiss, mica schist, hornblende schist and the
like; by sinking ourselves, as has been said, in the consideration of these first formations, we involuntarily get the
impression that something is here intended by all this; it is as
% --- p. 350 ---
though we were witnesses to a beginning execution of a grand plan of building; as though we had a glimpse into obscure
workshops where an invisible master is at work shaping the material, covering over the abysses, casting beams of stone
and building mountains. But it is only a fantasy; nothing is here prepared, for nothing is here intended; the whole is a
series of changes in which the laws are fulfilled: it all goes on much as when a covering of ice forms on water in a
gale.

If it is to be possible to recognise an earthly nature in progressive formations of the earth; if the globe is with right
to be said to have a history, then the real causality must be reflected in the ideal, object-power in self-power, the
objectivity of nature's being in the divine subjectivity of Spirit, and the conception must become teleological. If a
scientific thinking can make it possible for the physicist to determine the given facts in objective purity and to remove
every admixture of arbitrary opinions, why should a scientific thinking not also make it possible for the philosophy of
nature to keep away the finite representations of the Creator's merely supposed intentions and to recognise the
subjectivity that grounds everything in the purity of the concept?

When it is fully recognised that the globe, as a whole of nature, images in existence a concept of wholeness determining
from within, the teleological way of seeing is scientifically secured. For in this lies: 1) that the concept, peculiarly
determined and expressly willed by comprehending self-power, is with its whole content the indwelling constitution of the
rotating mass; 2) that the concept's successive actualisation, the constitution's gradual development, satisfies a demand
of selfhood, so that there can here with ground be talk of advance; 3) that every result of formation by which an advance
in the whole's development is attained must be determined as an intended goal, a relative end.

From this it cannot indeed be derived that the moon must
% --- p. 351 ---
have been formed out of a part of the mass of the earth; but it can be concluded with certainty that if the formation of
the moon has taken place in a way answering to the hypothesis, then this formation has been preformed in the globe's
constitution; the actual formation is then more than a physical result: it is a purpose attained. The material
presuppositions give only possibilities, and the mere possibilities, when they are actualised under the conditions that
obtain, give only results; but the constitution, which takes up presuppositions, possibilities and conditions into the
preforming unity of ideal causality, gives ends, for in the constitution there is something intended. It must above all
be held fast that it is not the outward events of nature but their grounding that is here in question. Let one represent
the formation of the moon as one will: the fundamental question about the event's inner nature remains constantly the
same. If the moon's coming-to-be be conceived solely as an effect of physical causes, then one may introduce as many
presuppositions, conditions and supervening circumstances as one pleases: the grounding is and remains formal,
superficial, imperfect, for the essential ground, the ground of nature itself, is wanting. The manifoldness of the
physical causes belongs to the parcelling out, and in the parcelling out accidentality rules. It then comes to look like
a pure accident that the earth has got a moon, while Venus, which is nearer the sun, and Mars, which is further from it,
are moonless. The immediate appeal to the accidentality of the natural facts has its answering extreme in an equally
immediate appeal to the Creator's arbitrariness; both ways of seeing are equally warranted and equally insufficient.
% The working copy bears here two pencil marks (before „Man“ and before „Kun“); the KB copy has
% no sign at either place. They are not text.
Only in the union of the two --- a union that is not a finite composition in the beholder's imagination but an inner
union in the constitution itself --- have we \emph{the sufficient ground} (\textit{ratio sufficiens}). In object-power,
and hence from the side of nature according to the outward, the moon's formation is a physically
% --- p. 352 ---
conditioned result of physically conditioning presuppositions, \emph{because} in self-power's sense, and hence from the
side of Spirit according to the inward, it is the realisation of a purpose, a teleological event that satisfies a demand
of selfhood.

The sufficient ground is doubly-sided; it is from the objective side real ground, from the subjective side motive: the
unity of motive and real ground is precisely the constitution. The conception of the globe's historical development can
accordingly be missed in two different ways. If one overlooks the real grounds while the motives are sought, then the
facts and the physical causes sink to a fantastic apparatus of a fleeting business: geological myths and home-made
theories of intention flourish. If one overlooks the motives while one digs everywhere in real grounds and rummages in
conditions, then one does indeed get an actually coherent whole; but in the whole's connexion there is no scientific
meaning, for meaning and motive are inseparable categories.

One opens the archives of the primeval world, considers the old stone monuments, speaks of picture-writing, of
hieroglyphs whose meaning it is a question of interpreting. What then is the meaning of these gigantic, cyclopean walls,
these enormous blocks of granite, this piling up of eurites, porphyries, gneisses, greywackes and whatever else they are
called? ``Do not ask about that, but investigate! Nature itself has no opinion about such things, and on men's opinion
nothing whatever depends.'' But if there is no
% --- p. 353 ---
meaning in the obscure writing, how then will one be able to interpret it? Investigate, investigate, for the key lies
hidden in the connexion! Why was the glowing earth wrapped as in a mantle of stone with many folds? Because the layers in
its surface were cooled by radiation of heat, and eruptive masses broke through the tilted strata. Why has it got a sea?
Because in the surrounding space of gas there were such great quantities of oxygen and hydrogen as entered into chemical
combination. Why is there in the earth's crust so extraordinary a diffusion of silicates, carbonates, alkalies and so
forth? Because the corresponding elementary substances have all been present beforehand in the rotating sphere of
vapour. It is the Why of the real grounds that constantly repeats itself. But the mere real grounds are only half
grounds, as is seen at once when one asks after the first essential presuppositions. Why has the mass of the earth
originally contained the many and differing substances? Here the grounding stands still. ``According to the Indian
primeval myth the earth is borne by an elephant; this is itself borne, that it may not fall, by a giant tortoise; but
whereon this tortoise rests, the believing Brahmins are forbidden to
% The printed mark stands between the quotation mark and the full stop: „spørge“*).
ask''\footnote{% printed mark: *)
Kosmos I. p.\ 251.}. In the same way the investigating empiricist is forbidden to ask wherein the elementary constituents
that lie at the ground of the geological formations are themselves grounded. --- Only when the result is bent back and
becomes the ground of the presupposition can the grounding be completed. There has then not merely formed a sea on the
earth because there was oxygen and hydrogen in the mass of the earth; but there has originally been oxygen and hydrogen
because a sea was to be formed. By this turn the result has become motive, and its presupposition has gone in under the
constitution.

In an interpretation of the motives of formation it may in many cases be difficult enough to distinguish between what is
means and what is end; for what in one respect is end may in many other and very different respects be means. It is
always of importance to hold to the facts and to penetrate into their connexion. But when the very most essential
categories of grounding --- constitution and motive, end and means --- are either overlooked or even slighted, it will be
impossible for the investigator to make a scientific distinction between side issue and main issue, between essential and
inessential; by mere real grounds everything becomes,
% --- p. 354 ---
for the sake of the connexion, equally essential and just therefore equally inessential. The moment the investigator
declares anything whatever --- be it a substratum of masses of stone, the formation of a sea, the shooting up of
mountains, the extension of mainland, the precipitation of lime, the purification of the atmosphere, whatever it may be
--- to be important and significant, he slips in a meaning and has a
% The printed mark stands BEFORE the full stop: „Motiv*).“
motive\footnote{% printed mark: *)
Cf. The Logic of the Fundamental Ideas II. p.\ 429 ff.}. The meaning is that something must be lower in order that
something else may be higher, something a means that something else may be an end, something a motive that something else
may be a real ground: without constitution and motives, without means and ends it is meaningless to speak of ``nature's
housekeeping'', impossible to interpret the obscure writing in the folded strata, impossible to find meaning in the
geognostic phenomena of the primeval world.

% Run-in head, printed in bold on the same line as the paragraph, read on the KB copy.
\runin{b) Intermediate Periods.} The development of the body of the earth is, like all development in nature, seen
originally, a development of constitution. What and how much such a constitution contains must be seen in facts from what
is developed. When then, in the course of the development, there are brought forth not merely strata of earth and masses
of water but also organic beings, then the conditions for the coming forth of the organic beings must also belong to the
constitution. If now, as we brought out above, the whole idea of the globe is enclosed in its constitution, then the
historical development shows that the idea of life must from the beginning be enclosed in the idea of the globe. The
first general answer to the great question how organic life has come forth on the earth must then be that it has come
forth from its peculiar constitution, precisely at that turning point when the earth's development had reached so far
that elementary nature could yield the first necessary conditions.

In what compass the organic activity has in its turn had an influence upon the articulation of certain strata of earth
and species of stone is established by Ehrenberg's
% --- p. 355 ---
% The Danish breaks mid-word at the foot of p. 354: „Op-“ / „dagelse“.
discovery of the first infusoria. These ``almost invisible organisms have existed already in the oldest periods of the
earth, ever since there was water on the earth . . . Whole mountains, nay the substrata of whole countries, consist
principally of the left corpses and skeletons of small extinct beings, of which billions can be contained in a cubic
inch, but which under favourable circumstances multiply so rapidly that their number is doubled in every hour that they
live.''\footnote{% printed mark: *) --- first of two notes on p. 355.
Cotta. Geol. Pictures. p.\ 103.}

The dissecting study of the animal and plant life of the fore-world has a double direction. ``The one is purely
morphological and is directed principally to
% The errata leaf corrects p. 355, line 11 from the top: „Organismens --- read Organismernes“.
% By the house rule the text stands AS PRINTED: Organismens (the singular). The reading is
% checked on the KB copy, which also prints Organismens. The working copy bears at the place a
% pencil mark (an angle after „er“ in the line above); it is not text. The English renders the
% sense, the misprint being noted here.
the organisms' natural description and physiology; by the perished formations it fills out the gaps in the series of the
more recent period. The other direction is purely geognostic and considers the fossil remains in their relation to the
superposition and relative
\renewcommand{\thefootnote}{**)}%
age of the sedimentary formation.''\footnote{% printed mark: **) --- second of two notes on p. 355.
Kosmos. I. p.\ 230.}%
\renewcommand{\thefootnote}{*)}

If we gather under a general survey the systems, groups and formations --- the greywacke group, the coal group, the
Zechstein group, the Trias group, the Jurassic, Cretaceous and Molasse groups --- which geology has for the present
ordered into great intermediate periods, then there are three marks that, in spite of the shifting of hypotheses and
doubts about particulars, raise themselves with unmistakable clarity and decisive definiteness. 1) Since the body of the
earth yields the conditions for organic life, the organisms in each peculiar period must necessarily presuppose a
peculiar stage of formation characteristic of the earth's development; land plants, for example, could not come forth
before masses of land had formed; lung-breathing, warm-blooded animals, which do not endure an atmosphere rich in
carbonic acid, could not come forth before the atmosphere had been purified, and so forth. 2) There is a characteristic
thoroughgoing difference between the plant and animal forms that
% --- p. 356 ---
essentially belong to the different periods of the earth. It begins with unicellular plants and cryptogams, fucoids,
algae in the greywacke's Devonian strata; calamites and lycopodiaceae, ferns, phanerogamous monocotyledons (grasses,
lilies, palms), conifers and cycads appear in the coal period, maples and poplars in the Tertiary period, and so on.
Still more striking are the animal world's typical changes of form. Upon the lowest sea animals, radiates, molluscs
(brachiopods, cephalopods) follow the articulates (trilobites), and already in the greywacke time the vertebrates begin
with the oldest fish forms (Cephalaspis, Pterichthys). In the Zechstein period the fishes, in the Trias group the
reptiles, and in the Jurassic period especially the crocodile-like animals (Ichthyosaurus, Plesiosaurus) and the flying
lizards have evidently the upper hand. Iguanodon shows itself at the beginning of the Cretaceous group. ``In the so-called
Tertiary time, \emph{the Molasse time}, the mammals already ruled on the earth; but as their forerunners, the marsupial
genera of the Jurassic time, belonged to a low-standing and, remarkably enough, declining order of animals, so now too
the more lowly organised herbivores appear first, and only after them the beasts of prey, which stand higher in life as
% The printed mark stands BEFORE the semicolon: „Rovdyr*);“.
in the system\footnote{% printed mark: *) --- first of two notes on p. 356.
There are, however, indications that the beasts of prey appeared very early.}; at the close of the Tertiary time, finally,
% The quotation mark stands BEFORE the full stop, and the note mark after it:
% „Aberne“.**)“ --- rendered as printed.
\renewcommand{\thefootnote}{**)}%
the apes''.\footnote{% printed mark: **) --- second of two notes on p. 356.
Cotta. Work cited. p.\ 250.}%
\renewcommand{\thefootnote}{*)}
{} 3) ``One believes,'' it is further said, ``that one has observed that the different classes, orders and families of
organic beings always begin with such forms as belong to none of the now existing subdivisions, but on the contrary often
require an extension of the main division's
% MISPRINT: „Charak-erer“, divided at the line-end, hence „Charakerer“ without the t. Both
% copies print it so (KB checked), so it is the compositor's fault and not the copy's.
% Rendered here by the sense; read „Charakterer“.
characters, as these have once for all been set up for the organisms now living. These fundamental forms often constitute
as it were the germ of a number of kindred series of forms which have developed out of it and have constantly attained
greater and greater
% --- p. 357 ---
% THE FOLIO IS MISPRINTED: the leaf bears „257“, not „357“. Its neighbours print 355, 356,
% [257], 358, 359, and the KB copy bears the same fault, so it is the compositor's and not the
% digitisation's. The page is marked here by its true place in the series.
independence, so that their outermost members are now very different. We see the fishes thus begin with some forms so
unlike the other kinds of fish that for a time they were taken for crabs or for large water beetles. Among the oldest
cold-blooded bony animals belong the labyrinthodonts, which were regarded now as lizards transformed into frogs, now as
frogs with lizard characters grafted on. The birds begin with a group of which traces of seven-foot strides have been
found in the sandstone of Massachusetts, according to which they must then have been twice as large as our largest
ostriches. Among the oldest remains of mammals belong, besides the marsupial remains found in the English Jurassic
formation, also the so-called Hydrarchos (Zeuglodon), a whale-like mammal of enormous length and most deviant form,
wherefore it was also long taken for a sea serpent or a crocodile-like animal. The class of corals begins in like manner
with the graptolites, that of the echinoderms with the cystideans and crinoids, that of the cuttlefishes with the
% MISPRINT: „Triboliterne“ with b before l; the book writes „Trilobiter“ on the preceding page
% (p. 356). Both copies print „Triboliterne“ (KB checked), so it is the compositor's
% transposition. Rendered here by the sense; read „Trilobiterne“.
orthoceratites, that of the crabs with the trilobites, that of the worms with the brachiopod genera.''

It will not, even from an abstractly objective geognostic standpoint, be easy to guard oneself against the impression
that in these periodic developments of plant and animal forms there is an end and an advance. By lower and higher
standing organisms one surely understands not merely less or more composite, less or more articulated forms; one
understands at the same time that specimens and individuals with richer organisation must stand at a higher stage,
because they satisfy in richer measure the fundamental demands of upward-striving life; in short, one understands the
teleological. The fishes stand lower than the reptiles, these in turn lower than the mammals; if now the fishes rule in
the Zechstein period (the Permian system), the reptiles in the Trias and
% The errata leaf corrects p. 357, line 2 from the foot: „Perioder --- read Periode“. By the
% house rule the text stands AS PRINTED: the plural. The reading is checked on the KB copy,
% which prints „Perioder“; in the working copy the final r is struck through in pencil, and
% tesseract therefore reads the word wrongly. The English reproduces the plural.
Jurassic periods, the mammals in the Tertiary periods (the Molasse group), then this says in other words that every
% --- p. 358 ---
succeeding period is, with regard to the end, on the whole higher than the preceding; only so can there with ground be
talk of ``the history of organic life on the earth'', and then also with ground of the globe's corresponding history of
development. None the less the geologists find themselves, as regards this point, in a certain embarrassment with their
own presuppositions; and it is not difficult to show whence the embarrassment comes.

It is a peculiar double gleam --- of existence and concept, of parcelled-out causes and ideal causality, of hypothetical
necessity and an overreaching freedom shining through plan and constitution --- that dazzles the eye. When the physicist,
under the whole impression of the mutual agreement of the manifold inorganic and organic formations, assumes that in all
this there is a plan, an end, a meaning, he does not notice that by such a way of considering he springs off from his
original standpoint and places himself at the centre of the idea. Only from the centre of the idea is the teleological to
be discovered. When the two diametrically opposed ways of seeing are run together, so that it looks now as if the lime
were precipitated in the sea for the organisms' sake, masses of land raised for the land animals' sake and the air
purified for the lungs' sake; now again as if the sea were sea, the mainland land and the air air without anything being
thereby intended, then one speaks with two tongues, and the two-tongued language often betrays in a quite naive way how
concepts that are equally warranted but of unlike origin clash with
% A NOTE OVER THREE LEAVES: the only note on p. 358 runs on at the foot of p. 359 and is
% completed at the foot of p. 360. The mark is not repeated on the continuations. By the house
% rule the note stands WHOLE here at its mark.
% The letterspacing in „warum“ (only the FIRST of the three occurrences), in „den absolute
% Nødvendighed“ and in the first „Warum“ is checked on the KB copy; the last two „warum“ are
% set in ordinary type.
one another.\footnote{% printed mark: *)
Examples of that kind of clash of concepts occur in almost every geology, but especially in half-popular writings. We
have in an earlier writing (Lectures on Philosophical Propaedeutic, 1860--61, p. 263) already singled out a treatise by
Dr. Burmeister, Vergangenheit und Gegenwart des Thierreichs (Geol. Bilder. 1ster Band. Leipzig 1851), as a grand model in
this respect. After the author has at the very beginning (p. 146), as regards the earth's development, referred to ``die
unwandelbare Nothwendigkeit'', and has further let the animals' fundamental types (p. 150) belong ``mit Nothwendigkeit
zur organischen Natur der Erde'', while these types' more or less perfect ``Darstellungsart'' is on the contrary ``den
Umständen preisgegeben'', he then again considers, by way of a closer understanding of what it really means to be ``den
Umständen preisgegeben'', on the same page ``die gesammten Phänomene der Genesis'' --- ``unter dem allein richtigen
Gesichtspunkte der absoluten Nothwendigkeit.'' Under this point of view the teleological way of considering nevertheless
raises itself now and then so involuntarily and with such insistence that it not seldom takes the power from simple
necessity; thus (p. 207): ``Alle beständigen Wasserbewohner, namentlich die im Meere lebenden, haben kurze Wirbelkörper
mit concaven Berührungsflächen (--- \emph{warum}?), weil eine solche Anlage der Wirbelsäule die schlängelnden Bewegungen
des Rumpfes beim Schwimmen am leichtesten ausführbar macht.'' Further (p. 213), where the author, after touching on the
probability that the fore-world's Iguanodon lived on plant food, continues: ``Demnächst ist es nicht unwahrscheinlich,
dass die plumpen Dinosaurier auf eine entsprechende Nahrung angewiesen waren (--- warum?), denn kaum möchte für so grosze
Geschöpfe eine ausreichende Fleischnahrung in jener Zeit vorhanden gewesen sein, wenn sie nicht von Fischen leben
sollten. Dagegen spricht aber ihr plumper Bau . . . . Grosse Raubthiere müssen zwar kräftig, aber nicht schwerfällig
gebaut sein, (--- warum?), weil zum Beutemachen Gewandtheit und Schnelligkeit ebenso sehr, wie Kraft, vonnöthen ist.'' It
might indeed seem as though the author with all this inclined more to a grounding by force of \emph{absolute necessity}
than by force of purposiveness's ``\emph{Warum}''; but the equilibrium is brought about in the finest way, in that after
a many-sided grounding and manifold answering of a manifold ``Warum'' he ends by breaking the staff over all scientific
grounding and confiscating without distinction every ``Warum''. For we are warned at the close in the most emphatic
manner against, as regards ``die Thatsachen'', ``die Gründe ihres Entstehens zu beurtheilen''; for it is declared outright
(pp. 243--44): ``Alle wissenschaftliche Forschung lehrt nichts mehr als den Gang kennen, welchen die organische
Entwickelung genommen hat; nicht das Warum desselben, nach dem so viele Neugierige und Unverständige täppisch fragen.''
--- ``Das Warum liegt überall ausser der Sphäre des Begreiflichen.''}

% --- p. 359 ---
% Run-in head, printed in bold on the same line as the paragraph, read on the KB copy. Only
% two lines of body text stand on p. 359; the rest of the page is taken up by the continuation
% of the note from p. 358, separated by the usual note rule.
\runin{c) The Concluding Period.} Whether it really will do to designate the stage at which the earth's development at
present
% --- p. 360 ---
finds itself as a concluding period; whether we must not rather regard it as a passing present-day period and be prepared
for this state too, after the lapse of millions of years, dissolving and sinking down into the series of the intermediate
periods --- that is a question which the philosophy of nature shall not be hasty in answering; it must in any case first
take counsel with research and listen to geology's pronouncements. ``The keystone in the whole building of the organic
kingdom is finally formed by man, who first appeared in the most recent geological period as creation's flower and the
present's ruler. How far back in time his origin lies is not known, but it certainly lies further back than history has
hitherto been willing to grant. How many millennia did not pass before the human race had come so far that it could think
of building temples and pyramids? Geology teaches us, however, that nothing on earth is lasting except the laws of
nature, and least of all the single organic individuals. The species die like the individuals; even the genera and
families disappear and yield their place to other conditions and to a new organic life. Is this a consequence of the age
of the species, the genus, the family, so that time as it were consumes them, or is it a consequence of altered
conditions of life on the earth's surface? --- Whatever the ground may be, we can expect similar changes in the future.
The species now living will also be replaced
% --- p. 361 ---
by new ones, probably more highly organised, and even man dare not claim to form an exception to this rule. Is he to
expect a more gifted successor, who after him shall be the earth's lord and form the last member of the organic series of
development? And when shall this happen? That lies altogether outside the bounds of our knowledge; but analogy allows us
to surmise that it will not happen within the next hundred thousand years.''\footnote{% printed mark: *)
Cotta. Geol. Pictures. pp.\ 251--52.}

So long as the fundamental relation between the empirically physical and the a priori teleological way of seeing still
stands in complete indeterminacy, it will be impossible to find out what really lies enclosed in the above consideration
and to what it leads. That nothing on earth is lasting but the laws of nature is certain: everything, the living even
more than the lifeless, is subject to vicissitude. The next thing we are referred to is analogy; but there is in this
something so indeterminate that judging subjectivity can get the most opposed results out of it. If one strictly keeps
away every accessory representation of plan and purpose, then the series of changes will be continued unceasingly without
regard to whether the development goes upward or downward. If one takes up a reflex of the teleological and holds fast
the difference between the \emph{lower} and the \emph{higher} organised as essential, then the explanation points in
several directions. Straightforwardly according to analogy one would surely assume that the whole advanced toward ever
greater perfection. With this it will then agree that not only the animal species now living but also the humanity now
existing must die away and be succeeded by newer, higher representatives. The ideal which the human race has by a mistake
placed at the beginning of its history and thereby conceived backwards, namely the ideal of paradise or of blessedness,
may then perhaps still come to honour and dignity, when

%  --- C. Jordlegemet som elementær Organisme  (printed pp. 362--374) ---
% --- p. 362 ---
% DIVISION CHANGE IN MID-PAGE. p. 362 opens FLUSH (without indentation) and continues without
% paragraph break the sentence running from p. 361; the word „det“ is whole, so there is no
% word division across the seam. Division B ends about two thirds down with „… Organismernes
% Væsen.“, whereupon follows a short centred rule, then „C.“ and the head; division C therefore
% begins on the same leaf. Read on the KB copy.
it is set in front and its actualisation on earth placed in an infinitely distant future, one not indeed reached by this
humanity, but still --- be that our comfort --- approximately by a succeeding one. The hopeful view becomes, however,
somewhat darkened when one looks back; for why does the movement advance into infinity? Why will the one humanity come to
be succeeded by the other? Because none of the many expected humanities reaches the goal; because Kronos constantly
continues to devour his own children; in other words, because the fantastic teleology cancels itself. If one does not let
the goal remove itself into infinity, but sets, according to the above indication, ``a more gifted successor'' of the
present humanity ``to be the earth's lord and form the last member of the organic series of development'', then analogy
is surely broken, if it ends with something after all coming to be at last that is ``lasting on the earth'', namely man's
``more gifted successor''. If on the other hand the last humanity is only to represent ``the last'' but also vanishing
member, in so far as it represents not the perfect but a more perfect that is nevertheless imperfect, why then can the
human race now existing, for all its imperfection, not form the last vanishing member? One sees that the geological
consideration, however many facts and analogies it may have in its favour, is wavering so long as one is not agreed with
oneself about the real goal of the history of the earth's development: the inner indwelling goal must be sought in the
essence of the organisms.

% A short centred rule between division B and division C, in the middle of p. 362. Checked on
% both copies.
\streg

% Danish: \afdeling{C.}{Jordlegemet som elementær Organisme.}   (printed pp. 362--374)
\afdeling{C.}{The Body of the Earth as Elemental Organism.}

In order to make clear what the in itself ambiguous concept \emph{an elemental organism} really has to signify, we first
direct
% --- p. 363 ---
% The Danish breaks mid-word here: p. 362 ends „hen-“, p. 363 begins „vende“.
attention to the finite opposition of unity of self and unity of connexion, which lets itself be seen particularly in the
representation of an \emph{organising of elements}. Next, since every whole organised out of elements in the form of
otherness involuntarily, on account of the ambiguity, makes the impression of a self-organising self-being, we note what
it means that the body of the earth is \emph{a selfless organism}. The contradictions that hereby come forth will then
resolve themselves into a physically critical separation of inorganic and organic, and the resolution will contribute to
setting \emph{the inorganic organism} in full light.

% The letterspacing on this page read on the KB copy: in „Forvirring og Uorden“ the „og“ is NOT
% letterspaced, whereas the „og“ in „Orden og Plan“ IS; hence two \emph in the first case and
% one in the second.
\runin{a) The Organising of Elements.} To organise a given material is, according to ordinary usage, to bring order into
an unordered material, an aggregate of unlike, accidentally mixed constituents. What is disorderly in the material
cannot, however, be determined abstractly and objectively as something confused in and by itself. \emph{Confusion} and
\emph{disorder} are negative categories and can be determined only by their express opposite, \emph{order and plan},
which must be fixed subjectively. Here there are accordingly presupposed: from the side of selfhood, the subjective
moment, a concept according to which ordering is done, and an ordering capacity for activity; from the side of otherness,
the objective moment, such a constitution in the existing material as makes it, on account of the possibilities dwelling
in the substances, let itself be transformed, formed and arranged according to the demand of selfhood uttered in the
plan. Where the organising is as outward as in an arbitrary mechanical production, the opposition between the determining
concept in the unity of self --- the manufacturer's plan --- and the form and combination of the parts --- the unity of
connexion --- is seen quite plainly in the manufactured article itself. As regards this kind of organising it is
incontestable that all characteristic determinations in the
% --- p. 364 ---
unity of connexion must derive from arbitrary determinations in the unity of self. It is otherwise when it is a whole of
nature that is organised. Here arbitrariness disappears, and with it the representation of a will ordering according to
purposes; here there can be no talk of an outside master workman who should force a predetermined form upon the material:
the forming takes place in that the constituents mutually order and form one another. Since the organising rests upon the
reciprocal action of the parts, it can then easily take on the appearance that the unity of self is superfluous, that the
ordering rests simply upon the connexion. Where there is connexion, one says, there is order, for there are always parts
which by the connexion constitute a whole. As the parts are, so is the whole; therein consists the eternal order of
things. How one-sided and false this way of seeing is may be seen from this, that in this way it becomes impossible to
make a difference between order and disorder. It is not even possible to represent a chaos to oneself without
representing some sort of reciprocal action and connexion between the wildly mixed parts. All actual order presupposes an
ordering principle; but an ordering principle is unthinkable without unity of self. Even if one went so far as to take
any chaotic reciprocal action of elements whatever for organising, one would still not be rid of the unity of self; for
reciprocal action, let us repeat it, rests upon relation, relation upon reflexion, reflexion upon reflecting selfhood.

% „Timæus“ stands in upright roman like the rest of the text, not in italic; read on the KB copy.
There is in this respect too something to be learned from Plato. When in his Timaeus he lets the world-builder consider
with himself how order may best be brought into the unordered, into chaos, lets him implant a soul in the world, stretch
out and divide the soul before he forms the elements and so forth, then Aristotle is not merely wrong in treating the
mythical images as philosophical expressions of concepts and applying the measure of the understanding's criticism to the
particulars
% --- p. 365 ---
% The Danish breaks mid-word here: p. 364 ends „Enkelt-“, p. 365 begins „hederne“.
of this organising; but he is besides, and that is the worst, far more wrong in so far as he overlooks the very most
essential point in Plato's great fundamental vision. What Plato wants to bring out is plainly the truth that the
immediately natural, the matter-like for itself, is without reason; that reason and order can not possibly come into the
dark, disorderly possibilities of the elementary ground unless it comes from the Godhead, the eternal self with the clear
wisdom that orders the whole according to ideas. This in itself fully valid presupposition Plato has not been in a
position to develop in rational form or to unite with the fundamental thought of an original, ordering activity in the
ground of nature itself. And what then has Aristotle done? After a rather easy victory over childishly pictorial
indications he has without more ado drawn a line through the principle of selfhood; instead of a cognition of the active,
ordering, organising Godhead he has delivered a metaphysical description of form and matter, and without more ado
presupposed the order of nature that he was to ground. According to Aristotle it is the form of nature, organising
nature, that organises nature; it is a tautology which at bottom says nothing else than that nature is organised because
it is organising, and organising because it is organised: \emph{nature organises itself.}

All organising activity is teleological; that is why it goes so easily with Aristotle, and after
him with all speculative philosophers of evolution, to make the whole of nature teleological. Even if one does not openly
profess the doctrine of a universal world-soul, which Aristotle does not do either, one nevertheless constantly slips in
the accessory representation that self-organising nature is a self-being. This mystification is in equal degree confusing
for concept and for fact. That the organising principle must proceed from a self-being is incontestable; but from that it
% --- p. 366 ---
by no means follows that what is organised must also be a self-being. The body of the earth is no self-being, and yet
the body of the earth is organised.

That an organised unity of connexion must always give a reflected gleam of unity of self is intelligible, since the
concept of organising proceeds from selfhood; but that this gleam can even be very strong without therefore dazzling the
eye, we learn from manufactured articles. It occurs to no one to make a pocket watch into a self-being because it keeps
an even course, because its parts co-operate toward a certain end, and the hand often knows better about the time than
the owner. That the watch has not, like a whole of nature, organised itself makes from this point of view nothing to the
matter. Nature's organising activity is just as certainly subject to an ordering, governing concept as are the formations
owed to human invention. In another way, to be sure. But the physical in the way is here by no means the decisive thing.
The striking difference between manufactured articles and products of nature has brought it about that philosophers have
all too easily sprung off from the unity of principle. Behind the manufactured article stands a manufacturer, therefore
the article is teleological; behind the product of nature stands no manufacturer; and with that one is finished. If the
product of nature, the organised whole, is none the less to be conceived teleologically, then one must either with Kant
declare the teleological principle a merely regulative and not a constitutive principle, make a critical separation
between the mechanical and the organic, and restrict the teleological to the living organisms; or else with the
philosophers of evolution make the teleological principle constitutive, break down the partition between the inorganic
and the organic, and regard elementary, self-organising nature as a self-being.

That the Kantian separation between the mechanical and the organic is arbitrary can easily be seen. How could one ground
the assertion that one part
% --- p. 367 ---
of nature should let itself be \emph{explained} according to a teleological principle and another part not? That would
surely be the same as to assert that one part of nature might be assumed to \emph{act} teleologically and another part
% MISPRINT: „teleologik“ stands thus printed, without the s, in BOTH copies (checked on the KB
% colour scan at the word itself); the errata leaf does not correct it. Rendered here by the
% sense; read „teleologisk“.
not, and that although nature as nature is in both cases subject to the laws and acts blindly. One of two things: either
the whole of nature is teleological, or else every representation of a teleology of nature must be an imagination; if
there is to be teleology in a cockchafer, then there must also be teleology in the solar system. If there can be teleology
neither in the solar system's nor in the globe's development, then neither can there be teleology in any animal or animal
organ whatever. But if we now, by force of the teleological, make the whole of nature, without difference between the
lifeless and the living, into a self-organising self-being, what follows from this? We shall see it in the body of the
earth.

\runin{b) A Selfless Organism.} It is characteristic of the immanent philosophy of nature, and especially of the Hegelian,
to identify the idea that comes to existence in the body of the earth with the universal idea of life; only that the idea
of life in the geological organism is still at its lowest stage. Hegel accordingly conceives it in this its immediacy as
non-life, and the body of the earth itself as ``a corpse of the life-process''. The earth, he says, is a whole, is the
system of life, but as crystal in the form of a structure of bones; it may be regarded as dead, since its limbs still
seem to have a formal subsistence for themselves, and its process falls outside it. While the powers in this process show
themselves independent over against their product, the animal, as process in itself, has its powers in itself; its limbs
are potencies of this process. The earth, on the other hand, is only this, that it has this place in the solar system,
occupies this position in the series of the planets. Since every limb, as regards the animal, has the whole in itself,
% --- p. 368 ---
space's outside-one-another is cancelled in the soul; the soul is indeed everywhere in the body, but undivided, not as an
outside-one-another. The limbs of the geological organism, on the contrary, are in actuality outside one another and for
that reason unensouled. Because the earth is the universal individual, such moments as magnetism, electricity and chemism
step freely apart from one another in the meteorological process; but the animal is no longer any magnetism, and
electricity is only something subordinate in it.

% The letterspacing in the citation ends at „Subject“: „for den meteorologiske Proces“ is set
% in ordinary roman. Read on the KB copy.
None the less the earth does have a life. ``The earth's organism lying dead, this crystal of life, which has its concept
outside itself in the sidereal connexion but its peculiar process in a presupposed past, \emph{is the immediate subject}
of the meteorological process.'' Its first determinate life is the atmosphere. The meteorological process is, however,
not the earth's life-process, although the earth is enlivened by it; for this enlivening is only the real possibility
that subjectivity essentially proceeds from it as something living. As pure movement, as ideal substance, the atmosphere
doubtless has the life of the heavenly spheres in its essence, since its changes hang together with the heavenly motions;
but it at the same time materialises these in its elements. It is the dissolved, purely tensed earth, the relation of
gravity and heat; it runs through the periods of the year, the month and the day and expresses them as changes of heat
and gravity. At the equator the air is in a constant ebb and flood. As the heavenly motions materialise themselves in the
air, so on the other side sea and earth reach into it and volatilise themselves in it: a processless, immediate
transition. The neutral earth, the sea, is likewise a movement of ebb and flood. As the air, the universal element,
borrows its tension from the earth, so the sea its neutrality. The earth evaporates toward the air as sea; but toward the
sea the earth is a crystal
% --- p. 369 ---
that rids itself of the superfluous water, which gathers into rivers. But this, as fresh water, is only abstract
neutrality; the sea, on the contrary, is the physical neutrality into which the earth-crystal passes over. The origin of
the inexhaustible springs may therefore not be presented in a mechanical and quite superficial way (as a seeping
through), any more than, on the other side, the arising of volcanoes and hot springs; but as the springs are the earth's
\emph{lungs} and secretory vessels for evaporation, so the volcanoes are its \emph{liver}, in which it presents its own
self-seething. At all events we see regions, especially strata of sandstone, that always give off moisture. ``I do not
see the mountains as collectors of rainwater that penetrates into them. The genuine springs that beget such streams as
the Ganges, the Rhone, the Rhine have an inner life, a longing and a drivingness like naiads; the earth rids itself of
the abstractly fresh water, which in these outpourings hastens toward its more concrete element of life, the sea.''

In so far as what is figurative in this presentation is only to serve to make intuitable how the earthly whole is in a
peculiar way active in its parts, there is nothing to object against it. But already the circumstance that the geognostic
demonstration of the formation of springs and the arising of volcanoes is regarded as mechanical and superficial, while
the explanation with liver and lungs and naiads is given out as scientifically thorough, looks questionable; still more
uncomfortable is the impression of the selfless self, the subjectless subject, the great living corpse with pulse-beat
and breathing, with liver and lungs in full activity.

% The letter in this head is read on the page itself (KB copy, high magnification): it stands
% „c)“, as the sequence requires. THE TABLE OF CONTENTS, ON THE OTHER HAND, PRINTS „e)“ at this
% line (p. IX). This was at first taken for a misreading in the contents' OCR; on transcribing
% the contents the leaf was examined at high magnification in BOTH copies, and the letter has a
% closed bowl and a horizontal cross-stroke: it is the compositor's wrong sort, not a
% misreading. Both forms are reproduced as printed, each in its place.
\runin{c) An Inorganic Organism.} That the body of the earth is at once lifeless and yet living, inorganic and yet
organic, is according to the philosophy of immanence a contradiction that merely waits for a breath of dialectic in order
that organic life, with the whole infinite riches of the forms of individuality, may
% --- p. 370 ---
% CITATION PRACTICE: Nielsen opens „ anew after the interpolation („siger Hegel“, „siger
% Schelling“); but here each piece is closed for itself („Havet selv“ … „Organisation“ …), so
% the quotation marks come out even on this page. No imbalance to record.
break forth. ``The sea itself,'' says Hegel, ``is in a higher state of living than the air; it is the subject of
bitterness and of neutrality and of dissolution --- a living process that is always on the point of breaking out into
life, but that always falls back again into water, because this contains all the moments of that process: the subject's
point, the neutrality, and the dissolution of that subject in this.''

The fundamental thought itself had already earlier been developed by Schelling.\footnote{% printed mark: *) --- first of
% two notes on p. 370.
% The title is letterspaced, not set in italic; read on the KB copy.
\emph{Von der Weltseele}, eine Hypothese der höheren Physik. 3te Auflage. Hamburg 1809.}
``Organisation,'' says Schelling, ``is on the whole nothing else than the arrested stream of causes and effects. Only
where nature has not checked this stream does it flow ever further (in a straight line). Where it checks it, it returns
(in a circular line) into itself .... Since the world is infinite only in its finitude, and an unrestricted mechanism
% The omission is printed with FOUR points, as here.
would disturb itself, the universal mechanism too must be checked in infinity, and there will be as many \emph{single,
particular worlds} as there are spheres within which the universal mechanism returns into itself; and so, when all comes
to all, the world is --- an organisation, and a universal organism is itself the condition (and so far the positive) of
mechanism.''

% The printed mark **) stands BEFORE the colon: „saaledes**):“. Rendered as printed.
The relation between the inorganic and the organic is then from this point of view more closely determined thus%
\renewcommand{\thefootnote}{**)}%
\footnote{% printed mark: **) --- second of two notes on p. 370.
Wk.\ cit. pp.\ 234--235. pp.\ 298, 299 ff.}%
\renewcommand{\thefootnote}{*)}%
: nature is in its blind bondage to law to be free, and conversely in its full freedom to be bound by law; in this union
alone lies the concept of organisation. To solve the problem of organisation Schelling refers to the concept
\emph{drive}. Instead of saying that nature in its formations acts at
% --- p. 371 ---
once by law and freely, we can say that in organic matter there works an original formative drive, in which it takes on,
maintains and constantly repeats a determinate shape. But, it is added, \emph{formative drive} is only an expression for
that original union of freedom and bondage to law in all formations of nature, and not a ground of explanation for this
union itself. In the concept of formative drive there lies that the formation is not blind, that is, does not take place
by forces proper to matter as such, but that, besides \emph{the necessary} lying in these forces, \emph{the accidental}
of a foreign influence comes in, which, by modifying matter's forming forces, at the same time compels them to produce a
determinate shape. In this peculiar shape, which matter left to itself does not take on, lies precisely what is
accidental in every organisation; it is really this accidental element of formation that is expressed by the concept
formative drive.
% LETTERSPACING ON PART OF THE WORD (cf. p. 314): the compositor letterspaces only the
% compound's last member, „kraft“ and „drift“, while „Dannelses“ stands close; checked on the
% KB copy at high magnification on both words.
Formative \emph{force} accordingly becomes formative \emph{drive} as soon as to the dead effect of the first there comes
something accidental, a sort of disturbing influence of a foreign principle. This foreign principle cannot in turn be a
force, for force in general is something dead; but this dead element that lies in mere forces is precisely what is here
to be excluded. The concept \emph{vital force} is of course an altogether empty concept. \emph{Organisation} and
\emph{life} express on the whole nothing subsisting in itself, but only a determinate form of being, something common to
several co-operating causes. The forces which, so long as life lasts, are in play are no \emph{peculiar} forces proper to
organic nature; but that which sets those forces of nature into the play whose result is \emph{life} must be a peculiar
principle. This principle, which is the cause of life, cannot in turn be a product of
life.\footnote{% printed mark: *) --- the mark stands AFTER the full stop: „Livet.*)“.
% THE NOTE RUNS OVER TWO LEAVES: it begins beneath the text on p. 371 and continues beneath the
% text on p. 372 („chemischen Verwandtschaft …“), where the mark is NOT repeated. By the house
% rule the whole note stands here at its mark.
% Greek, set in italic Greek type, with the ϑ-form of theta: ἄφθαρτον (lenis and acute over the
% alpha), read on the KB copy.
Dieses Princip, da es Ursache des Lebens ist, kann nicht als \emph{Bestandtheil} in den Lebensproces eingehen; keiner
chemischen Verwandtschaft unterworfen, ist es das Unveränderliche (ἄφθαρτον) in jedem Organisirten. Wk.\ cit. p.\ 302.}
It must accordingly stand
% --- p. 372 ---
in \emph{immediate} connexion with life's first organs. It must be universally diffused, although it works only where it
finds a determinate receptivity. Thus the cause of magnetism is everywhere present, and yet works only on few bodies. The
magnetic current finds the vanishingly small needle on the open free sea as well as in the closed chamber, and where it
finds it, it gives it its polar direction. Thus the current of life, from wherever it may come, meets the organs that are
receptive to it, and gives them, where it meets them, life's activity. Since now this
principle\footnote{% printed mark: *) --- the page's own note; it stands beneath the
% continuation of the note from p. 371.
Wk.\ cit. p.\ 305.}
maintains the continuity of the organic and the inorganic world and knits the whole of nature together into a universal
organism, we recognise in it anew that essence which the oldest philosophers, divining it, greeted as \emph{the universal
soul of nature}, and which some physicists of that time held to be one with the forming and shaping aether (the portion
of the noblest natures).

What is hovering, ambiguous and contradictory in this otherwise profoundly conceived presentation of the principle of
life is evident. The principle is to be both something peculiar in itself and yet something quite universal; is to be the
cause of the life-process without being able to enter into the development of life; its activity is to depend on
``meeting'' receptive organs, while the organs are precisely conditioned by the organising activity. Its intervention in
matter is to be determined by law and yet to utter itself as something accidental that exercises a certain disturbing
influence upon the dead forces;
% --- p. 373 ---
the principle is to be different from all forces, is both to hover above and yet to dwell in the world of the elementary
forces, is to be immaterial and yet a finer matter, one with the forming, shaping aether; more ambiguous a principle of
life can hardly be.

If organisation is essentially to be known by nature's arresting the otherwise continuous stream of causes and effects
and bending it round into a circular line, then such great totalities as the solar system and the globe must certainly be
taken for organisms; but then an express difference must also be made between lifeless and living organisms. The lifeless
organisms are \emph{inorganic}; that is to say, they are without organs. As elementary organism the globe is organised
thoroughly in another way than the solar system, the merely mechanical organism; none the less the body of the earth is
just as far from being alive as the solar system is. The inorganic organism is always recognisable by this, that the
essence in the selfless whole is realised as the being of matter. The being of matter in such a whole goes immediately
together with the parts, and therefore the parts can not possibly become organs. An organ is an instrument; such an
instrument presupposes an indwelling essence that forms and uses it, and hence an essence reflected in self-activity, a
self-being. However intimately the self-being may be grown together with the being of matter, there is nevertheless
between the two an essential difference. Such a difference is not there in the elemental whole, however thoroughly
organised it may be. The body of the earth, a being of matter articulated into nothing but forms of elementary otherness,
has no organs at all, can form none and use none. It is not merely a speculative childishness, it is a thoughtlessness of
principle to speak of the globe's organs, its liver, lungs, digestive apparatus and the like. There is in the
\emph{inorganic},
% --- p. 374 ---
% The letterspacing runs over the seam: „uorganiske“ closes p. 373, „organløse“ opens p. 374;
% the „d. v. s.“ between them is NOT letterspaced. Read on the KB copy.
that is, \emph{organless} organism not a spark of life, not a glimmer of living, neither in the parts, in the air or in
the sea, nor in the whole taken together. The inorganic organism belongs in each and all to elementary nature, but, be it
noted, so organised that it forms a real transition to the organic.

% THE CLOSE OF PART TWO. The text on p. 374 does NOT run to the foot: it ends after six lines,
% about a fifth of the way down, with „… til den organiske.“, whereupon follows blank setting
% and then a short centred rule about a third of the way down the leaf; the rest of the page is
% blank. The rule is, as at the close of Part One (p. 196), set DOUBLE (two thin parallel lines
% close together); by the house rule that difference is not reproduced --- \streg serves
% everywhere. Checked on both copies. p. 374 bears no sheet signature, and Part Three begins
% NOT on this leaf but on p. 375.
\streg

% ============================================================
%  Part Three. Organic Nature.
% ============================================================
% Danish: \deel{Tredie Deel.}{Den organiske Natur.}
\deel{Part Three.}{Organic Nature.}
% A short centred rule stands between the Deel head and the Kapitel head.
\streg

% Danish: \kapitel{Første Kapitel.}{Det organiske Livs Oprindelse.}   (printed pp. 375--444)
\kapitel{Chapter One.}{The Origin of Organic Life.}

% Danish: \afdeling{A.}{Liv og Organisme.}   (printed pp. 375--400)
\afdeling{A.}{Life and Organism.}

In opposition to the body of the earth, the elementary-selfless organism, the organic organisms, the living
individualities, are expressly to be conceived as particularly conditioned productions of a universal principle of
selfhood dwelling in the fundamental medium --- as self-organisms. Life in substantial universality is not an immediately
existing ``real'', but is nevertheless an essence subsisting in itself, a self-being formed into the ground of nature by
Spirit itself. How such a nature-bound self-being can step forth as an essence for itself and yet as an intermediate
essence between object-power in elementary nature and Spirit's own absolute self-power can become intelligible only when
the investigations are brought under categorically determined points of view. We fix our attention upon
% The letterspacing runs over the seam: „Almeenlivet som organiserende Selvvæsen“ closes p. 375,
% and the second letterspaced phrase continues across the turn of the leaf with „Stofvæsen“ on
% p. 376. The „og paa“ between the two is NOT letterspaced. Checked on the KB copy.
\emph{universal life as organising self-being} and upon \emph{the organic opposition of self-being and
% --- p. 376 ---
being of matter}, in order thereby to work our way forward to a deeper understanding of \emph{the self-being's
constitution in matter.}

% The run-in head is printed in bold, as the house rule found it on pp. 79, 101 and 107; the
% difference is not reproduced.
\runin{a) Life as Organising Self-Being.} The representation of a universal life streaming through the whole of nature is
false when it is taken in hylozoism's sense, as though there were life in all matter, as though all things were at bottom
living; but it contains a truth when the meaning is that actual nature, the mechanical just as well as the elementary, is
an integrating part of Spirit's own life, that Spirit lives --- in nature too.

If one asks how it is possible that Spirit can live in nature without making the whole of nature living, then the answer
lies in all that has gone before. From the fundamental opposition between self-power and power in another there follows a
corresponding opposition between ideal and real otherness, ideal and real activity, ideal and real motion. Out from the
absolutely universal centrality of the unity of self, Spirit moves through the actual world of things, the innumerable
realities and real determinations of the medium of objects, unceasingly back into itself.
% spacing.py reports a run on „tilbage i sig selv“; on the KB copy the line is ordinary,
% closely set roman. Rejected: narrow justification in the bitonal working copy.
This ideal circular motion is one with the infinite capacity by which the eternally free self-power reduces the whole
natural world to a moment for the revelation of its inexhaustible, content-laden dispositions of wisdom: thus Spirit
lives in nature. But object-power is, in consequence of the doubling of power, in turn so mighty that in absolute
dependence it preserves an actual independence. So there form themselves in the medium of objects mechanical and
elementary totalities, totalities of unity of connexion without unity of self, totalities of otherness's selfless
essence, that is, lifeless totalities. From this it is seen that a converse inference such as this --- Spirit lives in
nature, \textit{ergo} the whole of nature must be living --- is false.
% „ergo“ is set in italic against the surrounding roman; read on the KB copy.

By running together the elementary activity of the forces of nature with Spirit's living in the whole of nature one
neutralises
% --- p. 377 ---
nature and Spirit; the unclear result of the neutralising is the mystical representation of a world-soul. The moment of
truth in this representation is that Spirit's objectification of the elementary, selfless essence yields the
presupposition for its objectification of the universal essence of the life of nature, the organising self-being in the
physiological sense. What is peculiar to this essence is that as a being of nature it goes immediately together with the
being of matter, and yet as self-being is qualitatively different from matter, the selfless other-being.

In its unity with matter the organising self-being works as though everything were accomplished by the substances of
matter's own indwelling forces. Light and heat, electricity and chemism, indeed even mechanically moving force, belong so
essentially to the activity of life that the beholder gets an impression as though functions of selfhood and functions of
another ran out into one.

Light and heat, electricity and chemism present analogies with life; but these analogies can dazzle only those who
overlook the essential difference between selfhood and otherness. To confuse selfhood with otherness is, remarkably
enough, a blunder with which the physicists and philosophers of the present day seem in equal degree to have fallen in
love --- though from highly different motives. The philosophers set out to dissolve matter into points of force and to
make the points of force into living beings:
% LETTERSPACING ON PART OF THE SENTENCE (cf. pp. 314 and 371): the compositor letterspaces
% „Andetheden bliver“ and „et Skin af Selvheden.“ but lets „fra dette Synspunkt“ stand close.
% Exactly the same happens in the answering sentence below. Both places read on the KB copy at
% high magnification.
\emph{otherness becomes}, from this point of view, \emph{a semblance of selfhood.} From the opposite side the physicists
would have the activity of life set on the same footing as every other utterance of force proceeding from matter, so that
life itself must be regarded either as a peculiar force or as a resultant of co-operating elementary forces:
\emph{selfhood becomes}, from this point of view, \emph{a semblance of otherness.}

What peculiar kind of acuteness it is that would save life by annihilating matter, that is, by dissolving it
% --- p. 378 ---
into sense phenomena of active, space-filling forces, shall here --- in view of what we have striven to make clear about
the mechanical and elementary causes in the two preceding Parts --- be touched on only briefly. One presupposes an
outwardly actual space without any matter; that is a contradiction, for it is precisely actually extended matter that
makes space actual. One invents pure points of force into the space given beforehand; that is a thoughtlessness, for pure
points of force can as little as pure points of thought occupy any actually determinate place in space. What occupies a
place becomes in extension equal to that part of space which it fills; it becomes an extended thing. Actual determination
of place is impossible without an extension, however small and vanishing. Just as unreasonable is it to make the pure
points of force into living beings. If the points of force were to be living, they would surely have to be organic; each
of them would accordingly have to have otherness, the real opposition, the material element, in itself, although it is
precisely the material that is to be excluded. Finally there comes to this the wholly inconceivable question how the
whole swarm of living points is to be able to hover in space, how such points should be able to occupy each its
determinate place, co-operating as if according to plan and predetermined harmony, without there being any unity of cause
or determining fundamental unity. If there is no fundamental unity at all here, then we have blind accidentality, as
devoid of nature as of Spirit, to carry the whole. If there is to be a gathering, ordering unity here, then this unity
must have either otherness's or selfhood's form. If it has otherness's form, then materiality is at the same time posited
as what is really first, and the hypothesis of the living points' originality dissolves into nothing. If the fundamental
unity is itself to be a living unity, a comprehending, all-organising unity of self, then this unity of self, if it is to
count for more than an empty thought, must necessarily have
% --- p. 379 ---
real otherness, the material world, for an actually indwelling opposite.

Just as groundless and unscientific is the physical conception of life as \emph{a force}. What really is it that one
calls \emph{vital force}? An unknown cause! If no more is said, then it is much the same as nothing at all. If the force
is to be more closely determined, it must surely point to a property in some substance or other, a life-substance. Since
now no such life-substance can demonstrably be pointed out, it lies near to seek the unknown force in the series of the
other elementary forces. So one places vital force on the same footing as light, heat, electricity and so forth, only
with the difference that it works under conditions over which physics has no command. ``We can,'' says Liebig, ``compose
a crystal of alum out of its elements, out of sulphur, oxygen, potassium and aluminium, because we can to a certain
degree command their chemical affinity, the heat and therewith the ordering; but a part of sugar we cannot compose out of
its elements, for in the organising of the sugar atom's peculiar form the vital force co-operated, over which our will
cannot command in the same way as over heat, light, the force of gravity and so forth.''

One might now represent the matter as though the elementary forces sank to serving the vital force and thereby lost their
independence. But neither can physics acknowledge this assumption; it must in the most emphatic way maintain the
elementary forces' relative independence. ``Brain and muscle substance, the constituents of the blood, the milk, the bile
and so forth are compound atoms whose formation and subsistence rests upon the affinity active between their smallest
parts. It is affinity and no other force that effects their coming together. But light, heat, vital force, the force of
gravity exercise an altogether decisive influence upon the number of the single atoms that unite
% --- p. 380 ---
into a compound atom, as well as upon the manner in which they arrange themselves; they condition the form of the
properties, the peculiarity of the combinations, precisely because they have the capacity to impart motion to atoms at
rest and to annihilate motion by resistance.'' Vital force then sinks, according to this way of seeing, to becoming a
plain co-operating cause, a moment of cause among several, a \emph{condition.} ``Upon the elements' coming together into
a chemical combination the vital force has not the slightest influence; no element by itself alone is in a position to
serve for the nourishment and development of a plant or of the animal organism. All the substances that take part in the
life-process are lower groups of single atoms which by the vital force's influence come together into atoms of a higher
order. Under heat's dominion the chemical force conditions the form and properties of the least composite groups of
atoms; the vital force conditions the form and properties of the higher, the organised
atoms.''\footnote{% printed mark: *) --- the mark stands AFTER the quotation mark:
% „… Egenskaber.“*)“.
Liebig. Chem. Letters. 12th Letter.}

If vital force is now placed on the same footing as the other forces, then only a couple of questions remain: whence does
this force proceed? upon what does its activity rest? To let vital force rest, like the force of gravity, upon a
universal property of matter will not do; for the force of gravity works everywhere, vital force only in the living
organism. To let vital force proceed from a substance of its own, as the force of sulphur from sulphur, the force of
oxygen from oxygen, will not do either, since, as has been said, there is no peculiar life-substance. The so-called vital
force can accordingly not possibly be any peculiar force, but must be regarded as the resultant of the peculiar
co-operation of different forces --- of gravity, of heat, of magnetic, electric and chemical forces. Herewith selfhood is
denied and material otherness made one and all.

% --- p. 381 ---
It is now easy to see that the physicists' elementary conditions are just as little able to ground the organising
activity as the philosophers' points of force. Life's organising activity rests from the ground up upon a thoroughgoing
co-operation of selfhood and otherness. The differenceless unity, pure self-sameness, is an empty concept to which no
activity corresponds. That activity may set in, the like must take up a difference into itself, and the unity have an
opposite through which it can develop and confirm itself as actual unity of self. Since now the element of otherness goes
immediately in under the concept of power reflected in the self-being, the self-being becomes organising: in this sense
Spirit is the all-organising essence. But for that reason Spirit in its infinite self-power, its eternal freedom, is by
no means to be confused with an organic essence, a being of nature organising itself in matter.

The self-organising essence begins in matter in such a way that it forms itself out member by member by transforming and
articulating its material. It can, under the articulating development of form, become master over matter only in so far
as it enters upon those conditions of nature to which the transformation of the parts of matter is subject, and must for
that reason become dependent upon matter. This doubly-sidedness --- that the self-being in dependence upon matter governs
matter --- is in the decisive sense characteristic of the organising activity. The elementary parts work point for point
according to their own nature under otherness's laws, and yet --- remarkably! --- by this co-operation they come point
for point to work according to the self-being's demands and to satisfy selfhood's laws. It is seen in the formation of
the cell. To make the cell straightforwardly into the organic element and explain the organism by a growing together of
cells is just as superficial as to explain magnetism by elementary magnets. That a cell may form, its fundamental parts
must be just as originally subject to
% --- p. 382 ---
% The word „underlagte“ is divided over the turn of the leaf: „under-“ closes p. 381,
% „lagte,“ opens p. 382.
% In the working copy's left margin on p. 382 a pencilled vertical stroke runs along part of the
% setting. It is not text; the KB copy does not have it.
self-activity as is the swarm of already formed cells by whose transformation and fusing together the more composite
organism arises.

The formation of the cell and the formation of the organically conditioned self are in principle the same fundamental
formation. There is certainly a great distance and a thoroughgoing difference between the self-being in a unicellular
plant and the conscious self that man calls his I. None the less it becomes necessary, for physiology as well as for
psychology, to reflect upon the principle and to lead the separation in the outward back to the unity in the inward. The
organising self-being must at all stages of development maintain and preserve an essential likeness with itself; in the
contrary case it is no self-being. The physiologist who cannot explain the cell in such a way that therein are contained
presuppositions for an explanation of the conscious self cannot explain the cell at all, and has only a physically
superficial representation of organising activity. Conversely, the psychologist who does not know how to explain the
conscious self in such a way that the ideal formation of self casts light back upon the real formation of the cell cannot
explain the self at all. A remark such as that the cell has no self, and the self no cell, is an empty evasion.

To explain the formation of the cell we must begin with an immediate unity and yet an immediate opposition of self-being
and being of matter. That we in this way seem to begin with a contradiction says nothing, for the organising activity is
precisely the contradiction's solution.

There are two things that are here to be explained. The one is that protoplasm, nucleus, primordial membrane --- in
short, all a cell's parts --- order themselves in outwardness's matter-like form, so that elementary otherness is
recognisable at all points.
% Both OCR witnesses read here „Punkter,“ with a comma; the KB copy shows a full stop plainly.
% Set as on the leaf.
The other is that these parts, in motion and reciprocal action, everywhere answer, be it mediately or immediately, to
% --- p. 383 ---
demands of life determined from within; that the cell lives, that it is an organ for an essence living in it. To the
outer totality of determinations of otherness there answers accordingly an inner totality of determinations of selfhood;
herewith \emph{otherness is doubled.}

The doubling of otherness comes to light only one-sidedly in the formation of the cell; the ideal activity of life is
from the beginning an activity in matter and must for that reason look like an activity of matter. That the activity of
matter rests essentially upon a self-activity, and the self-activity in turn upon a development of content of the unity
of life's own indwelling moments of otherness, first becomes recognisable in proportion as organisms form themselves so
perfect that sensation, representation and voluntary motion come thereby to utterance. What indeed is sensation but a
result of organic activity --- a result, be it noted, in which the real has begun to take on the form of an ideal, a
result in which the self-being through its other (the affected organ) gets an immediate impression of itself. That the
impression of selfhood is under otherness's form is plainly uttered in representation. Representation --- the ideal other
--- is a production of the representing self, and yet the sensible representation is determined by the object by means of
the organ, the real other. Upon the unity of difference of the ideal and the real other is organically grounded that
opposition in which the organic self-being involuntarily grasps itself in its inner difference from the organs. Upon what
is involuntary in this self-determination rests voluntariness. Finally the turning point is reached, when the
representing, voluntarily self-determining self-being has got its circle of representation organised into so free an
ideal otherness that it can therethrough meet itself and become conscious of itself as a self.

From this it is seen that the inwardness which is set free at the highest stage of organising must already be secretly
present
% --- p. 384 ---
at the lowest. The same principle that organises the human body organises also the human soul and its self-consciousness.
But the principle of selfhood that organises the body must surely organise every cell in the body: the essence that is in
the self and the essence that is in the cell must in principle be the same essence.

% The division's second head stands some way down the leaf, separated from the preceding
% section by blank setting but WITHOUT a rule. This head too is printed in bold.
\runin{b) Organic Opposition of Self-Being and the Being of Matter.} The question how an immaterial essence can act upon
matter has already been answered in the doctrine of the elements. There is power in the ideal just as fully as, nay more
originally than in the immediately real. Life is capacity, ability: the self-being's unity with the being of matter must
be sought in the fundamental medium. None the less the relation between ideal and real is here another than in the merely
elementary. In the elementary formations every transition from possibility to actuality is a finitely conditioned
expression of power. Otherness's essence passes part by part into corresponding existence: what exists is the actualised
other. When oxygen and hydrogen are united into water, there is no peculiar water-essence, different from the essence of
oxygen and of hydrogen, that comes in to bring about the union in which it comes to existence; life, on the contrary, is
in its co-operation with matter an essence for itself, a self-being. Where life is to come to existence, it must work as
a particular principle. There is in the organising activity a double reciprocal action: 1) between the principle and the
material parts, and on this ground 2) between the parts themselves. This doubly-sided reciprocal action would be
impossible if there were not in the parts' own essence an ideality. The material parts' ideality shows itself in an inner
aptitude to be formed and transformed.

In how comprehensive a sense matter is disposed toward forming may be seen from the analyses of organic chemistry. ``The
general way in which organic bodies
% --- p. 385 ---
are composed is the cause of their likeness with regard to properties and chemical relations. Their general chemical
character is expressed by the ease with which they are decomposed, their great tendency to alter their composition under
the influence of vital force, heat, the strong acids and bases, the salt-formers and so forth, and thereby to convert
their elementary substances into new combinations. . . . In all organic bodies a part of the elementary substances must
be thought of as grouped into an element-like member of the combination, \emph{a compound radical} . . . The number of
the organic combinations seems to be unboundedly great, without our therefore daring to assume ``a similarly great number
of compound radicals''. It is ``highly probable that many organic bodies which seem to be single organic combinations
are in actuality composed by combinations of two or more different
radicals''.\footnote{% printed mark: *) --- the mark stands AFTER the full stop:
% „Radicaler“. *)“. --- The note is itself a citation; „parrede organiske Forbindelser“ is
% letterspaced in the note, while „Saadanne“ and „opstaae“ stand close. Read on the KB copy.
``Such \emph{paired organic combinations} arise frequently by the union of acids with certain other, either neutral or
acid, organic combinations, and that in such a peculiar way that the body combined with the acid (the pairing) does not
separate from it, but enters into the composition of the salt when, for example, the paired organic combination is
brought together with a base.'' F. Wöhler's Outline of Organic Chemistry. Prepared by Simon Groth. Copenhagen 1855.
Introduction.}
% UNBALANCED QUOTATION: the citation opens on p. 384 with „Den almindelige Maade …“ and is
% NEVER closed. The two inner citations on p. 385 („et lignende stort Antal …“ and „høist
% sandsynligt …“) are each opened and closed for themselves, so the pair of leaves pp. 384--385
% stands with one quotation mark in surplus. Printed as here; the house habit of reopening
% without closing is known from pp. 10, 11 and 16. Checked on the KB copy: there is no closing
% mark after „Forbindelser. . . .“ or after „Radical . . .“.

From an exclusively chemical standpoint it is naturally not possible to find out what and how much depends upon the
activity of life, and what is accomplished by the affinity of the substances alone --- so much the less as a whole number
of the products brought forth in the organism can also be presented by chemical art. No wonder that chemists and
physiologists are inclined to run the organic activity together with the chemical. That there is nevertheless a
difference here, and an essential difference, is
% --- p. 386 ---
% The word „umiskjendeligt“ is divided over the turn of the leaf: „umiskjende-“ closes p. 385,
% „ligt.“ opens p. 386.
% In the working copy's left margin on p. 386 runs a pencilled vertical serpentine stroke, and
% the word „Formvirksomheden“ in the second paragraph is underlined in pencil. Neither the
% stroke nor the underlining is text; the KB copy shows ordinary roman without underlining.
unmistakable. In merely chemical activity the producing and what is produced, the process and the product, must fall
apart from one another; in so far as the process is, the product is not yet; in so far as the product is, the process is
over.\footnote{% printed mark: *) --- the mark stands AFTER the full stop: „forbi.*)“.
% The leaf bears TWO notes, *) and **).
Cf. Lectures on ``Phil. Propaed.'' Copenhagen 1862. pp.\ 216--19.}
In the organising activity, on the contrary, the process is reflected in its product; what is produced is simultaneous
with its production: the vessels with which the organism operates are formed, shaped, repaired and maintained by the
operations themselves.

That the chemical is taken up as a moment in the activity of life appears especially from this, that the organically
composite formations are built upon nothing but formations of cells. It is here not merely a question of the quantities
of carbon, oxygen, hydrogen and nitrogen that go to such formations; the peculiarity lies in the activity of form. The
chemical does not fall outside or alongside the organic; but it is taken up and transformed therein: the chemical
activity is present in the activity of form, but this is more than chemical. However many organic acids, bases and
neutral substances chemistry may gradually learn to present, it never gets so far as to learn to present the
morphologically organising activity. The chemical processes outside the organism do not take place under the same
conditions, and accordingly \textit{in concreto} not in the same way as in the
organism%
% The leaf's second note; the mark **) stands BEFORE the semicolon: „Organismen**);“.
\renewcommand{\thefootnote}{**)}%
\footnote{% printed mark: **) --- second of two notes on p. 386.
That the fundamental chemical laws are the same goes without saying, but they are modified in their fulfilment by going
together with the laws of selfhood.}%
\renewcommand{\thefootnote}{*)}%
; but that they do not take place in the same way says in other words that they are not the same. According to life's
chemism the chemical goes together with the morphological; but in chemistry's chemism the organically morphological
falls outside the chemical: life's chemism is
% --- p. 387 ---
accordingly essentially different from chemistry's chemism. Not only outside but also within the organism the opposition
between the chemically living and the non-living is recognisable enough; there may be in a living organism many products
that do not take part in the activity of life.

What is adduced here contains nothing new; chemists and physiologists know very well that there is a difference between
the chemical in the activity of life and the chemical in an organic chemistry. The point about which the investigation
turns is wherein the difference really consists and how it is to be determined.

Those physiologists who assume that life cannot possibly be a self-being working in matter must see to helping themselves
by referring to some unknown condition or other. Whether this condition had best be called ``\emph{vital force}'' or
``\emph{molecular force}'' or a peculiar mechanical, electric or chemical ``\emph{motive force}'' is fairly indifferent;
the main thing is that the unknown is represented as a condition. From this point of view the dispute for and against
``\emph{vital force}'' can never end.
% The four letterspaced words all stand WITHIN the quotation marks, while the marks themselves
% are not letterspaced. spacing.py catches only two of them (it misses letterspacing inside
% quotation, cf. pp. 3 and 9); all four read on the KB copy.
Those who assume a vital force assume after all that this force must in a natural way work together with the other
forces; by this co-operation it is accordingly to be conceived as one of the organising cause's many subsidiary causes:
it becomes plainly and simply a condition. Those who reject vital force must put something in its place. If it is to be
molecular force, then it is this that is the condition. If one exchanges molecular force for motive force, then it is the
organ-forming motive force that is the condition. The fundamental contradiction is that the difference between the
organic and the inorganic, between self-being and the being of matter, is at once, in one breath, both posited and
effaced. It is posited, in that a peculiar condition decisive for life is assumed; it is effaced, in that the assumed
condition makes the cause of life of one kind with every other physically conditioned
% --- p. 388 ---
cause. When life, by help of the conditions, is made into a resultant of co-operating elementary forces, the
representation of a peculiar vital force breaks through anew. For it cannot be denied that the co-operation of forces
which constitutes life is quite peculiar; for the peculiar co-operation there must necessarily be a peculiar condition,
and what then is this condition but ``vital force''?

To conceal the embarrassment one sometimes has recourse to the expedient of conceiving life as ``a consequence of the
organism'', and quite forgets that the organism itself is a consequence of life. ``If we stop,'' it is said in a sketch
by Harting, ``the functions of life suddenly and in such a way that the form, that is to say, the organism's structure,
and at the same time the chemical composition, remains unaltered and by force of circumstances is obliged to remain
unaltered, then, as soon as the old state is again brought about, life will go its accustomed course; in other words: the
changes of matter and of form that constitute life's essence will begin again where they were for a time interrupted. A
seed that keeps its power of germination for centuries, a rotifer that lies whole years dried up on a glass plate, can
therefore be called neither living nor dead. They are only organic beings which under certain determinate conditions can
become living; in other words: they have vital capacity.'' The facts themselves no one will surely draw into
doubt\footnote{% printed mark: *) --- the mark stands BEFORE the semicolon: „Tvivl*);“.
% THE NOTE RUNS OVER TWO LEAVES: it begins beneath the text on p. 388 and continues beneath the
% text on p. 389 („bevise sig om, at deres Indvolde …“), where the mark is NOT repeated; the
% word „overbevise“ is divided over the seam. By the house rule the whole note stands here at
% its mark.
% „Ansichten über Natur- und Seelenleben“ is set in ITALIC and moreover LETTERSPACED; \textit
% is used and the letterspacing noted only here (cf. „Thema Coeli“, p. 142). „Tipula oleracea“
% is italic but NOT letterspaced. „Philosophie des Unbewussten“ stands in ordinary roman on
% p. 389 --- read on the KB copy --- and is therefore not italicised here.
In Autenrieth's \textit{Ansichten über Natur- und Seelenleben} there are notices of how cabbage caterpillars, pupae of
butterflies and larvae of \textit{Tipula oleracea} have been frozen into lumps of ice and called back to life again by
thawing; nay, it is even reported that fishes in the Arctic Sea, which had stiffened into so solid a mass of ice that one
could hew them in pieces with an axe and convince oneself that their entrails formed a solid, frozen-together lump, have
come to life by suitable warming. Cf. Ev. v. Hartmann. Philosophie des Unbewussten. Berlin. 1869. pp.\ 470--72. --- Even
if one deducts the story of the frozen fishes, it is nevertheless, as Hartmann also remarks, beyond doubt that an
organism in which every trace of life has disappeared can be revived.};
but everything depends on how they are interpreted.

% --- p. 389 ---
% p. 389 opens a NEW PARAGRAPH: the first line is indented in both copies, so it does NOT
% continue the sentence at the foot of p. 388. The foot of p. 389 carries the continuation of
% the footnote whose mark stands on p. 388; by the house rule that note is set whole at its
% mark above and is not repeated here.
Those who place the cause of life in a sum of conditions will naturally see in this a proof of their assertion. The
conditions for the activity of life can be subtracted and added quite as in elementary activity. A mixture of chlorine
and hydrogen is a necessary condition for the formation of hydrochloric acid; yet such a mixture may stand in the dark
without the combination setting in; a condition is still wanting, sunlight for example. If the mixture be set in the
sunlight, the combination takes place, and even with a report. In the same way with life. A grain of wheat that has lain
for some three thousand years in one of Egypt's tombs has in all that time preserved a compass of conditions of life and
yet been without life, because it has after all lacked some conditions --- moisture, soil and so forth. If one says that
in the germinable grain of wheat there is a slumbering life, then one misspeaks; for there are surely only some
conditions of life. If one says that life has departed, then one misspeaks too, for here after all is ``vital capacity'';
one might with the same right be able to say that motion has departed from a body that does not move. The whole
explanation is in the highest degree outward and superficial.

E. Hartmann's conception is more essential, but still unsatisfying. He sees very acutely that more is needed here than
the mere conditions, that life rests upon the action of an essence different from matter, and interprets the phenomenon
of the revival of frozen organisms in the following quite original way. Since an organism, so long as it is frozen, has
neither life nor soul, it follows from this that when life and
% --- p. 390 ---
soul return into it after the lapse of a certain time, this soul can no longer be held to be the \emph{same} as the one
that dwelt in the organism before the passage into the frozen state. Were two souls separated in time to be one and the
same, then it would be required for that that the activity of the first passed \emph{continuously} over into the activity
of the last; that the organism in question is the same, and the souls answering to it so far of \emph{like
constitution}, is not sufficient. It might, according to the current mode of representation, happen, when the old soul
had passed out at life's ceasing, that it was then \emph{another} soul that had passed in with life's return. ``What is
awry in the setting up of the question,'' the author adds, ``becomes, however, at once evident when one thinks of the
Unconscious's All-Unity and notes that old souls as well as new are activities, directed toward the same organism, of the
same essence of the All-Unity, which again \emph{seizes} and continues life in this organism so far as is possible
according to matter's laws.''\footnote{% printed mark: *) --- in the body, after the closing
% quotation mark.
Philosophie des Unbewussten. pp.\ 474--75.}
% EMPHASIS, p. 390: samme / continuerlig / lige Beskaffenhed / anden / griber all confirmed
% letterspaced on the KB copy. „Sjæl“ after „anden“, and „og fortsætter“ after „griber“, are
% PLAIN --- only the words marked here are spaced. The second „den samme“ is plain as well.

It follows strictly from the author's whole fundamental view that life, the essence of the All-Unity, must at all stages
of development be like itself: as life in the cell, as soul in the animal, as the conscious self in man. That the human
organism, so far as experience hitherto goes, cannot endure being frozen like a cabbage caterpillar or being shut in like
a frog does not concern life as life, nor the soul as soul. If we are now really to represent to ourselves that the soul
which ``passes into'' a frog that is revived after having lain shut in a piece of rock is not the soul that lived in the
frog before the shutting in, then we should also have to represent to ourselves that the soul which passed into a man who
came to life after having been not
% --- p. 391 ---
merely apparently dead but ``absolutely lifeless'' was not the old soul, the one the man had before, but a brand-new
soul, a soul that began the whole life, the life of consciousness too, over again. If one now finds that this is an
unreasonable, nay monstrous representation of the essence of the soul when one thinks of man's soul, then one must be
consistent and admit that it is just as monstrous when one thinks of the cabbage caterpillar's and the frog's soul. What
is monstrous lies in this, that the opposition of soul and organism is just as one-sidedly posited as it is absolutely
denied. From the one side the opposition is driven to such an extreme that one and the same organism can be ensouled
twice over without the one soul having anything to do with the other. From the other side the opposition has to such a
degree disappeared that what makes the soul a soul does not at all rest upon life's essence, which in unconscious
universality is common to all that lives, but solely and alone upon the organism's peculiarity: into the thawed cabbage
caterpillar there passes another soul than into the frog, because the caterpillar's organism is another than the frog's;
it is ``the essence of the All-Unity which \emph{seizes} and continues life'' in each of these organisms, ``so far as is
possible according to matter's laws''.
% „eensidig“ confirmed on the KB copy (the Google layer gives „ensidig“). EMPHASIS: „griber“
% letterspaced, „og fortsætter Livet“ plain --- KB.

% The run-in head is printed in bold on the same line as the paragraph it opens. Verified on
% the KB copy.
\runin{c) The Self-Being's Constitution in Matter.} The selfless totalities, the solar system and the globe, have
according to their constitution the original significance of unfolding otherness's reality, of actualising the content of
the concepts and thereby of the ideas of nature in that object-power through which Spirit confirms its absolute
self-power. The solar system, the mechanical totality, is otherness's hardest extreme, the immediate realm of power. That
this realm of power is none the less self-power's own moment comes especially to light in the globe's constitution. In
the body of the earth, the elementary totality, the selfless being of matter is worked through, articulated so richly and
% --- p. 392 ---
inwardly that the otherness fully objectified therein conditions and prepares the objectification of a new, peculiar
essence: life, the natural self-being. As the globe's
% „Selvvæsen.“ with a FULL STOP: both OCR witnesses read a comma, the KB copy shows a period.
constitution, though peculiar for itself, must --- as the real development shows --- be enclosed as a moment in the solar
system's, so the natural facts speak also for this, that universal life's constitution must, as regards the organisms on
the globe, be contained as a moment in the constitution of the earthly whole.

What lies in this conception, and how far it has validity, must be tested from the side of essence as well as from that
of actuality.

From the side of essence the doctrine of the self-being's constitution in matter is a straightforward and necessary
consequence of a view of the world in which all being and existence is brought under the two thoroughgoing principal
forms: selfhood and otherness.
% spacing.py flagged „Anlæg i“ here. REJECTED: the KB copy shows „Selvvæsenets Anlæg i Stoffet“
% in plain roman; the line is merely loosely set.
Everything possible, everything that is or is thought to be in nature's realm as well as in Spirit's, must necessarily be
brought either particularly under otherness's or particularly under selfhood's form, and on the whole be determined by
the reciprocity of both forms, since nothing can exist or be there in otherness's form without presupposing selfhood, and
no selfhood is possible without corresponding otherness. The opposition of Spirit's realm and nature's rests solely on
this, that otherness in Spirit as Spirit is ideal, in nature on the contrary real, that is, immediate and material. That
the opposition has its ground in an original unity is intelligible from this, that self-power and power in another is a
doubling of the same omnipotence, a confirmation of Spirit's absolutely overreaching sole power.\footnote{% printed mark:
% *) --- in the body, after the full stop.
% The title „Den absolute Magt.“ is LETTERSPACED, confirmed on the KB copy.
Fundamental Ideas' Logic in Brief. \emph{Absolute Power.} p.\ 37 ff.}
The cognition of life's essence as self-being now develops straightforwardly from
% --- p. 393 ---
the fundamental relation between nature and Spirit. In mechanical and elementary nature Spirit has posited its real
negation, its material other, other in otherness's form, the immediately existing other. Therefore the essence of things
is, with all its ideality, a selfless essence, and the causes of things selfless causes. Between the reality of otherness
in the outward and the free, inward, ideal otherness --- the concepts of intellect and ideas belonging to the medium of
self --- there now comes to light life's essence, a self-being which nevertheless in a peculiar way belongs to real
otherness. In that the self-being steps forth as a being of nature, Spirit has, in that circle of existence which is its
immediate counter-image, an essence posited in which it can meet itself by otherness's way and see its own reflection.
Religion teaches that man is created in God's image; the philosophy of nature climbs lower down and assures itself that
already the living cell is a materialised natural image of Spirit.

From the side of actuality the doctrine of the self-being's constitution in matter has so manifold points of connexion in
the phenomena and facts of life that investigating science will not easily be able to point to a single observation that
stands in conflict with it, nor to miss any essential element of grounding that it, fully developed and rightly
understood, should not contain.

In the altogether erroneous opinion that organic life could not possibly have arisen on the globe unless it had happened
by a miracle breaking the law-determined connexion of the natural causes, some investigators of nature have had recourse
to the desperate expedient of fetching germs of life down from distant regions of the heavens and transporting them by
help of meteoric stones. This exceedingly troublesome transport is quite superfluous when one considers that the act of
creation from which life proceeds is not essentially different from the one that lies at the ground of the formation of
the solar system and the globe
% --- p. 394 ---
% MISPRINTED FOLIO: this leaf prints „294“ for 394. Both copies agree, so it is the
% compositor's; the page is marked here by its true number.
--- an act of creation which, as we have seen, by no means breaks the laws of nature. The organic self-being has at
bottom the same supernatural origin and for that reason the same natural beginning as the selfless essence. Nor do we
need to patch up the old hypothesis, warmed over by more recent philosophers, of existing monads, reals, living points
of force and such like inconceivable single essences. If the self-being's constitution stands in inner, thorough
connexion with the other-being's, the globe's, constitution, then it is also herewith given that the one development of
constitution must follow the other, that organic life must break forth with law-determined necessity when the body of
the earth yields the conditions requisite for it.

Life has the capacity to form easily formable matter, but it has not, like Spirit, a self-power capable of giving itself
matter; the life that organises its matter receives the matter from without. That the matter comes from without by no
means signifies that the self-being has a monadic subsistence of its own over against, independent of, matter, any more
than it signifies that matter, existing matter, should be only a semblance formed by living points of force; it
signifies on the contrary that the organic essence is in the finite sense dependent upon matter, originally so bound to
matter that, although it is in truth a self-being, it must in the sensible world look like a being of matter.

The fact established by palaeontology, that the first, ``smallest'' life on earth must have come to breakthrough in
millions upon billions of organic points --- whether one will now explain the phenomenon by a simultaneous fulguration
of the vital capacity, or, appealing to the small organisms' immense power of multiplication, prefer to let the swarm of
organisms proceed from a single first organism --- can just as little be grounded by an assumption of invisibly existing
points of force, monadic single essences, as by
% --- p. 395 ---
% MISPRINTED FOLIO: this leaf prints „295“ for 395; both copies agree.
a reference to an indeterminate universal living in the earthly elementary essence. By assuming an indeterminable
manifoldness of single essences one misses the grounding unity. By referring to the elementary universal essence one
does indeed get something resembling a unity of principle; but how it comes about that this unity smouldering in the
ground of matter can at one stroke become punctual and break out in innumerable sparks is altogether inconceivable. Here
is a leap from the inorganic to the organic, and at the same time a leap from the universal to the single, which can be
motivated only by the self-being's own original fundamental division.

It has been shown in the doctrine of mechanical reality how the ideal unity of cause is multiplied by the parcelling out
of physical causes, in that it is divided without being parcelled out, so that it is present in every cause of a thing
both divided and undivided. What already holds of the selfless essence holds in a deeper sense of the self-being. Life
--- the universal constitution, enclosed within the constitution of the earth, which naturally presses forward to come
to development with the earthly whole --- begins its organising activity, its peculiar actualisation, with a radiating
out into innumerable points of selfhood. By this transition of essence from unity to manifoldness the self-being has by
no means become either dispersed or parcelled out: the whole universal self-being is part by part and yet undivided
active in every cell. By this fundamental division the living beings have their origin and beginning out of universal
life; upon this fundamental division rests the principle of life.

That the principle of life, like all principles of nature, proceeds from Spirit is certain; but Spirit as Spirit must
not for that reason be made one with the principle of life. The principle, the germ, is what is not yet developed, what
is to develop and for that reason needs a temporal self-development in development of another; the principle does not
point separately to the origin for itself any more
% --- p. 396 ---
than to the beginning: the principle (\textit{ἀρχή}) is in the natural sense the immediate unity of origin and beginning.
% The Greek is set in ITALIC against the surrounding roman; read off the KB copy at high
% magnification: smooth breathing on the alpha, acute on the eta.
As what is original, the principle is in itself a relative unconditioned that must begin by receiving conditions from
without; from this it is seen again how necessarily the organic self-being must be dependent upon matter. And yet life
makes itself, within the circle of the conditions, master over matter.

The peculiar way in which organic life, though dependent upon matter, governs matter is well suited to clear up the
relation between power and impotence. No elementary conditions or merely physical causes are in a position to use the
chemical affinities in so peculiar a way, or to bring them to actualise such individualised formations, as the principle
of life does. Life is an ideal power, a relative self-power, which reduces the material as its other to a subordinate,
serving element of transit, to a means. And yet from a sensible point of view it looks as if the real reality, the
actual power, were to be sought in matter. One can tear the plant apart, cut down the tree, kill the animal; by doing
the organism violence, or by placing it under such conditions that the circulation stops, one can bring life to a
standstill, nay even to disappear entirely. From ideality's point of view it is the power of matter, the real power,
that shows itself as essential impotence; from reality's point of view, on the contrary, it looks as if self-power were
an actual impotence. --- Essentially understood, selfhood, which in principle bears otherness in itself, is what is
original, the higher and so far the mightier; otherness, which has what integrates it outside itself, is on the contrary
what is derived, the lower, and so far the impotent, the weaker. But if now the emphasis falls, as actuality likewise
demands, upon the phenomenon, the immediacy striking in the fact, then the relation turns about; then otherness, which
has its strength precisely in the outwardly immediate,
% --- p. 397 ---
must necessarily show itself as what is original, the stronger, and that not merely \emph{although} it is, but
\emph{just because} it is, what according to its essence is derived, the weaker; likewise selfhood, whose validity rests
upon energy of reflexion, must necessarily present itself as what is derived, what is abstracted, a faint shading in
subjective reflexion, and that not merely \emph{although} it is, but just \emph{because} it is, what at bottom is itself
original, what in objective reflexion is eternal and inviolable.\footnote{% printed mark: *) --- in the body, after the
% full stop.
Lectures on ``Philos. Propaed.'' 1860--61. pp.\ 317--18.}
% EMPHASIS, p. 397, all confirmed on the KB copy. The two parallel clauses are NOT spaced
% alike: the first prints „u a g t e t … j u s t  f o r d i“, the second „u a g t e t … just
% f o r d i“ --- „just“ letterspaced in the first and plain roman in the second. Reproduced as
% printed; this is the partial-phrase letterspacing the book does elsewhere (pp. 314, 371), not
% a reading error.

It will in this connexion be possible to understand that the activity of life in an organism can be stopped without life
being therefore extinguished or annihilated. If life's essence were nothing but activity of life, in the same sense in
which a motion in space is nothing but change of place, then those would certainly be right who assert that life is
annihilated when the activity of life stops. But life's essence is precisely the cause of the activity. Although it lies
incontestably in the concept cause, and hence in the cause's essence, to have an effect, and likewise in the concept
effect, and hence in the effect's essence, to have a cause, logic must nevertheless already call attention to the fact
that the category of causality is a category of actuality, that there must consequently be not merely an actual unity
but also an actual opposition between cause and effect. The actual opposition now steps forth in existence in such a way
that the cause --- the original matter, the substantial essence --- can have the effect in itself as a possibility
waiting for conditions in order to be actualised, and that conversely, in so far as otherness rules, the existing effect
can continue after the cause has ceased to work. When one grants that the lifeless organism, the stiff-frozen larva, the
shut-in frog, can preserve its organic form although
% --- p. 398 ---
the organising cause has ceased to work, then one has surely ascribed to the effect a relative independence separated
from the cause; but then one ought also, for consistency's sake, to grant the cause a relative independence separated
from working and activity. To that one is indeed willing enough when the talk is of a cause of a thing such as oxygen,
carbon and the like. Why is one now not just as willing to make the same concession when the talk is of the cause of
life? The reason is easy to discover. The thing has properties; in the properties the forces can rest; accordingly the
thing can be cause in possibility (\textit{κατὰ δύναμιν}) without being so in actuality (\textit{κατ' ἐνέργειαν}). Those
% Both Greek phrases are set in ITALIC against the surrounding roman; the accents and
% breathings were read off the KB copy at high magnification.
who cannot bring themselves to hold life for an invisible thing, ``a real'', but place it in a peculiar mode of working
of the forces, are then compelled to let ``changes of matter and of form'' constitute life's essence and to declare life
annihilated when the organic activity has stopped. This mode of representation, with all the absurdities following from
it, falls away when life is referred to the fundamental medium and conceived as self-being.

The organic self-being has, like every actual being of nature, a threefold medium: in Spirit as ideal, in nature as
real, in the fundamental medium as ideal-real essence. Life can accordingly be bound to matter and rest in the inactive
organism, the real medium, although it is different from matter, because it essentially rests and has ideal subsistence
in Spirit's ideal medium; both determinations are united in this, that it rests in the fundamental medium. It is a law
of selfhood, a law of recollection, of inwardising, that is herewith fulfilled. Since all self-activity goes through
another back into itself, life's circular motion is to be conceived --- not as conscious memory --- but as an
inwardising in self-reflecting, a recollection. If this reflecting is existential in utterances of power, then the
recollection gives itself to be known in real activity, in the expression of visible effects. Life's essence
% --- p. 399 ---
is working. But the reflexion of inwardness can also take place in ideal, invisible form, as pure reflexion of essence;
in this sense we can say that life, even when it does not work outwardly, nevertheless reflects itself unconsciously, in
the reflex of another, as essence in itself; life, when it rests, is not a nothing, nor a dead existent, but an essence
reflecting itself in recollection.\footnote{% printed mark: *) --- in the body, after the full
% stop; single asterisk confirmed on the KB copy (the Google layer reads „**)“).
To declare all activity of life ceased because it does not let itself be observed empirically is unwarranted; an
activity may be vanishingly small without therefore being 0. Thus the motion of a pendulum with a small angle of
deflection, when cv is the air's resistance, $\vartheta$ the amplitude and t the running time, can be designated by
$\vartheta = A_0\,e^{-\frac{1}{2}ct}\left\{\cos q\,t + \frac{c}{2q}\,\sin q\,t\right\}$,
% The page prints an UPRIGHT theta-variant (ϑ) and upright roman variables throughout; per the
% house rule the default mathematical italic stands and is not forced to \mathrm.
a formula which shows that the oscillation, objectively reflected, is continued into infinity after it has, in the
empirical sense, long since ceased.}

In the organism life has a natural --- unconscious --- recollection of its own self-development. Life rests as
possibility in its own product, in the result of a development which is only checked by outward causes and which can be
continued when the checking causes disappear. That life without recognisable organic activity is bound to the organism
and held fast at a given stage of development in the material is, from ideality's side, an expression of the fact that
the determinate result of development attached to the organism is recollected by Spirit in perfect knowledge. By force
of Spirit's eternally clear recollection, the cause of life --- resting in hiddenness, that is, in inwardness's
invisibility, in the ideality peculiarly determined by the development of the constitution --- is still united with its
matter-like fundamental form, its material exterior: \emph{this} life with \emph{this} organism, \emph{this} soul with
\emph{this} body, waiting
% EMPHASIS confirmed on the KB copy: only the four demonstratives are letterspaced; Liv,
% Organisme, Sjæl and Legeme are plain roman.
for the conditions that may come in from without and make a continued development possible. It is the recollection of
Spirit that,
% --- p. 400 ---
in accordance with the nature of the case, binds the resting life to the resting organism through the fundamental
medium.

That ``the Philosophy of the Unconscious'', for all its analyses, as profound as they are acute, must let a new soul
begin the activity of life over again with every animal organism revived from ``absolute lifelessness'', has its ground
in this, that it mistakes, in the unconscious as well as in the conscious, the originality of the self-being and
therewith the fundamental law of recollection.\footnote{% printed mark: *) --- in the body, after
% the full stop.
% PRINTER'S DEFECT, p. 400: the footnote prints „det Ubevidstes Philosopie“, without the h,
% where the body four lines above prints „det Ubevidstes Philosophie“. Confirmed on the KB copy,
% so the fault is the compositor's and not this copy's. Rendered here by the sense; the errata
% leaf does not touch it.
According to ``the Philosophy of the Unconscious'' one would almost have to assume that life stood to the organism as
light to the prism. If a prism with a certain angle of refraction were in one way or another made opaque, that would
answer roughly to the frozen organism. If the cause of the opacity were removed, it would answer to the warming of the
organism. As unreasonable as it now is to want to ask whether it is the same colours, that is, colours of the same
aether waves, that come forth by the subsequent refraction as those that came forth by the preceding one, just as
unreasonable would it then be to want to ask whether the life that is in the thawed organism is the same as that which
disappeared in the frozen one. That in such comparisons a confusion of spheres is committed is easy to see. There is a
quite other relation between life and organism than between light and prism. Light and prism stand to one another as
other to other; but the organism is a product of the activity of life. The organism is life's property, and life is its
own owner. By producing its organism life has in inwardness appropriated a determinate exterior and has in this its
exterior won a subsistence of its own, an actual, individual shape. So long as there is in such an organism a
possibility of life, it is solely and alone the possibility of that life which has produced the organism.}

% A short centred rule stands at the foot of p. 400, after „Grundlov.*)“ and ABOVE the
% footnote's own separator rule. It closes division A. Verified on the KB copy as well as on
% the working scan. Division A therefore closes INSIDE p. 400, with a rule and no head: there
% is no B. head on this leaf.
\streg

% Danish: \afdeling{B.}{Organiske Grundformer.}   (printed pp. 401--423)
\afdeling{B.}{Basic Organic Forms.}

% „articulerede“ stands rightly in the print here (p. 401); the errata leaf's correction
% concerns the occurrence on p. 407, not this one.
% The letterspacing of „Individer og Arter“, „Arter, Slægter“, „Selvhedslove“ and „Lovene for
% Andet“ is confirmed on the KB copy. „mellem“ and „Familier o. s. v.“ are NOT letterspaced.
That the organising activity must be a matter-forming, unconscious activity of recollection is recognisable in the basic
organic forms. The development of a basic organic form is a peculiarly determined development of constitution answering
to an individualising principle of selfhood. The individualising principle comes forth by an inwardly as well as
outwardly conditioned fundamental division of the essence articulated in its real concept, the universal, content-rich
self-being. Upon this fundamental division rests the substantial relation between \emph{individuals and species},
between \emph{species, genera}, families and so forth. The universal self-being is everywhere seen to form the
substratum of the shapes of separate types of formation and living specimens that step forth in existence. Organic
nature stands under \emph{laws of selfhood} which are in their kind no less strict than the elementary laws of otherness
holding for matter; but in so far as the self-being's organising activity is bound to matter, the fulfilment of the laws
of selfhood must punctually go together with a corresponding fulfilment of \emph{the laws for another}.

% Run-in head, read on the page itself (KB copy): set in bold roman, on the same line as the
% paragraph it opens.
% „Vierchow“ stands thus in the print (Virchow); read on the KB copy. Rendered as printed.
% The citation is reopened with „ after the interpolation „siger Vierchow“ and closed only
% once, at „for sig.“ --- the book's house habit; the quotation marks balance here.
\runin{a) Individuals and Species.} That it is really the individuals and not the species that exist, or, what says the
same, that the species exist only in and by the individuals, is a matter on which all are agreed. It is, on the other
hand, not so easy to agree about what the living individual really is and where it is originally to be sought. This
uncertainty concerns most nearly the plant individual. ``For the one,'' says Vierchow, ``the whole plant counts as
individual, for the second the branch or the shoot, for the third the leaf or the bud, for the fourth the cell, and each
of these ways of seeing has weighty grounds in its favour.'' E.~Hartmann remarks on this occasion
% --- p. 402 ---
% The letterspacing of „Enten---Eller“ and „Baade---Og“ is confirmed on the KB copy; spacing.py
% overlooks them because the dash breaks the word.
that each of the four is right when he posits this or that as individual, but wrong when he disputes the others'
assertions, since what is here in question is not an \emph{either---or} but only a \emph{both---and}.%
% The note is marked *) in the print.
\footnote{Philosophie des Unbewussten. p.\ 434.}

There is in the sense of principle nothing absolutely individual but selfhood and self. Spirit, the almighty self, is
the one, eternal individual because it is the one, eternal universal. The mediation of the individual with the universal
takes place in infinite self-division. The self-division happens by means of otherness. Through its ideal otherness, the
rich wellspring of ideas, the inexhaustible system of content-filled concepts, the absolute self divides itself into an
inner infinite manifoldness of points of selfhood: in every concept, in every separate idea the universal self is in a
peculiar way. In and by object-power's real otherness the effects become real. The ideal self-division with its real
effects is so far from either breaking or weakening the original unity that it is on the contrary the unity's true and
complete confirmation: Spirit's activity is a self-activity individualising into infinity. What Spirit, the absolute
self-being, does with eternal freedom in eternal knowledge, that life --- the derived, matter-bound self-being --- is
disposed to imitate in unconscious form of conditioned necessity. Life is by its self-division a principle
individualising in matter, and the cell, the first organic individual, individualisation's first immediate
representative.

% MISPRINT, NOT CORRECTED BY NIELSEN. The second note on p. 402 prints „8.“ where „S.“ (= page)
% should stand. The figure 8 has two closed bowls and is on the KB copy easy to distinguish
% from the open S in the note above, which prints „S. 434“ --- the two notes stand one beneath
% the other and give each other's calibration. The working copy has the same. Rendered here by
% the sense; read „S. 439--40“.
% The citation is reopened with „ after „siger Lewes**)“ and runs on to p. 403.
The individualised life in the cell has in a peculiar degree held the physiologists' attention. ``The last twenty
years,'' says Lewes%
% The note is marked **) in the print; the mark is set locally and restored afterwards.
{\renewcommand{\thefootnote}{**)}\footnote{The Physiology of Common Life. p.\ 439--40.}}, ``will in history go under the
name of the cell period, so manifold have the investigations been and so striking the results with regard to the part
% --- p. 403 ---
% The points of omission stand in the print as three spaced points, here as on p. 404 and in
% the note on p. 410. Read on the KB copy.
the cells play. The cell is regarded as the true biological atom. Nothing lives but cells, or what can be led
straightforwardly back to cells . . . Every cell in the organism is independent: it is born, it grows, it reproduces
itself and it dies, as though it were a unicellular plant or a unicellular animal''. The individual life in the cell is
brought out with so decisive a one-sidedness that the eye for the unity of life penetrating the whole organism is
completely clouded over. ``An organ's growth and decay is like a nation's or a tree's growth and decay. The single cells
of which the organ is composed grow and die, as the single men or the single leaves grow and die. There is a certain
aggregate unity, but it consists of separate individuals. As a nation's or a tree's life is the compass of the life in
all their single parts, so the life in an organism is the compass of the life in its single cells. This is one of the
newer science's great revelations. One no longer speaks of a centre for the activity of life, for one knows that there
are myriads of centres.''

% LETTERSPACED RUN, confirmed on the KB copy: „Det menneskelige Legeme,“ is letterspaced, the
% interpolation „hedder det,“ is set close, and thereafter the whole period „er eet; … en
% Eenhed.“ is letterspaced again. On the working scan „hedder det,“ also looks bold; the KB
% copy decides.
We have in this utterance --- an utterance that by no means stands alone --- an example of how a superficial conception
of life's essence must lead even expert and acute men to plain misinterpretation of well-known facts. The assertion that
the organism is only an aggregate of independent cell individuals stands in glaring contradiction with morphology's most
unmistakable fundamental phenomena. \emph{The human body,} it is said, \emph{is one; all its parts are subordinated, all
are combined into a higher unity. Our consciousness teaches us that our life is a unity.} ``This argument,'' says Lewes%
% The note is marked *) in the print.
\footnote{Work cited. pp.\ 433--34.}, ``points to an important
% --- p. 404 ---
fact, but it is wrongly interpreted. There is unity, there is an agreement between all the organism's parts. But this
is, in my opinion, not to be ascribed to a principle of life existing independently of the organism; it is owing to the
organic subordination; all the parts stand in relation to one another, all work in union with one another by help of the
nervous system, as all the parts of an army co-operate by help of officers and discipline . . . One can divide a polyp
or a worm into several pieces, and all the pieces continue to live and to grow; but one cannot after all assume that one
has in such cases divided the principle of life into several principles. There is a life in the parts and a life
belonging to the whole organism. Every microscopic cell possesses an independent existence, runs through its own course
from its birth to its death, and the total sum of these single lives constitutes what one calls the animal's life. The
unity is an aggregate of forces, not a single governing force.''

What concept is one now, after such a mode of representation, to form of the unity of life in an animal organism? The
analogy borrowed from the army does not fit at all. If there is to be discipline in the army, then the soldiers cannot
be independent; if every soldier is to possess an independent existence and run through his own course on his own hand
from his own individual impulse, then all discipline is dissolved. It is not even enough to say that the organism's
single parts are subordinated to one another and work in union, for the whole's dominion over the parts is so
thoroughgoing that it is precisely the formation of the whole that determines the formation of the single parts. Whether
a cell is to come forth by division or by new formation, and how it is to develop, depends notoriously on the part
allotted to it in the organism. Vessels and organs are not formed merely in consequence of a number of cells coming by
chance to work together; rather, cells work together and fuse together because
% --- p. 405 ---
the organism expressly demands more than the ``total sum'' of the parts' life; it is an overreaching unity of end that
makes the single parts into means.

What hinders the physiologist from seeing these simple truths is his own outward, forbidding representation of the
principle of life. When one denies the principle of life because it cannot ``exist independently of the organism'', then
one may with the same right deny the law of gravity; for that this essence cannot exist except in and by matter is
certain%
% The note is marked *) in the print and stands BEFORE the full stop, as printed.
% The letterspacing of „være og at existere“ within the note is confirmed on the KB copy;
% „mellem at“ before it is not letterspaced.
\footnote{That gravity's essence falls together with matter's makes nothing to the matter in this connexion. That an
essence may come to existence, it must beforehand be as essence in the ideal medium, be it a being of otherness or a
self-being. The emphasis falls on the opposition between \emph{being and existing}. The assertion that the principle of
life, in order to work in the organism, would have to ``exist independently of the organism'' is just as false as the
assertion that gravity, in order to work in matter, would have to exist independently of matter; nay, just as false as
this: that a real concept, in order to come to existence in the thing, would have to exist independently of the thing.}.
If a principle of life could not be because a physiologist would otherwise cut the principle into pieces when he cut the
polyp into pieces, then the organism's real concept could not be either; one would surely have to fear that the
physiologist would go and boil the concept when he boiled the organism.

The fact that one can divide a polyp, and the pieces still continue to live and grow, is in good agreement with the
principle of life, if only one will reflect on the essential difference there is between parcelling out and fundamental
division. The material goes to pieces when it is divided; but the ideal does not go to pieces. So long as the ideal is
held in otherness's form, the fundamental division does not strike the eye; the selfless essence has in power's
expression immediate existence, and the emphasis rests on the division
% --- p. 406 ---
of existence. But the self-being is altogether inconceivable when one overlooks the fundamental division.

% The letterspacing of „Slutningseenhed“ and „Individer og Arter“ is confirmed on the KB copy.
The unity of life answering to the cell-forming fundamental division is more than a plain unity of connexion; it is in
principle a \emph{unity of conclusion}. This unity of conclusion the physiologists have now conceived so backwards that
what in the matter is original is reduced to something derived, and the derived without more ado made into the original.
We see it in clear facts when we rightly reflect on the relation between \emph{individuals and species}.

One cannot divide a man as one can divide a polyp. Organisms of the lowest species can be divided, but individuals of
higher species cannot: the degree of divisibility is determined by the species. The fundamental difference between lower
and higher species must be known by the relation between fundamental division and joining together. The lower the
species stands, the weaker is the joining together and the more immediately does the fundamental division step forth. In
the polyp organism the total unity is so weak that the parts can each for itself appropriate a separate unity; in the
% --- p. 407 ---
human organism, on the contrary, the parts' unity is so subordinate that the total unity rules. None the less the
principle of life is, seen generally, like itself in the polyp as well as in man. The life that is in the polyp organism
is not a sum of independent, accidentally heaped-together unities of life; rather, the parts' many single lives are
points of selfhood out from one and the same unity of the species; only that the unity of self at this low stage is so
lost in its fundamental division that the totalised unity of conclusion becomes almost unrecognisable. In the cutting
asunder of the organism the conditions of existence are altered. The unity of the species is not cut asunder; but this
unity of self, as yet little concentrated in its constitution, has its points of selfhood in the parts in so immediate a
way that it particularly asserts itself under the form of fundamental division, repeats itself part by part yet
undivided, and lives in the parts separated by the cutting asunder. If we consider the human organism besides, then its
ruling unity is by no means any abstract single unity, but a totalised unity of conclusion accomplished through
fundamental division. That this unity of conclusion has passed through its fundamental division is seen from this, that
% THE ERRATA LEAF corrects p. 407, line 6 from the top: „anticulerede --- read articulerede“.
% The print has „anticulerede“, read on the KB copy at high magnification; the n is
% unambiguous. Kept as printed there; the English renders the sense.
% THE WORKING COPY IS HERE CORRECTED IN PENCIL: the n is written over and there is a pencilled
% insertion wedge between the lines above, so tesseract reads „atticulerede“. The KB copy
% decides the reading.
the richly articulated organism constitutes a many-sidedly ramified system of organ systems. The reason why the human
organism cannot be parcelled out as the polyp organism can lies accordingly by no means in there being here a lack of
inner division. On the contrary! What physiology makes clear about man's partial living, with relative centres of life,
and the organism's partial death is correct. No, the reason is that the relative centres of life in man enter as moments
into the central unity that overreaches them by the joining together. Thus the higher can be explained from the lower,
and conversely. How it comes about that the different organ systems can each lead its own relative life we learn from
the polyp organism; how the polyp organism, though it seems to be an aggregate of independent individuals, nevertheless
presents a living unity of essence, an original unity of self, we learn from the human organism. If it is really to be
one of the newer science's ``great revelations'' that one no longer speaks of a ``centre for the activity of life''
because one knows that there are ``myriads of centres'', then it must be a revelation that one is no longer in a
position to see the human organism, because one has stared oneself blind on nothing but polyp organisms.
% Nielsen here renders Lewes's words with „Midtpunkt“/„Midtpunkter“, while the citation on
% p. 403 prints „Middelpunkt“/„Middelpunkter“. Both places as printed.

% The letterspacing of „virkelige“ and „væsenlige“ is confirmed on the KB copy.
The \emph{actual} individuality in the existing individuals rests upon the \emph{essential} individuality in the species
concerned. As individuality's real concept, a self-being peculiarly determined by its constitution, the species has
ideal subsistence in a reflected unity of demands of function belonging together; by these demands of function the
organism's typical form is preformed. Individuality is actualised in proportion as
% --- p. 408 ---
the functions come into force: the species comes to existence in the individuals. That one and the same species can
exist in several individuals is, in consequence of the fundamental division, just as explicable as that a polyp can be
divided and the parts again become polyps. The species is self-being; the self-being fundamentally divides itself under
given conditions of existence; it does not by the self-division become several species, nor does it go to pieces, but it
realises itself part by part --- yet whole and undivided --- in a plurality of existing self-beings. The individuals,
the existing self-beings, can then in turn present the species in a more or less complete way: the existential
differences between individuals of the same species may not be confused with the essential differences between
individuals of different species.

From the fundamental relation between individuals and species there is further understood the relation between origin
and propagation. The first coming forth rests upon the self-being's developing its constitution in a given matter
prepared in an elementary way. To the constitution's fundamental division belongs the sexual opposition; the species is
realised in individuals of opposite sex; that is propagation's presupposition. Propagation does not take place by simple
mixture or chemical reciprocal action of sexual substances; it takes place essentially by the species. As the species is
already beforehand in the sexually opposed individuals and has its real moment of selfhood in each of the sexual
substances, so the species is in turn originally co-operative in the reciprocal action of the sexual substances; the
fruitful joining together is built upon a fundamental division: in the new individual too the species itself will part
by part --- yet undivided --- come to existence as self-being.

% Run-in head, read on the page itself (KB copy): set in bold roman, on the same line as the
% paragraph it opens.
\runin{b) Species and Genera.} The species has provisionally been regarded as an intermediate essence between universal
life and the living individuals; but between the species and universal life there shows itself in turn a manifoldness of
intermediate essences: genera, families, orders, classes and so forth. None of these intermediate essences has a
particular existence for itself; their inner difference depends solely
% --- p. 409 ---
upon the way in which the universal reflects itself in the particular. When the species is the particular, the genus is
the universal; when the genus is the particular, the family is the universal; when the family is the particular, the
order is the universal. This may then be expressed thus: the genera are species when the family is genus, and the
families species when the order is genus, and so forth.
% MISPRINT, NOT CORRECTED BY NIELSEN. The print has „o. s v.“ --- the full stop after „s“ is
% wanting. Both copies show the same, and the book elsewhere prints „o. s. v.“, thus two lines
% above on p. 408 and on p. 414. Rendered here by the sense.
It is accordingly sufficient for the development of the concept to keep to the categories species and genus.

If universal life be posited as genus, then there must be at least two species of life that come to existence, namely
plant life and animal life; both species must as separate constitutions originally be enclosed in the genus's
comprehending fundamental constitution. This fundamental constitution accordingly includes in possibility's form two
organic systems of function, namely that which is necessary for plant life, and besides it that which, in unity with the
former, belongs particularly to animal life. Which of the two species is first to step out into existence cannot be
derived from this; but it can be seen that the one cannot by a gradual alteration become the other. The boundary in
principle between plant and animal is determined by systems of function of unlike kind. Those functions that belong
particularly to animal life are so peculiar that they can neither come forth out of the merely vegetative ones nor be
added to them by outward circumstances.

That the animal functions cannot come forth by continued development of the merely vegetative lies in the nature of the
case; of growing comes only growth, of nutritive activity only nutrition: the whole plant world is a copious example of
how many-sidedly the vegetative functions can be developed without the boundary being overstepped. That the animal
functions cannot come forth by finite addition to the vegetative either is provable from this, that the functions
peculiar to animal life are functions of selfhood of a higher order than the merely vegetative. If the vegetative
functions cannot be derived from outward circumstances, then the animal ones can
% --- p. 410 ---
still less. Animal life, with its peculiar union of vegetative and animal functions, accordingly begins originally out
from its own unity of self determined by the species:%
% The note is marked *) in the print and stands AFTER the colon, as printed. It fills the whole
% lower half of p. 410 and does not run over to p. 411.
% A SINGLE RAISED COMMA stands in the note between „Udrugningen“ and the reopening „. It is NOT
% a whole “ mark: the book's closing mark consists of two raised commas, and here there stands
% only one, with a clear space before the following low „. Both copies have the mark, so it is
% in the setting and not in the copy. Rendered as one raised comma (‘).
% The letterspacing of „animale“, „dyriske“, „vegetative“, „Sliimbladet“ and „Karbladet“ is
% confirmed on the KB copy; „eller“ is not letterspaced.
\footnote{``The changes'' that take place with a fertilised hen's egg during incubation ‘``in the germ the first day are
the following: It spreads itself somewhat more to all sides and separates itself into two layers, of which the upper is
more transparent and is called the \emph{animal} or \emph{animalian} layer, the lower is more milk-white and is called
the \emph{vegetative} or \emph{mucous layer}. That they are so called is because from the upper there develop . . . all
the organs for the utterances of life proper to animals: the nervous system, the organs of sense and of motion; from the
lower, on the contrary, the organs for the life of nutrition, namely the intestinal tube, the lungs and all the
glandular ducts. Accordingly the division of the body's different organs into those that represent the animal and those
that represent the vegetative life is shown us by nature itself in the germ's first forms. Soon, however, a third layer
forms between these two. It is called the \emph{vascular layer}, for therein develop the heart, the blood vessels and
the blood itself, and hence as it were a system common to both those classes of organs.'' Eschricht. Twelve Lectures.
pp.\ 4--5. Cf. Lectures on Phil. Propaed. 1860--61. p.\ 345 ff.} The separation of plant and animal life is
qualitatively determined by the self-being's, the genus's, fundamental division. How originally animal life begins its
own constitution of selfhood as animal is seen the more clearly the higher the animal organism stands; but what animal
life is in the higher organisms it must as animal life also be in the lower.

Might there now perhaps be able to form itself between the two species an intermediate species, a sort of transitional
species? Such a species would then have to be partly both animal and plant, partly neither animal nor plant. But in so
far as the organism is both animal and plant, it is either two organisms, each belonging to its own species, and hence
merely an outward growing together of two species; or, if it is originally one organism, then it must necessarily be an
animal
% --- p. 411 ---
organism. If the species is to be designated by the organism's being neither plant nor animal, then it becomes not
merely indeterminable but in itself impossible. It is quite another thing that between the existing individual forms,
which each on its own side plainly utter the two species separated in principle, there may be intermediate forms so
ambiguous that it is difficult to find out to which side they belong.

% „to Plantearter; Kryptogamer“ with a SEMICOLON, while the parallel „tre Dyrearter: Bløddyr“
% has a colon. Both read on the KB copy; reproduced as printed, not harmonised.
With plant life as genus we get two plant species; cryptogams and phanerogams; with animal life as genus three animal
species: molluscs, articulates and vertebrates. How definitely these species are separated, how thoroughly the
fundamental division is carried through, may be seen in the systems of function and the typical individual forms. That
the fundamental difference is more strongly uttered in the animal than in the plant goes without saying. The organism
with the diffuse nervous system has originally another constitution, another unity of self, than the organism with the
ventral-ganglion system, and the ventral-ganglion system's fundamental form is in turn different from that of the
spinal-cord system. But however sharply the typical differences may be stamped in such individual forms as must be
regarded as the most perfect representatives of opposed species, this fact by no means excludes the possibility that
within each of the species there may be intermediate forms by which a transition from the one species to the other is
prepared. From a completely stamped mollusc organism to a completely stamped articulate organism the transition can
happen only by a leap. Does the principle of fundamental division now require a sudden leap from species to species
without preparatory intermediate members? By no means! The species are borne by the genus dwelling in them. There may be
in the genus's fundamental constitution so intimate a connexion between the species' separate constitutions that the one
species can become a medium of transit for an appearance of the other. A mollusc organism not yet typically completed
may by the genus determining its species be
% --- p. 412 ---
determined in such a way that, instead of moving toward its typical completion, it on the contrary, under gradual
alterations in its system of function, approaches that boundary at which a new species with a new system comes to
breakthrough by its transformation. With the new species there is posited, in consequence of the fundamental division, a
new system of function; but since the new system of function sets in by a transformation in principle of the old, the
fundamental division is mediated by the intermediate form%
% The note is marked *) in the print and stands BEFORE the full stop, as printed.
\footnote{Thus, to take an example from mathematics, the ellipse can varyingly approach the parabola; none the less it
continues to be an ellipse until that boundary is set where the parabola begins; the approach does not efface the
qualitative boundary.}.

It will be possible to see from this that the knots on the thread of life's development are by no means all of equal
size. The fundamental division that separates plant from animal goes deeper than that which separates the species within
a given animal genus; for the self-being has in animal life a capacity of centralisation that cannot possibly be derived
from plant life. The fundamental division by which the animal main types are separated is thoroughgoing in another way
than that by which the species belonging under each of the types are separated. There is a more essential difference
between articulates and molluscs than between articulates or between molluscs themselves. The mediating transitional
forms, if such there be, do not efface the qualitative difference between the species; but they introduce the nodal
points, they prepare the change in principle. An intermediate form of mollusc continues for all modifications to express
the species mollusc until the boundary is reached where the change happens by the breakthrough of a new species, that
is, of a new type.

We shall now see how the higher stands to the lower. Between man and animal the gulf may well be said to be just as
great as between animal and plant;
% --- p. 413 ---
but no greater either. If the separation of plant and animal can have come forth by a fundamental division of universal
life, then the same must hold of the separation of animal and man. What is new and set in with animal life, namely the
heightened concentration of the self-being by which a sense-consciousness is formed, answers exactly to what is new and
set in with human life and concentrates itself into self-consciousness. Something similar holds of the transitional
forms. If it is possible that there can be transitional forms from molluscs to articulates and from these to vertebrates
without the boundary of selfhood characteristic of the types being effaced, then for the human organism too there may
well be a transitional form which by no means effaces but on the contrary presupposes the qualitative boundary between
animal and man.

It is not abstract concepts, nor actual facts, but real possibilities, relations of constitution in principle, ideas
alive with force, that we here have before us. Universal life is not an indeterminate living but the idea of life
itself, the infinitely concrete idea of selfhood. Its ideal fundamental divisions reflect themselves in Spirit's medium
and are determined by self-power's individualising will; its real self-divisions unfold themselves through the
fundamental medium into systems of constitutions of constitution in naturally connected series. Before we go on to a
consideration of the facts belonging here, however, we must see to securing the course of the investigation by a more
exact determination of the laws of development.

% Run-in head, read on the page itself (KB copy): set in bold roman, the letter „c)“ too.
% The letterspacing of „Som Væsenet er, saa ere dets Love;“ and „Forskjel i Væsen giver
% Forskjel i Love;“ is confirmed on the KB copy.
\runin{c) Laws of Selfhood and Laws for Another.} Essence and laws are inseparable in nature's realm as in Spirit's.
\emph{As the essence is, so are its laws;} this fundamental proposition, which is as unshakable in the philosophy of
nature as in physics, yields a trustworthy starting point for mutual understanding. But from this there follows another,
equally irrefutable proposition, namely that \emph{difference in essence gives difference in laws;} against this
fundamental proposition
% --- p. 414 ---
the physiologists sin unceasingly. As surely as there is an essential, actual difference between the living and the
lifeless, there must also be an essential, actual difference between the self-being and the selfless other-being, and in
consequence of this an essential difference between laws of selfhood and laws for another.

In elementary nature the laws of selfhood do not come to light; all physical laws, however different they may otherwise
be, are nothing but laws of otherness. With organic life the laws of selfhood come into force; but this happens, in the
nature of the case, in such a way that their fulfilment is by no means separated from, but on the contrary goes as one
with, the fulfilment of the laws for another. If physiology could now keep itself exclusively at a physical standpoint,
exactly observing the laws of otherness alone, then the results of its penetrating researches, so highly important for
science, would be raised above every objection. But that cannot be done. It is impossible for the physiologist to treat
the doctrine of motor and sensory organs without availing himself of categories such as sensibility, sensation,
consciousness, voluntariness and so forth. As his insight into the categories is, so must his interpretation of the
phenomena be.
% The citation is reopened with „ after the interpolation „siger Lewes“ and is not closed until
% p. 415; the page therefore stands with one quotation mark too many. That is not a fault but a
% citation running over the turn of the leaf.
% The letterspacing of „Sammentrækninger“ is confirmed on the KB copy; the word is divided over
% two lines in the print and letterspaced in both halves.
% THE ERRATA LEAF corrects p. 414, line 4 from the foot: „tl --- read til“. The print has „tl“,
% read on the KB copy; the i is wanting altogether. Kept as printed there; the English renders
% the sense.
% THE WORKING COPY IS HERE CORRECTED IN PENCIL: an i is drawn in between t and l.
``The unity of the nervous system through the whole animal kingdom has,'' says Lewes, ``become generally acknowledged;
but, remarkably enough, the unity of consciousness has not been derived from it. As we climb up through the series we
begin with animals that have no real locomotion, and in which we find only \emph{contractions}, and find gradually in
the higher animals ever greater variety and composition; but at the same time we find that no new property has been
acquired, but that the same primitive tissue with its contractility extends into all the new phenomena. None the less
one speaks, and that with right, of simple contractions, of voluntary and involuntary motions, of flexion and extension
and so forth. It seems to me one ought
% --- p. 415 ---
% The page continues the sentence from p. 414 --- no paragraph break at the turn of the leaf.
% HERE the Lewes citation opened on p. 414 is closed; the page's balance is therefore --1.
% The marks of omission are three points in both places, but they are NOT justified alike: the
% first (after „Tankebevidsthed“) is spaced, the second (after „Tilværelse“) closer. Rendered
% each in its own way; the second is a borderline case.
% The letterspacing of „Systembevidsthed“, „Sandsebevidsthed“, „Tankebevidsthed“ and of „føle“
% and „tænker“ is confirmed on the KB copy.
to do the same with regard to sensibility or consciousness, which is suitably divided into three sections: 1)
\emph{system-consciousness}; 2) \emph{sense-consciousness}; 3) \emph{thought-consciousness}\ .\ .\ .\ 
System-consciousness comprehends all those sensations that arise in the body taken as a whole, especially during the
organic processes, and constitute the more formless and confluent elements that form the feeling of existence .\,.\,.
Few writers can free themselves from the troublesome ambiguities of our everyday speech. The ambiguity leads them here to
assume that the concept consciousness, and even sensation, includes something of thinking. But although every animal
must \emph{feel}, it does not follow from that that it \emph{thinks}. Every animal must contract itself, but from that
it does not follow that it has locomotion. Thinking is the result of a more composite organisation than is given to
every animal; but consciousness is an attribute of every animal organism. However simple the forms of consciousness may
be in the lower animals, it must nevertheless be essentially akin to the more composite forms in the higher animals. If
one does not grant this, it is impossible to say where sensation first appears. If the mollusc does not sense, neither
does the crustacean. If the crab is a machine, so is the bee; if the bee is one, so are the beaver, the elephant, the
dog, the ape.''

The naive, sensualistically vigorous way in which Lewes everywhere brings out the identity of the organic and the mental
activity; the acute criticism with which he knows how to lay bare the inconsistencies of half-theories; his
comprehensive endeavour to demonstrate an organic substratum common to the differing utterances of the life of
consciousness --- all this deserves high appreciation. And yet it is evident that he is not in a position to interpret a
single organic phenomenon from the ground up, because he from first to last identifies the laws of life with the
physical laws of otherness. Whence
% --- p. 416 ---
% All letterspacing on this page is confirmed on the KB copy: „Attribut for en dyrisk
% Organisme“ is letterspaced THROUGHOUT (including „for“ and „en“), while in the following
% citation only parts of the phrase are spaced --- „en for“ and „almindelig“ stand solid,
% „Nervesystemet“ and „Egenskab“ letterspaced. The same partial letterspacing as on pp. 314 and
% 371; reproduced as printed.
% „a priori“ is printed in ordinary roman, neither italic nor letterspaced (read on the KB
% copy); hence no \textit here.
do we know, then, that consciousness cannot possibly be ``\emph{an attribute of an animal organism}'' or, as Lewes
elsewhere expresses it, ``a property of \emph{the nervous system} in general''? We know it from experience, and yet a
priori, that is, by force of a priori experience.

It goes with the laws of selfhood in rational physiology as with the laws of otherness in rational mechanics: they are
not abstracted from outward phenomena, they are seen to lie in the essence of the matter. In order to see what
consciousness is, the physiologist must necessarily begin by becoming conscious of his consciousness. But one who has
once clearly become conscious of consciousness's essence will at the same time have made the discovery that the form of
consciousness is an original form of selfhood, every function of consciousness a function of selfhood, and all the laws
of consciousness without exception laws of selfhood. Consciousness consists in a \emph{becoming conscious of oneself}.
Consciousness is not a property of a thing or of a being of matter; a thing, a being of matter, cannot become conscious
of itself; or: what can become conscious of itself may as little be taken for a being of matter as a circle for a
square. \emph{Sensibility}, on the other hand, is a property of the being of matter, namely of the organ --- precisely
that intermediate determination in which functions of otherness go together with functions of selfhood. That an essence
becomes conscious of itself means that it notices itself in an other, is \emph{sensing}; that by the self-sensation
answering to the sensation of another it finds itself well or ill, is \emph{feeling} and distinguishes itself from
another, is \emph{representing}; that in relation to feeling and representation it utters itself with peculiar
self-determination --- \emph{desiring}, \emph{resisting}, involuntarily-voluntarily \emph{willing}. By such functions we
know selfhood. And by what, then, do we know otherness? Precisely by this, that an existent which acts on another and is
acted on by another lacks that doubling of reflexion by which an essence can through another turn self-confirmingly back
into
% --- p. 417 ---
itself. The selfless essence exists as a thing, and the thing has its properties, has actuality as simple or compound
substance of matter with corresponding accidents; its otherness is everywhere recognisable by this, that it determines
another only in so far as it is determined by another, that in all reciprocal actions it is subject to the law of
inertia, indifferent to motion and rest, indifferent to separation and combination, dissolution and subsistence. Once
one has fully reflected on the difference in principle between self-being and other-being, between functions of selfhood
and functions of another, laws of selfhood and laws for another, one will be astonished at the confusion of concepts
that rules in physiology when the self-being is made into a nerve tissue, a thing, and consciousness into a property of
the thing.\footnote{% printed mark: *), set after the full stop.
With such an understanding for a substratum it is an easy matter to cut consciousness into pieces when one cuts the
organism into pieces. The worst of it is that one cuts the laws of life into pieces at the same time.}

In order to make the difference between the self-being and the selfless essence plain we have particularly brought out
consciousness; for only in consciousness's mirror can it be seen what a self-being is. By this it is not meant that it
is consciousness that makes an essence a self-being --- the organising essence begins its activity in matter without
consciousness ---; but it is hereby inculcated that consciousness, and especially self-consciousness, is the purest,
clearest, most transparent form in which the self-being comes to revelation.

``The Philosophy of the Unconscious'' has introduced two functionaries to serve in the interpretation of the
physiological phenomena, namely ``the blind will'' and ``the unconscious representation''. The effect of this is at many
points surprising. How far a blind will and an unconscious representation really belong among the concepts that can bear
to be seen in the light of knowledge
% --- p. 418 ---
% THE ERRATA LEAF corrects p. 418, line 4 from the top: „betegne --- read betegner“. The KB copy
% prints „betegne et væsenligt Fremskridt“, without the r; reproduced as printed, the English
% likewise dropping the verb's ending. THE WORKING COPY IS HERE CORRECTED IN PENCIL: an r is
% added. Another of the pencil's places; the correction is Nielsen's own.
and science is a question for itself. But already the circumstance that appeal is made to functions of selfhood which
have this in common with functions of another, that they are unconscious, signify an essential advance. In all
reciprocal determinations of conscious and unconscious functions --- in transitions, fusings and shadings --- the
unconscious plays a principal part. From the lower nerve centres (spinal cord and ganglia) there proceed utterances of
will of which the individual cannot become conscious, because they do not pass through the brain. The voluntary motions
governed by the conscious will cannot --- it is further said --- possibly take place without the co-operation of an
absolutely unconscious will. If the conscious will would move, for example, a finger, then this must happen by help of
an unconscious will that can move precisely that nerve-end in the brain which causes the finger's motion; and to find
this nerve-end's situation the unconscious will must in turn presuppose an unconscious representation. Since now the
representation and the will that form the necessary intermediate members between the conscious will and the organic
motion do not come to consciousness in the brain, where the process takes place, they can as little come to
consciousness in lower centres, where the process does not take place, and must consequently be absolutely unconscious.
--- All this can be listened to. The point is that the unconscious functions as a self-being on a par with the
conscious; the unconscious can everywhere operate from the background and yield conditions of essence for corresponding
utterances of consciousness: nay, consciousness itself is to be only a passing phenomenon of the unconscious essence.
But when this unconscious essence and the conditions of activity under which it comes to light in consciousness's form
are more closely investigated, then ambiguities and contradictions show themselves; then it is seen that the laws of
selfhood are broken, that the confusion of concepts is --- still the old one!

Representation and will are ascribed to the unconscious, but
% --- p. 419 ---
% The letterspacing of „grundløs“, „sin“ and --- within the German citation --- „Wesen“ is
% confirmed on the KB copy. The German and „gar Nichts“ are printed in ordinary roman, not
% italic.
neither will nor representation belongs to its essence: the blind will has in an altogether \emph{groundless} way
proceeded from the unconscious One and has taken representation with it. There is talk in the usual way of
brain-consciousness, spinal-cord consciousness, ganglion-consciousness; but no one will easily be able to find out what
kind of essence it is that may be held for consciousness's real subject. To make a ganglion knot into a subject that
becomes conscious of itself stands in conflict with the author's assumption that consciousness is an activity of the
unconscious reacting against the material conditions. One would then suppose that the unconscious itself would have to
be consciousness's subject; but that stands in plain conflict with the declaration that the activity in which the
unconscious comes to revelation has nothing to do with its essence. It is brought out that it is only
brain-consciousness, not consciousness in the lower centres, that man calls \emph{his} consciousness. While accordingly
in ganglia and spinal cord there is only one subject of consciousness answering to each of these organs --- for the rest
essenceless, neither material nor immaterial --- there would in the brain be two whole subjects, namely one that the
brain has, on a par with the lower centres, and besides that one which man calls his I; but how essenceless this last
subject is, is expressly uttered in the sentence: ``Was an mir \emph{Wesen} ist, bin ich nicht''. One sees from this
that the eternally unconscious is to such a degree raised above all laws that it may be regarded as a true metaphysical
Proteus, functioning now as other-being, revealing itself in matter, now as self-being, revealing itself in the living
organisms, now again neither as other-being nor as self-being, revealing what it at bottom is: the eternally blessed
``gar Nichts''.

It is not enough that physiology has respect for the laws of otherness; it must see to getting the laws of selfhood in
too, for only in the union of the two can the laws of life be cognised. The laws of life
% --- p. 420 ---
% The letterspacing of „Erindringens Grundlov“, „Ernæringsloven er en physiologisk
% Erindringslov“, „nærer sig selv“, „Formudviklingens Lov er en morphologisk Erindriugslov“,
% „Organ“ and „Organisme“ is confirmed on the KB copy.
% MISPRINT, p. 420: the print has „Erindriugslov“ with u for n. Read on the KB copy AND on the
% working scan at 600 ppi: the letter has its bow below and opens upward, unlike the n in
% „Formudviklingens“ in the line above; both copies agree, so it is a turned or wrong sort in
% the setting and not a scanning fault. Rendered here by the sense; the errata leaf does not
% mention the place. The right form „Erindringslov“ stands three times elsewhere on the page
% and on p. 421.
% MISPRINT, p. 420: a superfluous comma in „Omdannelsen er, en virkelig Opfyldelse af
% Selvhedsloven.“ Both copies have the comma; reproduced, the English taking the same comma.
let themselves be brought under a common law, \emph{the fundamental law of recollection}.

\emph{The law of nutrition is a physiological law of recollection}. It is seen in assimilation. The living essence makes
the nutriment, the matter-like real other, into its other, makes the selfless essence into an integrating moment in its
own essence, and shows itself precisely in and by this activity as self-being. Assimilation does not transform matter
into self-being, any more than conversely; but it brings it about that matter enters into so intimate a union with the
living essence that this therein actualises itself. Actualisation of selfhood is inwardising, inward appropriation of
material otherness; the living essence nourishes itself, \emph{feeds itself} by recollection of another. That matter is
recollected means that it loses its raw, immediate otherness, that it is transformed. The transformation is conditioned
by matter's properties; so far the law of otherness must be fulfilled; but the self-being that effects the
transformation forms in every part of this other a part of its own. What is transformed, though in the material sense a
chemical product of another, is in the ideal sense of inwardness, of inwardising, a product of selfhood; the
transformation is, an actual fulfilment of the law of selfhood. If it is now seen how the two laws of unlike kind are in
assimilation collateral determinations of one and the same law of life, then it will also be possible to understand that
the self-being must not merely seize matter but expressly seize itself in matter, and that this self-seizing is an
inwardising of another, a recollecting.

\emph{The law of the development of form is a morphological law of recollection}. The assimilating activity is
determined from within by life's fundamental demand. The fundamental demand is that the self-being make its assimilated
other into a means, an instrument, that is, into an \emph{organ}. The fundamental demand divides itself, ideally into a
plurality of demands, really into a plurality of functions; so there is formed a system of organs, an \emph{organism}.
% --- p. 421 ---
% p. 421 begins WITHOUT indentation, in the same paragraph as the close of p. 420. TWO notes on
% the leaf, printed *) and **); the second set with a local \renewcommand, restored after. The
% printed mark *) stands BEFORE the full stop. The note **) RUNS OVER THE TURN OF THE LEAF: it
% begins here on p. 421, fills the whole foot of p. 422 beneath thirteen lines of text, and ends
% there; the book does not repeat the mark on the continuation. By the house rule the whole note
% is set at its mark on p. 421. The French citation (Milne Edwards) and the title are printed in
% ordinary roman, not italic.
% MISPRINT in note **): the print has „da maattc slige Organer“ with c for e. Both copies agree.
% Rendered here by the sense; the errata leaf does not mention the place.
% The letterspacing of „Bevidsthedsloven er en psychologisk Erindringslov“ and of „Inderligheden
% kan frigjøres“ is confirmed on the KB copy.
Every organ has its own form and structure because it has its own office; therein the fundamental division is
accomplished. All organs co-operate toward one purpose; therein the fundamental division's undivided unity is realised.
The teleological is that the end is recollected in the means and the means in the end; the typical is that the organs'
difference is recollected in the shape's unity. The plant's fundamental leaves (cotyledons) are recollected in the stem
leaves, these in turn in the flower leaves: organic recollection lies at the ground of metamorphosis. Upon recollection
rests the correlation of the organs: the beast of prey's claws are recollected in its teeth, and conversely. To the
preformation in the constitution there answers recollection in the constitution's development. By means of recollection
the real unity of connexion passes into the ideal unity of self, and the law for another is fulfilled together with the
law of selfhood. Life, the in itself shapeless essence, then not merely gets an outward shape by help of another; it
recollects to itself its other, takes shape therein and gives itself shape: the recollected organism is the self-being's
own real shape\footnote{% printed mark: *), set BEFORE the full stop.
The law of heredity too is a law of recollection.}.

\emph{The law of consciousness is a psychological law of recollection}. Nutrition and form-giving are functions of the
self-being, but in all unconsciousness. The transition from the unconscious to the conscious is, on the presupposition
of the self-being's corresponding constitution, a straightforward, law-bound consequence of the organising activity.
There is in the organising activity a bound inwardness; therefore the activity is unconscious. The inwardness is not
matter's but the self-being's; in this lies that \emph{the inwardness can be set free}: the first faint indication of
freed inwardness is the first glimmer of dawning consciousness.%
\renewcommand{\thefootnote}{**)}%
\footnote{% printed mark: **); the note runs over the turn of the leaf and fills the foot of
% p. 422; the mark is not repeated on the continuation.
If consciousness were really nothing but a ``property'' of certain organs, then such organs would surely be able,
immediately by their properties, to become conscious of the corresponding
% --- the note's continuation stands on p. 422 ---
impressions. But the nerve as nerve becomes conscious of nothing, and the brain as brain no more. --- Dans la fonction
de la sensibilité, il y a donc une division du travail bien remarquable: les parties qui, par leur contact avec les
corps étrangers, sont susceptibles de donner naissance à des sensations, ne peuvent pas percevoir elles-mêmes ces
impressions, et, d'un autre côté, l'organe qui est le siége exclusif de la perception de ces impressions ne peut
lui-même en recevoir directement; il est insensible et ne peut être excité par les impressions qui lui sont transmises
par l'intermédiaire des nerfs. Milne Edwards. Zoologie. P. 139.\par
If one will now say that it is the reciprocal action between the nerve's and the brain's properties that under certain
conditions brings forth consciousness, as it is the reciprocal action between oxygen's and hydrogen's activity that
brings forth the flame, then one must not forget that the decisive condition for consciousness's coming forth is the
presence of the living soul. If the soul is gone, no reciprocal action between the organs' properties can bring forth
consciousness. The reciprocal action is a condition for the function, but not the function itself; consciousness is as
function of selfhood a function of the soul, but by means of another an assimilation, an ideal nutrition, a
recollection.}%
\renewcommand{\thefootnote}{*)}
Consciousness
% --- p. 422 ---
% The page bears only thirteen lines of text; the rest of the leaf is filled by the continuation
% of note **) from p. 421, which is set at its mark above.
is: that the organic recollection is itself recollected. The recollection is no knowledge of the organising activity,
but a reflexion in the organised self-being, a concentrating act of selfhood by which a self is formed. It is not a new
essence with new laws that comes in; it is the same essence that raises itself to a new stage, the same law of life that
gets a new and richer fulfilment. Sensation is conditioned by the affected organ and the sensing self being at once
reflected in immediate separation and yet in immediate unity. To investigate how the matter stands in fact with this
immediate and yet reflected unity, together with the corresponding fulfilment of the doubly-sided law, becomes a task
for psycho-
% --- p. 423 ---
% One footnote, printed mark *), set BEFORE the full stop. DIVISION B ENDS HERE. At the foot of
% the page stands a short centred rule --- read both on the working scan and on the KB copy,
% where it stands clearly above and separate from the longer, left-set note rule. The head
% „C. Grundformernes Systematik.“ stands NOT on this leaf but at the head of p. 424.
physics\footnote{% printed mark: *), set BEFORE the full stop.
Fechner. Elemente der Psychophysik. Leipzig 1860.}. So much is clear, that it is neither the ganglion nor the nerve but
the self-being that in sensation is what senses. The question how an action upon an organ, a being of matter, can awaken
sensation in a self-being is then not at all more difficult than the question how a self-being can work organisingly in
matter. In that consciousness raises itself on the ground of the unconscious, the self-being comes to live in two
spheres, a conscious and an unconscious; but the reciprocal determination of the two spheres is surely explicable in the
most natural way from this, that it is the same self-being that lives in both. Here, as regards the voluntary motions,
there are not two wills --- a conscious one that determines the motion alongside an unconscious one that goes to the
brain and affects the nerve-end; it is the same act of self-determination that works, consciously as will, unconsciously
as nerve-moving factor. Nor are there here two kinds of representation --- a conscious one belonging to the will
alongside a blind one that shows the way to the right point in the brain; a blind representation is an imaginary
magnitude, a phrase that says nothing.

``The Philosophy of the Unconscious'' has in a comprehensive and penetrating way established that every psychologically
conscious act of selfhood has an unconscious side and must in consequence stand in hidden connexion with a manifoldness
of unconscious acts of selfhood. To have made this clear is a real merit. But that categories such as ``blind will'',
``unconscious representation'', ``unconscious knowledge'' and the like cannot possibly obtain the freedom of the city in
science, but must, when they are more closely analysed, dissolve into ``gar Nichts'', may be seen from the laws.

% Printed rule at the foot of p. 423: it closes division B.
\streg

% Danish: \afdeling{C.}{Grundformernes Systematik.}   (printed pp. 424--444)
\afdeling{C.}{The Systematics of the Basic Forms.}

% The letterspacing of „det umiddelbart Typiske“, „Classification“, „det abstract
% Teleologiske“, „Darwinshypothesen“, „Eenhed af det Typiske og det Teleologiske“ and „System
% af naturbetingede Skabelser“ is confirmed on the KB copy. „faste, uforanderlige Grændser“ is
% NOT letterspaced.
The morphological law is decisive in the determination of the basic forms. That this law is at the same time
teleological we have indicated above. But the necessity of holding fast the teleological in unity of principle with the
morphological will first be fully seen when science itself has furnished sufficient proofs that it is precisely the
separation of the two principal sides, inseparable in the principle of selfhood, that makes a satisfying systematics of
the basic organic forms impossible. By one-sidedly keeping to \emph{the immediately typical} one remains at a schematic
\emph{classification} of species and genera with fixed, unalterable boundaries. By one-sidedly laying the whole emphasis
on life's universal struggle and striving, and hence on \emph{the abstractly teleological}, one dissolves the types
% MISPRINT (not entered on the errata leaf): after „Darwinshypothesen“ a full stop is wanting;
% the following sentence begins with a capital K in „Kun“. The space is wide, but there is no
% full stop --- neither on the working copy (600 ppi, bitonal, clean white) nor on the KB copy
% (checked pixel by pixel on the raw image). Both copies therefore agree, so the fault is the
% compositor's. Reproduced as printed, the English likewise taking no full stop.
into fluid variations --- \emph{Darwin's hypothesis} Only by recognising the \emph{unity of the typical and the
teleological} grounded in the self-being's nature does it become possible to discover the law-determined relation
between fixed and variable boundaries in a continuous \emph{system of naturally conditioned creations.}

% Run-in head a), printed in bold roman, with a colon after „Grundlag“ and a full stop at the
% end --- read on the head itself, not in the table of contents.
\runin{a) Typical Basis: Classification.} Even a grouping of plants and animals supported by tolerably trustworthy marks
offers unmistakable advantages for survey and
% The note's mark *) stands BEFORE the full stop. The note is French, set in ordinary roman (not
% italic), and RUNS OVER TWO LEAVES: it begins at the foot of p. 424 and is closed at the foot
% of p. 425. The book does not repeat the mark on the continuation; the whole note is set here
% at its mark. --- „régne“ stands thus in the print, with an acute for a grave; likewise
% „espêce“ with a circumflex. Both read on the KB copy and reproduced as printed.
% The letterspacing of „classification“ in the note's penultimate line is confirmed on the KB
% copy.
characterisation\footnote{% the mark *) stands before the full stop
Telle est effectivement la marche adoptée par les naturalistes. On divise d'abord le régne animal en un certain nombre
de groupes de premier degré, caractérisés chacun par certaines particularités de structure; puis on subdivise chacun de
ces groupes, et l'on caractérise de la même manière les groupes secondaires ainsi formés; ces derniers sont à leur tour
divisés de nouveau, et l'on multiplie ces sections successivement suivant les besoins, jusqu'à ce qu'on arrive enfin à
ne laisser dans le même groupe que les divers individus d'une même espêce. C'est cet échafaudage de divisions, dont les
supérieures contiennent les inférieures, qui constitue ce que les naturalistes appellent une \emph{classification}.
Milne Edwards. Zoologie. Sixième Ed. Paris 1852. p. 279.}. But the grouping acquires scientific significance only
% --- p. 425 ---
% The letterspacing of „naturlige“ and „konstige.“ is confirmed on the KB copy. The Latin
% generic names on this leaf --- Galanthus, Fraxinus, Salicornia, Agrimonia, Spiræa, Alchemilla,
% Poterium --- are set in ORDINARY ROMAN, not italic. (Cf. pp. 435--436, where Trifolium
% pratense and T. incarnatum are on the contrary italicised.)
in proportion as the system-forming division, the classification, is laid out so as to make the typical recognisable. In
this respect the \emph{natural} systems have an essential advantage over the \emph{artificial.} The marks of character
used in an artificial system must indeed be taken from what is given in nature and so far be natural; but since by an
abstraction they are torn out of their remaining organic connexion for merely heuristic use, the heuristic may rather be
said to veil than to bring out what is typically peculiar. The Linnaean system (the sexual system), which is built upon
the number and other relations of the organs of fertilisation, furnishes a grand example of this. According to this
system plants that are otherwise quite unlike one another may come to stand near one another when with regard to the
chosen mark they happen to be alike. Thus for example a barberry and a snowdrop (Galanthus) both belong to the sixth
class, an ash (Fraxinus) and a glasswort (Salicornia) both to the second class, although nature by a highly different
mode of organisation has parted the barberry from the snowdrop and the ash far from the glasswort. On the other hand
plants which in all else show themselves as mutually akin may be removed far from one another when it happens that they
% tesseract reads an em-dash before „Det Linneiske System“; there is none on the KB copy, only
% the usual sentence space after the full stop. The dash is rejected.
differ with regard to the mark of classification. The Linnaean system places agrimony (Agrimonia) in the eleventh,
meadowsweet (Spiræa) in the twelfth, lady's mantle (Alchemilla) in the fourth and burnet (Poterium) in the twenty-first
class, although they
% --- p. 426 ---
% The letterspacing of „blodførende“ and „blodløse“ is confirmed on the KB copy. Greek, set in
% italic Greek type: ἔναιμα and ἄναιμα (lenis and acute over the epsilon and the alpha
% respectively, read on the KB copy at high magnification). By this edition's practice the Greek
% is set upright.
% The leaf bears TWO notes, *) and **); the second is set with a local \renewcommand and the
% mark restored afterwards.
according to their natural peculiarity are all members of one and the same family, namely the rose family. --- Zoological
classification continued steadily to conform more to Aristotle than to nature until Linnaeus at last made an end of the
division into \emph{blood-bearing} (ἔναιμα) and \emph{bloodless} (ἄναιμα) animals. Linnaeus's predecessor and
contemporary, Klein, exerted himself in vehement opposition to his brilliant rival's growing authority to ground a
system built upon the number of feet, toes and similar organs (Animalia pedata --- A. apoda); but the system became in
its formal execution so artificial that its one-sidedness could not but strike the
eye\footnote{% the mark *) stands before the full stop
Cf. Oscar Schmidt. Lehrbuch der Zoologie. Vienna 1854. pp.\ 21--22.}.

% The letterspacing of „embranchement“ (three of the four occurrences), „L'unité“, „l'unité“ and
% „multiplicité“ (both occurrences) in note **) is confirmed on the KB copy; the second
% occurrence of „embranchement“ is NOT letterspaced. „decide“ and „difference“ stand without
% accent in the print and are reproduced as printed.
Only with Cuvier does zoological classification get a sure, firm scientific basis; this basis is in the very strictest
sense typical: the plan uttered in the nervous system's structure is what is decisive%
\renewcommand{\thefootnote}{**)}%
\footnote{% printed mark: **) --- stands before the full stop
La forme du système nerveux détermine donc la forme de tout l'animal, et la raison en est simple: c'est qu'au fond le
système nerveux est tout l'animal en effet, et que tous les autres systèmes ne sont là que pour le servir et
l'entretenir. Il n'est donc pas étonnant que, la forme de ce système restant la même pour chaque \emph{embranchement},
la forme générale de chaque embranchement reste la même, et que, cette forme changeant d'un \emph{embranchement} à
l'autre, la forme de chaque \emph{embranchement} change . . . \emph{L'unité}, la \emph{multiplicité} de forme du système
nerveux, voilà ce que decide de \emph{l'unité}, de la \emph{multiplicité} des formes du règne animal . . . et c'est
encore du système nerveux que les diverses types, comparés entre eux, tirent leur titre de distinction et de difference.
Cuvier. Histoire de ses travaux par P. Flourens. Paris 1845. p. 98--99.}%
\renewcommand{\thefootnote}{*)}%
. --- For the one-sided, unsymmetrical or radiate animals (Animalia mollusca) without real limbs the diffuse, irregular
nervous
% --- p. 427 ---
% The French citation is letterspaced in its entirety --- confirmed on the KB copy.
system is suited; the ventral-ganglion system for the one-sided, articulated animals (A. articulata) without inner
skeleton, with many limbs or none; the spinal-cord system for the one-sided, unarticulated animals (A. vertebrata),
whose bone structure encloses the spinal cord, the nervous system's axis, and supports four inwardly articulated
limbs\footnote{% the mark *) stands before the semicolon
Cf. Chr. Fr. Lütken. The Animal Kingdom. Copenhagen 1855.}; that is the governing thought running through the
systematics, however many alterations may be made in the schema. It is characteristic that it was precisely by bringing
out the typical and applying comparative anatomy that Cuvier was put in a position to found palaeontology: ``\emph{Une
seule dent m'a, pour ainsi dire, tout annoncé!}''

The classification laid out upon the immediately typical is, however, not consistent. An organic type cannot be grasped
straightforwardly as something outwardly given; its actual existence presupposes an actual becoming, a development
answering to the use of the organs. To develop and realise all a type's constitution is doubtless the goal set for the
individuals belonging under it, and the marks of character can so far have both a typical and a teleological stamp. But
the type does not always come to equal or many-sided development; there may on development's way set in changes,
deviations, from outward as well as inward grounds. When then some organs are checked while others are transformed or
trained in a one-sided direction, there are deposited marks which, although they have originally proceeded from one
type, nevertheless, to judge by the outward, form boundary lines between several sharply separated types. The
classification keeps straightforwardly to given marks without regarding the physiological conditions, that use of the
organs by which the type is realised: it builds upon the outward but takes no account of the inward, which is
nevertheless the essential thing, namely the types' origin and laws of development.

% --- p. 428 ---
% No letterspacing on this leaf; spacing.py's two results are the rows of points in the
% citation. Confirmed on the KB copy.
% The note's mark *) stands BEFORE the comma: „Forsker*),“.
The one-sidedness of classification --- to separate the typical from the physiological and to lump the inessential
together with the essential --- has been censured in the most emphatic way by Darwin. ``Let us see,'' says the famous
investigator\footnote{% the mark *) stands before the comma
On the Origin of Species. Trans. by J. P. Jacobsen. Copenhagen 1872. p.\ 506 ff.}, ``what rules have been followed for
classification, and the difficulties that thereby present themselves. . . . One might have believed (and so it was
believed in old days) that those parts of the structure which were decisive for the habits of life and for each single
being's place in nature's housekeeping must be of the greatest importance for classification. Nothing can be more false.
There is no one who thinks that the outward likeness between a mouse and a shrew, between a dugong and a whale, or
between a whale and a fish signifies anything. These likenesses, although they stand in the closest connexion with the
being's whole life, are reckoned only as adaptive or analogical characters. . . . How remarkable it is that plants'
organs of vegetation, on which their nutrition and life depend, have little to signify, while the organs of reproduction
with their seeds and germs have exceedingly great significance! . . . That an organ's bare physiological significance
decides nothing about its worth for classification is, so to speak, proved by the fact that in nearly related groups, in
which the same organ, as we have every possible reason to assume, has nearly the same physiological value, its
classificatory value is highly different. . . . To name an example among the insects: in one large division of the wasps
the antennae are, as Westwood has remarked, in a high degree constant in their structure; in the other division they
differ much from one another; but here the differences have only a very
% --- p. 429 ---
% No footnote and no letterspacing on this leaf; spacing.py's two results („for“, „selv“)
% rejected on the KB copy.
subordinate worth for classification, and yet there is surely no one who will say that the antennae in these two
divisions of the same order have unequal physiological significance. . . . Further, there is no one who will say that
rudimentary or stunted organs are of great physiological significance, and yet organs in this state are undoubtedly
highly valuable for classification. No one will object that the rudimentary teeth in the upper jaw of young ruminants
and certain rudimentary bones in the leg are in a high degree fitted to show the near kinship there is between rodents
and pachyderms. Robert Brown has most strongly maintained that the position of the rudimentary flowers is of the highest
importance for the classification of the grasses. --- The significance that inessential characters can have for
classification depends mainly on their standing in correlation with several other characters of more or less
significance. . . . The importance of a collection of characters that stand in connexion with one another, even if none
of them are important, can alone explain Linnaeus's aphorism that the characters do not make the genus, but the genus
gives the character; for this utterance seems to be grounded on weight having been laid upon many subordinate points of
likeness that are too small to be defined. . . . In practice, when the investigators of nature are at their work, they
do not trouble themselves to ask after the physiological value of the characters they use to define a group or to put
some species or other in its place. If they find a character that is nearly the same and common to a great number of
forms and not common to others, they use it as one that has great significance; if it is common to a somewhat smaller
number, they hold it for a character of subordinate value.''

These and similar objections, drawn from evident inconsistencies, gather at last into one: the natural
% --- p. 430 ---
% The letterspacing of „Deductionen.“ is confirmed on the KB copy. The book's habit of
% REOPENING „ after an interpolated clause shows itself here three times, but the balance on the
% leaf comes out even: three „ and three “.
system still lacks the very most essential thing, that which alone can give it a title to be called natural, namely a
demonstration of the transitions by which the mutual connexion of the different types, standing nearer to and further
from one another, is grounded --- in one word, \emph{the deduction.} ``What is meant,'' asks Darwin, ``by this'' (the
classification's natural) ``system? Some writers regard it merely as a schema by help of which one can get those beings
put together that most resemble one another, and separated those that are most unlike one another; or as an artificial
way of expressing general propositions in as few words as possible. Others mean something more by it; they believe that
in the system one sees the Creator's plan; but unless one gives detailed information as to whether one means by this
order in time and space or in both, it does not appear to me that our knowledge has been increased. . . . Nothing can be
more hopeless than to attempt to explain the likeness in structure in members of the same class by help of the theory of
utility or the doctrine of final causes. . . . According to the common way of considering the matter --- that every
single being is the result of a particular act of creation --- we can only say that so it is; that it has pleased the
Creator to arrange all animals and plants in each great class after a uniform plan; but this is not any scientific
explanation.''

% Run-in head b), printed in bold roman, with a colon after „Grundlag“ and a full stop at the
% end. The word is divided over two lines, „Darwinshypo-“ / „thesen.“, and thus stands as ONE
% word with the genitive s. Read on the head itself on the KB copy, not in the contents.
\runin{b) Abstractly Teleological Basis: Darwin's Hypothesis.} If the physiological be set above the typical, the
emphasis must fall on the development of constitution, and the basis of the systematics becomes so far teleological.
Although the teleological everywhere individualises itself in typical forms, the end of life may nevertheless be
conceived from so general a point of view
% THE ERRATA LEAF corrects p. 430, lines 2--3 from the foot: „Phænomer --- read Phænomener“. The
% print has „Phæno-“ / „mer“, divided at the line-end, read on the KB copy. Kept as printed
% there; the English renders the sense. THE WORKING COPY IS HERE CORRECTED IN PENCIL.
that the types look like temporary, shifting phenomena of the universal life that struggles in the development of form
and works its way forward by different means along different ways.

% --- p. 431 ---
% MISPRINT (not entered on the errata leaf), p. 431, line 4 from the top: the print has
% „Betydniug“ for „Betydning“ --- a turned or wrong n. In roman type u and n are unambiguously
% different, and the letter has a clear bowl below and a serif tail to the right. Read on the KB
% copy at high magnification; tesseract likewise reads „Betydniug“ on the working copy, so both
% copies agree and the fault is the compositor's. Rendered here by the sense.
% The letterspacing of „Arter og Varieteter.“ is confirmed on the KB copy. The Latin generic
% names Pelargonium, Fuchsia, Calceolaria, Petunia, Rhododendron are set in ordinary roman, not
% italic.
It is this fundamental thought that carries the Darwinian hypothesis, though it shows itself only indirectly in the way
in which it has been put forward. We shall in all brevity single out the points of view that have particular
significance for the philosophy of nature.

\emph{Species and varieties.} That the seed from a determinate animal species can fertilise the egg from another has
long been a known matter. But bastard formation is as a rule possible only between related species --- between horse and
ass, dog and wolf, canary and siskin --- and, a circumstance on which particular weight has been laid, the bastards are
in the sexual respect so imperfect that they readily die out with the first generation unless pairing takes place with
one of the original species. This has, according to the way of seeing hitherto prevailing, been conceived as a proof
that the boundaries of species are constant, that nature maintains them by restricting the varieties and making the
bastards infertile. Darwin urges that such a narrow boundary for the crossing of species is in actuality not fixed. It
is known in how highly composite a way one has, as regards plants, crossed the species of Pelargonium, Fuchsia,
Calceolaria, Petunia, Rhododendron and so forth, and yet there are many of these hybrids that set abundant seed and
multiply just as perfectly as if they were natural species. As regards animals, the different races of every kind of
domestic animal are quite fertile when crossed with one another, and yet they descend in many cases from two or more
wild species. Either --- he concludes --- the original parent species must in the beginning have brought forth fertile
hybrids, or the hybrids since brought forth must have become quite fertile. The degree of fertility both at the first
crossing and in the hybrids runs right from the zero point up to perfect fertility. It depends in the first place on the
systematic affinity, that is to say, on the likeness between the species in the structure of parts
% --- p. 432 ---
% No footnote and no letterspacing on this leaf; confirmed on the KB copy.
that are of great physiological significance and that differ but little in the nearly related species: between species
that are referred by the systematists to different families hybrids have never been brought forth. Next it depends on
the reproductive organs. It is a fact that when animals and plants are removed from their natural conditions they are
``exceedingly inclined to have their reproductive system seriously affected''. When two species are crossed, the one
sometimes has preponderant power to form the hybrids in its own image. It is assumed that the ass in this respect has
the preponderance over the horse, so that both the hinny and the mule have more likeness to the ass than to the horse.
But the hybrids' fertility is not dependent on the degree to which they outwardly resemble one of the parents. Sterility
can accordingly by no means be regarded as the barrier nature has set to the mixing of species. ``For why should
sterility be so highly different in degree when different species were crossed, since it must after all be equally
important for them all that they should not be mixed? Why should the degree of sterility be by nature variable in
individuals of the same species? Why should some species cross with ease and yet bring forth very sterile hybrids, and
other species cross with the greatest difficulty and yet bring forth very fertile hybrids? Why should there also be so
great a difference in the results to which a reciprocal crossing between the same two species leads? Why in all the
world, one might ask, has it finally been permitted that hybrids should be brought forth? To grant the species the
peculiar capacity of bringing forth hybrids, and then to stop their further propagation by different degrees of
sterility that stand in no definite relation to the greater or lesser degree of ease with which the first crossing
between their parents came about, seems to be a
% --- p. 433 ---
% The leaf bears TWO notes, *) and **). NOTE **) RUNS OVER THREE LEAVES: it begins at the foot
% of p. 433, fills almost the whole of p. 434 beneath two lines of text, and is closed on p. 435
% above the note *) there. The book does not repeat the mark on the continuations; the whole
% note is set here at its mark.
% The letterspacing of „Varieringens Love.“, „Qvalitetsvalg“ and --- within note **) --- the
% phrases marked below is confirmed on the KB copy. „(natural Selection)“ is NOT letterspaced
% and NOT italicised.
singular arrangement''\footnote{% the mark *) stands after the quotation mark and before the
% full stop
On the Origin of Species. p.\ 350.}. If there comes to this that, apart from the question of fertility and sterility,
there shows itself a general and very great likeness between the offspring of crossed species and of crossed varieties,
then we understand that there may be ground for declaring the old marks of distinction invalid: the variety is a
modified species, the species a pronounced, relatively constant variety.

\emph{The laws of variation.} The species vary, and every viable, fertile variety strives to win subsistence as species.
What is teleological in this striving is, however, in the most various ways obscured and veiled by the accidental
co-operation of outward conditions. There are a quantity of small differences frequently found in the offspring of the
same parents; these individual differences are of the highest importance, for they are often inherited and thus furnish
a material that nature can use in the so-called \emph{selection of quality} (natural Selection)%
\renewcommand{\thefootnote}{**)}%
\footnote{% printed mark: **) --- stands before the full stop
There are with regard to the selection of quality undeniably utterances in Darwin that seem to reject every teleological
presupposition. He declares (wk. cit. p. 157) that the circumstance that little or no modification has taken place since
the beginning of the ice age has weight only as an objection against those who believe in \emph{a necessary indwelling
law of development}, but is powerless against the doctrine of the selection of quality, which maintains only that
\emph{occasionally}, in \emph{a few species, there occur variations or individual differences of a fortunate nature,
which are then preserved.} But such utterances must be understood from the connexion. His rejection of \textit{an innate
and necessary law of development} is plainly directed only against the way of seeing according to which it would be a
law that made a variation of the species impossible; for the basis of the whole doctrine of the selection of quality is
from first to last teleological --- though to be sure \emph{abstractly teleological.} A substance of matter asserts
itself by reacting in accordance with its properties; but no one will say that carbon, for example, struggles for its
existence. The living being must struggle for its existence; but to struggle for one's existence is teleological. The
struggling individuals are more or less fortunately placed, and the selection of quality signifies not merely that
nature lets the stronger live, but it signifies also that nature, the inner driving force,
% THE PRINT'S QUOTATION MARKS CHANGE FORM IN MID-SENTENCE: „individuelle Forskjelligheder“ has
% the usual low opening mark, but ”nyttige” and ”heldige” have a HIGH opening mark, of the same
% cut as the closing mark (thus ”…” and not „…“). Read on the KB copy at high magnification and
% compared line by line with the low „ two lines above. Reproduced as printed; it is the first
% occurrence of ” in the book.
preserves such ``individual differences'' as are ``useful'', ``fortunate'' for the individuals. To preserve such
peculiarities is precisely teleological; they are preserved not merely \emph{so that}, but expressly \emph{in order
that} the individuals may be preserved. It goes so far that the organisms which most need the selection of quality also
vary most. They accordingly vary not merely \emph{so that}, but in consequence of the indwelling formative drive,
essentially \emph{in order that} a selection of quality may set in. Let one read, for example (pp. 520--23), the
remarkable description of the disguises by which butterflies of one species (\textit{Leptalis}) escape pursuers by
imitating another species (\textit{Ithomia}).
% The two generic names in parenthesis are ITALICISED here; the following occurrence of
% „Leptalis“ within the citation is on the contrary set in ordinary roman. Both read on the KB
% copy. The citation opens with the usual low „ but CLOSES with ” (high, of the same cut as the
% opening mark in ”nyttige”) --- not with the book's usual “. Reproduced as printed.
``Mr. Bates may almost be said straightforwardly to have been a witness to the process by which the imitators have come
to resemble the imitated so closely; for he found that some of the forms of Leptalis which imitate many other
butterflies varied in an exceedingly high degree.'' --- In every illustrating example, in every interesting elucidation
of the selection of quality, the teleological plays a principal part. ``That disguise is far more frequent in insects
than in other animals is,'' Darwin adds, ``probably a consequence of their small size; insects cannot defend themselves,
excepting the genera that sting, and I have never heard it said that these imitated other insects, though they are
imitated. Insects cannot by flight escape the larger animals; therefore they are, like most weak creatures, restricted
to deceit and dissimulation.'' We need not take the last words literally to get the impression that in such disguises
there is something intended. Here is not merely a circle of outward circumstances and accidental conditions under which
the struggle of life is carried on; here is in the living beings themselves a circle of inner conditions and
corresponding constitutions; here are not merely causes of otherness that can be straightforwardly sensed; here are also
causes of selfhood that must necessarily be thought; only in the peculiar co-operation of both kinds of cause have we
the whole causality: this causality is teleological. If one throws the teleological basis away and makes the selection
of quality into a merely physical result of colliding conditions, then one misunderstands Darwin and breaks ``the
fundamental laws of the life of nature''; for it lies expressly in his presupposition that the struggling life, whose
many-sided self-development is the struggle's goal, must in matter be an essence different from matter. The materialists
with their physical primordial generation and their organic-physical doctrine of perfection may well make use of the
Darwinian hypothesis, but they cannot find bottom in it.}%
\renewcommand{\thefootnote}{*)}%
. What the selection of
% --- p. 434 ---
% Only two lines of text on this leaf; the rest is taken up by the continuation of note **) from
% p. 433, which is set whole at its mark above.
quality is, and how it works, is made clear by examples. Let one think of a wolf that lived on different kinds of
% [text to be added: pp. 435--444]

% Danish: \kapitel{Andet Kapitel.}{Organiske Naturer.}   (printed pp. 444--509)
\kapitel{Chapter Two.}{Organic Natures.}

% Danish: \afdeling{A.}{Plantenaturen.}   (printed pp. 445--465)
\afdeling{A.}{Plant Nature.}

% [text to be added: pp. 445--455]
% [text to be added: pp. 456--465]

% Danish: \afdeling{B.}{Dyrenaturen.}   (printed pp. 466--490)
\afdeling{B.}{Animal Nature.}

% [text to be added: pp. 466--478]
% [text to be added: pp. 479--490]

% Danish: \afdeling{C.}{Menneskenaturen.}   (printed pp. 491--509)
\afdeling{C.}{Human Nature.}

% [text to be added: pp. 491--500]
% [text to be added: pp. 501--509]

% Danish: \kapitel{Tredie Kapitel.}{Slutning: Natur og Aand.}   (printed pp. 509--549)
\kapitel{Chapter Three.}{Conclusion: Nature and Spirit.}

% Chapter intro: prose stands between the chapter head and the A. head.
% [text to be added: pp. 509--510 --- chapter intro]

% Danish: \afdeling{A.}{Naturen skabes af Aanden.}   (printed pp. 510--522)
\afdeling{A.}{Nature is Created by Spirit.}

% [text to be added: pp. 510--522]

% Danish: \afdeling{B.}{Naturen bæres af Aanden.}   (printed pp. 523--534)
\afdeling{B.}{Nature is Sustained by Spirit.}

% [text to be added: pp. 523--534]

% Danish: \afdeling{C.}{Naturen forklares af Aanden.}   (printed pp. 535--549)
\afdeling{C.}{Nature is Explained by Spirit.}

% [text to be added: pp. 535--542]
% [text to be added: pp. 543--549]

% ============================================================
%  TRYKFEIL --- the book's own errata leaf, unfoliated, following p. 549.
%  Its entries correct Danish type, so the leaf stands in the transcription
%  only and is not translated. Nielsen's corrections are noted in % comments
%  at their sites in both files.
% ============================================================

\end{document}
